\documentclass[11pt,a4paper]{article}
\usepackage[utf8]{inputenc}
\usepackage[T1]{fontenc}
\usepackage[french]{babel}
\usepackage{lmodern}
\usepackage{amsmath,amssymb,amsthm,mathtools}
\usepackage{physics}
\usepackage{siunitx}
\usepackage{bm}
\usepackage{geometry}
\usepackage{booktabs}
\usepackage{array}
\usepackage{enumitem}
\usepackage{hyperref}
\usepackage{microtype}
\geometry{margin=2.2cm}
\hypersetup{colorlinks=true,linkcolor=blue,citecolor=blue,urlcolor=blue}
\setlength{\parskip}{0.5em}
\setlength{\parindent}{0pt}

\title{Méthodes SRC/MPCD quasi-incompressibles :\\ projection Q6/Q9, Q6-g-f force-aware, tension superficielle de Laplace, surface libre et couplage liquide--gaz explicite, CUDA résident}
\author{}
\date{\today}

\begin{document}
\maketitle



\begin{abstract}
Ce rapport rassemble l'état de développement de la méthode SRC/MPCD quasi-incompressible en distinguant explicitement deux générations de code. La première, historique, reposait sur une redistribution particulaire dure vers une cible d'occupation et sur une fermeture liquide virielle. Elle reste utile pour l'interprétation physique, mais ses résultats quantitatifs doivent être relus avec prudence car une partie des validations a été effectuée avant la correction du thermostat, avant la stabilisation des conditions limites par particules virtuelles et avant les premières vectorisations MATLAB. La seconde génération, désormais considérée comme la chaîne de référence, est la méthode Q6/Q9 : collision SRC/MPCD classique, projection de vitesse incompressible Q6, correction basse fréquence du flux de masse Q9, thermostat cellulaire corrigé et diagnostics conservatifs associés.

Le socle historique Q6/Q9 validait quatre points principaux. Premièrement, les cas Taylor--Green et Poiseuille montrent que Q6/Q9 réduit les modes compressifs de grande longueur d'onde sans détruire les structures hydrodynamiques moyennes. Deuxièmement, les campagnes Poiseuille multi-résolutions ont mis en évidence une dépendance apparente de la viscosité identifiée au nombre de mailles transverse ; une correction géométrique en \(N_y^2\) ramène les viscosités corrigées vers la référence bulk Taylor--Green, ce qui impose de distinguer viscosité brute de fit et viscosité corrigée par la hauteur discrète du canal. Troisièmement, le piston staircase avec fermeture virielle bulk fournit une équation d'état contrôlée, avec \(P_{\mathrm{vir}}=K_{\mathrm{virial}}(\rho-\rho_{\mathrm{EOSRef}})\), \(P_{\mathrm{tot}}=P_{\mathrm{kin}}+P_{\mathrm{vir}}\), \(P_{\mathrm{drive}}=P_{\mathrm{tot}}\) et des fits staircase de pente quasi parfaite sur les grandeurs bulk. Quatrièmement, le diagnostic mécanique de paroi avec particules virtuelles \emph{wallVP} est maintenant séparé en impulsions signées, pressions positives de compression et moyennes cumulées ; le test d'équilibre gaz entre deux murs valide l'ordre de grandeur de la pression mécanique cumulée, tandis que le staircase piston demande encore des fenêtres longues pour réduire la variance de l'estimateur impulsionnel.



Enfin, le chantier CUDA-TOPO ajoute une brique de pénalisation Darcy--Brinkman pilotée par un champ externe $\chi$. Ce champ permet de représenter des géométries complexes et de construire des proxies de coût, de portance et de traînée pour des canaux, obstacles et profils NACA. Le point d'étape validé est toutefois nuancé : Darcy--Brinkman est robuste comme pénalisation poreuse volumique, mais ne reproduit pas à lui seul une frontière solide impenetrable. La comparaison Von Karman montre qu'un cylindre Darcy peut annuler la vitesse dans la zone pénalisée sans produire le déficit de vitesse et le sillage du solide immergé codé ; un futur chantier devra donc dériver une vraie condition solide à partir de $\chi$.

Les développements 0418--0426 révisent ce statut. Le champ $\chi$ n'est plus seulement utilisé comme frein Darcy pur : il pilote désormais une chaîne plus complète combinant désactivation initiale des particules dans les zones $\chi\simeq0$, kick moyen Darcy--Brinkman, bain de réémission orienté par la normale $\nabla\chi/\|\nabla\chi\|$, et contribution de particules virtuelles de collision dérivées de $\chi$. Cette variante, notée dans les scripts \texttt{mean\_outward\_bath+chiVP}, permet d'utiliser $\chi$ comme représentation géométrique souple de formes solides quelconques, tout en conservant la distinction conceptuelle entre pénalisation poreuse et condition pariétale exacte. La correction 0426 des \emph{fastflags} CUDA résidents est capitale : elle fait repasser le cas Darcy/chiVP segmenté de l'ordre de $0.19$--$0.23\,\si{\second}$ par pas à environ $0.056\,\si{\second}$ par pas sur le backward step $960\times480$, soit un coût voisin du chemin solide spécialisé. Le goulet mesuré ne venait donc pas de Darcy/chiVP mais d'un aiguillage CUDA incomplet dans les scripts de démonstration.


Depuis les développements C++/OpenMP réalisés ensuite sur \texttt{SRC\_incompressible\_openMP}, une nouvelle étape peut être isolée. Les conditions limites inlet/outlet Q6/Q9 ont été combinées avec des solides immergés fixes à l'aide d'un masque cellule/face : les cellules fluides adjacentes au solide restent actives, les faces fluide--solide sont fermées, le halo Q9 autour du solide est désactivé comme stratégie nominale, et les flux normaux Q6/Q9 au travers du solide sont imposés nuls après fermeture géométrique. Les validations périodiques et inlet/outlet avec cylindre fixe montrent que cette formulation solide immergé est suffisamment robuste pour être considérée comme validée. En revanche, les tests ultérieurs sur poches de cellules pauvres ont mis en évidence une limite indépendante des solides : Q9, ses sécurités low-mass et les tentatives de garde virielle locale peuvent transformer une cellule pauvre en défaut eulérien actif, alors que le transport SRC classique tend plutôt à la convecter et la diffuser.

La réponse méthodologique récente est une troisième génération de prototype MATLAB, séparée de la branche solides : une formulation SRC/MPCD pondérée avec resampling local conservatif. Les particules portent désormais une masse variable $m_i$ ; les champs de cellule sont reconstruits par dépôt massique ; un pool préalloué, des rôles de particules \emph{fluide}/\emph{latente}/\emph{inactive}, un masque de cellules mouillées et une boucle extraction--insertion--remap maintiennent une population locale suffisante sans transformer les cellules pauvres en puits de pression. La masse de cellule, le moment et, lorsque nécessaire, l'énergie thermique locale sont restaurés par des corrections conservatrices. Les premiers tests démontrent la fermeture de la boucle sur poches vide/surpeuplée, sur hétérogénéités distribuées, sur remplissage d'un domaine initialement latent et, surtout, sur Taylor--Green forcé et Poiseuille avec particules virtuelles de paroi. Le cas Poiseuille wallVP indique que Q6+resampling conserve une viscosité effective proche de la référence SRC wallVP pure tout en contrôlant la masse cellulaire au niveau du roundoff.


Depuis le portage C++/OpenMP réalisé dans \texttt{SRC\_incompressible\_openMP}, cette stratégie n'est plus seulement un prototype MATLAB. L'infrastructure particulaire pondérée, les rôles \texttt{Fluid}/\texttt{Inactive}/\texttt{Latent}, les fichiers d'état \texttt{.smpcd} enrichis, le dépôt pondéré, la collision SRC pondérée, Q6 sur champs pondérés, le pool de recyclage, le remap local, la renormalisation thermique, le \texttt{mass guard} et, surtout, le garde de population $N_{\min}/N_{\mathrm{target}}/N_{\max}$ ont été traduits dans le code OpenMP. Les validations récentes montrent que la chaîne OpenMP conserve la dynamique Taylor--Green et Poiseuille, impose la masse cellulaire au roundoff dans les cas resamplés, accélère fortement le remplissage d'un domaine initialement inactif et contrôle le support particulaire dans les cellules mouillées. La différence essentielle par rapport au premier portage OpenMP est que le resampling est maintenant redevenu \emph{population-driven}, comme dans MATLAB : on corrige d'abord le nombre de particules actives par cellule, puis seulement ensuite la masse et l'énergie locales.


Le cycle 0493w5--0493w8 complète maintenant la projection Q6 multi-espèces sur la branche \texttt{surf}, jusqu'au commit \texttt{caa5b6b}. Le mode historique \texttt{weighted} reste une projection barycentrique du mélange et est conservé pour reproductibilité. Un nouveau mode \texttt{independent\_masked} construit au contraire un support d'occupation propre à chaque espèce, résout un problème Q6 masqué distinct pour chaque espèce déclarée incompressible et applique la correction uniquement aux particules du type correspondant. Le contrat strict est \texttt{q6StrengthDeclared=0} $\Rightarrow$ aucune correction Q6 directe. Le chemin reste CUDA résident, fonctionne sur les familles périodique, canal à parois, entrée/sortie pleine face, entrée/sortie segmentée et Darcy--Brinkman déjà acceptées par Q6 résident, et conserve les exclusions générales du backend pour les domaines mobiles, les masques de projection de solides immergés non pris en charge et la réponse de capacité fermée.

Les validations associées séparent désormais la convergence du flux de face projeté et la divergence du champ cellulaire réellement obtenu après application aux particules. La matrice multi-conditions-limites passe $5/5$. Une qualification Taylor--Green à trois graines établit, à la précision numérique, la non-régression monoespèce du nouveau mode par rapport au Q6 historique et la neutralité du simple découpage d'un même fluide entre deux identifiants de type sous SRC. Lorsque les deux types identiques sont projetés indépendamment, la viscosité effective moyenne reste à $0.15\,\%$ de la moyenne monoespèce, les écarts graine par graine restent inférieurs à $1.7\,\%$, et la divergence est fortement réduite. Cette qualification ne vaut pas encore fermeture physique complète d'une interface liquide--gaz : la condition interne de correction de pression reste une condition de Dirichlet simple et aucune tension superficielle, pression gazeuse interfaciale ou loi de mouillage n'est encore imposée.


Le cycle 0493x constitue un changement de statut supplémentaire. Le cas de référence 
\emph{dam-break}, longtemps hors de portée d'une qualification fine à cause des coûts de calcul et des défauts successifs de support ou de fermeture, a servi de test discriminant pour la chaîne CUDA résidente. L'observation d'un liquide entièrement rempli sous gravité a montré que la séquence historique force--advection--collision--Q6 corrigeait la divergence trop tard : une particule était déjà déplacée de $g\Delta t^2$ avant que la projection ne puisse agir. La famille Q6-g issue de 0493x3--0493x4b projette donc la vitesse tentative incluant la force \emph{avant} le streaming ; la variante \texttt{prestream\_single\_fused} conserve un seul solveur Q6 par pas et fusionne le dépôt de la vitesse tentative et l'application force+Q6 dans le chemin CUDA résident. Les étapes 0493x5--0493x6 séparent ensuite le support numérique Q6 de l'interface physique $\alpha=0.5$, préparent un stencil interfacial résident et imposent enfin la condition physiquement motivée $p_l|_\Gamma=p_g$.

La génération courante est désignée dans ce rapport par \textbf{Q6-g-f}. Elle ajoute à Q6-g une reconstruction face--particule affine compatible avec le flux projeté et, surtout, une restauration de densité intégrée directement au second membre de la projection. Dans le bulk liquide, la contrainte devient $\nabla\!\cdot\bm u_{\rm proj}=(f-1)/\tau_\rho$ avec $f=M_l/M_{l,\rm ref}$ et une constante de temps physique qualifiée $\tau_\rho=0.25$ ; l'interface conserve le stencil $\alpha=0.5$ et peut simultanément recevoir la condition $p_l|_\Gamma=p_g$. Les runs longs $300\times150$, $\Delta t=0.005$, $5000$ pas et $600\times300$, $\Delta t=0.0025$, $10000$ pas atteignent le même temps $t=25$ sans dérive volumique séculaire : l'excès géométrique final reste respectivement proche de $5.02\,\%$ et $5.86\,\%$ avec la pression gazeuse active. Une calibration SRC séparée montre que ces deux mailles ne conservent pas le même Reynolds effectif ($Re\simeq726$ contre $Re\simeq1670$) ; la comparaison est donc une validation de robustesse de l'opérateur et de la fermeture volumique, non une preuve de convergence hydrodynamique complète. La même infrastructure est conçue pour recevoir ensuite le saut capillaire $p_l-p_g=\sigma\kappa$ sans modifier la matrice CG.


Les développements 0493x7f--0493x7q complètent enfin cette chaîne sur trois plans qui se sont révélés indissociables en pratique. Le domaine admissible de Q6-g-f est étendu aux topologies statiques déjà supportées par le backend CUDA résident, puis au Darcy--Brinkman en tant que source déterministe de vitesse tentative. La restauration de densité devient signée et cohérente : une branche de compression et une branche de traction/déplétion sont admises seulement lorsque le dépassement du seuil est confirmé par au moins un voisin direct, ce qui évite d'asservir le solveur au bruit d'occupation MPCD. Le CG Q6-g-f est ensuite rendu entièrement résident dans un kernel coopératif 0493x7j, tandis que les diagnostics et réductions hôte sont retirés du chemin chaud hors cadence de sortie.

La correction 0493x7q traite un défaut plus fondamental découvert par Poiseuille : dans un domaine monophase périodique, la reconstruction affine B1/RT0 était neutre en moment au niveau des centres de cellules, mais pas exactement au niveau des particules, car le barycentre d'un échantillon MPCD fini n'est pas confondu avec le centre géométrique de chaque cellule. Le résidu ainsi créé contaminait le mode spatial uniforme $k=0$ et pouvait annuler une force volumique moyenne malgré un solveur elliptique convergé. x7q mesure désormais le moment réellement appliqué par B1 aux particules et retire sur GPU le résidu uniforme exact, uniquement dans les directions périodiques et uniquement pour le chemin B1 monophase \texttt{fullDomain}. Le dam-break partiel conserve donc son kernel historique. Le probe périodique force seul retrouve alors $\bar u_x=a_xt$ et le moment théorique au niveau du roundoff. La campagne finale \texttt{src}/\texttt{src-q6}/\texttt{src-q6-g-f} sur Taylor--Green, Poiseuille, bend-pipe et same-face IO, complétée par un dam-break long, valide la robustesse multi-CL du chemin Q6-g-f signé avec tolérance de projection $10^{-5}$.

Le document corrige donc le rapport précédent en reclassant les résultats historiques, en documentant le thermostat corrigé, les conditions limites virtuelles, le solveur \texttt{general\_bc} avec fixation du mode constant, les diagnostics de conservation du moment, les corrections de viscosité liées à la discrétisation transverse, les premières stratégies de parallélisation/vectorisation MATLAB et la nouvelle stratégie de resampling pondéré. Le portage C++/OpenMP reste une étape naturelle, mais le socle MATLAB fournit maintenant un protocole de validation plus fiable : Taylor--Green pour la préservation des structures et la viscosité bulk de référence, Poiseuille pour la viscosité effective corrigée et les parois, piston staircase pour l'EOS bulk, gaz à l'équilibre entre murs pour la pression mécanique wallVP, et tests de population contrôlée pour la robustesse locale du support particulaire.


Le dernier cycle de développement OpenMP ajoute quatre éléments qui modifient fortement l'architecture pratique du code. Premièrement, les conditions limites ouvertes ne sont plus seulement définies par face entière : des segments \texttt{inlet}/\texttt{outlet} multiples, en coordonnées relatives, peuvent être placés sur une même face solide, les portions non couvertes restant solides. Deuxièmement, l'ancien aiguillage \texttt{method=classic/q6} a été supprimé au profit de deux interrupteurs explicites, \texttt{projectionEnable} et \texttt{resamplingEnable}, qui définissent les quatre cas d'ablation \texttt{classic}, \texttt{classic\_ resampling}, \texttt{q6} et \texttt{q6\_ resampling}. Troisièmement, un mécanisme de réponse de capacité fermée a été introduit : lorsque la masse contenue dans un domaine fluide fermé dépasse la masse nominale, l'intensité effective de Q6 est diminuée par modulation de \texttt{q6ProjectionStrength}, le module viriel effectif augmente, et le remap de masse suit une cible cellulaire effective qui conserve la surmasse globale au lieu de l'effacer. Quatrièmement, la pression bulk reconstruite, $P_{\mathrm{kin}}+P_{\mathrm{vir}}$, est maintenant projetée en charge normale sur les portions solides des parois ; le test statique de boîte uniformément surcompressée vérifie que les quatre faces reçoivent la même pression et que la force pariétale nette est nulle à la précision numérique.

Enfin, la branche \texttt{SRC\_GPU} ajoute un jalon d'architecture parallèle : le SRC classic, défini ici comme advection/streaming, décalage de grille, collision/rotation SRC et thermostat, dispose désormais d'un chemin CUDA résident validé sur les conditions limites périodiques, wall-simple, solide rectangle, piston/mobile wall, inlet/outlet full-face et inlet/outlet segmenté. Les fermetures restent des modules activables séparément. Q6 est maintenant CUDA résident et, avec \texttt{independent\_masked}, qualifié en mono-espèce et en multi-espèces sur les topologies périodique, canal à parois, entrée/sortie pleine face, entrée/sortie segmentée et Darcy--Brinkman ; le resampling et le viriel conservent leurs périmètres de qualification propres.



Le cycle CUDA 0294--0303 complète ce statut : un module de resampling CUDA post-SRC, non destructif et conservatif, a été ajouté au chemin résident. Il comprend un survey passif du support particulaire, un reconditionnement conservatif des masses, un garde de population local par split/merge de particules représentatives, une restauration de l'énergie cinétique relative et un filtrage boundary-aware. Les validations longues sur marche arrière montrent que le réglage nominal provisoire $N_{\min}=12$, $N_{\mathrm{target}}=20$, $N_{\mathrm{max}}=32$ supprime les cellules humides vides derrière la marche pour $U_{\mathrm{in}}=0.30$, $0.45$ et $0.60$, tout en conservant masse, impulsion et énergie relative au bruit numérique.



Le chantier 0490a--0490p ajoute enfin un traitement multi-espèces au resampling pondéré. Le champ \texttt{type} devient une identité physique transportée, distincte du \texttt{role} algorithmique \texttt{Fluid}/\texttt{Inactive}/\texttt{Latent}. Les dépôts par cellule et par espèce, la fermeture de masse, le garde de population, le remplissage des cellules temporairement vides et les transferts donneur--receveur sont désormais conservatifs par espèce. Le chemin validé 0490p garde ces opérations sur l'état CUDA partagé, y compris la politique wet/poor/rich et le pool de slots, et termine l'audit avec \texttt{remaining\_cpu\_scope=none} dans le sous-ensemble périodique actuellement qualifié.

Le cycle 0491 complète cette infrastructure par un Q6 multi-espèces CUDA résident. Une seule projection elliptique reste appliquée au mouvement barycentrique ; sa correction finale est ensuite distribuée selon les fractions massiques et le coefficient \texttt{q6StrengthDeclared} de chaque espèce, avec une identité de recomposition qui conserve exactement la correction barycentrique. Le mode commun reproduit le Q6 historique, tandis que le mode pondéré rend les vitesses relatives sensibles au \texttt{type} sans modifier masse, population ou identité particulaire. Les validations bulk sont encourageantes ; le resampling demeure cependant optionnel en raison de son coût et des questions encore ouvertes sur l'autodiffusion.



Le cycle 0493x8 qualifie les conditions limites ouvertes segmentées nécessaires aux écoulements externes confinés. L'inlet impose un profil de Poiseuille local au segment de façon cohérente aux niveaux particulaire et Q6-g-f ; l'outlet \texttt{neumann} est une sortie passive cinétique--pression avec déflation des modes longitudinaux lents et centrage de la cible de relaxation de densité. Après un premier calcul en domaine axial réduit, la comparaison bibliographique a été reprise sur le domaine exact de Zovatto--Pedrizzetti, $15D$ en amont et $40D$ en aval, sur une grille $4400\times400$ exécutée sur GPU A100 de compute capability \texttt{sm\_80} au CCUB et prolongée jusqu'à $75000$ pas par restart. L'inspection des champs montre que l'allée n'est pleinement développée sur toute la zone aval qu'au cours du troisième segment ; l'analyse bibliographique définitive est donc restreinte à la fenêtre globale $59001\ldots74901$, soit $160$ frames et $3.085$ cycles POD. Elle donne $f_{\rm POD}=0.0969335$ et $f_{\rm probe}=0.0974010$, soit $0.481\,\%$ d'écart, avec $R^2_{\rm phase}=0.999923$, $R^2_{\rm probe}=0.9416$ et une capture de la paire POD de $73.1\,\%$. La vitesse bulk mesurée à $-3D$ du cylindre vaut $0.120095201$, soit $-0.6158\,\%$ de la cible, et les dispersions amont de $Q_v$, $\bar\rho$ et $J_\rho$ restent respectivement de $0.2238\,\%$, $0.0223\,\%$ et $0.2162\,\%$. La réserve historique sur le Reynolds est fermée par une calibration Taylor--Green directement exécutée sous le chemin Q6-g-f de production : huit graines donnent $\nu_{\rm Q6GF}=6.651267455\times10^{-4}$, $\sigma_\nu=3.258\times10^{-5}$ et $CV=4.90\,\%$. Avec la vitesse bulk mesurée, le run tardif correspond à $Re_H=282.125$, très proche du point $Re_H=280$ de Zovatto--Pedrizzetti. La période mesurée vaut désormais $T^*=0.791017$ contre $0.83$, soit $-4.70\,\%$ ; l'écart résiduel est donc de l'ordre de $4$--$5\,\%$, et non plus $6\,\%$. Il ne peut plus être expliqué par une erreur grossière de viscosité, par un mauvais débit amont ou par un sillage encore en phase de croissance ; la principale réserve quantitative restante est la durée statistique de la fenêtre strictement pleinement développée, limitée à un peu plus de trois cycles.


Le cycle 0493x9 ajoute enfin une fermeture capillaire explicite au chemin Q6-g-f. La géométrie physique de surface libre reste définie par le crossing $\alpha=0.5$, mais la courbure est estimée sur un champ auxiliaire destiné uniquement à la géométrie : trois passes binomiales $3\times3$, gradient de Scharr, normale $\bm n_{AB}=-\nabla\alpha_K/\lVert\nabla\alpha_K\rVert$, puis $\kappa=\nabla\cdot\bm n_{AB}$. Le saut de Laplace $p_A-p_B=\sigma\kappa$ n'est pas appliqué comme une force CSF ni comme un kick particulaire direct ; il est introduit dans la condition de Dirichlet de la projection Q6-g-f, $\phi_\Gamma=\Delta t\,(p_B+\sigma\kappa-p_{\rm ref})/\rho_{A,\rm ref}$, puis transmis aux particules par la reconstruction B1 déjà qualifiée. Les tests de goutte et d'ellipse montrent une réponse de Laplace quasi linéaire, une circularisation monotone et des oscillations capillaires amorties ; les runs de dripping produisent goutte pendante, pincement, chute, impact, flaque puis impact d'une goutte suivante. Une régularisation de résolution limite uniquement la courbure utilisée dans $\sigma\kappa$ lorsque le rayon estimé tombe sous quelques mailles. Le statut courant distingue le défaut du code, qui reste $R_{\min}=0$, les calibrateurs/splash x12 et Taylor--Culick qui utilisent $R_{\min}=4h$, et x12cal qui conserve $3h$ dans son régime de faible courbure. Le mouillage reste moins mûr : la construction x9m reproduit correctement des géométries statiques sur paroi dans la plage qualifiée, mais l'angle de contact dynamique et la courbure au voisinage d'une ligne triple ne sont pas encore quantitativement fermés.

Les cycles 0493x10--0493x12 ont depuis profondément révisé le traitement cinétique de la surface libre. Les premières formulations par barrière, réflexion spéculaire ou réaction collective ont été conservées comme ablations mais ne constituent plus la chaîne de travail. Le chemin courant sépare désormais la contrainte capillaire macroscopique --- toujours portée par le saut de Laplace Q6-g-f --- de la conservation du support particulaire au voisinage de l'interface. Il combine la paroi thermique x10o, un champ $\alpha$ cinétique obtenu par dépôt CIC distinct du champ x6c de Q6, une reconstruction biquadratique tensorielle Q2, la relocalisation one-for-one x10u, l'échange local conservatif de vitesses x10v, la résolution des recouvrements initiaux x10p/x10q et le refroidissement thermique local x12a. Le dépôt CIC est fusionné avec un passage particulaire déjà nécessaire ; le vrai Q2 n'ajoute ni grille de coefficients ni passe globale ; x10v réutilise les buffers cellulaires x10m et n'ajoute qu'un marqueur d'un octet par particule ; x12a ajoute un champ d'un float par cellule et fusionne son mapping thermostat avec le dépôt de moments existant.

Cette chaîne n'est pas présentée comme définitivement fermée. Les ablations d'onde capillaire montrent que la relocalisation x10u est nécessaire pour supprimer les pertes de support mais peut introduire du chattering, tandis que le swap vectoriel complet x10v ajoute un amortissement/ralentissement systématique des modes interfacials courts. La branche x10w de redistribution thermique pairwise, exactement conservative en quantité de mouvement et énergie cinétique, a été explorée puis laissée désactivée ; les variantes cinématiques x10r, x10s et x10t restent disponibles pour ablation mais sont également désactivées dans les runners x12. x12a réduit localement la température effective et l'épaisseur cinétique dans les ligaments ou petites gouttes selon
\[
  \frac{k_BT_{\rm eff}}{k_BT}
  =\min\!\left[1,\left(\frac{L_{\rm loc}}{R_c}\right)^2\right],
  \qquad R_c/h=8\sqrt{10},
\]
où $L_{\rm loc}$ est la demi-distance vers l'interface opposée mesurée le long de la normale Q2 intérieure.

La qualification a également été séparée en deux objets. Le calibrateur dynamique x12cal mesure la dispersion des ondes résolues par
\[
  \omega^2=\frac{\sigma_{\rm dyn}}{\rho}k^3\tanh(kH)
\]
et met en évidence une dépendance au mode : les modes $n=2$ et $n=3$ donnent des gains bruts voisins de $0.78$, tandis que $n=4$ est nettement plus dispersif et ne permet pas une calibration scalaire robuste à cette résolution. Le calibrateur mécanique x12yl utilise au contraire des gouttes statiques appariées à rayon et graine identiques et définit
\[
  \Delta p_{\rm cap}=p(\sigma)-p(0)
  =\sigma_{\rm eff}\,\langle\kappa\rangle_{\rm active}.
\]
La campagne finale à $\sigma_{\rm declared}=945$ donne $\sigma_{\rm eff}=940.965$, soit $G_\sigma=0.99573$, avec $R^2=0.999978$. La tension superficielle mécanique est donc proche de la valeur déclarée ; le ralentissement des ondes est traité comme une erreur dynamique/dispersive distincte et non comme une renormalisation générale de $\sigma$.

Une requalification ultérieure sur le fluide économique x13h et sur la chaîne de production complète modifie toutefois le statut de la dynamique résolue. Les gouttes oscillantes bidimensionnelles $n=2,3,4$ donnent des gains de fréquence tardifs proches de l'unité, respectivement $G_\omega\simeq0.992$, $0.999$ et $0.987$, avec amortissement positif et fits compatibles avec les références retenues. En revanche, le benchmark de rétraction de nappe Taylor--Culick révèle un défaut spécifique au régime de rim mobile et collecteur de masse : pour une nappe de hauteur totale $H=0.25$, $\rho=524288$ et $R_{\min}=4h$, le cas long $L/H=12$ donne $G_{\rm TC}=0.8079$ à $\sigma=10000$ et $0.7004$ à $\sigma=2500$ sur la fenêtre principale. Le prolongement temporel confirme un régime de front presque linéaire sans retour tardif vers $G_{\rm TC}=1$. Une reconstruction offline validée de la géométrie x6c/x6e, de la courbure p3 et du cutoff x9r donne cependant une résultante capillaire locale du rim compatible avec $2\sigma$ à quelques pourcents. Le défaut Taylor--Culick ne peut donc pas être résumé par une simple renormalisation de $\sigma$ ni par un facteur temporel global ; il reste un défaut du couplage dynamique complet interface--Q6--particules à localiser par bilan de quantité de mouvement.

Enfin, les campagnes splash et le benchmark JFM 524 étendent la qualification à des changements de topologie violents. Les impacts appariés wall/puddle avec $R/h=40$, $k_BT=0.125$, $\sigma=945$ et $v_y=-0.35$ montrent que le passage de \texttt{surfaceTensionMinRadiusCells}=3 à 4 ne supprime ni lamelle, ni couronne, ni fragmentation. Le runner x12d transpose ensuite le cas de Josserand, Lemoyne, Troeger et Zaleski \cite{JosserandLemoyneTroegerZaleski2005} avec $D/h=320$, $We\simeq514$, un obstacle de hauteur $H/D\simeq0.01875$ placé à $d/D\simeq0.540625$, et explicite le principal écart encore ouvert : le fluide SRC actuellement qualifié donne $Re\simeq649$, loin du $Re\simeq12200$ expérimental. Le benchmark JFM est donc à ce stade un cas de mesure et de développement, pas encore une reproduction quantitative complète de l'expérience.




Le dernier cycle expérimental de surface libre a ensuite cherché à corriger le déficit Taylor--Culick par des fermetures thermiques et de rétention plus agressives. Les branches x13r--x13x (refroidissement de cellule d'interface, refroidissement anisotrope, loi progressive x13t, rétention one-for-one, absorption/reseed x13w et rétention probabiliste) ont permis d'augmenter fortement $G_{\rm TC}$, jusqu'à des valeurs proches de l'unité sur une grille affinée. La validation croisée par gouttes oscillantes a cependant montré que la fermeture x13t+x13w contracte artificiellement le support de phase et détruit la dynamique capillaire, surtout sur la grille fine. Ces développements sont donc conservés comme résultats négatifs documentés, mais exclus du chemin de production. Le point de référence surface libre est figé au commit \texttt{7655b81b1b2fd16eecefa8d8b3bebac4cd9f87f1}, tag \texttt{surf-tension-qualified-x13h-20260831}, avec la chaîne x10o+CIC+Q2+x10p/q+x10u+x10v full-vector swap+x12a. Une revalidation après retour exact au dépôt donne $G_{\omega,n=2}=1.00151$ et $G_{\rm TC}=0.79457$ à $\sigma=10000$.


Le cycle 0493x14 ajoute ensuite le couplage avec une phase gazeuse particulaire explicite sans modifier la chaîne liquide de référence. Les collisions SRC restent communes au mélange, tandis qu'un thermostat par espèce permet des masses et températures distinctes. La géométrie cinétique x10u est étendue bilatéralement au couple liquide--gaz ; le gaz reçoit une réflexion spéculaire normale x14l qui conserve sa vitesse tangentielle. La pression gazeuse x6g est corrigée par une fraction de volume accessible dérivée de l'alpha Q6 (x14s) et continue de se composer avec le saut capillaire $p_L=p_G+\sigma\kappa$. x14v transfère au liquide l'impulsion cinétique normale excédentaire des réflexions gazeuses après soustraction de la traction thermodynamique déjà représentée par x6g. La qualification des interfaces courbes a ensuite séparé la distribution locale de traction et sa résultante globale : x14ad conserve les normales x10n et échantillonne localement la pression de jauge sur les faces x6g représentées, tandis que x14ai-fix1 impose au seul mode résultant $n=1$ la quantité de mouvement Q6 réellement appliquée après B1 et correction périodique. Sur la goutte oscillante $n=2$, le moment total reste au niveau de l'arrondi numérique et $G_\omega=0.97988$ avec $R^2=0.9971$ ; le benchmark de traînée x14ah conserve une migration aval nette, $\Delta x/R=0.799$, tout en réduisant d'environ un facteur 17 la dérive transverse observée avant la fermeture. Le coût x14ai reste sous la résolution du chronométrage apparié. Cette fermeture est retenue comme candidate pour une composante liquide fermée et isolée des frontières Q6 externes ; son extension à un liquide touchant une paroi ou une ouverture n'est pas qualifiée. Le Couette biphasique indique parallèlement que le SRC commun transmet déjà la contrainte tangentielle sans terme de frottement artificiel, avec $a_L/a_G=0.23555$ contre $\mu_G/\mu_L\simeq0.23769$.

\end{abstract}


\tableofcontents


\section*{Note de version et statut des résultats}
\addcontentsline{toc}{section}{Note de version et statut des résultats}

Ce document complète une version antérieure du rapport. Les sections consacrées à la redistribution dure, au zonage et aux premières validations Q6/Q9 sont conservées parce qu'elles expliquent la construction progressive de la méthode. Elles ne doivent toutefois plus être lues comme des validations quantitatives définitives lorsque les résultats provenaient des scripts historiques : ces scripts ne contenaient pas encore le thermostat cellulaire corrigé, n'utilisaient pas les conditions limites par particules virtuelles dans les canaux, et n'étaient pas encore structurés pour les premiers gains de vectorisation/parallélisation MATLAB.



Depuis les patchs post-0143, le code OpenMP utilise en outre une configuration plus explicite : \texttt{projectionEnable} contrôle seul Q6, \texttt{resamplingEnable} contrôle seul le garde de population et le remap pondéré, et les ouvertures partielles sont décrites par des segments relatifs. Les développements de capacité fermée, de viriel faiblement compressible et de charge pariétale ne remplacent pas les validations précédentes ; ils constituent un enrichissement activé par paramètres, conçu pour les domaines fermés sur-remplis et pour le futur couplage à des parois solides déformables.


Depuis le cycle 0294--0303, la branche \texttt{SRC\_GPU} inclut en outre un resampling CUDA post-SRC destiné au contrôle du support particulaire. Ce module est volontairement séparé de Q6 CUDA : Q6 CPU/OpenMP reste réactivable dans l'architecture, mais la validation initiale du resampling CUDA concerne le chemin SRC classic résident suivi d'une opération de re-échantillonnage conservatif. Les listings de paramètres 0292 joints au présent rapport restent la référence pour le socle CUDA classic, inlet/outlet, thermostat et fermeture liquide OpenMP ; les options 0295--0303 sont documentées ci-dessous comme complément spécifique au resampling CUDA.


Le cycle 0490a--0490p étend ce resampling aux mélanges de plusieurs espèces enregistrées. La conservation n'est plus contrôlée seulement sur la masse totale : les opérations discrètes, le remplissage, la fermeture de masse et les transferts sont audités séparément pour chaque valeur physique de \texttt{type}. Le checkpoint de référence de cette section est le commit \texttt{36abd23} de la branche \texttt{surf}. Les inventaires CSV joints ont été reconstruits à partir des appels \texttt{getenv} et du parseur \texttt{params.kv} de cet état du code ; ils remplacent les inventaires consolidés 0436b pour les options ajoutées ou omises depuis cette date.

Le cycle 0491 rend ensuite l'application Q6 sensible aux espèces sans créer de solveur elliptique par phase. La section 33 distingue désormais explicitement la formulation 0490 du resampling et la formulation 0491 de Q6, documente leurs paramètres, leur ordre d'exécution, leurs invariants et leur niveau de qualification.

Le présent addendum 0493x documente la branche force-aware et interface liquide--gaz développée après 0493w8. Il introduit le paramètre \texttt{q6ForceProjectionMode}, étend \texttt{speciesQ6Mode} avec \texttt{free\_surface\_masked}, puis ajoute les étapes x6a--x6g de géométrie/interfacialité et la reconstruction face--particule B1 de x6h. Les développements 0493x7c--x7q complètent désormais cette architecture par la restauration de densité dans le RHS, sa paramétrisation par un temps physique, le gate cohérent signé compression/traction, l'élargissement multi-CL et Darcy, le CG Q6-g-f coopératif entièrement CUDA résident, la symétrisation des reconstructions de faces et la fermeture exacte du moment périodique au niveau particulaire. Le chemin complet actuellement qualifié est nommé \textbf{Q6-g-f} dans la suite du rapport.

Les inventaires CSV ont été réaudités sur le snapshot local du 26 août 2026 et prolongés jusqu'à l'état x12 courant. Ils documentent maintenant, en plus de la physique capillaire x9, l'ensemble des contrôles runtime x10/x12 réellement lus par le code : x10o, CIC cinétique, reconstruction Q2, x10u, x10v, x10p, les ablations x10r/x10s/x10t, la branche x10w et le refroidissement x12a. Cette révision corrige plusieurs ambiguïtés historiques. Le défaut du champ \texttt{phaseInterfaceKineticReflectionFraction} dans \texttt{SimulationParams} est $0$, tandis que les runners x12 de production imposent explicitement $r=1$. \texttt{surfaceTensionMinRadiusCells} n'a pas de valeur universelle : le défaut du struct reste $0$, les benchmarks JFM et x12yl utilisent actuellement 4, tandis que x12cal conserve 3 dans un régime de faible courbure où ce cutoff n'est pas censé être actif. Les alias parseur \texttt{kineticReflectionFraction} et \texttt{evaporationTargetType} sont également acceptés. Les variables d'environnement x10/x12 servent à sélectionner des branches algorithmiques et des ablations ; les paramètres physiques pérennes restent portés par \texttt{params.kv} autant que possible.


\paragraph{Extension x14 au 4 septembre 2026.}
Le rapport est prolongé au-delà du point de référence du 31 août par le snapshot de travail \nolinkurl{snap_040926.zip}. Cette extension ne modifie pas rétroactivement le statut du tag \nolinkurl{surf-tension-qualified-x13h-20260831}; elle documente la couche de couplage liquide--gaz construite par-dessus ce socle et les deux campagnes de qualification d'interface poursuivies jusqu'à x14ai-fix1. Après le thermostat par espèce x14a--x14g, la géométrie bilatérale x14k, la réflexion spéculaire gazeuse x14l, la pression accessible-volume x14s, le transfert cinétique excédentaire x14v et le Couette x14w, les jalons x14y--x14ad explorent différentes géométries de soustraction de la traction thermodynamique. x14ae et x14af sont des diagnostics causaux temporaires ; x14ai introduit la fermeture de résultante Q6 réellement appliquée et x14ai-fix1 corrige sa cible pour inclure exactement la correction périodique B1. Les inventaires de paramètres et de variables d'environnement sont réaudités sur ce snapshot ; aucune hash Git n'est attribuée à l'archive fournie, qui ne contient pas les métadonnées \texttt{.git}.


Le cycle x8 n'introduit pas une seconde famille de paramètres physiques pour l'inlet/outlet : il affine les sémantiques existantes de \texttt{inletVelocitySpatialProfile} et \texttt{openBoundaryOutletMode}. Pour la surface libre, une distinction d'architecture doit être conservée : le champ CIC x10cic est bien dédié à l'interface cinétique et ne remplace pas le champ x6c de Q6/capillarité, mais son orchestration reste, dans le snapshot courant, située dans le chemin Q6 \texttt{apply\_independent\_masked\_species\_q6\_0493w5}. Il n'existe donc pas encore de chemin CIC interface SRC-only indépendant de Q6.


Pour les sections historiques Q6/Q9 et piston, la chaîne de référence reste celle de la branche \texttt{q9-piston-case-virial-kick}. L'addendum 0493x décrit séparément le chemin CUDA force-aware de la branche \texttt{surf}. Sauf indication contraire dans ces sections historiques, on appelle \textbf{Q6/Q9 complet} la séquence suivante : étape SRC/MPCD classique, projection de vitesse Q6, projection basse fréquence du flux de masse Q9, thermostat corrigé, correction globale exacte du moment pour les kicks additionnels, puis diagnostics bulk/paroi. Les résultats historiques restent utiles pour comparer les idées, mais les chiffres de calibration doivent être repris avec les scripts corrigés.

\section{Architecture générale de la méthode}

Le code simule un fluide particulaire en deux dimensions dans un domaine rectangulaire de taille $L_x \times L_y$, discrétisé en un maillage cartésien régulier de $N_x \times N_y$ cellules de taille
\begin{equation}
 a_0 = \Delta x = \frac{L_x}{N_x},
 \qquad
 \Delta y = \frac{L_y}{N_y},
 \qquad
 V_c = \Delta x\,\Delta y,
\end{equation}
où $V_c$ désigne l'aire d'une cellule dans la simulation bidimensionnelle.

Les variables particulaires sont la position $\bm{x}_i \in \mathbb{R}^2$ et la vitesse $\bm{v}_i \in \mathbb{R}^2$ de la particule $i$, pour $i=1,\dots,N_p$, où $N_p$ est le nombre total de particules. Le pas de temps est noté $\Delta t$.

La séquence temporelle de référence des validations Q6/Q9 historiques n'est plus la chaîne avec redistribution dure. Elle s'écrit maintenant, au niveau d'un pas de temps,
\begin{enumerate}[label=\arabic*)]
  \item advection particulaire avec les forces imposées et application des conditions limites géométriques ;
  \item collision SRC/MPCD avec décalage aléatoire de grille, rotation SRD et thermostat cellulaire corrigé ;
  \item reconstruction particules--grille des champs d'occupation, de vitesse moyenne, de température et de pression cinétique ;
  \item projection de vitesse Q6, c'est-à-dire correction du champ moyen pour réduire \(\nabla_h\cdot u\) ;
  \item projection basse fréquence du flux de masse Q9, c'est-à-dire relaxation de \(\nabla_h\cdot(Nu)\) sur les modes cohérents de grande longueur d'onde ;
  \item interpolation des corrections grille vers les particules avec correction globale exacte de la quantité de mouvement lorsque des kicks non strictement conservatifs sont appliqués ;
  \item diagnostics de densité, divergence, flux de masse, température, pression bulk, pression mécanique de paroi et, selon le cas, vorticité ou profils moyens.
\end{enumerate}

Cette liste décrit le chemin Q6/Q9 historique dans lequel la projection intervient après transport/collision. Le cycle 0493x a montré qu'il n'est pas adapté tel quel à un liquide incompressible soumis à une force volumique persistante : la position est déjà advectée avec la vitesse accélérée avant que Q6 ne corrige cette vitesse. Pour les cas force-aware qualifiés sur la branche \texttt{surf}, le chemin \texttt{prestream\_single\_fused} utilise au contraire
\begin{equation}
  \text{dépôt de }(\bm v+\bm a\Delta t)
  \;\longrightarrow\; Q6
  \;\longrightarrow\; (\bm v\leftarrow\bm v+\bm a\Delta t+\delta\bm u_{Q6})
  \;\longrightarrow\; \text{streaming}
  \;\longrightarrow\; \text{collision SRC}
  \;\longrightarrow\; \text{thermostat}.
\end{equation}
Ce chemin n'ajoute qu'un solveur Q6 par pas forcé et est détaillé à la section~\ref{sec:q6g-dambreak}. Les chemins généraux Q6/Q9, ouverts, Darcy ou resampling conservent leurs propres contrats tant qu'ils n'ont pas été requalifiés avec ce séquencement.

La chaîne historique, fondée sur redistribution particulaire puis fermeture liquide par pression virielle, est décrite plus loin parce qu'elle a joué un rôle important dans la construction de la méthode. Elle n'est toutefois plus le chemin principal pour les validations Q6/Q9 actuelles. Dans les cas piston récents, le terme viriel n'est pas utilisé pour déplacer durement les particules vers une cible de densité ; il sert à construire une fermeture bulk explicite,
\begin{equation}
  P_{\mathrm{vir}} = K_{\mathrm{virial}}(\rho-\rho_{\mathrm{EOSRef}}),
  \qquad
  P_{\mathrm{tot}}=P_{\mathrm{kin}}+P_{\mathrm{vir}},
  \qquad
  P_{\mathrm{drive}}=P_{\mathrm{tot}}.
\end{equation}

On distinguera donc dans la suite :
\begin{itemize}
  \item l'\textbf{état MPCD classique}, obtenu après transport, conditions limites, collision et thermostat ;
  \item l'\textbf{état Q6}, obtenu après projection de vitesse ;
  \item l'\textbf{état Q9}, obtenu après relaxation basse fréquence du flux de masse ;
  \item l'\textbf{état piston/viriel}, dans lequel une pression bulk additionnelle est reconstruite pour tester une réponse EOS quasi incompressible ;
  \item les \textbf{diagnostics de paroi}, qui mesurent un flux d'impulsion et ne doivent pas être confondus avec la pression thermodynamique reconstruite dans le bulk.
\end{itemize}

\section{Méthode MPCD : principes, advection, rotation, décalage de grille}

\subsection{Principe général}

La méthode MPCD représente un fluide par un ensemble de particules ponctuelles qui alternent deux sous-étapes : une sous-étape de transport balistique (ou advection libre) et une sous-étape de collision collective, dans laquelle les vitesses sont mélangées cellule par cellule.

La variable conservative fondamentale au niveau microscopique est la quantité de mouvement particulaire. La dissipation, la diffusion de quantité de mouvement et la température émergent de l'alternance entre transport et collisions.

\subsection{Advection avec forçage volumique}

À l'instant $t^n = n\Delta t$, la vitesse de la particule $i$ est d'abord modifiée par les forces volumiques imposées. Le code utilise actuellement une force volumique horizontale $f_x$ et une accélération verticale uniforme $g$. L'étape d'advection s'écrit donc
\begin{equation}
 v_{i,x}^{\star} = v_{i,x}^{n} + f_x\,\Delta t,
 \qquad
 v_{i,y}^{\star} = v_{i,y}^{n} + g\,\Delta t,
\end{equation}
puis
\begin{equation}
 \bm{x}_i^{\star} = \bm{x}_i^{n} + \Delta t\,\bm{v}_i^{\star}.
\end{equation}

Ici $v_{i,x}$ et $v_{i,y}$ sont les composantes de la vitesse $\bm{v}_i$, et l'exposant $\star$ désigne l'état pré-collision après transport et forçage.

\subsection{Découpage en cellules et décalage aléatoire de grille}

L'étape de collision est effectuée dans des cellules cartésiennes de taille $a_0$. Pour limiter les artefacts liés à la grille et restaurer l'invariance galiléenne de la méthode, le code applique à chaque pas un décalage aléatoire uniforme
\begin{equation}
 \bm{s} \in \left[-\frac{a_0}{2},\frac{a_0}{2}\right]^2,
\end{equation}
puis assigne les particules aux cellules de la grille décalée. Si la cellule d'indice $c$ contient $N_c$ particules, la vitesse moyenne cellulaire est
\begin{equation}
 \bm{u}_c = \frac{1}{N_c}\sum_{i\in c} \bm{v}_i.
\end{equation}

Le code distingue proprement les axes périodiques et non périodiques au moment de cette affectation, ce qui est essentiel lorsque les bords physiques ou un piston mobile rendent le domaine non périodique dans une direction.

\subsection{Collision SRD par rotation des vitesses relatives}

Dans la collision SRD, on décompose la vitesse particulaire sous la forme
\begin{equation}
 \bm{v}_i = \bm{u}_c + \bm{w}_i,
 \qquad i\in c,
\end{equation}
où $\bm{w}_i$ est la fluctuation de vitesse de la particule $i$ relativement à la vitesse moyenne de sa cellule.

Ces fluctuations sont ensuite tournées d'un angle $\alpha$ ou $-\alpha$ aléatoire, identique pour toutes les particules d'une même cellule. Si $\mathbf{R}(\theta_c)$ désigne la matrice de rotation d'angle $\theta_c\in\{+\alpha,-\alpha\}$ dans la cellule $c$, la collision s'écrit
\begin{equation}
 \bm{w}_i' = \mathbf{R}(\theta_c)\,\bm{w}_i,
 \qquad
 \bm{v}_i' = \bm{u}_c + \bm{w}_i'.
\end{equation}

Cette étape conserve exactement la quantité de mouvement de chaque cellule, puisque
\begin{equation}
 \sum_{i\in c} \bm{w}_i = \bm{0}
 \quad \Longrightarrow \quad
 \sum_{i\in c} \bm{v}_i' = \sum_{i\in c} \bm{v}_i.
\end{equation}

\section{Liens entre unités MPCD et grandeurs physiques}

\subsection{Choix d'unités réduites}

Le code est écrit en unités réduites MPCD. Les grandeurs indépendantes choisies comme unités de base sont généralement :
\begin{itemize}
  \item la longueur de maille $a_0$ ;
  \item le pas de temps $\Delta t$ ;
  \item la masse particulaire $m$.
\end{itemize}

À partir de ces trois échelles, on déduit :
\begin{equation}
 u_0 = \frac{a_0}{\Delta t}
\end{equation}
comme échelle de vitesse,
\begin{equation}
 \varepsilon_0 = m u_0^2 = m\frac{a_0^2}{\Delta t^2}
\end{equation}
comme échelle d'énergie, et
\begin{equation}
 \nu_0 = \frac{a_0^2}{\Delta t}
\end{equation}
comme échelle de viscosité cinématique.

Dans le code actuel, les diagnostics et la fermeture liquide utilisent la variable nommée \verb|rho| pour désigner en pratique une \textbf{densité de nombre surfacique}
\begin{equation}
 n = \frac{N_c}{V_c},
\end{equation}
et non une masse volumique au sens strict. La masse volumique physique en 2D s'obtiendrait par
\begin{equation}
 \rho_m = m n.
\end{equation}

Cette distinction est importante : les quantités reconstruites et les pressions utilisées dans la fermeture s'expriment dans le code à partir de $n$, avec le choix implicite $m=1$ pour les diagnostics de fermeture.

\subsection{Température et énergie thermique}

La quantité reconstruite dans le code est un champ local que l'on peut noter $k_B T_c$, défini à partir de l'énergie cinétique des fluctuations autour de la vitesse moyenne de cellule. Si $d=2$ désigne la dimension spatiale et si $N_c\ge 2$, alors
\begin{equation}
 k_B T_c = \frac{1}{d(N_c-1)} \sum_{i\in c} m\,\abs{\bm{v}_i-\bm{u}_c}^2.
\end{equation}

Dans l'implémentation courante, $m=1$ et la grandeur stockée sous le nom \verb|kBT| ou \verb|T| correspond directement à $k_B T_c$ dans les unités réduites choisies.

\subsection{Pression cinétique}

La pression cinétique locale reconstruite en post-traitement et utilisée dans la fermeture liquide est
\begin{equation}
 P_{\mathrm{kin},c} = n_c\,k_B T_c,
\end{equation}
où $n_c=N_c/V_c$ est la densité de nombre de la cellule $c$.

Cette pression est une pression locale de type gaz parfait dans les unités réduites du modèle. Elle ne doit pas être confondue avec la pression de paroi obtenue par flux d'impulsion aux frontières.

\subsection{Viscosité et nombre de Reynolds}

La viscosité cinématique $\nu$ peut être exprimée sous forme réduite, puis convertie physiquement par
\begin{equation}
 \nu_{\mathrm{phys}} = \nu^*\,\frac{a_0^2}{\Delta t}.
\end{equation}

Le nombre de Reynolds associé à une vitesse caractéristique $U$, une longueur caractéristique $L$ et une viscosité cinématique $\nu$ vaut
\begin{equation}
 \mathrm{Re} = \frac{U L}{\nu}.
\end{equation}

Dans l'état actuel du code, la viscosité n'est pas seulement pensée comme une grandeur théorique MPCD ; elle est aussi \textbf{identifiée numériquement} via des diagnostics de type Poiseuille. Lorsque l'écoulement est piloté par une force volumique horizontale $f_x$, le profil parabolique de référence s'écrit
\begin{equation}
 u_x(y) = \frac{f_x}{2\nu}\,y(H-y),
\end{equation}
où $H=L_y$ est la hauteur du canal. Le simulateur en déduit une estimation de $\nu$ soit par courbure du profil, soit par les contraintes de paroi. Cette stratégie est particulièrement adaptée lorsque la redistribution et la fermeture liquide modifient l'écoulement au-delà du cas MPCD théorique pur.

\section{Thermostat MPCD}

Après la rotation SRD des vitesses relatives, le code applique un thermostat cellulaire optionnel. Son principe est de conserver la vitesse moyenne de cellule tout en renormalisant la norme quadratique des fluctuations pour imposer une température cible $k_B T$.

Pour une cellule $c$, on calcule d'abord l'énergie cinétique relative après rotation,
\begin{equation}
 E_{\mathrm{rel},c}' = \sum_{i\in c} \abs{\bm{w}_i'}^2,
\end{equation}
puis le nombre effectif de degrés de liberté
\begin{equation}
 \mathrm{dof}_c = d(N_c-1).
\end{equation}

Le thermostat introduit un facteur d'échelle
\begin{equation}
 \lambda_c = \sqrt{\frac{\mathrm{dof}_c\,k_B T}{E_{\mathrm{rel},c}'}}
\end{equation}
quand $N_c>1$ et $E_{\mathrm{rel},c}'>0$, puis remplace
\begin{equation}
 \bm{w}_i' \leftarrow \lambda_c\,\bm{w}_i'.
\end{equation}

L'état final après collision et thermostat est donc
\begin{equation}
 \bm{v}_i^{n+1/2} = \bm{u}_c + \lambda_c\,\mathbf{R}(\theta_c)\,\bigl(\bm{v}_i^{\star}-\bm{u}_c\bigr).
\end{equation}

Le thermostat conserve la quantité de mouvement de chaque cellule, mais ne conserve pas son énergie cinétique totale ; il impose une énergie thermique relative cible. Dans le code actuel, il constitue le principal mécanisme explicite de contrôle de la température au niveau MPCD.


\subsection{Correction du bug historique et thermostat vectorisé}

Les premières campagnes Q6/Q9 avaient été réalisées avec une version historique du thermostat. La correction récente a deux conséquences importantes pour l'interprétation des résultats. D'une part, le diagnostic \(k_BT_c\) doit être calculé à partir des fluctuations relatives autour de la vitesse moyenne cellulaire, avec \(d(N_c-1)\) degrés de liberté effectifs lorsque \(N_c>1\). D'autre part, la renormalisation doit être appliquée cellule par cellule en conservant exactement la vitesse moyenne de collision. Une erreur dans cette étape peut modifier la viscosité effective, la pression cinétique reconstruite et la comparaison entre cas classic, Q6 et Q9.

La version corrigée est aussi plus favorable à MATLAB : les sommes par cellule, les énergies relatives et les facteurs \(\lambda_c\) sont calculés par opérations groupées de type \texttt{accumarray}/indexation vectorisée lorsque cela est possible. Cette vectorisation ne change pas la méthode mathématique, mais elle réduit le coût de la boucle de collision et rend les longues validations piston ou Poiseuille plus praticables. Les résultats historiques antérieurs à cette correction doivent donc être conservés comme résultats de développement, non comme valeurs finales de calibration.

\section{Conditions limites}

Le code gère des conditions limites différentes sur chacun des quatre côtés du domaine. Les modes disponibles peuvent être résumés physiquement ainsi.

\subsection{Périodicité}

Quand un bord est périodique, une particule qui sort par ce bord rentre par le bord opposé avec la même vitesse. Si l'axe $x$ est périodique, cela se traduit par
\begin{equation}
 x_i \leftarrow x_i \pm L_x,
\end{equation}
sans modification de $\bm{v}_i$.

\subsection{Réflexion spéculaire}

Sur un bord solide spéculaire, la composante normale de la vitesse change de signe, tandis que la composante tangentielle est inchangée. Par exemple sur la paroi basse $y=0$,
\begin{equation}
 v_{i,y} \leftarrow -v_{i,y},
\qquad
 v_{i,x} \leftarrow v_{i,x}.
\end{equation}

Ce choix réalise une paroi sans pénétration et sans friction tangentielle.

\subsection{Rebond intégral (\emph{bounce-back})}

Dans le mode \emph{bounce-back}, la composante tangentielle est aussi réfléchie relativement à une vitesse de paroi imposée. Si la paroi inférieure se déplace à la vitesse tangentielle $U_{\mathrm{bottom}}$, on écrit
\begin{equation}
 v_{i,x} \leftarrow 2U_{\mathrm{bottom}} - v_{i,x},
\qquad
 v_{i,y} \leftarrow -v_{i,y}.
\end{equation}

Ce mode mime un non-glissement plus fort qu'une réflexion spéculaire pure.

\subsection{Thermalisation aux parois}

Dans le mode \emph{thermalize}, la composante normale sortante est remplacée par une variable aléatoire orientée vers l'intérieur, et la composante tangentielle est tirée autour de la vitesse de paroi. Pour la paroi inférieure, on a schématiquement
\begin{equation}
 v_{i,y} \sim \abs{\mathcal{N}(0,\sigma_w^2)},
\qquad
 v_{i,x} \sim U_{\mathrm{bottom}} + \mathcal{N}(0,\sigma_w^2),
\end{equation}
où $\sigma_w$ est l'écart-type thermique imposé par la paroi.

Ce mécanisme agit comme un thermostat de bord et injecte ou retire de l'énergie selon l'état incident du fluide.

\subsection{Paroi piston}

Le mode \emph{piston} impose sur le bord supérieur une position $y_{\mathrm{p}}(t)$ et une vitesse normale
\begin{equation}
 U_p(t) = \dot y_p(t).
\end{equation}
Lors d'un impact sur le piston,
\begin{equation}
 v_{i,y} \leftarrow 2U_p - v_{i,y}.
\end{equation}

Le piston réduit la \emph{hauteur active} du domaine,
\begin{equation}
 H_{\mathrm{actif}}(t)=y_{\mathrm{p}}(t)-m_p,
\end{equation}
où $m_p$ est une marge géométrique optionnelle (\verb|pistonActiveMargin|) destinée à définir la zone effectivement disponible sous le piston. Le maillage reste fixe : ce n'est pas la taille des cellules qui change, mais le nombre effectif de cellules actives, éventuellement fractionnaires sur la dernière rangée coupée par le piston.

Si la $j$-ième rangée a pour bornes verticales $y_j^-=(j-1)\Delta y$ et $y_j^+=j\Delta y$, la fraction active de cette rangée vaut
\begin{equation}
 \phi_j(t)=\frac{\max\!\bigl(0,\min(H_{\mathrm{actif}}(t),y_j^+)-y_j^-\bigr)}{\Delta y},
 \qquad 0\le \phi_j\le 1.
\end{equation}
Cette fraction est recopiée sur toutes les colonnes de la rangée, ce qui définit une carte de fraction active \(\phi_c\) au niveau cellulaire.

Dans la version piston \emph{incompressible} du code, la cible de redistribution n'est plus adaptée au volume actif courant : la cible d'occupation de référence \(\gamma_0\) est conservée, et seule la géométrie active change. Le piston introduit alors deux quantités distinctes :
\begin{equation}
 \gamma_{\mathrm{target}} = \gamma_0,
 \qquad
 \gamma_{\mathrm{geom}}(t)=\frac{N_p}{\sum_c \phi_c(t)}.
\end{equation}
La première est la cible incompressible effectivement demandée à la redistribution ; la seconde est l'occupation moyenne qui serait géométriquement requise pour loger toutes les particules dans le volume actif sous le piston. Lorsque \(\gamma_{\mathrm{geom}}>\gamma_{\mathrm{target}}\), le piston pousse le système vers un régime de sur-occupation qui ne peut être compensé que par une forte montée de pression, puis par un blocage progressif de la redistribution lorsque la capacité disponible devient insuffisante.

Si l'option correspondante est activée, le code compacte aussi géométriquement le fluide sous le piston par un repositionnement affine vertical. Ce repositionnement n'est pas une modification de densité cible : il constitue uniquement un transport géométrique imposé, auquel la fermeture liquide réagit ensuite par la montée de la pression virielle.

\subsection{Pression de paroi}

L'impulsion échangée avec les parois est accumulée pour construire des diagnostics de pression instantanée de type flux d'impulsion. Si $\Delta P_y^{\mathrm{bot}}$ est la quantité de mouvement verticale échangée avec la paroi inférieure pendant un pas de temps, la pression instantanée associée est estimée par
\begin{equation}
 p_{\mathrm{bot}}^{\mathrm{inst}} = \frac{\Delta P_y^{\mathrm{bot}}}{L_x\,\Delta t}.
\end{equation}

Cette pression de paroi est une grandeur différente de la pression cinétique locale $P_{\mathrm{kin}}$ ou de la pression motrice $P_{\mathrm{drive}}$ utilisée dans la fermeture liquide.


\subsection{Particules virtuelles de paroi : modèle \texttt{wallVP}}

Les validations récentes de Poiseuille et de piston utilisent une condition limite enrichie par particules virtuelles de paroi. Le principe est de compléter les cellules coupées par une paroi solide avec une population virtuelle représentant la partie solide de la cellule. Cette population participe à la collision SRD afin de transmettre une contrainte tangentielle et normale plus réaliste que la seule réflexion géométrique des particules incidentes.

Pour une cellule de collision contenant \(N_f\) particules fluides et \(N_v\) particules virtuelles, on note
\begin{equation}
  \bm P_f = \sum_{i=1}^{N_f}\bm v_i,
  \qquad
  \bm P_v = \sum_{j=1}^{N_v}\bm v_j^{(v)},
  \qquad
  \bm U = \frac{\bm P_f+\bm P_v}{N_f+N_v}.
\end{equation}
La collision virtuelle reconstruit alors l'impulsion post-collision des particules virtuelles sous la forme
\begin{equation}
  \bm P_v' = N_v\bm U + \mathbf R(\theta)\bigl(\bm P_v-N_v\bm U\bigr),
\end{equation}
avec le même angle de rotation SRD que dans la cellule. L'impulsion transmise par la partie virtuelle est
\begin{equation}
  \Delta \bm P_v = \bm P_v' - \bm P_v.
\end{equation}
Ce terme doit être traité comme une impulsion signée. Pour la paroi supérieure, une compression positive correspond à une impulsion verticale positive dans le fluide ; pour la paroi inférieure, elle correspond à une impulsion verticale négative. Le diagnostic récent sépare donc explicitement : impulsion signée, pression positive de compression, contribution impact, contribution \texttt{wallVP}, total impact+VP, et moyennes cumulées sur fenêtre.

La pression mécanique instantanée est normalisée par
\begin{equation}
  p_{\mathrm{wall}}^{\mathrm{inst}} = \frac{\Delta p_n}{\Delta t\,L_x},
\end{equation}
la simulation étant bidimensionnelle avec profondeur unité. Cette grandeur est très bruitée, car elle divise une impulsion de collision par un pas de temps court. Le diagnostic physiquement exploitable est donc la pression moyenne obtenue par différence d'impulsions cumulées sur une fenêtre suffisamment longue,
\begin{equation}
  \overline p_{\mathrm{wall}}[t_a,t_b]
  = \frac{P_{\mathrm{wall}}^{\mathrm{cum}}(t_b)-P_{\mathrm{wall}}^{\mathrm{cum}}(t_a)}{(t_b-t_a)L_x}.
\end{equation}
Le test gaz à l'équilibre entre deux murs montre que cette moyenne cumulée retrouve la pression cinétique bulk, alors que les valeurs instantanées peuvent osciller de plusieurs milliers d'unités réduites autour d'une moyenne de l'ordre de \(2.1\times10^2\).

\section{Méthode de redistribution : cible, étapes et interprétation}

\subsection{Cible d'occupation et de densité}

La redistribution vise à contraindre l'occupation cellulaire vers une cible locale $N_c^{\star}$. Dans le cas le plus simple sans piston, cette cible vaut
\begin{equation}
 N_c^{\star} = \gamma_0,
\end{equation}
où $\gamma_0$ est le nombre moyen visé de particules par cellule pleinement active.

En présence d'une fraction de volume active $\phi_c\in[0,1]$ dans la cellule $c$, la cible locale devient
\begin{equation}
 N_c^{\star} = \gamma_{\mathrm{target}}\,\phi_c.
\end{equation}
Dans la version \emph{incompressible} actuelle du traitement du piston, on impose
\begin{equation}
 \gamma_{\mathrm{target}} = \gamma_0,
\end{equation}
indépendamment de la position du piston. Ainsi, la densité cible de nombre vaut
\begin{equation}
 n_c^{\star} = \frac{\gamma_0\,\phi_c}{V_c}.
\end{equation}

Il est utile d'introduire en parallèle l'occupation géométriquement requise
\begin{equation}
 \gamma_{\mathrm{geom}} = \frac{N_p}{\sum_c \phi_c},
\end{equation}
qui mesure combien de particules par cellule active seraient nécessaires si la totalité des particules restait confinée dans le volume disponible sous le piston. Le rapport
\begin{equation}
 \mathcal{R}_{\mathrm{cap}} = \frac{\gamma_{\mathrm{geom}}}{\gamma_{\mathrm{target}}}
\end{equation}
sert alors de diagnostic de compression. Lorsque \(\mathcal{R}_{\mathrm{cap}}\) dépasse l'unité, la compression géométrique demandée par le piston excède la capacité incompressible nominale du maillage : la redistribution doit forcer le système au voisinage de sa borne de capacité, et la réaction physique attendue est une forte montée de la pression virielle plutôt qu'une adaptation automatique de la densité cible.

Cette distinction entre \(\gamma_{\mathrm{target}}\) et \(\gamma_{\mathrm{geom}}\) est centrale. Dans l'ancienne logique quasi-compressible, la cible elle-même augmentait avec la diminution du volume actif ; dans la version actuelle, on garde au contraire une densité cible constante et on laisse la pression porter la réaction à la compression.

\subsection{Bande admissible}

La redistribution n'impose pas strictement $N_c=N_c^{\star}$ à chaque pas ; elle cherche plutôt à maintenir $N_c$ dans une bande admissible autour de la cible. Une écriture physique simple est
\begin{equation}
 N_c^{\mathrm{low}} \approx N_c^{\star}(1-\eta),
\qquad
 N_c^{\mathrm{high}} \approx N_c^{\star}(1+\eta),
\end{equation}
où $\eta$ est un paramètre de tolérance relié au champ \verb|coef| dans le code. L'implémentation utilise ensuite des arrondis, des planchers et, dans certaines configurations, des seuils locaux modifiés par la topologie active.

\subsection{Étapes de la redistribution}

La redistribution agit comme un opérateur de projection géométrique sur le champ d'occupation. Elle n'intègre pas une équation d'évolution de type Navier--Stokes ; elle déplace directement des particules pour rapprocher le champ discret $N_c$ de la cible locale $N_c^{\star}$, tout en préservant au mieux la structure spatiale du fluide.

Pour fixer les notations, on introduit pour chaque cellule $c$ :
\begin{itemize}
  \item l'occupation actuelle $N_c$ ;
  \item la cible locale $N_c^{\star}$ ;
  \item les seuils admissibles $N_c^{\mathrm{low}}$ et $N_c^{\mathrm{high}}$ ;
  \item le surplus
  \begin{equation}
    \Delta N_c^{+} = \max\bigl(N_c - N_c^{\mathrm{high}}, 0\bigr),
  \end{equation}
  \item le déficit
  \begin{equation}
    \Delta N_c^{-} = \max\bigl(N_c^{\mathrm{low}} - N_c, 0\bigr).
  \end{equation}
\end{itemize}

Une cellule est dite \emph{trop dense} si $\Delta N_c^{+}>0$, et \emph{trop pauvre} si $\Delta N_c^{-}>0$. Le processus de redistribution est alors organisé en plusieurs sous-opérations.

\paragraph{a) Identification topologique des cellules.}
Le code distingue les cellules de \emph{bulk}, entourées de voisins occupés, des cellules \emph{interfaciales} ou \emph{adjacentes à une paroi}. Cette distinction est importante car la politique de redistribution ne doit pas être la même dans un cœur de fluide homogène, près d'une interface libre, ou au voisinage d'un bord mouillant. La topologie locale sert donc à définir l'espace des cibles admissibles et la direction privilégiée des transferts.

Du point de vue algorithmique, une cellule $c$ n'est déclarée \emph{bulk} que si deux conditions sont satisfaites simultanément :
\begin{enumerate}[label=(\roman*)]
  \item la cellule n'est adjacente à aucune paroi physique ;
  \item tous les voisins du stencil de Moore d'ordre~1, noté $\mathcal{N}_1(c)$, sont occupés.
\end{enumerate}
Si l'on note $\chi_c\in\{0,1\}$ l'indicatrice d'occupation de la cellule $c$, la logique peut s'écrire sous la forme
\begin{equation}
  c\in\mathcal{B}
  \quad \Longleftrightarrow \quad
  c\notin\partial\Omega_{\mathrm{phys}}
  \ \text{et}\ 
  \chi_{c'}=1 \ \text{pour tout}\ c'\in\mathcal{N}_1(c),
\end{equation}
où $\mathcal{B}$ désigne l'ensemble des cellules de bulk et $\partial\Omega_{\mathrm{phys}}$ l'ensemble des cellules collées à une frontière non périodique.

Cette définition est volontairement stricte : dès qu'un voisin orthogonal ou diagonal est vide, ou dès qu'une cellule touche une paroi réelle, la cellule est sortie du bulk et traitée comme cellule interfaciale. Cette convention est cohérente avec la logique physique du code : le bulk est la région où la redistribution peut être effectuée comme un transport quasi isotrope, tandis que les cellules non-bulk doivent respecter la topologie de l'interface, le mouillage ou la géométrie active imposée par un piston.

\paragraph{b) Traitement des cellules trop denses.}
Si $N_c > N_c^{\mathrm{high}}$, la cellule doit céder $\Delta N_c^{+}$ particules. Le transfert ne s'effectue pas vers des cellules arbitraires, mais vers un ensemble de cellules réceptrices $\mathcal{A}_c$ satisfaisant deux contraintes :
\begin{equation}
  N_{c'} < N_{c'}^{\mathrm{target}},
  \qquad c'\in\mathcal{A}_c,
\end{equation}
où $N_{c'}^{\mathrm{target}}$ est soit le seuil bas, soit le seuil haut selon l'étape du processus.

Le code procède d'abord par décharge vers les cellules manifestement sous-peuplées, puis, s'il subsiste un excès, vers des cellules simplement capables d'absorber du fluide sans dépasser leur borne haute. Formellement, si $\mathcal{A}_c^{\mathrm{poor}}$ désigne l'ensemble des cellules sous-denses et $\mathcal{A}_c^{\mathrm{relief}}$ les cellules capables d'absorber un surplus, la redistribution réalise schématiquement
\begin{equation}
  c \longrightarrow \mathcal{A}_c^{\mathrm{poor}} \subseteq \mathcal{A}_c^{\mathrm{relief}}.
\end{equation}

Le classement des receveurs n'est pas isotrope pur : il combine la capacité disponible, la distance, et une direction de transfert privilégiée induite par le gradient local d'occupation ou par la normale intérieure à la paroi. Si $\bm{n}_c$ désigne cette direction privilégiée et $\bm{d}_{cc'}$ le vecteur reliant les centres des cellules $c$ et $c'$, un score typique est de la forme
\begin{equation}
  S_{cc'} = w_1\,\mathcal{C}_{cc'} + w_2\,\frac{\bm{d}_{cc'}\cdot \bm{n}_c}{\norm{\bm{d}_{cc'}}} - w_3\,\norm{\bm{d}_{cc'}},
\end{equation}
où $\mathcal{C}_{cc'}$ représente la capacité d'absorption de la cellule réceptrice, et $w_1,w_2,w_3$ des pondérations implicites dans l'algorithme. L'idée physique est simple : on favorise les transferts vers des régions moins denses et dirigées vers l'intérieur du fluide.

La direction $\bm{n}_c$ n'est pas calculée sur le champ brut $N_c$, mais sur un champ d'occupation lissé $\widetilde{N}_c$ obtenu par convolution locale de type
\begin{equation}
  \widetilde{N}_{i,j} = \frac{1}{16}\Bigl(N_{i-1,j-1}+2N_{i-1,j}+N_{i-1,j+1}+2N_{i,j-1}+4N_{i,j}+2N_{i,j+1}+N_{i+1,j-1}+2N_{i+1,j}+N_{i+1,j+1}\Bigr),
\end{equation}
à renormalisation près lorsque le stencil est tronqué par une frontière non périodique. Le gradient utilisé pour orienter la redistribution s'écrit alors
\begin{equation}
  \nabla \widetilde{N}_c \approx \left(\partial_x \widetilde{N}_c,\, \partial_y \widetilde{N}_c\right),
\end{equation}
avec des différences finies centrées à l'intérieur du domaine et décentrées au premier ordre sur les bords non périodiques. Si l'axe considéré est périodique, la différence finie correspondante est évaluée avec raccordement circulaire des indices ; sinon on utilise une dérivée unilatérale. En notation continue, pour une cellule centrée en $(x_i,y_j)$,
\begin{equation}
  \partial_x \widetilde{N}_{i,j} \approx
  \begin{cases}
    \dfrac{\widetilde{N}_{i+1,j}-\widetilde{N}_{i-1,j}}{2\Delta x}, & 2\le i\le N_x-1,\\[1ex]
    \dfrac{\widetilde{N}_{2,j}-\widetilde{N}_{1,j}}{\Delta x}, & i=1 \text{ et bord non périodique},\\[1ex]
    \dfrac{\widetilde{N}_{N_x,j}-\widetilde{N}_{N_x-1,j}}{\Delta x}, & i=N_x \text{ et bord non périodique},
  \end{cases}
\end{equation}
et de même pour $\partial_y \widetilde{N}_{i,j}$. L'orientation de la redistribution dépend donc explicitement du choix des conditions aux limites.

Près d'une paroi physique, cette direction n'est plus laissée au seul gradient d'occupation : elle est remplacée ou combinée avec la normale intérieure à la paroi, ce qui impose que les particules en excès soient réinjectées vers le bulk plutôt que parallèlement au mur.

\paragraph{c) Traitement des cellules trop pauvres.}
Si $0 < N_c < N_c^{\mathrm{low}}$, la cellule doit recevoir des particules. Le code sélectionne alors des cellules donneuses voisines $d$ telles que
\begin{equation}
  N_d > N_d^{\mathrm{low}},
\end{equation}
avec une préférence pour les cellules du bulk, afin de ne pas déstabiliser l'interface. Si $\mathcal{D}_c$ désigne l'ensemble des donneurs admissibles, le nombre prélevé sur la cellule $d\in\mathcal{D}_c$ est borné par
\begin{equation}
  n_{d\to c} \le \min\bigl(N_d - N_d^{\mathrm{low}},\, N_c^{\mathrm{low}}-N_c\bigr).
\end{equation}

On voit ainsi que la redistribution n'a pas pour but d'égaliser instantanément tous les $N_c$ à $N_c^{\star}$, mais de maintenir chaque cellule dans une bande admissible, avec priorité donnée à la stabilité topologique et à la connectivité du fluide.

\paragraph{d) Redistribution à portée finie.}
Le voisinage principal est d'ordre~1, mais le code ouvre aussi une portée~2 si les cibles immédiates sont insuffisantes. Cette extension peut être vue comme une relaxation du principe local, nécessaire lorsque le fluide se compacte fortement ou lorsqu'une paroi impose une géométrie active réduite. L'opérateur reste cependant fortement local : on ne transporte jamais des particules sur de longues distances en un seul pas.

En pratique, l'algorithme utilise d'abord le voisinage de Moore d'ordre~1,
\begin{equation}
  \mathcal{N}_1(c) = \{c' \neq c\ ;\ \max(|\delta i|,|\delta j|)=1\},
\end{equation}
puis, si l'équilibrage échoue, le voisinage d'ordre~2,
\begin{equation}
  \mathcal{N}_2(c) = \{c' \neq c\ ;\ \max(|\delta i|,|\delta j|)\le 2\}.
\end{equation}
Les candidats sont ensuite ordonnés de manière hiérarchique : d'abord par capacité disponible, puis par alignement avec la direction privilégiée, enfin par distance géométrique. Le code exploite pour cela le pas minimal entre cellules en tenant compte de la périodicité éventuelle sur chaque axe, de sorte que la notion de ``plus proche'' dépend bien de la topologie réelle du domaine.

\paragraph{e) Overflow bulk-interface et saturation.}
Lorsque le cœur du fluide est trop dense et que les voisins proches ne suffisent plus, le code autorise un report du surplus du bulk vers les cellules géométriquement interfaciales les plus proches. Cela correspond à une politique d'\emph{overflow} : l'excès de masse numérique est repoussé vers les régions déjà identifiées comme interface ou bord libre. Si même cette étape ne suffit pas, le code signale une \emph{saturation} de redistribution. Physiquement, cela signifie que la contrainte de quasi-incompressibilité devient plus difficile à satisfaire avec les seules relocalisations locales permises par le maillage et la topologie instantanée.

\paragraph{f) Interprétation physique.}
Du point de vue continu, on peut interpréter la redistribution comme une étape de projection approximative vers l'hypersurface discrète
\begin{equation}
  \mathcal{M}_{\rho^{\star}} = \left\{ \{N_c\}_{c=1}^{N_xN_y} \; ; \; N_c \approx N_c^{\star} \text{ pour tout } c \right\},
\end{equation}
qui représente un régime à compressibilité très faible. L'opérateur ne résout pas une équation de pression au sens analytique ; il agit comme une correction de densité imposée sur le champ particulaire lui-même.

\subsection{Correction locale de quantité de mouvement dans la redistribution}

Si $\bm{P}_c^{\mathrm{before}}$ et $\bm{P}_c^{\mathrm{after}}$ sont les quantités de mouvement de la cellule $c$ avant et après les transferts, le code calcule
\begin{equation}
 \delta \bm{P}_c = \bm{P}_c^{\mathrm{before}} - \bm{P}_c^{\mathrm{after}},
\end{equation}
et applique, quand cela est permis, une correction uniforme par particule
\begin{equation}
 \delta \bm{u}_c = \frac{\delta \bm{P}_c}{N_c},
\end{equation}
qui redonne à la cellule une quantité de mouvement proche de sa valeur précédente. Cette correction reste locale et ne constitue pas encore la fermeture liquide elle-même ; elle sert d'abord à éviter que la redistribution ne détruise inutilement le champ moyen.


\subsection{Traitement de surface, mouillage et réorientation interfaciale}

Le code distingue soigneusement deux familles de mécanismes de surface :
\begin{enumerate}[label=\alph*)]
  \item un \textbf{traitement géométrique de surface} intégré à la redistribution, destiné à préserver une topologie cohérente de l'interface et, si nécessaire, une couche mouillante près des parois ;
  \item une \textbf{réorientation cinématique interfaciale} appliquée après la fermeture liquide, destinée à modéliser une cohésion tangentielle ou une tension superficielle effective.
\end{enumerate}

Ces deux mécanismes ne jouent pas le même rôle.

\paragraph{a) Surface et mouillage dans la redistribution.}
Lorsqu'une paroi porte une condition de type mouillage, la redistribution peut imposer localement une occupation minimale dans les cellules voisines de cette paroi. Si l'on note $N_c^{\mathrm{wet}}$ la cible locale de mouillage, la contrainte s'écrit schématiquement
\begin{equation}
  N_c \gtrsim N_c^{\mathrm{wet}}
  \qquad \text{pour } c \in \mathcal{W},
\end{equation}
où $\mathcal{W}$ désigne l'ensemble des cellules de la couche mouillante.

Le but n'est pas d'introduire une force de surface explicite, mais de maintenir une population suffisante à proximité des bords physiques pour éviter que l'interface ne se décroche numériquement de la paroi. Cette logique est particulièrement utile lorsque le fluide est soumis à la gravité, à une compression par piston, ou à des redistributions répétées.

\paragraph{b) Direction privilégiée et normales de surface.}
Dans les cellules interfaciales, la redistribution utilise une normale intérieure approchée $\bm{n}_c$ obtenue à partir du gradient local d'occupation lissée. Si $\widetilde{N}(x,y)$ désigne un champ d'occupation discrètement lissé, la normale intérieure s'écrit, à normalisation près,
\begin{equation}
  \bm{n}_c \propto \nabla \widetilde{N}_c.
\end{equation}
Le signe est choisi de façon à pointer vers le bulk. Cette normale intervient à la fois dans le classement des cellules receveuses lors de la redistribution et, plus tard, dans la réorientation des vitesses interfaciales.

Du point de vue algorithmique, cette normale est obtenue par les mêmes opérateurs de gradient que ceux utilisés pour l'orientation de redistribution. Le champ lissé $\widetilde N$ est d'abord construit, puis les composantes
\begin{equation}
  g_x = \partial_x \widetilde N,
  \qquad
  g_y = \partial_y \widetilde N
\end{equation}
sont évaluées cellule par cellule. La direction intérieure associée à une cellule interfaciale est alors
\begin{equation}
  \widehat{\bm n}_c = \frac{(g_x,g_y)}{\sqrt{g_x^2+g_y^2}+\varepsilon_n},
\end{equation}
où $\varepsilon_n>0$ évite les divisions par zéro lorsque le gradient local est trop faible. Si la cellule est directement au contact d'une paroi physique, cette normale est remplacée par la normale géométrique intérieure au mur ; si plusieurs parois se rencontrent, par exemple dans un coin, les normales associées sont sommées puis renormalisées. Cette règle garantit que la direction privilégiée reste compatible avec la géométrie du domaine même lorsque le gradient d'occupation local devient ambigu.

\paragraph{c) Réorientation interfaciale.}
Après redistribution, correction de champ et kick liquide, le code peut appliquer une réorientation partielle des vitesses dans les cellules identifiées comme interfaciales. Soit $\bm{v}_i$ la vitesse d'une particule d'interface, et $\widehat{\bm{n}}_c$ la normale intérieure unitaire associée à la cellule $c$. La réorientation construit une vitesse cible
\begin{equation}
  \bm{v}_i^{\mathrm{tar}} = \abs{\bm{v}_i}\,\widehat{\bm{n}}_c,
\end{equation}
puis applique une interpolation linéaire
\begin{equation}
  \bm{v}_i^{\mathrm{new}} = (1-\beta_c)\,\bm{v}_i + \beta_c\,\bm{v}_i^{\mathrm{tar}},
\end{equation}
où $\beta_c\in[0,1]$ est un coefficient local de réorientation.

Le code ne réoriente qu'une fraction $\varphi_c\in[0,1]$ des particules de la cellule, ce qui revient à imposer une correction cinématique partielle. Les coefficients $\beta_c$ et $\varphi_c$ peuvent être renforcés en fonction d'un indicateur de courbure ou de désalignement local. On peut formaliser cela par
\begin{equation}
  \beta_c = \beta_0\,\mathcal{F}_\beta(\kappa_c),
  \qquad
  \varphi_c = \varphi_0\,\mathcal{F}_\varphi(\kappa_c),
\end{equation}
où $\kappa_c$ désigne un indicateur local de courbure ou de désalignement, et $\beta_0$, $\varphi_0$ des paramètres de base.

Dans le code, l'indicateur $\kappa_c$ n'est pas une courbure au sens différentiel strict $\nabla\cdot\widehat{\bm n}$, mais un défaut d'alignement des normales locales calculé à partir du voisinage de la cellule interfaciale. L'effet est le suivant : plus les normales voisines varient fortement, plus l'interface est localement coudée ou irrégulière, et plus les paramètres de réorientation sont augmentés. La réorientation agit donc davantage dans les régions de forte courbure apparente ou de surface mal alignée.

\paragraph{d) Signification physique.}
La réorientation n'est pas une tension superficielle au sens strict d'une force de courbure écrite comme $\sigma \kappa \bm{n}$. Elle agit plutôt comme une régularisation cinématique qui tend à réaligner les vitesses interfaciales vers le bulk, donc à limiter l'étalement tangent et la dispersion de l'interface. On peut la voir comme un analogue discret d'une dissipation interfaciale orientée ou d'une cohésion de surface effective.

Cette construction est conservée ici comme historique de développement mais elle ne doit pas être confondue avec la tension superficielle de production du cycle 0493x9. Le chemin x9 n'oriente pas directement les vitesses particulaires : il calcule une courbure géométrique et impose le saut normal de Laplace dans la condition de pression de Q6-g-f, comme détaillé dans la section consacrée au cycle 0493x9.

Le choix de la placer \emph{après} la fermeture liquide est cohérent avec l'architecture actuelle : la réorientation admet localement des cellules de surface plus pauvres, alors que redistribution, correction et kick cherchent au contraire à maintenir une occupation contrôlée. Il est donc préférable de ne pas inclure la réorientation dans l'état de référence utilisé pour la fermeture.

\section{Correction des champs : vitesse, température, énergie}

\subsection{Référence, état redistribué et défaut à corriger}

Soit $\bm{u}_c^{\mathrm{ref}}$ la vitesse moyenne de la cellule $c$ dans l'état de référence pré-redistribution, et $\bm{u}_c^{\mathrm{red}}$ la vitesse moyenne de cette même cellule après redistribution. On définit le défaut de quantité de mouvement
\begin{equation}
 \delta \bm{P}_c^{\mathrm{rep}} = \bm{P}_c^{\mathrm{ref}} - \bm{P}_c^{\mathrm{red}}.
\end{equation}

La correction externe de fermeture liquide introduit alors un champ de réparation uniforme dans la cellule,
\begin{equation}
 \delta \bm{u}_c^{\mathrm{rep}} = \frac{\delta \bm{P}_c^{\mathrm{rep}}}{N_c^{\star}},
\end{equation}
où $N_c^{\star}$ est la cible locale d'occupation.

Le code multiplie ensuite ce champ par un scalaire $\beta_{\mathrm{rep}}\ge 0$ choisi soit manuellement, soit par minimisation quadratique d'un défaut global sur les vitesses de cellule. La vitesse réparée s'écrit donc
\begin{equation}
 \bm{u}_c^{\mathrm{rep}} = \bm{u}_c^{\mathrm{red}} + \beta_{\mathrm{rep}}\,\delta \bm{u}_c^{\mathrm{rep}}.
\end{equation}

Dans le code, le coefficient $\beta_{\mathrm{rep}}$ peut être choisi par minimisation quadratique de l'écart entre champ redistribué et champ de référence. Si l'on note
\begin{equation}
  \bm e_c = \bm u_c^{\mathrm{red}} - \bm u_c^{\mathrm{ref}},
  \qquad
  \bm a_c = \delta \bm u_c^{\mathrm{rep}},
\end{equation}
alors le problème traité est de minimiser
\begin{equation}
  J(\beta) = \sum_c \abs{\bm e_c + \beta \, \bm a_c}^2,
\end{equation}
d'où l'optimum formel
\begin{equation}
  \beta_{\mathrm{rep}}^{\star} = -\frac{\sum_c \bm e_c\cdot \bm a_c}{\sum_c \abs{\bm a_c}^2},
\end{equation}
ensuite tronqué si nécessaire pour respecter les contraintes de l'utilisateur. Ainsi, la correction de champ n'est pas seulement une translation imposée ; elle est calibrée pour rapprocher au mieux le champ redistribué du champ de référence dans une norme quadratique discrète.

\subsection{Température : ce qui est corrigé et ce qui ne l'est pas}

Dans l'état actuel du couplage intégré, la correction externe agit explicitement sur la \textbf{vitesse moyenne} de chaque cellule, pas sur la température. Cette absence d'une correction thermique explicite n'est pas incohérente : une correction uniforme de type
\begin{equation}
 \bm{v}_i \leftarrow \bm{v}_i + \delta \bm{u}_c,
 \qquad i\in c,
\end{equation}
ne modifie pas les fluctuations relatives $\bm{v}_i-\bm{u}_c$ à appartenance cellulaire fixée, donc ne modifie pas directement $k_B T_c$.

En revanche, la redistribution elle-même déplace des particules d'une cellule à l'autre et peut donc modifier la température locale reconstruite. Dans le code actuel, la régularisation thermique est surtout assurée par le thermostat MPCD des collisions SRD et, éventuellement, par les parois thermalisantes.

\subsection{Effet énergétique d'un kick uniforme par cellule}

Considérons une cellule contenant $N_c$ particules et soumise à une correction uniforme $\delta \bm{u}_c$. Si l'on note $\bm{P}_c = \sum_{i\in c} \bm{v}_i$ sa quantité de mouvement avant correction, alors la variation d'énergie cinétique totale de la cellule est
\begin{equation}
 \Delta E_c = \bm{P}_c\cdot \delta \bm{u}_c + \frac{N_c}{2}\abs{\delta \bm{u}_c}^2.
\end{equation}

Cette variation affecte l'énergie macroscopique liée au mouvement moyen de la cellule. L'énergie interne relative,
\begin{equation}
 E_{\mathrm{int},c} = \frac{1}{2}\sum_{i\in c}\abs{\bm{v}_i-\bm{u}_c}^2,
\end{equation}
reste inchangée par un \emph{kick} uniforme à appartenance cellulaire fixée.

Ainsi, dans l'état actuel du code, la correction et le \emph{kick} incompressible modifient la répartition entre énergie moyenne et énergie potentiellement injectée/retraitée par le forçage, mais ne constituent pas un thermostat supplémentaire.

\section{Kick incompressible et fermeture liquide}

\subsection{Pression motrice}

La fermeture liquide actuelle repose sur une pression motrice construite à partir de l'état de référence pré-redistribution,
\begin{equation}
 P_{\mathrm{drive}} = P_{\mathrm{kin}} + P_{\mathrm{vir}}.
\end{equation}

La partie cinétique est
\begin{equation}
 P_{\mathrm{kin},c} = n_c^{\mathrm{ref}}\,k_B T_c^{\mathrm{ref}},
\end{equation}
où $n_c^{\mathrm{ref}}$ est la densité de nombre de la cellule $c$ dans l'état de référence, et $T_c^{\mathrm{ref}}$ la température locale reconstruite dans ce même état.

La partie virielle est définie par
\begin{equation}
 P_{\mathrm{vir},c} = K_{\mathrm{virial}}\,\bigl(n_c^{\mathrm{ref}} - n_c^{\star}\bigr),
\end{equation}
où $K_{\mathrm{virial}}$ est le coefficient de raideur virielle et $n_c^{\star}$ la densité cible issue de la redistribution.

Ainsi,
\begin{equation}
 P_{\mathrm{drive},c} = n_c^{\mathrm{ref}} k_B T_c^{\mathrm{ref}} + K_{\mathrm{virial}}\,\bigl(n_c^{\mathrm{ref}} - n_c^{\star}\bigr).
\end{equation}

\subsection{Interprétation physique}

Le rôle de $P_{\mathrm{vir}}$ n'est pas de remplacer la pression cinétique, mais d'ajouter une raideur effective au fluide lorsque la redistribution maintient $n_c$ près d'une cible donnée. Dans un régime quasi incompressible, il est naturel que la pente hydrostatique moyenne soit portée principalement par $P_{\mathrm{vir}}$ plutôt que par $P_{\mathrm{kin}}$.

Si la redistribution est très serrée, le gradient de densité accessible au fluide est lui-même borné ; dans ce cas, modifier $K_{\mathrm{virial}}$ change surtout l'amplitude de la contre-pression et la charge algorithmique supportée par la redistribution, sans nécessairement changer fortement l'état stationnaire moyen.


\subsection{Compression par piston et module volumique effectif}

Dans le cas du piston incompressible, la densité cible n'est plus adaptée géométriquement : la compression du volume actif doit donc se traduire par une montée de pression. Si l'on note \(V_{\mathrm{actif}}(t)=L_x H_{\mathrm{actif}}(t)\) l'aire active sous le piston et \(P_{\mathrm{tot,actif}}(t)\) la pression totale moyenne du modèle sur ce domaine actif, on définit le module volumique effectif
\begin{equation}
 B_{\mathrm{eff}}(t)
 =
 -\,V_{\mathrm{actif}}(0)\,
 \frac{P_{\mathrm{tot,actif}}(t)-P_{\mathrm{tot,actif}}(0)}
 {V_{\mathrm{actif}}(t)-V_{\mathrm{actif}}(0)}
 \;\approx\;
 \frac{\Delta P_{\mathrm{tot,actif}}(t)}{-\Delta V_{\mathrm{actif}}(t)/V_{\mathrm{actif}}(0)}.
\end{equation}
Sa réciproque
\begin{equation}
 \kappa_{\mathrm{eff}}(t)=\frac{1}{B_{\mathrm{eff}}(t)}
\end{equation}
joue le rôle d'une compressibilité effective.

Le code calcule aussi un module local apparent dérivé de l'équation d'état discrète,
\begin{equation}
 \left(\frac{\mathrm d P_{\mathrm{tot}}}{\mathrm d \rho}\right)_{\mathrm{eff}}
 \approx
 \frac{P_{\mathrm{tot,actif}}(t)-P_{\mathrm{tot,actif}}(0)}
 {\rho_{\mathrm{actif}}(t)-\rho_{\mathrm{actif}}(0)},
\end{equation}
ainsi que la quantité
\begin{equation}
 \rho_{\mathrm{actif}}(t)\,
 \left(\frac{\mathrm d P_{\mathrm{tot}}}{\mathrm d \rho}\right)_{\mathrm{eff}},
\end{equation}
qui doit être comparable à \(B_{\mathrm{eff}}\) si la montée de pression est effectivement portée par la variation de densité moyenne dans le domaine comprimé. Ces diagnostics permettent de quantifier explicitement la faible compressibilité du modèle plutôt que de la déduire uniquement de la forme des profils.

\subsection{Gradient de pression et \emph{kick} de vitesse}

Le code calcule le gradient de $P_{\mathrm{drive}}$ avec un opérateur compatible avec les conditions aux limites de chaque axe. Le \emph{kick} incompressible de fermeture s'écrit ensuite
\begin{equation}
 \delta \bm{u}_c^{\mathrm{EOS}} = -\beta_{\mathrm{EOS}}\,\frac{\Delta t}{n_c^{\star}}\,\nabla P_{\mathrm{drive},c},
\end{equation}
où $\beta_{\mathrm{EOS}}$ est un facteur scalaire de réglage et $n_c^{\star}$ la densité cible locale.

Dans la version actuelle du couplage, les dérivées spatiales sont évaluées sur la grille d'analyse par différences finies. Si $P_{i,j}$ désigne la pression motrice dans la cellule centrée en $(x_i,y_j)$, alors les composantes du gradient sont approchées par
\begin{equation}
  \partial_x P_{i,j} \approx
  \begin{cases}
    \dfrac{P_{i+1,j}-P_{i-1,j}}{2\Delta x}, & 2\le i\le N_x-1,\\[1ex]
    \dfrac{P_{2,j}-P_{1,j}}{\Delta x}, & i=1,\\[1ex]
    \dfrac{P_{N_x,j}-P_{N_x-1,j}}{\Delta x}, & i=N_x,
  \end{cases}
\end{equation}
\begin{equation}
  \partial_y P_{i,j} \approx
  \begin{cases}
    \dfrac{P_{i,j+1}-P_{i,j-1}}{2\Delta y}, & 2\le j\le N_y-1,\\[1ex]
    \dfrac{P_{i,2}-P_{i,1}}{\Delta y}, & j=1,\\[1ex]
    \dfrac{P_{i,N_y}-P_{i,N_y-1}}{\Delta y}, & j=N_y.
  \end{cases}
\end{equation}
Cette écriture explicite la logique du code : dérivées centrées dans le bulk, décentrées au premier ordre au voisinage des bords. Le choix est cohérent avec le fait que l'état de fermeture n'est pas construit comme un champ périodique abstrait, mais comme un champ défini dans la géométrie physique du canal ou du domaine comprimé par piston.

Le signe est celui attendu pour une force de pression. Si la direction verticale $y$ est croissante vers le haut et si la gravité vaut $g<0$, alors un équilibre hydrostatique moyen correspond à
\begin{equation}
 \frac{\mathrm{d}P}{\mathrm{d}y} \approx n^{\star} g
\end{equation}
dans les unités du code, ou plus généralement à
\begin{equation}
 \frac{\mathrm{d}P}{\mathrm{d}y} \approx \rho_m g
\end{equation}
si l'on réintroduit explicitement la masse particulaire.

\subsection{Résidu hydrostatique}

Le post-traitement reconstruit les profils moyens selon $x$ et calcule le résidu hydrostatique
\begin{equation}
 R_h(y) = \frac{\mathrm{d}P_{\mathrm{tot}}}{\mathrm{d}y} - n(y)g
\end{equation}
ou, selon la convention choisie,
\begin{equation}
 R_h^{\star}(y) = \frac{\mathrm{d}P_{\mathrm{tot}}}{\mathrm{d}y} - n^{\star}(y)g.
\end{equation}

Dans une simulation MPCD bruitée, un résidu instantané local non nul n'est pas en soi rédhibitoire. La quantité significative est plutôt le résidu moyen ou lissé, et la comparaison entre la pente moyenne de $P_{\mathrm{tot}}$ et le forçage $n^{\star}g$ dans le \emph{bulk} du domaine.

\section{Bilan de la chaîne redistribution/liquid-closure : état historique}

L'état actuel du code peut être résumé de la manière suivante.

\begin{enumerate}[label=\arabic*)]
  \item La méthode MPCD proprement dite est cohérente : advection, décalage aléatoire de grille, rotation SRD et thermostat cellulaire sont bien en place.
  \item Les conditions limites sont suffisamment riches pour traiter des parois spéculaires, non glissantes, thermalisantes et un piston mobile.
  \item La redistribution agit comme un opérateur de projection vers un régime faiblement compressible contrôlé par une cible d'occupation locale. Dans la version piston actuelle, cette cible peut être maintenue constante afin de rendre la compression essentiellement virielle.
  \item La correction externe actuellement couplée agit d'abord sur le \emph{champ moyen de vitesse}, puis applique un \emph{kick} de pression construit à partir de la pression motrice \(P_{\mathrm{kin}}+P_{\mathrm{vir}}\) évaluée sur la référence.
  \item Il n'existe pas encore, dans le couplage intégré, de correction thermique explicite séparée ; la température reste principalement contrôlée par le thermostat MPCD et par les conditions limites thermalisantes.
  \item Les diagnostics récents distinguent maintenant les états instantanés, les champs moyens temporels, les diagnostics Poiseuille et les diagnostics piston, ce qui permet d'analyser séparément le bulk, la paroi et la réaction de la fermeture liquide.
\end{enumerate}

\section{Résultats historiques de la redistribution et de la fermeture liquide}

\subsection{Cas Poiseuille}

Le cas Poiseuille a servi de banc d'essai principal pour qualifier à la fois le bulk MPCD, les conditions limites et l'effet de la redistribution/fermeture liquide.

\paragraph{Effet de la résolution et des conditions limites.}
Les premiers essais ont montré un glissement important aux parois lorsque la direction transverse était insuffisamment résolue ou lorsque les parois utilisaient une thermalisation. Le raffinement transverse \((N_y=100)\) et l'emploi de conditions limites de type \emph{bounce-back} ont nettement amélioré la situation : le profil longitudinal \(u_x(y)\) devient quasi parfaitement parabolique, avec un coefficient de détermination du fit quadratique proche de l'unité. Cette observation montre que la qualité des parois et la résolution transverse sont déterminantes avant même d'analyser la fermeture incompressible.

\paragraph{Comparaison compressible / incompressible.}
En présence de redistribution et de fermeture liquide, les fluctuations d'occupation sont fortement réduites, typiquement avec \(\mathrm{std}(N)/\gamma \approx 0.13\) et \(\mathrm{outBand}=0\), alors que le cas compressible laisse réapparaître des fluctuations plus larges et des cellules hors bande. Le profil de Poiseuille reste cependant très propre dans les deux cas une fois la résolution et les parois convenablement choisies. La redistribution et la fermeture liquide agissent alors surtout comme un resserrement du système : elles réduisent la compressibilité numérique, diminuent le glissement effectif et stabilisent l'identification de la viscosité.

\paragraph{Convergence à long temps.}
Sur des temps de simulation suffisamment longs, les estimations de viscosité par courbure du profil et par contrainte pariétale deviennent cohérentes. Dans les runs de référence longs, on observe typiquement
\begin{equation}
 \nu_{\mathrm{curve}} \approx \nu_{\tau},
\end{equation}
avec un fit quadratique excellent (\(R^2 \simeq 0.999\)). Cela indique que le bulk et les parois racontent enfin la même physique, et que le régime de Poiseuille est effectivement établi. Il reste un glissement absolu non nul aux parois, qui doit être interprété comme une propriété effective de la condition limite \emph{bounce-back} à résolution finie, plutôt que comme une rupture de la structure parabolique du profil.

\subsection{Cas \texttt{compressibility} : hydrostatique moyenne et fermeture liquide}

Le cas \texttt{compressibility} a été utilisé pour qualifier la capacité de la fermeture liquide à produire une hydrostatique moyenne. Les diagnostics instantanés sont naturellement bruités, comme attendu en MPCD, mais les champs moyens temporels montrent un comportement très clair :
\begin{itemize}
  \item la température moyenne reste presque uniforme ;
  \item la densité moyenne reste proche de la cible, avec une stratification faible mais non nulle ;
  \item la pression totale moyenne
  \begin{equation}
   \overline{P}_{\mathrm{tot}}(y)=\overline{P}_{\mathrm{kin}}(y)+\overline{P}_{\mathrm{vir}}(y)
  \end{equation}
  suit une loi affine raisonnable dans le bulk ;
  \item le résidu hydrostatique moyen
  \begin{equation}
   \overline{R}_h(y)=\frac{\mathrm d \overline{P}_{\mathrm{tot}}}{\mathrm d y}-\overline{n}(y)\,g
  \end{equation}
  reste faible au centre du domaine et ne se dégrade nettement qu'au voisinage des bords.
\end{itemize}

L'interprétation principale est que l'hydrostatique moyenne est portée majoritairement par le terme viriel. Ce résultat est cohérent avec l'objectif poursuivi : pour un liquide quasi incompressible, on n'attend pas que la seule pression cinétique porte la réaction mécanique, mais qu'une pression de fermeture très raide prenne le relais. La comparaison de différentes valeurs de \(K_{\mathrm{virial}}\) montre par ailleurs qu'une redistribution très serrée borne déjà fortement le défaut de densité accessible ; dans ce régime, augmenter \(K_{\mathrm{virial}}\) modifie surtout l'amplitude de la pression virielle et la charge algorithmique de redistribution, sans changer drastiquement l'état moyen.

\subsection{Cas piston : compression incompressible et montée de pression}

Le traitement piston a été réorienté vers une logique véritablement quasi incompressible. Le piston réduit le volume actif en diminuant le nombre effectif de cellules disponibles, mais sans adapter la cible \(\gamma_{\mathrm{target}}\) à cette réduction. Par conséquent :
\begin{itemize}
  \item l'occupation géométriquement requise \(\gamma_{\mathrm{geom}}\) augmente avec la compression ;
  \item la cible incompressible \(\gamma_{\mathrm{target}}=\gamma_0\) reste constante ;
  \item la redistribution ne peut plus ``absorber'' la compression en augmentant artificiellement la densité cible ;
  \item la réaction attendue est donc une forte montée de \(P_{\mathrm{vir}}\), puis de \(P_{\mathrm{tot}}\).
\end{itemize}

Les diagnostics temporels montrent effectivement ce comportement. Pour des compressions géométriques modérées de l'ordre de quelques pourcents,
\begin{equation}
 \frac{\gamma_{\mathrm{geom}}}{\gamma_{\mathrm{target}}} > 1,
\end{equation}
tandis que la pression cinétique active reste presque constante et que la pression virielle augmente fortement, entraînant avec elle la pression totale active. Dans les exemples de référence, une compression géométrique de l'ordre de \(4\%\) conduit à une augmentation très marquée de \(P_{\mathrm{vir}}\), alors que \(P_{\mathrm{kin}}\) varie peu. Le piston est donc bien le cas où la fermeture liquide révèle le plus clairement sa faible compressibilité.

Le calcul du module volumique effectif
\begin{equation}
 B_{\mathrm{eff}} \approx \frac{\Delta P_{\mathrm{tot,actif}}}{-\Delta V_{\mathrm{actif}}/V_{\mathrm{actif}}(0)}
\end{equation}
et de la compressibilité effective \(\kappa_{\mathrm{eff}}=1/B_{\mathrm{eff}}\) fournit enfin un diagnostic direct de cette raideur. Les valeurs obtenues dans les tests de référence sont cohérentes avec la pente \(\mathrm d P_{\mathrm{tot}}/\mathrm d \rho\) reconstruite sur le domaine actif et confirment que le modèle se comporte comme un fluide très peu compressible. À mesure que la compression progresse, on s'attend naturellement à l'apparition d'un régime de saturation de redistribution : ce n'est pas un artefact, mais l'expression numérique du caractère bloquant d'un liquide quasi incompressible confiné dans un volume trop petit.


\section{Redistribution par zones : principe, découpage et double passe}

\subsection{Objectif de la branche par zones}

La redistribution par zones a été introduite comme une étape intermédiaire entre la redistribution de référence, entièrement séquentielle sur le domaine complet, et une future implémentation parallèle plus ambitieuse. L'idée directrice est de conserver l'opérateur physique actuel de redistribution, mais de l'appliquer localement sur des sous-domaines rectangulaires du maillage, de manière à tester la sensibilité de la physique numérique à un traitement ``simili-parallèle'' avant tout portage compilé.

Dans la branche actuelle, l'ordre général du \emph{step} n'est pas modifié : l'advection, les conditions limites, la collision SRD, la capture de l'état de référence, la redistribution, la fermeture liquide et les diagnostics conservent la même séquence que dans la version de référence. Seul l'opérateur de redistribution peut être remplacé par une variante ``par zones'', commandée par les paramètres \verb|useZoneRedistribution|, \verb|zoneTileNx|, \verb|zoneTileNy|, \verb|zoneUseShiftedSecondPass|, \verb|zoneStoreDiagnostics|, \verb|zonePassthroughMode| et \verb|zoneSingleFullDomainMode|.

Cette branche a été utilisée comme outil d'évaluation : il ne s'agit pas encore d'une redistribution parallèle complète, mais d'un cadre permettant de mesurer ce qui est préservé ou perdu lorsque l'on remplace un traitement global par une suite d'opérateurs locaux.

\subsection{Construction du découpage en zones}

Soit un maillage cartésien de taille $N_x \times N_y$. On choisit deux tailles de zones entières
\begin{equation}
N_x^{(z)} = \texttt{zoneTileNx},
\qquad
N_y^{(z)} = \texttt{zoneTileNy}.
\end{equation}
Le domaine est alors découpé en rectangles alignés sur la grille. Pour la passe dite \emph{de base}, les indices de départ sont
\begin{equation}
 i_x^{(0)} \in \{1,\,1+N_x^{(z)},\,1+2N_x^{(z)},\dots\},
 \qquad
 i_y^{(0)} \in \{1,\,1+N_y^{(z)},\,1+2N_y^{(z)},\dots\},
\end{equation}
chaque zone étant ensuite tronquée au bord du domaine si nécessaire. Si l'on note
\begin{equation}
\mathcal Z_{pq}^{\mathrm{base}} = \{(i,j)\, :\, i_{x,p}^{(0)} \le i \le i_{x,p}^{(1)},\; i_{y,q}^{(0)} \le j \le i_{y,q}^{(1)}\},
\end{equation}
alors chaque cellule appartient à une zone de cette famille de base.

Pour un domaine $50\times 50$ et des zones $25\times 25$, la passe de base produit quatre sous-domaines :
\begin{equation}
 x = 1{:}25 \;\text{ou}\; 26{:}50,
 \qquad
 y = 1{:}25 \;\text{ou}\; 26{:}50.
\end{equation}
L'interface principale du pavage est alors située entre les colonnes $25$ et $26$, et entre les lignes $25$ et $26$. Sur un domaine physique $[0,10]\times[0,10]$, avec $50$ cellules dans chaque direction, cela correspond à $x=5$ et $y=5$.

\subsection{Seconde passe décalée}

Lorsque \verb|zoneUseShiftedSecondPass=true|, une seconde famille de zones est construite avec un décalage d'environ une demi-zone dans chaque direction. Le décalage est pris entier :
\begin{equation}
 o_x = \left\lfloor \frac{N_x^{(z)}}{2} \right\rfloor,
 \qquad
 o_y = \left\lfloor \frac{N_y^{(z)}}{2} \right\rfloor.
\end{equation}
Les points de départ de la famille décalée deviennent alors
\begin{equation}
 i_x^{(s)} \in \{1\} \cup \{1+o_x,\,1+o_x+N_x^{(z)},\,1+o_x+2N_x^{(z)},\dots\},
\end{equation}
\begin{equation}
 i_y^{(s)} \in \{1\} \cup \{1+o_y,\,1+o_y+N_y^{(z)},\,1+o_y+2N_y^{(z)},\dots\}.
\end{equation}
La présence systématique du point de départ $1$ garantit une couverture correcte du bord même lorsque le décalage ne divise pas exactement le domaine.

Pour $N_x^{(z)}=N_y^{(z)}=25$ sur un domaine $50\times 50$, on obtient
\begin{equation}
 o_x=o_y=12,
\end{equation}
et les départs de la passe décalée sont
\begin{equation}
 x : 1,\;13,\;38,
 \qquad
 y : 1,\;13,\;38,
\end{equation}
ce qui produit des zones du type
\begin{equation}
1{:}25,
\qquad
13{:}37,
\qquad
38{:}50.
\end{equation}
La deuxième passe ne constitue donc pas un simple balayage dans l'ordre inverse ; elle reconstruit une nouvelle famille géométrique de zones, toujours rectangulaires et alignées sur la grille, mais décalées d'une demi-zone.

\subsection{Transferts strictement intra-zone}

Le point structurel le plus important est que chaque famille de zones est traitée séquentiellement, zone par zone, avec des transferts strictement intra-zone. Si l'on note $\mathcal Z$ la zone courante, les cellules sources denses et les cellules cibles admissibles sont filtrées par le masque de $\mathcal Z$ ; par conséquent, une zone ne peut pas, au cours d'une même passe, décharger directement un surplus dans une cellule située dans une autre zone, ni recevoir directement des particules d'une autre zone.

La seconde passe n'introduit donc pas une compensation interzone globale ; elle fournit seulement une deuxième occasion de redistribuer les particules dans une autre partition du domaine. Une cellule qui se trouvait près d'une frontière de la famille de base peut ainsi se retrouver au milieu d'une zone de la famille décalée, et être corrigée à ce moment-là. Cette correction reste cependant indirecte.

\subsection{Modes de contrôle et validation}

Deux modes auxiliaires ont été utiles pendant la mise au point :
\begin{itemize}
  \item \verb|zonePassthroughMode=true| : la branche zones est traversée, mais délègue entièrement à la redistribution de référence ; elle sert à vérifier que l'orchestrateur zones ne modifie pas la physique lorsqu'il est neutre.
  \item \verb|zoneSingleFullDomainMode=true| : une seule zone recouvre tout le domaine ; le chemin ``zones'' est alors exercé dans une configuration contrôlée, sans introduire de véritable découpage géométrique.
\end{itemize}
Ces deux étapes ont permis de vérifier que la structure \emph{runner -> simulateur -> step -> redistribution par zones} était correctement branchée avant de tester le vrai mode \verb|multi_zone_active|.

\section{Lissage interzone de $P_{\mathrm{drive}}$}

\subsection{Motivation}

Les premiers diagnostics sur le cas \texttt{compressibility} ont montré que la redistribution par zones ne détruisait pas fortement les champs moyens juste après la redistribution, mais qu'un défaut local persistait ensuite après le \emph{post-kick}. L'analyse détaillée du \emph{step} a montré que la fermeture liquide applique, après une correction de quantité de mouvement, un \emph{kick} de vitesse fondé sur le gradient de la pression motrice
\begin{equation}
 P_{\mathrm{drive}} = P_{\mathrm{kin}} + P_{\mathrm{vir}},
\end{equation}
avec
\begin{equation}
 P_{\mathrm{vir}} = K_{\mathrm{virial}}\,(\rho^{\mathrm{ref}} - \rho^{\star}).
\end{equation}
Le \emph{kick} vertical vaut alors
\begin{equation}
 \delta u_y^{\mathrm{EOS}} = -\beta_{\mathrm{EOS}}\,\frac{\Delta t}{\rho^{\star}}\,\partial_y P_{\mathrm{drive}}.
\end{equation}
En présence d'une petite discontinuité de $P_{\mathrm{drive}}$ à une interface de zones, le gradient discret peut amplifier localement cette discontinuité et créer un saut visible sur $u_y$.

\subsection{Principe du lissage}

Pour limiter cette amplification, un lissage local de $P_{\mathrm{drive}}$ a été introduit avant le calcul du gradient. L'idée est de ne lisser que dans une bande centrée sur les interfaces géométriques du zonage. Trois paramètres pilotent cette correction :
\begin{itemize}
  \item \verb|smoothPdriveInterzoneWidthCells| : demi-largeur de la bande de lissage, exprimée en nombre de cellules ;
  \item \verb|smoothPdriveInterzoneBlend| : coefficient de mélange entre le champ brut et le champ lissé ;
  \item \verb|smoothPdriveIncludeShiftedLayouts| : choix d'inclure ou non les interfaces de la famille décalée dans le masque de lissage.
\end{itemize}
Si l'on note $P_{\mathrm{drive}}^{\mathrm{raw}}$ le champ brut et $\widetilde P_{\mathrm{drive}}$ sa version lissée localement, la correction utilisée est du type
\begin{equation}
 P_{\mathrm{drive}}^{\mathrm{corr}}
 =
 (1-\lambda) P_{\mathrm{drive}}^{\mathrm{raw}}
 + \lambda \, \widetilde P_{\mathrm{drive}},
\end{equation}
avec
\begin{equation}
 \lambda = \texttt{smoothPdriveInterzoneBlend}.
\end{equation}
Le bulk éloigné des interfaces est laissé inchangé.

\subsection{Comportement observé}

Sur le cas \texttt{compressibility}, ce lissage améliore effectivement la régularité du saut interzone de $u_y$, mais ne le fait pas disparaître totalement. Les essais montrent un comportement non monotone :
\begin{itemize}
  \item une largeur de lissage de $1$ cellule améliore légèrement le défaut ;
  \item une largeur de $2$ cellules constitue le meilleur compromis testé ;
  \item une largeur de $3$ cellules sur-lisse le champ et dégrade de nouveau les profils moyens.
\end{itemize}
Le réglage le plus satisfaisant parmi ceux testés est donc apparu être, pour le cas hydrostatique, un lissage de largeur $2$ cellules avec un mélange proche de l'unité. Ce point doit être lu comme une correction ciblée des interfaces, et non comme un nouveau modèle physique autonome.

\section{Bilan des essais de la redistribution par zones}

\subsection{Cas \texttt{compressibility} : cas sévère pour le zonage}

Le cas \texttt{compressibility} a été l'outil principal pour révéler les limites de l'approche par zones. Il combine en effet trois ingrédients défavorables :
\begin{itemize}
  \item un gradient de pression essentiellement vertical ;
  \item une pression virielle très raide, donc sensible aux petites erreurs de densité ;
  \item un \emph{kick} vertical directement proportionnel à $\partial_y P_{\mathrm{drive}}$.
\end{itemize}
Dans ce contexte, la moindre discontinuité de redistribution aux interfaces horizontales du pavage se répercute dans le gradient de pression et apparaît comme une cassure de $u_y$ ou comme un pic dans les diagnostics hydrostatiques.

Les tests ont montré que :
\begin{enumerate}[label=\alph*)]
  \item avec de grandes zones ($25\times25$ sur $50\times50$), la seconde passe décalée atténue partiellement les défauts, mais ne supprime pas la cicatrice interzone ;
  \item avec de petites zones ($10\times10$), la situation se dégrade fortement : la combinaison du pavage de base et du pavage décalé introduit alors des interfaces horizontales tous les $5$ points environ, ce qui revient pratiquement à traiter l'hydrostatique par sous-couches horizontales ;
  \item le lissage de $P_{\mathrm{drive}}$ améliore la situation, mais n'efface pas totalement la stratification résiduelle.
\end{enumerate}
Il faut donc considérer \texttt{compressibility} comme un cas très sévère pour l'approche par zones : il est excellent pour révéler les défauts, mais probablement trop exigeant pour juger à lui seul de la pertinence de la méthode dans des cas d'usage plus courants.

\subsection{Cas Poiseuille : robustesse sur un écoulement cisaillé}

Le cas Poiseuille s'est révélé beaucoup plus favorable. Les profils moyens $U_x(y)$ obtenus avec la redistribution par zones restent très proches de la référence, et les quantités intégrées telles que le débit $Q$, la viscosité identifiée par courbure ou la viscosité de paroi restent dans la variabilité naturelle du bruit MPCD. Des tests répétés sur plusieurs graines aléatoires ont montré que l'écart entre référence et redistribution par zones, avec lissage interzone, est du même ordre que la dispersion statistique observée entre deux réalisations indépendantes de la méthode MPCD.

Ce résultat est important : il indique que, pour un écoulement de type canal où le \emph{kick} de pression n'est pas dominé par une hydrostatique très raide, la redistribution par zones ne modifie pas de manière manifeste les observables principales du problème.

\subsection{Cas piston : compression quasi incompressible en volume actif variable}

Le cas piston est lui aussi nettement plus rassurant que \texttt{compressibility}. La cinématique imposée par le piston, la compression géométrique, la densité active requise et les pressions moyennes du modèle restent très proches entre la référence et la version par zones. Les diagnostics de pression active, de module effectif et de redistribution évoluent dans la même enveloppe, ce qui suggère que la compression quasi incompressible elle-même n'est pas disqualifiée par le zonage.

En pratique, le cas piston confirme que le principal verrou de l'approche par zones n'est pas la seule présence d'une pression virielle, mais bien la combinaison particulière ``gradient hydrostatique vertical fort + interfaces horizontales répétées + dérivation de $P_{\mathrm{drive}}$'' qui caractérise le cas \texttt{compressibility}.

\section{Discussion sur l'intérêt de l'approche par zones pour une parallélisation future}

\subsection{Pourquoi l'approche par zones reste intéressante}

L'approche par zones a été conçue pour préparer une future parallélisation de l'opérateur de redistribution, sans toucher à l'advection, aux collisions, aux conditions limites, à la fermeture liquide ou à la réorientation interfaciale. De ce point de vue, elle remplit déjà plusieurs objectifs :
\begin{itemize}
  \item elle isole clairement l'opérateur de redistribution dans une branche contrôlée ;
  \item elle permet de comparer directement une redistribution globale de référence et une redistribution locale ``simili-parallèle'' ;
  \item elle fournit déjà les paramètres géométriques essentiels d'un futur découpage parallèle : taille de sous-domaines, familles de zones, diagnostics de passe.
\end{itemize}
Pour les cas usuels sans surface libre tels que \texttt{poiseuille}, \texttt{piston} et vraisemblablement \texttt{diffusion}, les résultats obtenus jusqu'à présent suggèrent que le coût de développement peut être justifié : les écarts avec la référence semblent du même ordre que la variabilité statistique naturelle de la méthode dans plusieurs observables pertinentes.

\subsection{Limites actuelles}

Trois limites majeures demeurent.

Premièrement, la compensation entre zones reste indirecte : une zone ne peut pas, au cours d'une même passe, être alimentée ou déchargée par une autre zone ; seule la seconde passe décalée offre une nouvelle occasion de corriger les défauts dans un autre pavage.

Deuxièmement, la seconde passe actuelle garde la même orientation géométrique que la première ; elle est simplement décalée d'une demi-zone. Pour des problèmes à fort gradient vertical, cela peut donc préserver, voire multiplier, les interfaces les plus nuisibles. Une perspective naturelle serait d'explorer une seconde passe réellement orientée différemment, par exemple sous forme de bandes verticales plutôt que de blocs rectangulaires de même orientation.

Troisièmement, le noyau multi-zones ne supporte actuellement que les cas sans surface libre et sans mouillage. Les géométries à interface libre (ellipse, barrage, mouillage explicite) exigeraient de réintroduire dans le noyau zones la logique bulk/interface déjà présente dans la redistribution de référence, ce qui constitue un chantier distinct.

\subsection{Jugement provisoire}

Le bilan actuel conduit à la lecture suivante :
\begin{itemize}
  \item l'approche par zones n'est pas encore suffisamment robuste pour servir de référence générale sur des cas hydrostatiques sévères de type \texttt{compressibility} ;
  \item en revanche, elle paraît suffisamment prometteuse sur des cas d'écoulement et de compression plus usuels pour justifier la poursuite du développement ;
  \item la suite la plus raisonnable n'est pas de raffiner indéfiniment le cas \texttt{compressibility}, mais d'évaluer l'intérêt du zonage sur les cas sans surface libre qui correspondent aux usages visés du code.
\end{itemize}


\section{Méthodes de projection incompressible Q6/Q9}

La chaîne la plus récente du dépôt ne repose plus sur la redistribution dure comme opérateur principal d'incompressibilité. Elle introduit une famille de méthodes particules--grille dont l'objectif est de conserver le caractère mésoscopique SRC/MPCD tout en supprimant les modes compressifs les plus nuisibles. Les deux briques essentielles sont :
\begin{itemize}
  \item Q6, qui projette le champ de vitesse cellulaire sur un champ discret quasi solénoïdal ;
  \item Q9, qui complète Q6 par une projection basse fréquence du flux de masse $Nu$, afin de réduire les grandes structures cohérentes de densité.
\end{itemize}

Ces méthodes sont moins intrusives que la redistribution : elles ne déplacent pas directement les particules d'une cellule vers une autre pour imposer une occupation cible. Elles modifient les vitesses particulaires après dépôt sur grille, résolution d'un problème de projection et interpolation de la correction vers les particules. Cette structure est aussi plus favorable à une future parallélisation.

\subsection{Dépôt particules--grille}

À chaque pas, après l'étape SRC/MPCD classique, les particules sont déposées sur une grille cartésienne. Pour une cellule $c$, on définit l'occupation
\begin{equation}
  N_c = \#\{i\in c\},
\end{equation}
la densité de nombre discrète
\begin{equation}
  n_c = \frac{N_c}{V_c},
\end{equation}
et la vitesse moyenne cellulaire
\begin{equation}
  \bm{u}_c = \frac{1}{N_c}\sum_{i\in c}\bm{v}_i,
  \qquad N_c>0.
\end{equation}
Dans les cellules vides, la vitesse est laissée nulle ou reconstruite par les conventions de l'opérateur de projection. Pour les géométries avec obstacles ou régions solides, un masque fluide/solide peut neutraliser les cellules non fluides dans les projections.

La grandeur de densité manipulée dans Q9 est généralement l'occupation $N_c$ plutôt que la masse volumique au sens strict. Lorsque $m=1$ et $V_c$ est constant, utiliser $N_c$, $n_c$ ou $\rho_c$ ne change que des facteurs de normalisation. Dans les expressions ci-dessous, on écrit donc indifféremment le flux de masse discret sous la forme
\begin{equation}
  \bm{M}_c = N_c\bm{u}_c
  \qquad \text{ou} \qquad
  \bm{M}_c = \rho_c\bm{u}_c,
\end{equation}
selon la convention de normalisation retenue.

\subsection{Projection de vitesse Q6 : contrainte $\nabla\cdot u\simeq 0$}

Q6 applique une projection de Helmholtz--Hodge discrète au champ cellulaire $\bm{u}^{\star}$ issu de l'étape SRC/MPCD. On cherche un potentiel scalaire $\phi$ tel que
\begin{equation}
  \bm{u}^{Q6} = \bm{u}^{\star} - \nabla_h \phi
\end{equation}
vérifie
\begin{equation}
  \nabla_h\cdot \bm{u}^{Q6} \simeq 0.
\end{equation}
On obtient donc l'équation de Poisson discrète
\begin{equation}
  \Delta_h \phi = \nabla_h\cdot \bm{u}^{\star}.
\end{equation}

Selon la géométrie, cette équation est résolue par une variante périodique FFT ou par un opérateur périodique en $x$ et borné en $y$. Le cas canal utilise une périodicité longitudinale et une condition compatible avec l'absence de flux normal aux parois. Le cas Taylor--Green utilise une projection entièrement périodique.

La correction de vitesse cellulaire est
\begin{equation}
  \delta\bm{u}_c^{Q6}=\bm{u}_c^{Q6}-\bm{u}_c^{\star},
\end{equation}
puis elle est interpolée vers les particules :
\begin{equation}
  \bm{v}_i \leftarrow \bm{v}_i + \lambda_u\,\delta\bm{u}^{Q6}(\bm{x}_i),
\end{equation}
où $\lambda_u=\verb|projectionStrength|$ est le paramètre de force de projection. Dans les validations de référence, on utilise généralement
\begin{equation}
  \lambda_u=1.
\end{equation}

Q6 doit être interprétée comme une correction cinématique. Elle supprime efficacement la divergence reconstruite de $\bm{u}$, mais elle ne garantit pas à elle seule que les modes basse fréquence de densité ne se développent pas, car la densité particulaire reste fluctuante.


\subsection{Opérateur elliptique \texttt{general\_bc} et fixation du mode constant}

La projection Q6 repose sur la résolution d'une équation de Poisson discrète ou, plus généralement, d'un problème elliptique compatible avec les conditions aux limites. La version récente utilise un opérateur \texttt{general\_bc} destiné à traiter dans un même cadre les directions périodiques, les murs et les géométries piston. Pour une correction de pression scalaire \(\phi\), on écrit schématiquement
\begin{equation}
  \Delta_h \phi = \frac{1}{\Delta t}\nabla_h\cdot\bm u^\star,
  \qquad
  \bm u^{n+1}=\bm u^\star-\Delta t\,\nabla_h\phi.
\end{equation}

Lorsque le problème possède des conditions de Neumann ou des directions périodiques, le mode constant de \(\phi\) n'est pas déterminé. La correction récente impose explicitement un \emph{gauge-fix} du mode constant nul. Cette opération ne change pas le gradient \(\nabla_h\phi\), donc ne change pas la correction de vitesse, mais elle rend le solveur bien posé numériquement et évite les singularités ou mauvais conditionnements artificiels lors des runs piston et canal.

\subsection{Projection du flux de masse Q9 : contrainte sur $\nabla\cdot(Nu)$}

Q9 ajoute à Q6 une correction du flux de masse. La motivation est que la contrainte
\begin{equation}
  \nabla\cdot \bm{u}=0
\end{equation}
reste insuffisante dans une méthode particulaire à occupation fluctuante : le transport de population est contrôlé par
\begin{equation}
  \nabla\cdot (N\bm{u}),
\end{equation}
qui peut rester compressif même si $\nabla\cdot\bm{u}$ est faible.

On définit donc le flux discret
\begin{equation}
  \bm{M}_c^{\star}=N_c\bm{u}_c^{Q6}.
\end{equation}
Q9 cherche une correction de type gradient
\begin{equation}
  \bm{M}_c^{Q9}=\bm{M}_c^{\star}-\nabla_h \psi_c
\end{equation}
telle que la divergence du flux relaxe les modes de densité visés :
\begin{equation}
  \nabla_h\cdot \bm{M}^{Q9}=S_N.
\end{equation}
Le second membre $S_N$ n'est pas une contrainte de densité locale dure. Il est construit à partir du défaut d'occupation filtré basse fréquence. Si
\begin{equation}
  \delta N_c = N_c-\gamma,
\end{equation}
et si $\mathcal{P}_{\mathrm{low}k}$ désigne le filtre spectral retenu, alors
\begin{equation}
  \delta N_c^{\mathrm{low}k}=\mathcal{P}_{\mathrm{low}k}(\delta N)_c.
\end{equation}
La cible de relaxation peut s'écrire, à un facteur de normalisation absorbé dans le paramètre numérique,
\begin{equation}
  S_{N,c} \simeq -\beta_N\,\delta N_c^{\mathrm{low}k},
\end{equation}
où
\begin{equation}
  \beta_N = \verb|massFluxDensityRelaxationBeta|.
\end{equation}
On obtient l'équation de Poisson
\begin{equation}
  \Delta_h\psi = \nabla_h\cdot \bm{M}^{\star}-S_N.
\end{equation}

Une fois le flux corrigé, la correction de vitesse associée est reconstruite par
\begin{equation}
  \delta\bm{u}_c^{Q9}=\frac{\bm{M}_c^{Q9}-\bm{M}_c^{\star}}{\max(N_c,N_{\min})},
\end{equation}
puis appliquée aux particules :
\begin{equation}
  \bm{v}_i\leftarrow \bm{v}_i + \lambda_M\,\delta\bm{u}^{Q9}(\bm{x}_i),
\end{equation}
avec
\begin{equation}
  \lambda_M=\verb|massFluxProjectionStrength|.
\end{equation}

Le choix central de Q9 est de ne corriger que les grandes longueurs d'onde de densité. Le filtre courant est
\begin{equation}
  \verb|massFluxTargetFilter|=\verb|lowpass_fft|,
  \qquad
  \verb|massFluxLowKMaxIndex|=2.
\end{equation}
Les modes locaux, assimilables au bruit cellulaire naturel du MPCD, ne sont donc pas forcés brutalement.

\subsection{Projection finale de nettoyage}

Après interpolation de la correction Q9 vers les particules, la reconstruction cellulaire peut réintroduire une divergence faible mais non nulle, en particulier lorsque l'interpolation et le dépôt ne sont pas des opérateurs adjoints exacts. Un nettoyage optionnel est donc appliqué :
\begin{equation}
  \verb|massFluxFinalVelocityProjectionCleanup|=\verb|true|.
\end{equation}
Il consiste à appliquer une nouvelle projection de vitesse de force réduite
\begin{equation}
  \lambda_{\mathrm{clean}}=\verb|massFluxFinalVelocityProjectionStrength|.
\end{equation}
Le compromis retenu dans les cas validés est généralement
\begin{equation}
  \lambda_{\mathrm{clean}}=0.5.
\end{equation}
Un nettoyage trop fort rend $\nabla\cdot u$ très petit, mais peut réduire l'effet du contrôle low-$k$ ; un nettoyage trop faible préserve mieux la correction de flux de masse, mais laisse une divergence reconstruite plus élevée.

\subsection{Paramètres de référence}

Les paramètres Q9 de référence, avant adaptation à certaines géométries, sont :
\begin{center}
\begin{tabular}{ll}
\toprule
Paramètre & Valeur typique \\
\midrule
\verb|projectionStrength| & $1.0$ \\
\verb|massFluxProjectionMode| & \verb|relax_to_uniform_lowk| \\
\verb|massFluxProjectionStrength| & $1.0$ \\
\verb|massFluxDensityRelaxationBeta| & $2\times10^{-3}$ \\
\verb|massFluxApplyAfterVelocityProjection| & \verb|true| \\
\verb|massFluxTargetFilter| & \verb|lowpass_fft| \\
\verb|massFluxLowKMaxIndex| & $2$ \\
\verb|massFluxFinalVelocityProjectionCleanup| & \verb|true| \\
\verb|massFluxFinalVelocityProjectionStrength| & $0.5$ \\
\bottomrule
\end{tabular}
\end{center}

Pour la marche descendante, le paramètre de relaxation a dû être réduit à
\begin{equation}
  \texttt{massFluxDensityRelaxationBeta}=5\times10^{-4},
\end{equation}
afin d'éviter une sur-correction au voisinage du coin solide. Cette adaptation est importante : Q9 possède un paramètre de raideur effectif, mais sa valeur optimale dépend de la géométrie et du niveau de bruit mécanique.

\section{Diagnostics utilisés pour Q6/Q9}

\subsection{Densité locale et densité basse fréquence}

La densité locale est suivie par l'écart-type de l'occupation cellulaire,
\begin{equation}
  \mathrm{std}(N)=\sqrt{\langle (N_c-\langle N\rangle)^2\rangle_c},
\end{equation}
et par la fraction de cellules hors bande
\begin{equation}
  N_c\notin \gamma(1\pm \varepsilon_N),
\end{equation}
où $\varepsilon_N$ vaut typiquement $0.2$ dans les validations.

Le diagnostic le plus spécifique de Q9 est l'énergie basse fréquence du défaut relatif de densité. En posant
\begin{equation}
  r_c = \frac{N_c}{\gamma}-1,
\end{equation}
et en notant $\widehat r_{k_x,k_y}$ sa transformée de Fourier discrète, on mesure une quantité du type
\begin{equation}
  E_{\mathrm{low}k}=\sum_{0<|k_x|,|k_y|\le K}\abs{\widehat r_{k_x,k_y}}^2,
\end{equation}
avec $K=\verb|lowKMaxIndex|$. La normalisation exacte dépend du script de diagnostic, mais les comparaisons Q6/Q9 utilisent toujours la même convention.

\subsection{Divergence et flux de masse}

Après projection, on reconstruit
\begin{equation}
  D_u = \left\|\nabla_h\cdot\bm{u}\right\|_{\mathrm{RMS}}
\end{equation}
et, pour Q9,
\begin{equation}
  D_M = \left\|\nabla_h\cdot(N\bm{u})-S_N\right\|_{\mathrm{RMS}}.
\end{equation}
On distingue le résidu de projection du flux de masse, calculé immédiatement au niveau grille, et le défaut reconstruit après retour aux particules. Cette distinction est essentielle : le solveur de projection peut être très précis sur grille, tandis que l'interpolation particulaire réintroduit une erreur effective.

\subsection{Diagnostics de structures}

Pour les cas vorticitaires, on utilise la vorticité discrète
\begin{equation}
  \omega = \partial_x u_y - \partial_y u_x
\end{equation}
et l'enstrophie moyenne
\begin{equation}
  \mathcal{E}_{\omega}=\langle \omega^2\rangle.
\end{equation}

Dans le cas Taylor--Green, on projette le champ de vitesse sur le mode analytique initial afin d'obtenir une amplitude de mode, une énergie modale et une cohérence. Dans le cas marche, on suit la longueur de recirculation, l'aire de recirculation, le RMS de vorticité dans la couche de cisaillement et des sondes aval en $\omega$ et $u_y$.

\subsection{Pression effective au piston}

Pour l'évaluation d'une équation d'état effective, trois pressions sont distinguées :
\begin{align}
  P_{\mathrm{kin}} &= \langle n_c k_B T_c\rangle,\\
  P_{\mathrm{wall}} &= \frac{\Delta p_y^{\mathrm{top}}}{\Delta t_{\mathrm{sample}} L_x},\\
  P_{\mathrm{excess}} &= P_{\mathrm{wall}}-P_{\mathrm{kin}}.
\end{align}
La pression mécanique au mur est estimée par l'impulsion cumulée échangée avec le piston pendant la fenêtre de mesure retenue. C'est cette pression qui permet de définir une raideur mécanique effective.

En posant
\begin{equation}
  \chi=\frac{\rho}{\rho_0}=\frac{y_{\mathrm{top},0}}{y_{\mathrm{top}}},
\end{equation}
on écrit localement
\begin{equation}
  P(\chi)=P_0+K_{\mathrm{eff}}\ln \chi,
\end{equation}
ou, pour de faibles compressions,
\begin{equation}
  \Delta P \simeq K_{\mathrm{eff}}(\chi-1).
\end{equation}


\section{Validations Q6/Q9 corrigées}

\subsection{Statut des anciennes validations}

Les résultats Q6/Q9 historiques du rapport précédent sont conservés comme jalons méthodologiques, mais ils ne doivent plus être utilisés seuls pour calibrer la méthode. Trois éléments ont changé : le thermostat cellulaire a été corrigé, les conditions limites de canal ont été renforcées par particules virtuelles \texttt{wallVP}, et plusieurs routines MATLAB ont été réorganisées/vectorisées. Les comparaisons qualitatives restent utiles : Q9 réduit les modes compressifs basse fréquence et ne détruit pas les structures moyennes principales. En revanche, les valeurs de viscosité effective, de raideur piston et de pression de paroi doivent être réévaluées avec les scripts récents.

\subsection{Taylor--Green périodique : préservation des structures lisses}

Le cas Taylor--Green reste le test le plus propre du bulk périodique, car il élimine les complications de paroi. La configuration de référence récente utilise typiquement un domaine \(64\times64\), une occupation \(\gamma=20\), \(\Delta t=10^{-3}\), \(k_BT=10^{-2}\) et une amplitude initiale de l'ordre de \(0.1\). Les runs forcés avec Q6/Q9 complet donnent une viscosité effective bulk de l'ordre de
\begin{equation}
  \nu_{\mathrm{eff,TG}} \simeq 2.0\times10^{-2},
\end{equation}
avec une température stabilisée par le thermostat corrigé et une diminution des composantes basse fréquence du champ de densité.

Le rôle de ce test n'est pas d'identifier une loi d'état, mais de vérifier que les corrections Q6/Q9 ne projettent pas abusivement la dynamique vorticitaire. Le résultat important est que l'amplitude du mode Taylor--Green, l'énergie modale et la cohérence de phase restent exploitables : Q9 agit sur les défauts compressifs cohérents plutôt que comme un filtre destructeur de vorticité.

\subsection{Poiseuille : viscosité effective, murs et particules virtuelles}

Poiseuille reste le test principal de canal. Dans les versions anciennes, les profils pouvaient être bons en apparence mais dépendaient fortement du traitement de paroi et du thermostat. La version récente impose donc une lecture plus stricte : un run Poiseuille valide doit fournir simultanément un profil moyen parabolique, une température contrôlée, une densité basse fréquence réduite par Q9, une estimation de viscosité robuste après exclusion des cellules proches des murs, et une réponse pariétale compatible avec la pression cinétique.

L'introduction des particules virtuelles de paroi modifie sensiblement l'interprétation. Le \emph{bounce-back} géométrique impose une contrainte d'imperméabilité et une forme de non-glissement, mais il ne suffit pas toujours à reconstruire correctement la contrainte collisionnelle au mur. Le modèle \texttt{wallVP} ajoute une population solide effective dans les cellules coupées ; il améliore le couplage fluide--paroi, mais impose de diagnostiquer la pression par impulsions cumulées, et non par échantillons instantanés isolés.

\subsection{Poiseuille multi-résolutions : viscosité brute et correction en taille de grille}

Une difficulté apparue dans les runs Poiseuille récents est la variation apparente de la viscosité identifiée lorsque l'on modifie le nombre de mailles dans la direction transverse. À paramètres MPCD identiques, la viscosité mesurée par fit quadratique du profil peut changer lorsque l'on diminue ou augmente \(N_y\), même lorsque le profil reste visuellement parabolique. Cette observation ne doit pas être interprétée immédiatement comme une modification physique du fluide par Q6/Q9 : une part importante vient de la manière dont le profil discret est ajusté sur une hauteur de canal représentée par un nombre fini de cellules.

Le point de départ est le profil continu de Poiseuille,
\begin{equation}
  u_x(y)=\frac{f_x}{2\nu} y(H-y),
\end{equation}
où \(H\) est la hauteur utile du canal. Si le fit est effectué dans une coordonnée cellulaire normalisée de manière implicite par \(N_y\), la courbure discrète porte un facteur d'échelle en \(H^2\), donc en première approximation en \(N_y^2\) lorsque \(L_y\) et le domaine physique de référence sont conservés. Une viscosité brute extraite directement de la courbure discrète peut donc suivre une loi apparente
\begin{equation}
  \nu_{\mathrm{raw}}(N_y) \propto \frac{1}{N_y^2},
\end{equation}
si l'on ne réintroduit pas correctement l'échelle de hauteur du canal dans l'identification.

La correction utilisée dans les campagnes récentes consiste à rapporter les viscosités à une résolution transverse de référence \(N_{y,\mathrm{ref}}=64\), selon
\begin{equation}
  \nu_{\mathrm{corr}}
  =
  \nu_{\mathrm{raw}}
  \left(\frac{N_y}{N_{y,\mathrm{ref}}}\right)^2.
  \label{eq:nu-corr-ny}
\end{equation}
Cette écriture est volontairement pratique : elle ne prétend pas remplacer une dérivation analytique complète de la viscosité MPCD avec parois, mais elle corrige le biais de lecture introduit par la discrétisation transverse du profil.

Les campagnes comparant plusieurs grilles, notamment des canaux de type \(32\times48\), \(48\times64\) et \(64\times96\), ont montré que les viscosités brutes issues du fit pouvaient varier fortement avec \(N_y\), alors que les viscosités corrigées par \eqref{eq:nu-corr-ny} se regroupaient autour de la viscosité bulk de référence obtenue sur Taylor--Green, typiquement
\begin{equation}
  \nu_{\mathrm{eff,TG}} \simeq 0.02.
\end{equation}
Dans ces runs, les valeurs corrigées se plaçaient typiquement dans une bande de l'ordre de \(0.020\)--\(0.023\) pour Q9 et de \(0.020\)--\(0.021\) pour le cas classic, hors cas manifestement insuffisamment convergé ou mal ajusté. Ce regroupement est un résultat important : la diminution de \(N_y\) ne signifie pas nécessairement que la méthode change de fluide ; elle peut surtout changer l'échelle effective de lecture du profil de Poiseuille.

La conséquence méthodologique est que tout tableau Poiseuille doit désormais distinguer explicitement :
\begin{enumerate}[label=\roman*)]
  \item la viscosité brute de fit, \(\nu_{\mathrm{raw}}\), directement extraite du profil discret ;
  \item la résolution transverse \(N_y\), la hauteur utile \(H\), et le nombre de cellules exclues près des murs ;
  \item la viscosité corrigée \(\nu_{\mathrm{corr}}\), calculée avec \eqref{eq:nu-corr-ny} ou, plus généralement, avec le facteur \((H/H_{\mathrm{ref}})^2\) si la hauteur physique effective varie ;
  \item la comparaison avec la référence bulk périodique Taylor--Green, qui sert de contrôle indépendant des effets de paroi.
\end{enumerate}

Cette correction ne dispense pas de traiter correctement les conditions limites. Les particules virtuelles \texttt{wallVP}, l'exclusion de quelques cellules adjacentes aux murs et la durée du moyennage restent nécessaires pour obtenir un profil exploitable. La correction en \(N_y^2\) répond à un problème différent : elle évite de confondre une erreur d'échelle géométrique dans l'identification de \(\nu\) avec une vraie modification de transport due à Q6/Q9. Les prochains rapports de validation devront donc toujours présenter les deux colonnes \(\nu_{\mathrm{raw}}\) et \(\nu_{\mathrm{corr}}\), et ne comparer à Taylor--Green que la viscosité corrigée.

\subsection{Piston staircase : EOS bulk Q6/Q9 avec fermeture virielle}

Le piston staircase récent constitue le résultat bulk le plus net. Dans le run \texttt{Q9\_FULL} à \(N_p=20480\), grille \(32\times32\), \(\gamma=20\), \(\Delta t=10^{-3}\), \(k_BT=10^{-2}\), compression finale de \(5\%\), la densité physique finale est \(2.1558\times10^4\), la température cellule finale vaut \(0.01\), la densité basse fréquence finale est de l'ordre de \(10^{-7}\), et le limiteur viriel n'est pas déclenché dans le régime doux testé.

Le résultat essentiel est l'EOS bulk. Les fits staircase donnent
\begin{align}
  \frac{\mathrm d P_{\mathrm{kin}}}{\mathrm d\rho} &\simeq 0.010534, & R^2&=1,\\
  \frac{\mathrm d P_{\mathrm{vir}}}{\mathrm d\rho} &\simeq 0.05, & R^2&=1,\\
  \frac{\mathrm d P_{\mathrm{tot}}}{\mathrm d\rho} &\simeq 0.060534, & R^2&=1,\\
  \frac{\mathrm d P_{\mathrm{drive}}}{\mathrm d\rho} &\simeq 0.060534, & R^2&=1.
\end{align}
La valeur finale correspond à
\begin{equation}
  P_{\mathrm{kin}}\simeq 226.95,
  \qquad
  P_{\mathrm{vir}}\simeq 53.89,
  \qquad
  P_{\mathrm{tot}}\simeq 280.84.
\end{equation}
Ce résultat valide la fermeture virielle bulk : le code reproduit exactement la pente imposée par \(K_{\mathrm{virial}}\), et la pression totale suit la somme \(P_{\mathrm{kin}}+P_{\mathrm{vir}}\). Le résidu global de quantité de mouvement du kick viriel reste au niveau du bruit machine et le limiteur de kick n'est pas activé dans ce cas.

\subsection{Validation séparée du diagnostic mécanique \texttt{wallVP}}

Le même run piston a montré que les pressions instantanées de paroi pouvaient être algébriquement cohérentes mais physiquement inutilisables sur des plateaux courts : les résidus de reconstruction \(P_{\mathrm{total}}=P_{\mathrm{impact}}+P_{\mathrm{VP}}\) étaient nuls, mais les moyennes par plateau oscillaient fortement. Cette observation a conduit à isoler un test plus simple : gaz à l'équilibre entre deux murs, sans piston et sans terme viriel.

Dans ce test d'équilibre \(32\times32\), \(N_p=20480\), \(k_BT=0.01\), la pression idéale nominale vaut \(\rho k_BT=204.8\), tandis que la pression cinétique bulk mesurée vaut environ \(215.6\). Sans \texttt{wallVP}, les pressions cumulées haut/bas valent environ \(208.48\) et \(207.53\). Avec \texttt{wallVP}, les contributions impact seules montent à environ \(226.92\) et \(223.67\), les contributions VP signées valent environ \(-12.28\) et \(-9.07\), et les totaux cumulés reviennent à \(214.64\) et \(214.60\). L'écart relatif top/bottom tombe à environ \(2\times10^{-4}\).

Ce test valide la convention de signe, la normalisation \(\Delta p/(\Delta t L_x)\) et l'usage des impulsions cumulées. Il montre aussi que la contribution VP ne doit pas être rendue positive par construction : elle peut corriger à la baisse une contribution impact trop grande. La pression physique de paroi est le total signé impact+VP, moyenné sur une fenêtre longue.

\subsection{Retour au staircase : prudence sur les fenêtres courtes}

Les premiers retours staircase après correction du diagnostic montrent encore des moyennes de fenêtre instables lorsque les plateaux sont courts. Ce résultat n'invalide pas \texttt{wallVP} : le test d'équilibre a précisément montré que les pressions instantanées peuvent fluctuer entre des valeurs de plusieurs milliers d'unités autour d'une moyenne de l'ordre de \(2.1\times10^2\). La règle de validation actuelle est donc la suivante : comparer \(P_{\mathrm{tot,bulk}}\) aux moyennes de pression pariétale obtenues par différence d'impulsions cumulées sur des fenêtres longues, et non aux moyennes d'échantillons instantanés peu nombreux.


\section{Équation d'état bulk et statut mécanique de paroi}

\subsection{EOS bulk actuellement validée}

L'état validé côté piston doit être formulé de manière plus stricte que dans les versions précédentes du rapport. La loi d'état actuellement fiable est une loi bulk mesurée sur les cellules actives, pas encore une loi mécanique pariétale complète. Elle repose sur
\begin{equation}
  P_{\mathrm{vir}} = K_{\mathrm{virial}}(\rho-\rho_{\mathrm{EOSRef}}),
  \qquad
  P_{\mathrm{tot}} = P_{\mathrm{kin}}+P_{\mathrm{vir}},
  \qquad
  P_{\mathrm{drive}}=P_{\mathrm{tot}}.
\end{equation}

Le protocole staircase impose une suite de compressions à nombre de particules constant et hauteur active variable. Pour chaque plateau, on élimine le transitoire de déplacement, puis on moyenne les grandeurs bulk sur la partie finale du plateau. Les fits récents montrent que \(P_{\mathrm{kin}}\), \(P_{\mathrm{vir}}\), \(P_{\mathrm{tot}}\) et \(P_{\mathrm{drive}}\) sont linéaires en densité moyenne avec \(R^2\simeq1\). La pente de \(P_{\mathrm{vir}}\) retrouve directement \(K_{\mathrm{virial}}\), tandis que la pente de \(P_{\mathrm{tot}}\) est la somme de la pente cinétique effective et de la pente virielle imposée.

Cette validation est importante parce qu'elle dissocie deux questions. La première est : le modèle bulk produit-il une réponse pression--densité contrôlée ? La réponse est oui dans le régime doux testé. La seconde est : la pression mécanique mesurée sur le piston restitue-t-elle immédiatement cette même loi ? La réponse est plus prudente : pas sur des plateaux courts, car l'estimateur de paroi est fortement bruité.

\subsection{Kick viriel et conservation de la quantité de mouvement}

Le kick viriel actif est appliqué sous une forme faible et limitée. La version actuelle diagnostique le RMS de correction de vitesse, le rapport à la vitesse thermique, la fraction de cellules limitées, le maximum autorisé et le résidu de quantité de mouvement après correction globale. Dans le run de référence récent, le limiteur n'est pas déclenché, le RMS du kick reste de l'ordre de \(10^{-5}\), et le résidu global de moment reste proche de \(10^{-15}\). Ces valeurs indiquent que le terme viriel modifie la réponse EOS sans introduire de dérive globale détectable de la quantité de mouvement.

\subsection{Statut du diagnostic piston mécanique}

Le diagnostic mécanique piston n'est pas encore considéré comme validé pour les runs staircase courts. Le point acquis est algébrique : les champs de diagnostic reconstruisent correctement
\begin{equation}
  P_{\mathrm{wall,total}} = P_{\mathrm{wall,impact}} + P_{\mathrm{wall,VP}},
\end{equation}
avec un résidu de cohérence nul. Le point encore en cours de validation est statistique et physique : sur une fenêtre de plateau courte, le total instantané peut changer de signe ou prendre des valeurs trop grandes, alors que le bulk suit parfaitement l'EOS. La correction adoptée consiste à stocker les impulsions cumulées et à utiliser des moyennes de fenêtre longues.

La lecture correcte des prochains runs staircase sera donc :
\begin{enumerate}[label=\arabic*)]
  \item vérifier d'abord \(P_{\mathrm{kin}}\), \(P_{\mathrm{vir}}\), \(P_{\mathrm{tot}}\) et \(P_{\mathrm{drive}}\) dans le bulk ;
  \item vérifier que top et bottom donnent des pressions mécaniques cumulées compatibles à compression nulle ;
  \item n'interpréter les plateaux comprimés qu'après des temps de maintien longs ;
  \item comparer séparément la pression mécanique cumulée à \(P_{\mathrm{kin}}\) et à \(P_{\mathrm{tot}}\), afin de déterminer si le terme viriel bulk est effectivement transmis mécaniquement au piston ou seulement utilisé comme pression motrice interne.
\end{enumerate}

\section{Parallélisation et vectorisation MATLAB}

\subsection{Objectif}

La motivation initiale de Q6/Q9 est aussi algorithmique : remplacer une redistribution particulaire dure, difficile à paralléliser proprement, par des opérations plus régulières. Les briques principales de Q6/Q9 sont favorables à une exécution parallèle ou vectorisée : dépôt particules--grille, sommes par cellule, reconstruction de champs, gradients, projections elliptique/FFT, interpolation grille--particules et diagnostics.

\subsection{Ce qui est déjà favorable dans MATLAB}

La version MATLAB récente a commencé à réduire le coût de plusieurs boucles critiques par des opérations groupées. Le thermostat corrigé et les reconstructions cellulaires utilisent des index de cellule et des accumulations vectorisées lorsque c'est possible. Les diagnostics de densité, température, flux de masse, moyenne de profil, pression bulk et impulsions cumulées sont eux aussi structurés comme opérations sur tableaux.

Le gain n'est pas seulement un gain de temps. Une écriture vectorisée rend les diagnostics plus reproductibles : le même champ cellulaire sert à calculer la température, la pression, la divergence, les corrections Q6/Q9 et les métriques de conservation. Cela limite les divergences entre scripts de validation.

\subsection{Limites de MATLAB et direction C++/OpenMP}

MATLAB reste toutefois limitant pour les longs runs piston, les sweeps multi-graines et les géométries plus sévères. Les opérations particulaires avec conditions limites complexes, particules virtuelles, interpolation et stockage d'historiques deviennent coûteuses. La stratégie rationnelle est donc : conserver MATLAB comme banc de validation et de prototypage, puis porter les briques stabilisées vers C++/OpenMP, voire MPI/OpenMP ou GPU.

Les briques prioritaires pour un portage sont :
\begin{enumerate}[label=\arabic*)]
  \item collision SRC/MPCD avec thermostat corrigé ;
  \item conditions limites bounce-back, piston et \texttt{wallVP} ;
  \item projection Q6 avec opérateur \texttt{general\_bc} ;
  \item correction Q9 basse fréquence du flux de masse ;
  \item diagnostics cumulés de moment et de pression de paroi ;
  \item sorties légères pour profils moyens, vorticité, densité basse fréquence et EOS piston.
\end{enumerate}


\section{Synthèse actuelle et conséquences pour la parallélisation}

Le bilan actuel est plus clair que dans la version précédente du rapport. La chaîne historique par redistribution dure a démontré qu'il était possible d'imposer une faible compressibilité effective, mais elle est intrusive et difficile à paralléliser. La chaîne Q6/Q9 corrigée est moins intrusive : elle agit sur les champs projetés, réduit les modes compressifs cohérents et préserve mieux la structure particulaire naturelle du MPCD.

Les points considérés comme validés sont :
\begin{itemize}
  \item la projection Q6 avec opérateur compatible conditions limites et mode constant fixé ;
  \item la correction Q9 basse fréquence du flux de masse ;
  \item le thermostat cellulaire corrigé ;
  \item la stabilité thermique dans les cas de référence ;
  \item la réduction des modes basse fréquence de densité ;
  \item la préservation qualitative et quantitative des structures Taylor--Green ;
  \item la cohérence bulk de l'EOS piston virielle avec \(R^2\simeq1\) ;
  \item la cohérence du diagnostic \texttt{wallVP} en gaz d'équilibre lorsqu'il est évalué par impulsions cumulées.
\end{itemize}

Les points encore ouverts sont :
\begin{itemize}
  \item la statistique des pressions mécaniques de paroi dans les runs staircase comprimés ;
  \item la calibration multi-graines de la viscosité effective en canal avec \texttt{wallVP}, en reportant systématiquement \(\nu_{\mathrm{raw}}\) et \(\nu_{\mathrm{corr}}=\nu_{\mathrm{raw}}(N_y/N_{y,\mathrm{ref}})^2\) ;
  \item la transmission mécanique exacte de la pression virielle bulk vers le piston ;
  \item la robustesse sur écoulements plus complexes, notamment von Kármán et géométries à interfaces libres ;
  \item le coût des longs runs MATLAB, qui justifie un portage parallèle.
\end{itemize}

La prochaine étape scientifique n'est donc pas de modifier à nouveau le bulk Q6/Q9, mais de consolider la statistique de paroi : équilibre gaz long, Poiseuille long avec \texttt{wallVP}, puis staircase piston à plateaux longs. La prochaine étape numérique est de factoriser les briques validées pour préparer un portage C++/OpenMP.




\section{Addendum C++/OpenMP : inlet/outlet, solides immergés et cellules pauvres}

Les développements ultérieurs au socle MATLAB ont été poursuivis dans le dépôt C++/OpenMP \texttt{SRC\_incompressible\_openMP}. Ils ont d'abord concerné les conditions limites inlet/outlet Q6/Q9, puis le traitement des solides immergés fixes. Cette section met à jour le statut méthodologique du rapport : les solides immergés sont désormais séparés du problème plus général des cellules pauvres.

\subsection{Conditions limites inlet/outlet retenues}

La politique inlet/outlet retenue pour les cas ouverts est la suivante. L'inlet est alimenté par un réservoir de type \texttt{hard\_cell\_density}, de manière à imposer une occupation cible. L'outlet peut fonctionner en mode \texttt{balanced\_flux}, \texttt{neumann} ou \texttt{hybrid}. Le mode \texttt{hybrid} est devenu le choix nominal pour les cas d'injection plus difficiles : il part d'un flux local de type Neumann, ajoute éventuellement un faible mélange vers un profil équilibré, et applique une correction de balance globale faible au niveau de l'outlet.

Un point important des validations inlet/outlet est l'abandon des exclusions volumétriques près des frontières ouvertes comme stratégie nominale. Les réglages validés sont
\[
  \texttt{q9OpenBoundaryExclusionCells}=0,\qquad
  \texttt{virialOpenBoundaryExclusionCells}=0.
\]
Les exclusions non nulles créaient une interface active/inactive se comportant comme une paroi numérique près de l'outlet. La rampe d'inlet doit en outre être appliquée de façon cohérente à l'injection particulaire, au flux Q6 et au flux massique Q9.

Les apertures segmentées inlet/outlet sont maintenant interprétées comme des géométries physiques, par exemple une fente ou une buse, et non comme un correctif numérique destiné à stabiliser un canal canonique. Les diagnostics Q9 spatialisés sont écrits sous forme de sidecars binaires \texttt{q9\_diagnostics\_step\_*.q9bin}, lisibles par les outils MATLAB de visualisation.

\subsection{Formulation cellule/face pour les solides immergés}

Les essais précédents avaient montré que désactiver Q9 dans un halo autour d'un obstacle produisait une couche morte : les cellules adjacentes au solide devenaient incapables de produire une correction de pression ou de flux massique, précisément là où les accumulations de densité apparaissent. La formulation retenue est donc une formulation cellule/face.

On distingue :
\begin{itemize}
  \item les cellules solides, inactives pour Q6, Q9 et le viriel ;
  \item les cellules fluides, actives même lorsqu'elles sont adjacentes au solide ;
  \item les cellules coupées, actives si leur fraction fluide dépasse le seuil fixé ;
  \item les faces fluide--fluide, ouvertes au flux correctif ;
  \item les faces fluide--solide ou coupées fermées, sur lesquelles le flux normal correctif est imposé nul.
\end{itemize}

Le réglage nominal est donc
\[
  \texttt{projectionImmersedSolidMaskEnable}=\texttt{true},\quad
  \texttt{projectionImmersedSolidCloseCutFaces}=\texttt{true},\quad
  \texttt{q9ImmersedSolidHaloCells}=0.
\]
Le paramètre de halo Q9 reste éventuellement disponible comme mode conservateur ou de débogage, mais il ne doit plus être considéré comme une stratégie physique.

\subsection{Fermeture géométrique des flux Q6/Q9 aux faces solides}

Les premiers tests avec cercle immergé ont révélé que l'opérateur elliptique fermait bien les faces solides, mais que le diagnostic de fuite Q9 pouvait rester non nul lorsque la face fermée était portée par une cellule propriétaire solide. Une fermeture géométrique explicite a donc été ajoutée après reconstruction des flux :
\[
  F_{\mathrm{proj}}^{Q9}(\text{face solide})=0,\qquad
  F_{\mathrm{corr}}^{Q9}(\text{face solide})=-F_{\mathrm{base}}(\text{face solide}).
\]
La même convention a ensuite été appliquée au diagnostic Q6. Les grandeurs de fuite effectives
\[
  \texttt{q6ImmersedSolidLeakProjectedFluxRms},\qquad
  \texttt{q9ImmersedSolidLeakMassFluxRms}
\]
désignent désormais les flux après fermeture géométrique. Les diagnostics avant fermeture restent disponibles pour vérifier l'amplitude de la correction imposée.

Cette correction a été validée sur cercle périodique, cylindre périodique et cylindre en canal ouvert. Les résultats caractéristiques sont :
\[
  \texttt{q6ImmersedSolidLeakProjectedFluxRms}=0,\qquad
  \texttt{q9ImmersedSolidLeakMassFluxRms}=0,\qquad
  \texttt{q9ImmersedHaloExcludedCells}=0.
\]
Les cellules adjacentes et coupées restent actives, et les sidecars Q9 confirment l'absence de couronne inactive autour du solide.

\subsection{Viriel près des solides}

Le kick viriel a été modifié pour ne pas pousser artificiellement dans le solide. Dans les cellules fluides adjacentes au solide, le viriel reste actif, mais la composante normale du kick orientée vers une face solide fermée est neutralisée. Les composantes tangentielles sont conservées. Les diagnostics ajoutés comptent notamment les cellules adjacentes actives, le nombre de composantes normales supprimées et les normes du kick supprimé.

La validation périodique avec cylindre montre que le viriel reste actif au voisinage du solide sans créer de fuite de flux normal. La formulation actuelle est donc :
\begin{itemize}
  \item viriel actif en cellule fluide adjacente ;
  \item composante normale entrante dans le solide supprimée ;
  \item pas de halo viriel mort autour du solide ;
  \item cellules coupées traitées par fraction fluide.
\end{itemize}

\subsection{Validation visuelle des solides immergés}

La validation s'est faite par paliers afin de limiter les coûts :
\begin{enumerate}
  \item smoke-tests courts avec cercle ou cylindre périodique ;
  \item visualisation du démarrage sur quelques centaines de pas ;
  \item run périodique plus long avec forçage modéré ;
  \item premier cas inlet/outlet avec cylindre fixe.
\end{enumerate}

Les visualisations périodiques montrent un sillage cohérent, un déficit de vitesse en aval du cylindre et l'absence de couche inactive autour de l'obstacle. En inlet/outlet à vitesse modérée, les diagnostics restent propres : fuite Q6/Q9 solide nulle, halo Q9 nul, cellules adjacentes actives, limiteur Q9 peu ou pas sollicité. À ce stade, le traitement des solides immergés fixes peut donc être considéré comme validé pour les besoins du développement suivant.

\section{Limite identifiée : cellules pauvres et défauts de population}

Les essais plus sévères ont mis en évidence une difficulté indépendante de la présence d'un solide. Dans des cas ouverts ou avec poche initiale de cellules vides, certaines cellules pauvres peuvent se comporter comme des défauts persistants de population. L'effet est particulièrement visible lorsque l'on introduit une poche vide \(6\times6\) dans un canal ouvert sans solide.

Le témoin SRC classique convecte et diffuse progressivement la poche. En revanche, les variantes Q9/viriel enrichies par une garde locale de population ont tendance à faire l'inverse : la poche devient un défaut eulérien actif, dévie le champ de vitesse et se comporte comme un obstacle mou ou un puits fixé à la grille.

\subsection{Essais de garde locale de population}

Une première idée a consisté à écrire
\[
  P_{\mathrm{vir}} = P_{\mathrm{bulk}}(\rho) + P_{\mathrm{guard}}\left(\frac{N}{N_\star}\right),
  \qquad
  N_\star=\gamma_{\mathrm{ref}}\,\theta_f,
\]
où \(\theta_f\) est la fraction fluide de la cellule. Le terme \(P_{\mathrm{guard}}\) devait être nul dans une bande \([r_{\mathrm{low}},r_{\mathrm{high}}]\), positif dans les cellules surpeuplées et négatif dans les cellules sous-peuplées. L'objectif était de créer une pression locale dense--vers--pauvre sans redistribution dure de particules.

Cette approche a été testée sous plusieurs formes :
\begin{itemize}
  \item mode \texttt{pressure\_gradient}, injectant \(P_{\mathrm{guard}}\) dans le gradient cell-centered existant ;
  \item mode \texttt{face\_donor}, appliquant le kick depuis la cellule donneuse vers la cellule receveuse face par face ;
  \item mode \texttt{face\_donor\_surplus}, n'autorisant que les cellules réellement surpeuplées à donner ;
  \item mode \texttt{face\_donor\_capacity}, autorisant les cellules normales à donner jusqu'à un plancher de population.
\end{itemize}

Ces essais ont clarifié le mécanisme mais n'ont pas fourni de solution robuste. Le mode \texttt{face\_donor} simple peut transformer une cellule vide en puits aspirant ses voisines. Le mode \texttt{face\_donor\_surplus} évite cette aspiration mais laisse la poche vide agir comme un obstacle mou. Le mode \texttt{face\_donor\_capacity}, intermédiaire, n'a pas supprimé ce comportement et a même accentué l'effet dans certains cas.

\subsection{Interprétation}

Le résultat principal est le suivant : une cellule pauvre ne doit pas devenir une source de pression attractive. Une accumulation peut légitimement produire une pression répulsive, mais un vide de particules ne doit pas imposer un champ de vitesse eulérien fort. Dans la méthode particulaire, un trou local doit plutôt être convecté, diffusé et progressivement réparé par le transport SRC, l'inlet ou les voisins ; le transformer en objet actif le rend plus persistant.

La difficulté vient du couplage entre trois faits :
\begin{enumerate}
  \item Q9 corrige un flux de masse mais ne crée pas de porteurs dans une cellule vide ;
  \item les sécurités low-mass protègent les cellules peu peuplées, mais peuvent les rendre passives ;
  \item une garde attractive locale peut fixer un défaut de population sur la grille.
\end{enumerate}

La conclusion provisoire est donc de ne plus poursuivre les gardes locales attractives dense--vers--pauvre sous cette forme. Le problème des cellules pauvres doit être repris dans une branche dédiée, séparée de la branche des solides immergés, en partant du constat que le SRC classique transporte correctement les poches pauvres alors que certaines corrections Q9/viriel les eulérianisent.

\subsection{Pistes pour la branche suivante}

Les pistes les plus raisonnables pour la branche suivante sont :
\begin{itemize}
  \item rendre les cellules low-mass transparentes pour Q9, afin qu'elles ne deviennent pas des murs numériques ;
  \item éviter tout terme de pression négative attaché aux cellules pauvres ;
  \item tester une pression de garde positive seulement, agissant uniquement sur les accumulations ;
  \item séparer correction des accumulations et transport des trous ;
  \item examiner si la logique \texttt{ramp\_floor} de Q9 doit être remplacée par un mode \texttt{transparent\_floor} ;
  \item conserver la chaîne classic comme témoin systématique pour toute poche de cellules pauvres.
\end{itemize}

Cette difficulté ne remet pas en cause la validation des solides immergés. Elle identifie un problème plus général de robustesse locale de population dans les cas ouverts ou fortement hétérogènes.


\section{Branche resampling pondéré : réponse aux cellules pauvres et domaine fluide évolutif}

Les derniers développements MATLAB ont repris le problème des cellules pauvres sous un angle différent. Les essais Q9/viriel/guard avaient pour point commun d'essayer de réparer un défaut de population par une correction du champ de vitesse, du flux de masse ou de la pression. Or les tests avec poches vides ont montré que cette approche devient instable conceptuellement lorsque la cellule contient trop peu de particules : le champ hydrodynamique local n'est plus fiable, mais la correction eulérienne continue à l'interpréter comme un degré de liberté robuste. La poche pauvre peut alors devenir un objet actif de grille, au lieu d'être simplement transportée et diffusée par la dynamique particulaire.

La nouvelle branche, créée à partir de la base MATLAB Q6/Q9 intégrée, introduit une formulation \textbf{SRC/MPCD pondérée avec resampling local conservatif}. L'objectif n'est plus de forcer une cellule pauvre par une pression attractive, mais de maintenir un support particulaire local suffisant pour que les champs $\rho$, $u$ et $\rho u$ restent représentatifs. Cette section décrit le principe, les briques algorithmiques ajoutées et les premiers résultats obtenus.

\subsection{Difficulté fondamentale des cellules pauvres}

La projection Q6 porte sur la contrainte cinématique
\begin{equation}
  \nabla\cdot u \simeq 0,
\end{equation}
et la correction Q9 sur une contrainte de flux massique
\begin{equation}
  \nabla\cdot(\rho u) \simeq 0.
\end{equation}
Ces contraintes ont un sens hydrodynamique fort lorsque les champs sont reconstruits à partir d'une population statistiquement suffisante. En revanche, dans une cellule vide ou presque vide, la vitesse moyenne
\begin{equation}
  u_c = \frac{\sum_{i\in c} m_i v_i}{\sum_{i\in c} m_i}
\end{equation}
devient mal définie ou dominée par très peu de particules. La difficulté n'est donc pas seulement numérique : l'opérateur agit alors sur une variable hydrodynamique qui n'est plus correctement représentée.

Les essais de garde locale ont clarifié la contradiction. Une correction assez faible laisse le trou persister. Une correction assez forte pour remplir rapidement le trou agit comme une force interne artificielle. Les variantes \texttt{pressure\_gradient}, \texttt{face\_donor}, \texttt{face\_donor\_surplus} et \texttt{face\_donor\_capacity} ont donc produit soit des oscillations, soit une aspiration des voisins, soit un obstacle mou. Le principe retenu est dès lors le suivant :
\begin{quote}
Une cellule pauvre ne doit pas être corrigée par une attraction eulérienne. Elle doit être traitée comme une insuffisance locale du support particulaire.
\end{quote}

\subsection{Passage à des particules pondérées}

La branche resampling introduit un état particulaire de la forme
\begin{equation}
  \{x_i, v_i, m_i\}_{i=1}^{N_{\mathrm{cap}}},
\end{equation}
où $m_i$ est désormais une masse particulaire variable et $N_{\mathrm{cap}}$ la capacité mémoire préallouée. Le cas $m_i=m_0$ pour toutes les particules doit reproduire le comportement SRC/MPCD classique pondéré de façon neutre.

Les dépôts particules--grille sont remplacés par des dépôts massiques :
\begin{equation}
  M_c = \sum_{i\in c} m_i,
  \qquad
  P_c = \sum_{i\in c} m_i v_i,
  \qquad
  U_c = \frac{P_c}{M_c}.
\end{equation}
La collision SRC pondérée se fait autour du centre de masse massique $U_c$ :
\begin{equation}
  v_i' = U_c + R(\theta_c)(v_i-U_c),
\end{equation}
ce qui conserve exactement le moment pondéré de la cellule tant que la rotation est appliquée aux particules fluides de cette cellule.

La projection Q6 reste une projection de vitesse. Elle est appliquée au champ $U_c$ reconstruit par dépôt pondéré, puis la correction $\delta U$ est interpolée aux particules fluides. Une correction globale exacte de moment pondéré est appliquée lorsque nécessaire, afin d'éviter une dérive du moment total.

\subsection{Pool préalloué et rôles de particules}

Pour rester compatible avec un futur portage OpenMP, l'ajout de particules n'est pas traité par une croissance dynamique non bornée des tableaux. Le code utilise un pool préalloué :
\begin{equation}
  N_{\mathrm{cap}} = \alpha_{\mathrm{cap}}\,\gamma N_xN_y,
\end{equation}
avec typiquement $\alpha_{\mathrm{cap}}=2$. Les particules possèdent un rôle logique :
\begin{itemize}
  \item rôle 0 : particule inactive, disponible dans le pool libre ;
  \item rôle 1 : particule fluide réelle, participant au dépôt, à la collision, au thermostat, à Q6 et au remap ;
  \item rôle 2 : particule latente, stockée dans le domaine mais ne représentant pas encore du fluide.
\end{itemize}
La fonction de masque actif ne sélectionne que les particules de rôle fluide. Les particules inactives ou latentes sont exclues des dépôts, de la collision, du thermostat, de la projection et des diagnostics hydrodynamiques.

Cette distinction est importante pour les futurs solides mobiles, fronts libres ou domaines partiellement remplis : une particule peut exister dans les tableaux sans être un porteur de masse fluide.

\subsection{Extraction, insertion et remap local conservatif}

La population cible d'une cellule mouillée est définie par trois seuils :
\begin{equation}
  N_{\min},\qquad N_{\mathrm{target}}=\gamma,
  \qquad N_{\max}.
\end{equation}
Si $N_c<N_{\min}$, des particules sont insérées depuis le pool libre jusqu'à $N_{\mathrm{target}}$. Si $N_c>N_{\max}$, des particules sont extraites et replacées dans le pool libre. Le choix nominal des particules extraites est \texttt{closest\_to\_cell\_mean}, c'est-à-dire que l'on retire d'abord les particules dont la vitesse est la plus proche de la vitesse moyenne locale, afin de minimiser la perturbation du moment avant remap.

Le remap local impose ensuite, dans chaque cellule mouillée non vide,
\begin{equation}
  \sum_{i\in c} m_i^{\mathrm{new}} = M_c^{\mathrm{target}},
  \qquad
  \sum_{i\in c} m_i^{\mathrm{new}} v_i^{\mathrm{new}} = M_c^{\mathrm{target}} U_c^{\mathrm{target}}.
\end{equation}
Dans le mode simple \texttt{scale\_preserve\_velocity}, les masses d'une cellule sont d'abord multipliées par un facteur commun. Ce mode impose exactement $M_c^{\mathrm{target}}$ et conserve la vitesse moyenne pondérée lorsque la cible de moment est la vitesse avant remap. Un solveur \texttt{min\_change} existe également, mais le mode de mise à l'échelle s'est montré plus robuste pour les premiers tests.

L'ordre nominal après SRC/Q6 est le suivant :
\begin{enumerate}[label=\arabic*)]
  \item reconstruire le champ cible $U_c$ avant édition de population ;
  \item extraire les particules des cellules surpeuplées ;
  \item insérer les particules dans les cellules sous-peuplées ;
  \item remapper les masses pour imposer masse et moment cibles ;
  \item appliquer éventuellement un thermostat pondéré post-remap.
\end{enumerate}
Ainsi, l'extraction et l'insertion modifient le support particulaire, tandis que le remap restaure les grandeurs hydrodynamiques de cellule.

\subsection{Masque mouillé/non mouillé et domaine fluide évolutif}

La branche ajoute un masque de cellule
\begin{equation}
  \chi_c \in \{0,1\},
\end{equation}
où $\chi_c=1$ désigne une cellule mouillée, donc active pour le remap de masse, et $\chi_c=0$ une cellule non mouillée. La masse cible devient alors conceptuellement
\begin{equation}
  M_c^{\mathrm{target}} = \chi_c\,\rho_0 V_c.
\end{equation}
Dans la version actuelle, les cellules non mouillées ne sont pas forcées vers une masse cible et ne déclenchent ni insertion, ni extraction, ni remap de masse fluide. Le masque peut être fixe ou mis à jour par hystérésis de masse, ce qui prépare le traitement de domaines fluides évolutifs, d'interfaces libres ou de solides mobiles.

La couche latente complète ce formalisme. Un domaine peut être initialisé avec des particules latentes de masse et vitesse nulles ; une cellule source peut ensuite activer des particules fluides de masse non nulle. Le remap n'est appliqué qu'aux cellules mouillées, de sorte qu'un domaine vide n'est plus rempli artificiellement au premier pas.

Cette structure sera également utile pour les solides immergés : les cellules solides ou non mouillées auront $M_c^{\mathrm{target}}=0$, tandis que les cellules fluides ou partiellement fluides auront une cible ajustée à leur fraction fluide. Les particules virtuelles de paroi restent, elles, des contributions collisionnelles temporaires et non des particules fluides transportées.

\subsection{Sécurité de masse et renormalisation conservatrice}

Les premiers tests hétérogènes ont montré que le remap masse/moment peut satisfaire les contraintes algébriques en créant des particules très lourdes ou très légères lorsque l'état initial est extrêmement déséquilibré. Une sécurité conditionnelle a donc été ajoutée : si les masses candidates sortent d'une bande admissible, par exemple
\begin{equation}
  m_i \notin [0.25m_0,4m_0],
\end{equation}
la cellule bascule vers un mode \texttt{uniform\_mass\_velocity\_shift}. Les masses sont remises à
\begin{equation}
  m_i = \frac{M_c^{\mathrm{target}}}{N_c},
\end{equation}
puis les vitesses sont décalées uniformément pour restaurer le moment cible.

Une renormalisation plus douce a ensuite été conçue pour limiter la dérive long terme de la distribution des masses. Pour chaque cellule, on mémorise
\begin{equation}
  M_c,
  \qquad
  U_c,
  \qquad
  E_{\mathrm{th},c}=\frac12\sum_{i\in c}m_i\abs{v_i-U_c}^2.
\end{equation}
Les masses sont relaxées vers $M_c/N_c$, puis renormalisées pour conserver $M_c$. Les vitesses sont ensuite décalées pour restaurer $U_c$, et les fluctuations sont rescalées pour restaurer $E_{\mathrm{th},c}$. La correction préserve donc localement
\begin{equation}
  M_c,
  \qquad
  U_c,
  \qquad
  E_{\mathrm{th},c},
\end{equation}
tout en réduisant l'hétérogénéité des masses.

Deux validations récentes ont précisé le statut de cette correction. D'abord, un smoke Poiseuille court de 500 pas avec une renormalisation unique au pas 300 a montré une chute nette de la dispersion globale des masses, sans rupture des autres contraintes. À pas communs, la dispersion relative finale passe typiquement de
\begin{equation}
  \frac{\sigma_m}{\langle m\rangle}\simeq 0.079
  \quad\text{sans renormalisation}
  \qquad \text{à} \qquad
  \frac{\sigma_m}{\langle m\rangle}\simeq 0.059
  \quad\text{avec renormalisation},
\end{equation}
à 500 pas. Au pas de renormalisation, les résidus de moment et d'énergie thermique locale restent au niveau du roundoff, ce qui confirme que la correction ne se comporte pas comme une force volumique cachée.

Ensuite, le protocole long de référence a consisté à appliquer une seule renormalisation au pas 3000 d'un Poiseuille Q6+resampling wallVP de 5000 pas, avec tous les autres paramètres inchangés. Le résultat final est
\begin{equation}
  \frac{\sigma_m}{\langle m\rangle}=0.200,
  \qquad
  M_{\mathrm{relRMS}}\simeq 10^{-16},
  \qquad
  k_BT\simeq 0.00987,
\end{equation}
contre environ $0.234$ de dispersion des masses dans le run sans renormalisation. Le profil Poiseuille reste excellent, avec $R^2\simeq0.9978$ et $\nu_{\mathrm{eff}}\simeq18.0$, à comparer à $R^2\simeq0.9979$ et $\nu_{\mathrm{eff}}\simeq17.65$ sans renormalisation. La correction est donc validée comme sécurité numérique ponctuelle ou conditionnelle : elle réduit l'hétérogénéité des masses sans dégrader visiblement la masse cellulaire, la température ou le profil hydrodynamique.

\subsection{Tests fonctionnels de population}

Le premier test conçu pour faire tourner explicitement la boucle complète est le cas \texttt{paired\_pockets}. Il crée simultanément une poche vide et une poche surpeuplée en déplaçant les particules de l'une vers l'autre. Un résultat typique est :
\begin{itemize}
  \item au pas 0 : $N_{\min}=0$, $N_{\max}=40$, $M_{\mathrm{relRMS}}\simeq 0.221$ ;
  \item au pas 1 : 500 particules extraites, 500 insérées, $N\in[19,21]$, $N_{\mathrm{active}}$ constant ;
  \item après 300 pas : $N\in[17,23]$, $M_{\mathrm{relRMS}}\simeq 10^{-16}$ et les compteurs d'extraction/insertion restent fixés à 500.
\end{itemize}
Ce test démontre que le pool fonctionne comme un tampon de recyclage et non comme une source non bornée de particules.

Un second protocole impose une hétérogénéité distribuée de population. Les cas avec écart-type initial modéré sont ramenés vers la bande $N\in[14,26]$ avec $M_{\mathrm{relRMS}}$ au roundoff. Pour des hétérogénéités extrêmes, la sécurité de masse empêche les masses aberrantes observées avant son introduction. Dans un cas sévère représentatif, on observe après 300 pas
\begin{equation}
  N\in[14,26],
  \qquad
  M_{\mathrm{relRMS}}\simeq 10^{-16},
  \qquad
  \frac{\sigma_m}{\langle m\rangle}\simeq 0.13,
\end{equation}
avec des masses restant dans une plage contrôlée.

Le test d'injection/remplissage part d'un domaine entièrement latent, puis active progressivement des particules dans une cellule source. Dans un run accéléré, le nombre de cellules mouillées croît de 0 à 2048, la fraction mouillée atteint 1 vers le pas 105, et la masse des cellules mouillées reste contrôlée avec $M_{\mathrm{relRMS}}\simeq 10^{-14}$. Ce test valide la séparation entre particules latentes et particules fluides. Il montre aussi qu'un remplissage brutal est exigeant pour la distribution des masses : dans le run analysé, $\sigma_m/\langle m\rangle$ atteint environ 0.40 avec des masses proches de la borne de sécurité supérieure. Ce cas doit donc être lu comme un stress-test du mécanisme, non comme un réglage physique nominal.

\subsection{Validation hydrodynamique préliminaire}

Une fois la boucle de population fonctionnelle, la priorité est de vérifier que la méthode produit encore un comportement fluide crédible. Deux cas ont été retenus : Taylor--Green forcé en domaine périodique et Poiseuille forcé avec parois et particules virtuelles de paroi.

Sur Taylor--Green, le cas Q6 resamplé donne typiquement
\begin{equation}
  \mathrm{rmsDiv}\simeq 4\times 10^{-17},
  \qquad
  M_{\mathrm{relRMS}}\simeq 10^{-16},
  \qquad
  k_BT\simeq 0.0113,
\end{equation}
avec $N\in[14,26]$. La cohérence du mode Taylor--Green est meilleure que dans le témoin classic, avec une valeur représentative d'environ 0.893 contre 0.747, et l'enstrophie bruitée est réduite. Ce résultat suggère que Q6+resampling ne détruit pas les structures cohérentes de bulk.

Sur Poiseuille, les premiers essais sans particules virtuelles de paroi n'étaient pas exploitables : le profil était trop influencé par le glissement pariétal. Les particules virtuelles \texttt{wallVP} ont donc été réintroduites en version pondérée. Elles sont ajoutées seulement comme contributions agrégées à la collision près des murs :
\begin{equation}
  U_c^{\mathrm{coll}}=
  \frac{P_{\mathrm{real}}+P_{\mathrm{virt}}}{M_{\mathrm{real}}+M_{\mathrm{virt}}},
\end{equation}
et seules les particules fluides réelles sont mises à jour. Les particules virtuelles ne participent ni au transport, ni au pool, ni au remap.

Le run Poiseuille wallVP à 5000 pas fournit une comparaison directe :
\begin{center}
\begin{tabular}{lccc}
\toprule
Cas & $R^2$ parabolique & $\nu_{\mathrm{eff}}$ & $M_{\mathrm{relRMS}}$ \\
\midrule
SRC classic wallVP pur & $0.9959$ & $16.99$ & $0.15$ \\
Q6+resampling wallVP & $0.9979$ & $17.65$ & $\simeq10^{-16}$ \\
Q6+resampling wallVP + renorm. ponctuelle & $0.9978$ & $18.01$ & $\simeq10^{-16}$ \\
\bottomrule
\end{tabular}
\end{center}
Le résultat important n'est pas la valeur absolue de la viscosité, mais la proximité entre le témoin SRC wallVP pur et le cas Q6+resampling. La méthode resamplée conserve donc la physique de transport du cas wallVP de référence, avec un écart d'environ 4 \%, tout en supprimant les fluctuations de masse cellulaire. Le coût du run sans renormalisation est une dispersion des masses particulaires d'environ
\begin{equation}
  \frac{\sigma_m}{\langle m\rangle}\simeq 0.23.
\end{equation}
La renormalisation ponctuelle au pas 3000 abaisse cette valeur à environ $0.20$ à 5000 pas, tout en conservant un fit parabolique de qualité comparable et une viscosité effective très proche. Cette comparaison montre que la dispersion des masses peut être contrôlée sans disqualifier la dynamique Poiseuille.

\subsection{Interprétation et statut de la branche resampling}

La nouvelle méthode ne doit pas être interprétée comme une simple variante de Q9. Elle remplace le traitement des trous de population par une régularisation du support particulaire. Le principe devient :
\begin{equation}
\begin{aligned}
  &\text{support particulaire contrôlé}
  + \text{masse cellulaire imposée}
  + \text{projection Q6} \\
  &\hspace{3em}\Longrightarrow
  \text{fluide quasi incompressible pondéré}.
\end{aligned}
\end{equation}
Q9 peut rester utile comme diagnostic ou correction basse fréquence, mais il n'est plus nécessairement le mécanisme principal de contrôle de la densité lorsque le remap impose $M_c$ cellule par cellule.

Le statut actuel est donc le suivant :
\begin{itemize}
  \item la boucle extraction--insertion--remap est fonctionnelle ;
  \item les cellules non mouillées et les particules latentes permettent un domaine fluide évolutif ;
  \item les tests de population démontrent la réparation rapide des défauts de support sans force attractive ;
  \item Taylor--Green valide qualitativement le comportement bulk ;
  \item Poiseuille wallVP valide un premier comportement canal/paroi avec viscosité effective proche du SRC de référence ;
  \item la distribution des masses reste le principal point de vigilance, mais la renormalisation locale préservant $M,u,E$ est maintenant validée comme sécurité ponctuelle ou conditionnelle.
\end{itemize}
La branche resampling devient donc la voie la plus prometteuse pour traiter les cellules pauvres, les fronts de remplissage, les solides immergés et, à terme, les solides mobiles ou interfaces libres. Elle doit encore être validée statistiquement sur des runs plus longs et plusieurs graines, puis portée vers C++/OpenMP une fois la formulation stabilisée.




\section{Portage C++/OpenMP du resampling pondéré}

Les développements précédents ont stabilisé la formulation MATLAB du resampling pondéré. La dernière phase du projet a consisté à en traduire le cœur dans le dépôt \texttt{SRC\_incompressible\_openMP}, sans transformer le code C++ en générateur de cas physiques. La règle de workflow retenue est volontairement stricte : les états initiaux \texttt{.smpcd} sont préparés par des scripts MATLAB \texttt{prepare\_\dots.m} dans \texttt{init/}, les scripts \texttt{.sh} écrivent les fichiers \texttt{params*.kv} et lancent les exécutables OpenMP dans \texttt{runs/}, puis les scripts MATLAB \texttt{analyze\_\dots.m} post-traitent les sorties. Le noyau C++ ne fabrique donc pas les cas physiques : il lit un état initial, applique la dynamique SRC/Q6/resampling et écrit des diagnostics et dumps reproductibles.

Cette séparation est importante. Elle évite de multiplier les chaînes de génération, conserve MATLAB comme référence pour les scénarios de validation et permet de comparer directement les mêmes configurations entre prototype et code parallèle. Le C++ reste focalisé sur les opérations à fort coût : streaming, conditions limites, collision, projection elliptique, resampling, dépôt pondéré, diagnostics de masse et écriture d'états.

\subsection{Infrastructure OpenMP ajoutée}

Le portage a été introduit par jalons courts et testables. L'état particulaire OpenMP a été étendu avec :
\begin{itemize}
  \item une masse individuelle $m_i$ ;
  \item un rôle logique \texttt{Fluid}, \texttt{Inactive} ou \texttt{Latent} ;
  \item une distinction explicite entre \texttt{role} et \texttt{type}, afin de conserver la possibilité future d'espèces fluides multiples ;
  \item un format \texttt{.smpcd} V2 persistant les masses et les rôles ;
  \item un dépôt pondéré ne comptant que les particules \texttt{Fluid} réelles ;
  \item des diagnostics de masse, population, pool, remap, sécurité de masse et garde de population.
\end{itemize}

La séparation rôle/type est une différence structurante par rapport à certains prototypes MATLAB. Le rôle indique si une particule participe à la dynamique fluide courante. Le type est réservé à la nature physique future de la particule : espèce, matériau, traceur ou famille de paroi. Ainsi, une particule peut être présente dans les tableaux mais inactive, latente ou fluide. Les particules \texttt{Inactive} et \texttt{Latent} sont exclues du dépôt, de la collision, de Q6, du thermostat, du remap, du \texttt{mass guard} et des diagnostics hydrodynamiques.

Le dépôt pondéré utilisé par la collision et les diagnostics OpenMP est
\begin{equation}
  N_c = \sum_{i\in c,\;r_i=\mathrm{Fluid}} 1,\qquad
  M_c = \sum_{i\in c,\;r_i=\mathrm{Fluid}} m_i,
\end{equation}
\begin{equation}
  P_c = \sum_{i\in c,\;r_i=\mathrm{Fluid}} m_i v_i,
  \qquad
  U_c = \frac{P_c}{M_c}.
\end{equation}
Les particules virtuelles de paroi \texttt{wallVP} sont conservées dans leur rôle collisionnel : elles modifient le centre de collision près des murs, mais ne sont pas versées dans le pool, ne sont pas remappées et ne sont pas interprétées comme des particules fluides transportées.

\subsection{Correction d'une première divergence avec MATLAB}

Une étape importante du portage a été la mise en évidence d'une divergence conceptuelle. La première version OpenMP du resampling était surtout \emph{mass-driven} : les cellules pauvres ou riches étaient détectées à partir de seuils sur $M_c$. Cette version pouvait donc imposer
\begin{equation}
  M_c \simeq M_{\mathrm{target}}
\end{equation}
tout en laissant une cellule mouillée avec un nombre très faible de particules actives. Le champ de masse devenait alors parfait, mais le support particulaire $N_c$ pouvait rester insuffisant. Ce défaut est apparu dans le cas d'obstacle en canal ouvert : la masse était ramenée au roundoff, mais le domaine aval se dépeuplait localement et les masses individuelles compensaient trop fortement le manque de particules.

Ce comportement n'était pas conforme au prototype MATLAB. Dans la méthode de référence, le resampling est d'abord un contrôle de population :
\begin{equation}
  N_c < N_{\min} \quad\Longrightarrow\quad \text{repeupler la cellule},
\end{equation}
\begin{equation}
  N_c > N_{\max} \quad\Longrightarrow\quad \text{dépeupler la cellule},
\end{equation}
puis seulement ensuite une correction de masse/moment :
\begin{equation}
  M_c \rightarrow M_{\mathrm{target}},\qquad
  P_c \rightarrow M_{\mathrm{target}}U_c.
\end{equation}

Le patch OpenMP 0140 a donc réintroduit le garde de population $N_{\min}/N_{\mathrm{target}}/N_{\max}$ comme mécanisme principal. Pour $\gamma=20$, les scripts de validation utilisent typiquement
\begin{equation}
  N_{\min}=14,\qquad N_{\mathrm{target}}=20,\qquad N_{\max}=26.
\end{equation}
La logique devient conforme au développement MATLAB : on maintient d'abord un support particulaire suffisant, puis on utilise les masses pondérées pour fermer la contrainte de masse.

\subsection{Variante hybride efficace pour OpenMP}

Le mécanisme retenu est une variante hybride pensée pour éviter un appariement global donneur--receveur coûteux. Trois familles d'opérations sont distinguées.

Premièrement, si une cellule mouillée est sous-peuplée mais non vide, le code effectue un \emph{split local} de particules fluides lourdes. Une particule de masse $m$ et vitesse $v$ est remplacée par deux particules fluides de masse $m/2$ et vitesse $v$ :
\begin{equation}
  (m,v)\quad\longrightarrow\quad
  \left(\frac{m}{2},v\right)+
  \left(\frac{m}{2},v\right).
\end{equation}
Cette opération augmente $N_c$ tout en conservant localement masse, moment et énergie cinétique de translation. Elle utilise un slot \texttt{Inactive} du pool, mais ne crée pas de masse nouvelle.

Deuxièmement, si une cellule est surpeuplée, le code applique une extraction/merge locale qui renvoie des slots vers le pool inactif et évite que $N_c$ ne dérive trop au-dessus de $N_{\max}$. Cette opération reste locale à la cellule ; elle évite les tris globaux et les recherches donneur--receveur à coût quadratique.

Troisièmement, les mécanismes d'activation latente, d'extraction et d'insertion préexistants restent disponibles pour les cas de front mouillé, d'inlet, de cellules vides ou de pool déjà préparé. La logique générale est donc :
\begin{enumerate}[label=\arabic*)]
  \item dépôt pondéré des particules \texttt{Fluid} ;
  \item classification wet/dry et classification $N_c$ ;
  \item garde de population $N_{\min}/N_{\mathrm{target}}/N_{\max}$ ;
  \item activation latente, extraction et insertion si activées ;
  \item remap masse/moment tous les $K$ pas ;
  \item \texttt{mass guard} avec la même cadence que le remap de masse ;
  \item renormalisation thermique à chaque pas lorsque le resampling est actif.
\end{enumerate}

Ce choix est plus favorable à OpenMP que le transfert global donneur--receveur. Les décisions sont essentiellement locales à la cellule ; l'accès au pool peut être organisé par listes compactes et préallocations ; le coût principal reste linéaire en nombre de particules et de cellules. Les optimisations déjà intégrées évitent aussi les reconstructions coûteuses de plans de mutation à chaque dépôt et suppriment un diagnostic thermique quadratique qui provoquait un ralentissement massif sur Poiseuille.

\subsection{Fréquence des opérations dans la version OpenMP}

Le portage OpenMP distingue explicitement les fréquences des opérations. Le pas SRC classique et Q6, lorsqu'il est activé, restent exécutés à chaque pas. Le resampling ajoute les opérations suivantes :
\begin{center}
\small
\begin{tabular}{p{0.33\linewidth}p{0.36\linewidth}p{0.22\linewidth}}
\toprule
Opération & Rôle & Fréquence / paramètre \\
\midrule
\texttt{resamplingEnable} & verrou global du chemin mutateur resampling & lu à chaque pas \\
Dépôt pondéré & calcule $N_c,M_c,P_c,U_c$ sur \texttt{Fluid} & chaque pas \\
Classification wet/dry & détermine les cellules mouillées & chaque pas, \texttt{resamplingWetMaskMode} \\
Garde de population & impose la bande $N_{\min}/N_{\mathrm{target}}/N_{\max}$ & chaque pas, \texttt{resamplingPopulationGuardEnable} \\
Activation latente & \texttt{Latent}$\rightarrow$\texttt{Fluid} & chaque pas si activée \\
Extraction/insertion & basculement \texttt{Fluid}$\leftrightarrow$\texttt{Inactive} & chaque pas si activé \\
Remap de masse & impose $M_c\rightarrow M_{\mathrm{target}}$ & tous les $K$ pas, \texttt{resamplingMassRenormalizationPeriod} \\
\texttt{mass guard} & borne $m_i$ dans $[m_{\min},m_{\max}]$ & même cadence que le remap \\
Renormalisation thermique & restaure/cohère $U_c$ et $E_{\mathrm{th},c}$ & chaque pas si activée \\
\bottomrule
\end{tabular}
\end{center}
La convention pour $K$ est :
\begin{equation}
  K=1 : \text{remap à chaque pas},\qquad
  K>1 : \text{remap tous les }K\text{ pas},\qquad
  K=0 : \text{remap désactivé}.
\end{equation}
Pour les validateurs post-0140, le réglage nominal est
\begin{equation}
  K=10,\qquad m_{\min}=0.5,\qquad m_{\max}=2.0.
\end{equation}

\subsection{Différences pratiques entre MATLAB et OpenMP}

La version OpenMP ne doit pas être vue comme une simple copie ligne à ligne du prototype MATLAB. Elle en conserve la logique physique, mais avec des choix adaptés au coût parallèle :
\begin{itemize}
  \item le C++ lit des états \texttt{.smpcd} générés par MATLAB ; il ne génère pas les cas physiques ;
  \item les rôles sont persistés dans le fichier d'état, ce qui rend possible un redémarrage exact avec pool, particules latentes et inactives ;
  \item les opérations locales de garde de population privilégient split/merge cellulaire plutôt qu'un appariement global donneur--receveur ;
  \item la renormalisation de masse est cadencée par un paramètre, alors que la garde de population reste à chaque pas ;
  \item les diagnostics runtime sont plus nombreux que dans MATLAB : masses, population wet, pool, guard, solid leak, divergence Q6, opérations d'extraction/insertion et coût mural ;
  \item les restrictions de sécurité du noyau C++ sont plus fortes : par exemple, le cas \texttt{Q6 + inlet/outlet + solide immergé} n'autorise encore que certains solides rectangulaires masqués, afin d'éviter un désaccord entre le masque solide de projection et la dynamique particulaire.
\end{itemize}

Cette dernière restriction explique pourquoi le vrai cylindre en canal ouvert n'a pas encore été activé avec Q6/resampling. Les cylindres périodiques et les obstacles rectangulaires ouverts sont validables avec le noyau actuel ; le cylindre inlet/outlet demandera un futur chantier C++ spécifique sur les cellules coupées et les faces solide--fluide dans l'opérateur elliptique.

\subsection{Validations OpenMP actuellement obtenues}

Les validations post-portage montrent que les briques essentielles fonctionnent dans plusieurs régimes. Le tableau suivant résume le statut, en mettant l'accent sur le gain par rapport au témoin \texttt{classic}.

\begin{center}
\small
\begin{tabular}{p{0.23\linewidth}p{0.32\linewidth}p{0.34\linewidth}}
\toprule
Cas & Observation principale & Gain par rapport à \texttt{classic} \\
\midrule
Taylor--Green périodique et forcé & Q6 et Q6+resampling conservent le motif cohérent ; le cas forcé maintient une amplitude stationnaire lorsque le forçage est calibré. & Q6 réduit la divergence de grille ; Q6+resampling impose la masse cellulaire au roundoff et maintient une population bornée. \\
Population aléatoire distribuée & Le cas déclenche un désordre initial de population comparable au prototype MATLAB. & Le resampling ramène le support vers la bande $[14,26]$ et supprime l'erreur de masse, alors que \texttt{classic}/Q6 gardent des fluctuations massiques. \\
Poiseuille \texttt{wallVP} & Les trois cas donnent un profil parabolique cohérent ; Q6+resampling ne détruit pas le transport. & La masse cellule est ramenée au roundoff et la population reste bornée. Après optimisation 0132--0133, le surcoût court observé était d'environ $6.4\,\mathrm{s}$ contre $4.4\,\mathrm{s}$ pour Q6 seul sur 1000 pas, soit environ $1.45\times$ au-dessus de Q6. \\
Cylindre périodique immergé & Le solide circulaire périodique reste propre et le sillage reste stable. & Le \texttt{solidLeakMass} représentatif passe d'environ 76 en \texttt{classic} à 22 en Q6 et 18 en Q6+resampling, avec masse cellulaire contrôlée. \\
Backward step & Le cas inlet/outlet avec solide rectangulaire reste crédible ; la recirculation n'est pas encore très marquée, mais l'écoulement ne présente pas de pathologie. & Q6+resampling stabilise la masse et le support sans casser les conditions ouvertes ni le masque solide. \\
Injection/fill depuis domaine inactif & Le domaine part vide/inactif puis se remplit progressivement par l'inlet localisé. & \texttt{classic} et Q6 seuls accumulent localement de très fortes masses ; Q6+resampling propage le front mouillé et maintient la masse des cellules mouillées au niveau cible. \\
Obstacle rectangulaire en canal ouvert & Le cas a révélé le défaut du premier portage : masse correcte mais support particulaire appauvri. & Après 0140, ce cas devient le test décisif du garde $N_{\min}/N_{\mathrm{target}}/N_{\max}$ et doit remplacer les conclusions de la variante mass-driven. \\
\bottomrule
\end{tabular}
\end{center}

Les gains ne sont donc pas seulement quantitatifs sur un indicateur. Ils sont surtout structurels : le cas resamplé maintient simultanément un champ de masse quasi incompressible, un support particulaire suffisant et la possibilité de domaines mouillés évolutifs. C'est ce point qui rend la méthode SRC/MPCD pondérée beaucoup plus puissante que les corrections eulériennes de cellules pauvres testées auparavant.

\subsection{Points de vigilance restants}

Le portage OpenMP est désormais plus conforme à la méthode MATLAB, mais plusieurs points restent ouverts :
\begin{itemize}
  \item la campagne post-0140 doit confirmer que le garde de population supprime le dépeuplement observé dans les canaux ouverts ;
  \item la fréquence optimale du remap de masse $K$ doit être balayée selon les cas ;
  \item les bornes $m_{\min},m_{\max}$ doivent être interprétées comme des sécurités numériques, non comme des paramètres physiques ;
  \item le vrai cylindre inlet/outlet avec Q6 exigera un masque circulaire compatible projection, faces coupées et diagnostics de fuite de flux projeté ;
  \item le coût de Q6 elliptique dans les géométries ouvertes reste un chantier d'optimisation séparé ; l'essai de warm-start a été abandonné pour ne pas polluer les validations physiques.
\end{itemize}

Le prochain objectif scientifique est donc de terminer la revalidation post-0140 sur les cas Taylor--Green, Poiseuille, remplissage, obstacle ouvert et marche, avant de chercher une allée tourbillonnaire plus nette. Le prochain objectif algorithmique, séparé, sera l'optimisation du solveur Q6 et la généralisation propre du masque solide courbe en inlet/outlet.




\section{Développements OpenMP post-0143 : segments ouverts, configuration explicite et capacité fermée}

Les développements postérieurs au portage OpenMP du resampling pondéré ont poursuivi deux objectifs complémentaires. Le premier est logiciel : réduire les états de configuration ambigus et permettre des conditions limites ouvertes plus générales, notamment plusieurs segments \texttt{inlet}/\texttt{outlet} sur une même face. Le second est physique : traiter le cas limite d'un domaine fermé déjà plein soumis à un flux entrant net, situation incompatible avec une fermeture incompressible stricte mais naturellement interprétable comme une compression faible, puis forte, d'un liquide dans une enceinte rigide.

Cette section documente les patchs 0143 à 0154. Elle ne remplace pas les sections précédentes : les cas Taylor--Green, Poiseuille, solides immergés et remplissage restent les validations de référence du couple Q6+resampling. Les nouveaux développements ajoutent un chemin d'étude pour les domaines fermés surpressurisés et pour la future interaction avec une paroi déformable.

\subsection{Conditions limites ouvertes segmentées en coordonnées relatives}

La première simplification concerne les conditions limites ouvertes. L'ancien formalisme associait un mode à une face entière, par exemple \texttt{bcLeft=inlet} ou \texttt{bcRight=outlet}, avec une option d'ouverture partielle limitée à un unique intervalle. Ce formalisme ne permettait pas d'écrire proprement un cas du type : entrée sur la partie haute de la face gauche, sortie sur la partie basse de la même face, et paroi solide ailleurs.

Le nouveau formalisme introduit une liste de segments ouverts :
\begin{verbatim}
openBoundarySegmentsEnable = true
openBoundarySegmentCount = N
openBoundarySegmentK = face mode sMin sMax ux uy type mass
\end{verbatim}
où \texttt{face}\(\in\{\texttt{left},\texttt{right},\texttt{bottom},\texttt{top}\}\), \texttt{mode}\(\in\{\texttt{inlet},\texttt{outlet}\}\), et \(s\in[0,1]\) est la coordonnée tangentielle relative à la face. Sur les faces gauche/droite, \(s\) désigne la coordonnée verticale relative ; sur les faces bas/haut, il désigne la coordonnée horizontale relative. Les champs \texttt{ux}, \texttt{uy}, \texttt{type} et \texttt{mass} servent à l'injection ; ils sont présents dans la syntaxe des \texttt{outlet} pour garder un parseur à longueur fixe.

La convention retenue est volontairement stricte : les segments ouverts sont posés sur une face solide par défaut, et les portions non couvertes restent solides. Par exemple,
\begin{verbatim}
bcLeft   = solid
bcRight  = solid
bcBottom = solid
bcTop    = solid

openBoundarySegmentsEnable = true
openBoundarySegmentCount = 2
openBoundarySegment0 = left inlet  0.65 0.90  0.10 0.00  0 1.0
openBoundarySegment1 = left outlet 0.10 0.35  0.00 0.00  0 1.0
\end{verbatim}
représente une face gauche solide percée localement d'un inlet haut et d'un outlet bas. Les anciens paramètres d'ouverture partielle de type \texttt{leftOpenYMin}/\texttt{leftOpenYMax} ont été supprimés afin d'éviter que de vieux scripts produisent des comportements implicites non maîtrisés.

Du point de vue algorithmique, cette évolution a trois conséquences. Premièrement, le traitement lagrangien des frontières est effectué par face et par segment : une particule sortant par une portion solide est réfléchie, tandis qu'une particule sortant par un segment \texttt{outlet} est convertie en particule inactive. Deuxièmement, le réservoir \texttt{inlet} parcourt les segments d'injection et non plus une face unique. Troisièmement, Q6 reconstruit des profils de flux ouverts par face à partir des segments ; le solveur elliptique n'a pas besoin d'être modifié, car les profils discrets de flux étaient déjà stockés par rangée ou colonne de face.

Cette abstraction prépare aussi l'injection multi-type. Le champ \texttt{type} peut d'abord être interprété comme un traceur passif de provenance, par exemple pour distinguer un jet central et une gaine coaxiale. Une vraie physique multi-espèces exigerait toutefois des collisions, thermostats, dépôts massiques et remaps dépendants du type ; cette extension reste séparée.

\subsection{Modernisation de la configuration : suppression de \texttt{method}}

Le second nettoyage supprime l'ancien aiguillage \texttt{method=classic/q6}. Dans les versions précédentes, \texttt{method=q6} activait implicitement la projection, tandis que \texttt{projectionEnable=true} l'activait aussi. Les deux options étaient donc redondantes, et des configurations comme
\begin{verbatim}
method = q6
projectionEnable = false
\end{verbatim}
étaient sémantiquement ambiguës. Or Q6 n'est pas un modèle séparé de la projection : Q6 est précisément l'opérateur de projection de vitesse.

La configuration moderne repose donc sur deux interrupteurs principaux :
\begin{equation}
  \texttt{projectionEnable}\in\{\texttt{false},\texttt{true}\},
  \qquad
  \texttt{resamplingEnable}\in\{\texttt{false},\texttt{true}\}.
\end{equation}
Ils définissent quatre cas d'ablation :
\begin{center}
\begin{tabular}{lll}
\toprule
\texttt{projectionEnable} & \texttt{resamplingEnable} & Interprétation \\
\midrule
\texttt{false} & \texttt{false} & SRC/MPCD classique \\
\texttt{false} & \texttt{true}  & SRC classique avec garde numérique de support \\
\texttt{true}  & \texttt{false} & SRC avec projection Q6, sans resampling \\
\texttt{true}  & \texttt{true}  & Q6+resampling pondéré \\
\bottomrule
\end{tabular}
\end{center}

Le paramètre \texttt{resamplingPopulationGuardEnable} a également été supprimé. Lorsque \texttt{resamplingEnable=true}, le garde de population est considéré comme le coeur du resampling, et son comportement est piloté par les bornes
\begin{equation}
  N_{
m min} < N_{
m target} < N_{
m max}.
\end{equation}
Cette simplification est importante pour l'interprétation physique : le resampling peut être utilisé soit avec Q6, comme élément de la chaîne quasi-incompressible, soit sans Q6, comme outil purement numérique de maintien du support particulaire SRC classique.

\subsection{Masque mouillé pour domaines pleins et correction du dépeuplement pariétal}

Les essais en cuve pleine avec inlet/outlet localisés ont mis en évidence un défaut du mode \texttt{resamplingWetMaskMode=occupied}. Dans ce mode, une cellule est considérée mouillée uniquement si elle contient déjà de la masse. Près d'une paroi solide, une fluctuation peut vider temporairement une cellule ; celle-ci devient alors \emph{dry}, sort du champ du garde de population, et peut rester durablement vide. Le problème ne vient pas d'une fuite solide directe, mais d'un effet de cliquet wet/dry.

Pour une cuve initialement pleine, le bon choix est \texttt{resamplingWetMaskMode=active\_domain}. La cellule reste mouillée tant qu'elle appartient au domaine fluide actif, même si son occupation instantanée tombe à zéro. Le garde de population et l'insertion locale peuvent alors la réparer. Le passage à \texttt{active\_domain} a supprimé les cellules vides persistantes près des parois dans les runs de cuve pleine.

La règle pratique devient donc :
\begin{itemize}
  \item pour un remplissage depuis domaine vide ou latent, \texttt{occupied} reste adapté, car les cellules non atteintes par le front doivent rester sèches ;
  \item pour un domaine plein ou une cuve déjà mouillée, \texttt{active\_domain} est nécessaire pour empêcher le dépeuplement pariétal irréversible.
\end{itemize}

\section{Réponse de capacité fermée et viriel faiblement compressible}

\subsection{Motivation physique}

Un domaine fermé, plein, à parois infiniment rigides et soumis à un inlet sans outlet n'est pas compatible avec une incompressibilité stricte. Pour un fluide incompressible dans un volume fixe,
\begin{equation}
  \nabla\cdot u = 0,
  \qquad
  \int_{\partial\Omega} u\cdot n\,\mathrm{d}S = 0.
\end{equation}
Un flux entrant net viole cette compatibilité. Un liquide réel répondrait par sa compressibilité finie : la masse volumique augmenterait, la pression monterait, et la température pourrait monter si la compression est suffisamment adiabatique. La limite incompressible stricte correspondrait à une pression infinie pour imposer une quantité finie de liquide supplémentaire dans un volume rigide.

L'objectif n'est donc pas de stabiliser artificiellement ce cas, mais de permettre une transition continue : tant que le domaine reste sous sa capacité nominale, le comportement quasi-incompressible est conservé ; si un flux entrant maintenu produit une surmasse, l'intensité effective de Q6 diminue, la fermeture virielle se raidit, et le remap de masse cesse d'effacer la compression globale.

\subsection{Mesure de surcapacité}

On définit une masse de référence fermée
\begin{equation}
  M_{\rm ref} = N_{\rm ref}\,M_{\rm cell,ref},
\end{equation}
où \(N_{\rm ref}\) est le nombre de cellules du domaine fluide de référence et \(M_{\rm cell,ref}\) la masse cellulaire nominale, typiquement \(\gamma m_0\). La surcapacité relative est
\begin{equation}
  \eta_M = \max\left(0,\frac{M_{\rm tot}-M_{\rm ref}}{M_{\rm ref}}\right).
\end{equation}
Toutes les réponses de capacité sont continues en \(\eta_M\). Si \(\eta_M=0\), les cas validés précédemment restent strictement inchangés.

\subsection{Modulation de Q6 par \texttt{q6ProjectionStrength}}

La projection Q6 n'est pas remplacée par un nouveau solveur. Le coefficient d'application de la projection est simplement modulé :
\begin{equation}
  \alpha_{Q6}^{\rm eff}
  = \alpha_{Q6}^{0}\,
  \exp\left[-\left(\frac{\eta_M}{\eta_{Q6}}\right)^{p_{Q6}}\right],
\end{equation}
où \(\alpha_{Q6}^{0}\) est \texttt{q6ProjectionStrength}. Lorsque la cuve n'est pas sur-remplie, \(\alpha_{Q6}^{\rm eff}=\alpha_{Q6}^{0}\). Lorsque la surmasse augmente, Q6 s'affaiblit progressivement. Cette écriture évite un basculement brutal entre un solveur incompressible et un solveur compressible.

Dans les diagnostics OpenMP, on suit séparément la force nominale de Q6, le facteur de capacité, et la force effective réellement appliquée. Les premiers tests montrent que, dans un cas inlet seul maintenu, \(\alpha_{Q6}^{\rm eff}\) chute rapidement vers zéro lorsque la surcapacité atteint quelques pourcents : le système quitte alors volontairement le régime quasi-incompressible.

\subsection{Renforcement viriel continu}

La fermeture virielle est renforcée selon une loi de capacité, par exemple
\begin{equation}
  K_{\rm vir}^{\rm eff}
  = K_{\rm vir}^{0}\left[1+G_{\rm vir}
  \left(\frac{\eta_M}{\eta_{\rm vir}}\right)^{p_{\rm vir}}\right].
\end{equation}
La pression virielle cellulaire est reconstruite relativement à la masse cellulaire nominale :
\begin{equation}
  P_{{\rm vir},c}
  = K_{\rm vir}^{\rm eff}\left(\frac{M_c}{M_{\rm cell,ref}}-1\right).
\end{equation}
La pression totale utilisée pour les diagnostics bulk et pariétaux est
\begin{equation}
  P_{{\rm tot},c}=P_{{\rm kin},c}+P_{{\rm vir},c}.
\end{equation}

Un kick viriel optionnel applique une correction de vitesse proportionnelle à \(-\nabla P_{\rm vir}/\rho\), avec correction globale de moment si elle est activée. Dans un état uniformément pressurisé, ce gradient est faible ; la pression peut donc monter sans générer de kick volumique important. En revanche, un inlet localisé crée des gradients viriels et peut produire des accélérations violentes, voire une rupture numérique, si le flux est maintenu.

\subsection{Cible effective du remap de masse}

Le premier prototype de réponse fermée montrait un verrou : le remap de masse ramenait chaque cellule vers \(M_{\rm cell,ref}\), donc il effaçait presque toute surmasse avant que le viriel ne puisse vraiment monter. Le correctif consiste à conserver la surmasse globale dans la cible de remap :
\begin{equation}
  M_{\rm cell,target}^{\rm eff}
  = M_{\rm cell,ref} + \frac{\max(0,M_{\rm tot}-M_{\rm ref})}{N_{\rm ref}}.
\end{equation}
Ainsi, le resampling conserve son rôle numérique : homogénéiser la masse, le moment et le support particulaire. Mais il ne détruit plus la compression globale imposée par un flux entrant net. La masse totale peut alors augmenter, \(\eta_M\) croît, Q6 s'affaiblit et \(K_{\rm vir}^{\rm eff}\) augmente.

Cette modification préserve les validations antérieures car elle n'est active que lorsque la réponse de capacité est explicitement activée et que \(M_{\rm tot}>M_{\rm ref}\). Sinon, le remap conserve sa cible fixe historique.

\subsection{Injection additive et diagnostics de moment}

Pour qu'une cuve pleine puisse réellement se pressuriser, l'inlet ne doit pas seulement remplacer localement les particules du réservoir. Une option d'injection additive associe au flux normal imposé un budget de masse positif. Si le domaine est fermé, cette masse reste dans le domaine et augmente \(M_{\rm tot}\). Les diagnostics suivent alors : masse totale, surcapacité, force Q6 effective, module viriel effectif, pression virielle moyenne, kick viriel et activité du resampling.

Les runs de cuve fermée avec inlet seul montrent une montée de masse jusqu'à plusieurs pourcents de surcapacité, une extinction rapide de Q6, une croissance forte de \(K_{\rm vir}^{\rm eff}\) et une montée de l'énergie. Ils montrent aussi un budget de quantité de mouvement encore incomplet : les contributions Q6 et resampling ne suffisent pas à expliquer les sauts tardifs du moment global. Les contributions de paroi, d'inlet et de viriel doivent donc être instrumentées plus finement avant d'interpréter complètement la dynamique. Cette limite ne remet pas en cause la validation statique de la charge pariétale, mais elle encadre l'usage dynamique du modèle.

\section{Pression pariétale virielle et objectif parois déformables}

\subsection{Pourquoi la pression bulk ne suffit pas}

Dans un continuum, une pression uniforme exerce une traction normale sur une paroi, même si \(\nabla p=0\) dans le bulk. Dans le code particulaire, une fermeture virielle utilisée seulement comme kick volumique \(-\nabla P_{\rm vir}\) ne transmet pas automatiquement une pression uniforme à une paroi solide. Pour préparer un couplage avec une paroi déformable ou un modèle de rupture, il faut donc construire explicitement une charge pariétale.

Le diagnostic ajouté convertit la pression cellulaire reconstruite,
\begin{equation}
  P_{{\rm tot},c}=P_{{\rm kin},c}+P_{{\rm vir},c},
\end{equation}
en charge normale sur les faces solides. Pour une face solide adjacente à une cellule \(c\), la contribution élémentaire est
\begin{equation}
  \Delta \bm{F}_{\rm wall} = P_{{\rm tot},c}\,A_f\,\bm{n}_{\rm wall},
\end{equation}
où \(A_f\) vaut \(\Delta y\) pour les faces verticales et \(\Delta x\) pour les faces horizontales en 2D de profondeur unité. Les segments \texttt{inlet}/\texttt{outlet} sont exclus de cette intégration : seule la portion solide de la paroi porte la charge.

Les diagnostics OpenMP séparent la pression cinétique, la pression virielle, la pression totale par face, les longueurs solides par face, et les forces nettes \(F_x,F_y\). Cette séparation est essentielle pour distinguer une pression hydrostatique uniforme d'un chargement fortement asymétrique induit par un inlet localisé.

\subsection{Validation statique : boîte uniformément surcompressée}

Le test de validation 0152 initialise une boîte fermée, sans inlet, avec une masse uniforme
\begin{equation}
  M = \lambda M_{\rm ref},
  \qquad \lambda\in\{1.00,1.02,1.05,1.10\}.
\end{equation}
La dynamique est presque statique et virielle pure : pas de thermostat de paroi, pas de \texttt{wallVP} thermique et pas d'inlet. Le résultat attendu est
\begin{equation}
  p_{\rm wall,left}\simeq p_{\rm wall,right}
  \simeq p_{\rm wall,bottom}\simeq p_{\rm wall,top},
  \qquad
  \bm{F}_{\rm wall}^{\rm net}\simeq 0.
\end{equation}

Les résultats obtenus valident ce chaînage. La pression moyenne pariétale suit la loi attendue \(K_{\rm vir}^{\rm eff}\eta_M\), la pression par face est confondue pour les quatre parois, et la force nette relative est au niveau du bruit numérique. Ce test démontre que la pression bulk virielle est bien convertible en charge normale pariétale dans le cas de référence uniformément surcompressé.

\subsection{Cas dynamique inlet seul et limite actuelle}

Dans un cas inlet seul, la pression pariétale augmente fortement et devient non uniforme : la face proche de l'injection et la face supérieure peuvent recevoir une charge plus forte que les autres. Cette non-uniformité n'est pas forcément un défaut ; elle est attendue dans un régime transitoire de cuve fermée alimentée localement. En revanche, le budget de quantité de mouvement doit encore être complété. Pour un futur couplage solide, il faudra exposer dans les fichiers \texttt{summary\_runtime.csv} les impulsions associées aux parois, à l'inlet, au kick viriel brut, au kick viriel corrigé, au thermostat et aux corrections de projection.

Le statut actuel est donc le suivant : la conversion \(P_{\rm kin}+P_{\rm vir}\) vers une charge normale pariétale est validée en statique uniforme ; l'usage dynamique dans un domaine fermé alimenté par un inlet est prometteur, mais demande encore un budget d'impulsion complet.

\section{Positionnement bibliographique : MPCD quasi-incompressible, projection et resampling}

\subsection{Contexte MPCD/SRD}

La méthode MPCD/SRD a été introduite comme modèle mésoscopique de solvant par Malevanets et Kapral \cite{MalevanetsKapral1999}. La revue de Gompper, Ihle, Kroll et Winkler \cite{GompperIhleKrollWinkler2009} synthétise les variantes principales, l'interprétation hydrodynamique, le rôle du streaming/collision, les coefficients de transport et les extensions vers fluides complexes. Dans ce cadre, le fluide standard est un solvant fluctuant, localement conservatif en masse et moment, mais il n'est pas strictement incompressible. En pratique, on se place souvent à bas Mach ou à densité particulaire suffisante pour réduire les effets compressibles.

Les travaux sur MPCD non idéal traitent explicitement cette limite. Ihle, Tüzel et Kroll ont proposé un algorithme particulaire avec équation d'état non idéale fondé sur des collisions multiparticulaires biaisées dépendant de la densité locale \cite{IhleTuzelKroll2006}. Plus récemment, Zantop et Stark ont introduit une règle de collision MPCD modifiée produisant une équation d'état non idéale et une compressibilité fortement diminuée \cite{ZantopStark2021a,ZantopStark2021b}. Ces travaux sont proches dans l'objectif général -- réduire la compressibilité MPCD -- mais ils agissent principalement sur la collision et l'équation d'état.

\subsection{Projection et méthodes particulaires incompressibles}

Le choix Q6 s'inscrit dans la tradition des méthodes de projection de pression pour les équations incompressibles, initiée par Chorin \cite{Chorin1968}. Dans les méthodes particulaires, la même tension entre incompressibilité et compressibilité artificielle se retrouve en SPH. PCISPH impose l'incompressibilité par une procédure prédictive-corrective de pression \cite{SolenthalerPajarola2009}, tandis que les travaux récents sur l'ACSPH et le \(\delta\)-SPH mettent en évidence les liens entre correction pression/vitesse et compressibilité artificielle \cite{DeCourcy2024}. Des solveurs Navier--Stokes faiblement compressibles à projection de pression évolutive explorent également l'idée de ne pas imposer strictement \(\nabla\cdot u=0\), mais d'utiliser une projection pour amortir les ondes acoustiques et contrôler la pression \cite{Yang2021}.

La contribution Q6+resampling est différente. La projection Q6 traite la partie cinématique : elle retire la composante compressive du champ de vitesse lorsque le bilan global de flux est compatible avec l'incompressibilité. Le resampling pondéré traite une difficulté distincte : la qualité du support particulaire. Une cellule pauvre ou vide près d'une paroi ne doit pas être interprétée comme une dépression physique ; elle doit d'abord être réparée comme défaut de représentation. Le garde \(N_{
m min}/N_{
m target}/N_{
m max}\), les rôles \texttt{Fluid}/\texttt{Latent}/\texttt{Inactive} et le remap conservatif masse--moment--énergie séparent donc clairement
\begin{equation}
  \text{incompressibilité cinématique}
  \quad \text{et} \quad
  \text{support particulaire local}.
\end{equation}

Cet avantage est particulièrement net par rapport à deux stratégies concurrentes : augmenter fortement \(\gamma\), ce qui réduit le bruit mais augmente le coût, ou modifier la collision pour obtenir une équation d'état non idéale. L'approche développée ici conserve le coeur SRC/MPCD classique, permet des ablations propres \texttt{classic}/\texttt{q6}/\texttt{classic\_resampling}/\texttt{q6\_resampling}, et laisse la fermeture de pression comme un module séparé. Cette modularité est importante pour la parallélisation OpenMP et, à terme, GPU.

\section{Positionnement bibliographique : viriel faiblement compressible et pression pariétale}

\subsection{Contraintes, tenseurs de stress et forces pariétales en MPCD}

Les contraintes et tenseurs de stress MPCD ont été analysés par Winkler et Huang \cite{WinklerHuang2009}. Ils distinguent notamment des formulations internes et externes du tenseur de contrainte, et considèrent des fluides confinés. Imperio, Padding et Briels ont proposé une méthode de calcul des forces sur parois et particules incluses en MPCD avec conditions de non-glissement, en évitant de reconstruire explicitement le profil de vitesse près de la surface \cite{ImperioPaddingBriels2011}. Ces travaux constituent les références les plus proches pour la question des forces mécaniques exercées sur un solide dans MPCD.

La nouveauté de la fermeture actuelle ne réside pas dans l'existence d'une force de paroi MPCD, mais dans le fait que la pression bulk reconstruite,
\begin{equation}
  P_{\rm tot}=P_{\rm kin}+P_{\rm vir},
\end{equation}
est explicitement projetée en charge normale sur les portions solides des parois dans un domaine fermé surcompressé. Le test statique 0152 montre que cette charge donne une pression identique sur les quatre faces et une force nette nulle lorsque la surpression est uniforme. Cette étape est indispensable pour envisager un couplage avec un solide déformable : une pression uniforme de bulk doit charger une paroi même lorsque \(\nabla P=0\) dans le domaine.

\subsection{Parenté avec SPH faible compressibilité et FSI}

Les méthodes SPH faiblement compressibles utilisent depuis longtemps une équation d'état raide pour convertir une variation de densité en pression. Les modèles de paroi WCSPH et la qualité des pressions pariétales ont été étudiés par Valizadeh et Monaghan \cite{ValizadehMonaghan2015}. Côté fluide--structure et rupture, des cadres SPH récents couplent WCSPH à des modèles de solides déformables ou fracturables, par exemple le cadre d'Islam pour les grandes déformations et la fracture \cite{Islam2024}, ou des modèles 3D combinant WCSPH et pseudo-ressorts pour la rupture fragile sous chargement hydrodynamique \cite{Singh2025}. Ces références montrent que la chaîne \emph{pression fluide -- charge pariétale -- déformation/rupture} est une voie active dans les méthodes particulaires.

La différence est que la méthode présente n'est pas une formulation SPH de pression. Elle conserve le coeur SRC/MPCD, ajoute une projection de vitesse modulable, maintient le support particulaire par resampling pondéré, et introduit une pression virielle qui devient dominante seulement lorsque la capacité fermée est dépassée. À notre connaissance, cette combinaison -- projection MPCD affaiblie par surcapacité, remap conservant la surmasse globale, renforcement viriel continu et projection pariétale de \(P_{\rm kin}+P_{\rm vir}\) -- n'a pas été rapportée sous cette forme dans la littérature MPCD.

La revendication doit rester prudente : les ingrédients individuels existent dans des familles voisines, notamment MPCD non idéal, projection incompressible, artificial compressibility, stress MPCD et SPH-FSI. L'originalité se situe dans leur assemblage modulaire pour passer continûment d'un SRC/MPCD quasi-incompressible validé à un régime faiblement compressible capable de charger des parois fermées.



\section{Optimisation OpenMP, profilage contrôlé et branche de production allégée}

Les développements précédents ont établi la formulation physique et numérique de la chaîne SRC/Q6/resampling/viriel. Une étape séparée a ensuite été menée sur le dépôt C++/OpenMP afin de réduire le temps d'exécution sans modifier les fonctionnalités ni la trajectoire numérique. Cette étape a suivi une règle stricte : toute optimisation modifiant un chemin algorithmique devait être validée par une comparaison discriminante entre la version de référence et la version optimisée. Les cas utilisés pour cette comparaison étaient les quatre cas mono-configuration déjà employés dans le dépôt : Taylor--Green périodique, Poiseuille avec parois, obstacle rectangulaire ouvert, et piston avec fermeture virielle. Le critère d'acceptation était l'égalité, à tolérance fixée, des métriques physiques, numériques et de resampling sur l'ensemble des cas.

\subsection{Méthode de travail et séparation des branches}

La phase d'optimisation a conduit à distinguer trois usages du code OpenMP. La branche \texttt{SRC\_openMP\_optimized} est conservée comme branche de développement et de préparation GPU : elle garde les diagnostics internes fins, les scripts de profilage et les outils de comparaison. La branche \texttt{clean/openmp-light} dérive de cette version optimisée mais désactive les diagnostics internes par défaut afin de fournir une version OpenMP portable, plus proche d'un code de production. Enfin, une future branche \texttt{clean/openmp-production-stripped} pourra retirer complètement le code de profilage interne lorsque la version allégée aura été suffisamment stabilisée.

Cette séparation évite de confondre deux objectifs. D'une part, la branche instrumentée reste nécessaire pour identifier les noyaux dominants, préparer un futur backend GPU et comprendre les régressions. D'autre part, la branche light vise un exécutable qui ne crée pas de fichiers de profilage internes en usage normal et qui laisse les analyses de cas au post-traitement externe, conformément à l'architecture initiale du code.

\subsection{Instrumentation initiale et identification des points chauds}

Une première série de patches a ajouté des profils internes contrôlés : temps par phase du pas SRC/MPCD, profil Q6/CG, profil des dépôts particules--cellules, profil des gardes de population et de masse, et diagnostics post-guard. Ces profils ont été utiles pour localiser précisément les verrous. Ils ont montré que le coût n'était pas concentré dans un seul module, mais dans plusieurs chemins distincts : projection elliptique Q6, dépôts pondérés répétés, garde de population, garde de masse, et rafraîchissement thermique. Ils ont aussi permis d'écarter certaines optimisations intuitives mais non pertinentes, par exemple la seule réutilisation des buffers temporaires du \texttt{mass\_guard}, dont le gain s'est révélé marginal.

Le profilage a également révélé plusieurs coûts anormaux. Le plus net concernait la réponse de capacité fermée : lorsque \texttt{closedCapacityResponseEnable=false}, le code effectuait encore un dépôt de masse cellulaire avant de constater que la fermeture était désactivée. Un retour anticipé au début du chemin viriel a supprimé ce coût parasite. Un second coût anormal provenait du changement de rôle dans le \texttt{population\_guard} : les fonctions de mutation appelaient indirectement des vérifications globales de rôles à chaque activation ou inactivation locale. Le profil profond des mutations a montré que ces appels dominaient presque totalement le coût du \texttt{population\_guard}.

\subsection{Optimisations OpenMP validées}

Les optimisations finalement conservées sont celles qui ont passé les validations discriminantes. Elles sont résumées ci-dessous.

\begin{center}
\small
\begin{tabular}{p{0.24\linewidth}p{0.44\linewidth}p{0.24\linewidth}}
\toprule
Zone du code & Modification retenue & Effet principal \\
\midrule
Réponse virielle fermée & Retour immédiat si la fermeture de capacité est désactivée, avant tout dépôt masse/cellule. & Suppression d'un coût inutile dans les cas sans viriel. \\
Projection Q6/CG & Réorganisation du noyau elliptique, pré-calculs et fusion de certaines opérations vectorielles. & Réduction du coût du solveur elliptique tout en conservant les résidus Q6. \\
Resampling sans édition de rôles & Suppression des reconstructions de pool après remap ou renormalisation thermique lorsque seules les masses ou vitesses changent. & Moins de scans particulaires inutiles. \\
\texttt{population\_guard} & Listes compactes de cellules candidates overfull/underfull, conservant l'ordre croissant des cellules. & Évite de parcourir toutes les cellules lorsque seules quelques cellules sont candidates. \\
Mutations de rôles & Écriture locale préconditionnée des rôles \texttt{Fluid}/\texttt{Inactive} dans les mutations, sans revalidation globale répétée. & Suppression du verrou principal du \texttt{population\_guard}. \\
Dépôt post-thermal & Rafraîchissement spécialisé des moments et vitesses après renormalisation thermique, sans reconstruire les listes candidates ni les plans. & Forte baisse du coût \texttt{post\_thermal\_deposit}. \\
Dépôt post-guard & Réutilisation stricte des \texttt{cellId} après \texttt{population\_guard}, avec mise à jour locale des particules activées/inactivées, mais dépôt complet conservé. & Réduction du coût de \texttt{post\_guard\_deposit} sans changer les listes/plans. \\
Diagnostics internes & Mode light par défaut, profils réactivables par \texttt{MPCD\_INTERNAL\_PROFILES=1}, puis suppression du coût runtime des profils lorsque désactivés. & Version portable sans fichiers de profilage internes en usage normal. \\
\bottomrule
\end{tabular}
\end{center}

Deux optimisations ont été explicitement rejetées malgré des gains apparents. Le fast path du \texttt{mass\_guard}, qui évitait la projection de masse lorsque la cellule paraissait déjà admissible, a modifié la trajectoire discrète du resampling dans le test discriminant. De même, une première tentative de refresh local complet du dépôt post-guard a accéléré le code mais a changé certaines décisions de resampling sur Taylor--Green, obstacle ouvert et piston viriel. Ces rejets sont importants : ils montrent que le critère de validation n'était pas seulement temporel, mais bien fonctionnel.

\subsection{Résultats de validation et gains mesurés}

La validation discriminante finale du jalon OpenMP optimisé a comparé la version de référence et la version optimisée sur les quatre cas standard, avec \(64\times64\) cellules, \(\gamma=20\), \(1000\) pas et \(8\) threads. Les quatre cas ont passé la comparaison avec zéro métrique échouée sur 76 métriques par cas, soit 304 métriques comparées sans divergence. Les speedups mesurés sur cette campagne étaient de l'ordre de \(1.45\times\) pour l'obstacle ouvert, \(1.52\times\) pour le piston viriel, \(1.96\times\) pour Poiseuille et \(1.47\times\) pour Taylor--Green. Ces valeurs dépendent naturellement de la machine, du bruit d'ordonnancement OpenMP et des scripts utilisés, mais elles donnent l'ordre de grandeur du gain obtenu sans changer la logique physique.

Sur le profil \texttt{q6\_resampling} court, le verrou \texttt{population\_guard} a été presque éliminé comme poste dominant. Après l'écriture locale des rôles, son coût est devenu secondaire devant Q6, \texttt{mass\_guard}, collision et dépôts. Le dépôt post-thermal a également cessé d'être un verrou significatif après spécialisation. Le dépôt post-guard a été réduit sans recourir au refresh local non équivalent : la réutilisation stricte des \texttt{cellId} conserve le dépôt complet et donc les classifications/listes/plans, mais évite une partie substantielle du coût de calcul cellulaire.

\subsection{Réorganisation des diagnostics}

Une seconde étape a consisté à séparer les diagnostics de développement des diagnostics fluides de production. Les fichiers de profilage internes ont été désactivés par défaut dans la branche \texttt{clean/openmp-light}. Sont notamment concernés :
\begin{itemize}[nosep]
  \item \texttt{phase\_profile} et \texttt{q6\_cg\_profile} ;
  \item \texttt{deposit\_profile} et \texttt{resampling\_guard\_profile} ;
  \item \texttt{post\_guard\_profile} et les diagnostics d'équivalence post-guard.
\end{itemize}
 Les sorties physiques et numériques générales restent conservées : masse, nombres de particules fluides/inactives, moment global, température effective, statistiques de population, résidus Q6 utiles, compteurs globaux de resampling et diagnostics viriels lorsque la fermeture est active.

La variable d'environnement \texttt{MPCD\_INTERNAL\_PROFILES=1} permet encore de réactiver les profils internes dans la branche light. Cela préserve une possibilité de diagnostic ponctuel sans encombrer les runs de production. Un patch complémentaire a ensuite évité le coût runtime des timers et accumulateurs internes lorsque cette variable n'est pas active. Des scripts de smoke test vérifient les deux comportements : aucun fichier de profilage interne en mode light, et production effective de profils lorsque \texttt{MPCD\_INTERNAL\_PROFILES=1} est explicitement défini.

Enfin, les scripts de profilage historiques, les rapports intermédiaires et les CSV de comparaison ont été déplacés vers \texttt{dev\_history/}. La branche light garde donc les scripts de build, de validation générique et de smoke test, mais n'expose plus les nombreux scripts de développement liés aux jalons d'optimisation. Une validation finale \texttt{clean/openmp-light} contre \texttt{SRC\_openMP\_optimized} a confirmé l'absence de modification fonctionnelle : les quatre cas standard ont passé la comparaison avec zéro métrique échouée.

\subsection{Conséquences pour la suite GPU}

Cette phase d'optimisation OpenMP fournit aussi une base utile pour le port GPU. Les profils conservés dans la branche instrumentée ont identifié les noyaux qui deviendront prioritaires sur de grandes grilles : projection Q6, dépôt particules--cellules, collision par cellule et éventuellement construction d'index cellule--particules. À l'échelle \(1000\times1000\) avec \(\gamma\simeq20\), le nombre de particules atteint l'ordre de \(2\times10^7\), et les gains purement OpenMP auront probablement des rendements décroissants. La branche optimisée/instrumentée doit donc rester le support du développement GPU, tandis que la branche \texttt{clean/openmp-light} fournit une référence CPU portable et validée.



\section{Portage CUDA du SRC classic, architecture résidente et thermostat boundary-aware}

\subsection{Définition du périmètre : SRC classic et fermeture liquide}

Le chantier GPU a été mené sur la branche \texttt{SRC\_GPU}, dérivée de la branche OpenMP allégée. Il est essentiel de fixer la terminologie, car elle conditionne l'architecture du code. Dans cette branche, on appelle \textbf{SRC classic} le pas physique complet
\begin{equation}
  \text{advection/streaming}
  + \text{décalage aléatoire de grille}
  + \text{rotation/collision SRC}
  + \text{thermostat}.
\end{equation}
La fermeture liquide incompressible ou faiblement compressible est un module additionnel :
\begin{equation}
  \text{SRC classic}
  + \text{Q6}
  + \text{resampling pondéré}
  + \text{fermeture virielle/capacité}.
\end{equation}
Le résultat CUDA obtenu à ce stade concerne donc le SRC classic. Il ne doit pas être confondu avec une migration CUDA complète de Q6, du resampling ou du viriel. Cette séparation a été conservée volontairement afin de ne pas verrouiller les chantiers futurs, en particulier la migration Q6 CUDA.

Le choix de CUDA, plutôt qu'un backend abstrait OpenCL/CUDA, a été retenu pour ce cycle de développement afin de réduire le coût d'intégration et de concentrer l'effort sur la résidence mémoire, les transferts et les chemins physiques. Le backend OpenMP reste disponible explicitement, ce qui permet de comparer les trajectoires et de conserver une référence CPU robuste.

\subsection{Principe d'architecture : état persistant et chemins spécialisés}

La difficulté principale du portage n'est pas seulement d'écrire des noyaux CUDA isolés. Un port naïf qui transfère toutes les particules vers le GPU à chaque sous-étape, puis les redescend immédiatement vers l'hôte, reproduit le calcul mais perd l'essentiel du gain. La stratégie retenue est donc une architecture \emph{résidente} : les particules et les espaces de travail cellule restent autant que possible côté GPU pendant la séquence SRC classic.

On distingue plusieurs couches.
\begin{itemize}[nosep]
  \item L'état particulaire persistant contient positions, vitesses, masses, rôles et métadonnées nécessaires aux noyaux de streaming, frontière, collision et thermostat.
  \item L'espace de travail cellule persistant contient les dépôts, moments, vitesses moyennes, énergies thermiques et compteurs de population nécessaires à la collision SRC et au thermostat.
  \item Les chemins de streaming/frontière sont spécialisés par famille géométrique : périodique, parois simples, obstacle rectangle, piston/mobile wall, inlet/outlet full-face et inlet/outlet segmenté.
  \item Le noyau de collision SRC applique la rotation par cellule en utilisant les dépôts cellule, les signes aléatoires et les informations de rôle/matériau.
  \item Le thermostat CUDA reconstruit l'énergie thermique relative et applique le rescale cellulaire lorsque le chemin physique le permet.
\end{itemize}

Cette organisation explique la différence entre deux types de chemins. Pour un cas SRC classic-only, le chemin cible est entièrement fusionné côté GPU :
\begin{equation}
  \text{frontière/streaming CUDA}
  \longrightarrow
  \text{collision SRC CUDA}
  \longrightarrow
  \text{thermostat CUDA}.
\end{equation}
Pour un cas où Q6, resampling, viriel ou réponse de capacité modifient les vitesses ou les masses après la collision SRC, le thermostat fusionné n'est pas correct. L'ordre conservé est alors
\begin{equation}
  \text{collision SRC CUDA}
  \longrightarrow
  \text{étapes CPU/OpenMP de fermeture liquide}
  \longrightarrow
  \text{thermostat CUDA post-CPU}.
\end{equation}
Ce second chemin est notamment celui validé pour le piston/mobile wall. Il préserve la cohérence physique et la possibilité de réactiver Q6 CPU/OpenMP, resampling CPU/CUDA et viriel.

\subsection{Étapes du portage CUDA}

Le développement CUDA a progressé de manière incrémentale. Les premiers jalons ont validé les briques élémentaires : détection de toolchain CUDA, backend de projection expérimental, dépôts cellule, thermostat isolé, collision SRC isolée, puis état particulaire persistant. Les jalons suivants ont introduit les espaces cellule persistants, les chemins streaming/frontières, les obstacles, l'inlet/outlet et les variantes résidentes.

La séquence validée peut être résumée ainsi.
\begin{center}
\begin{tabular}{p{0.18\linewidth}p{0.72\linewidth}}
\toprule
Jalons & Rôle \\
\midrule
0186--0208 & Mise en place du backend CUDA, tests Taylor--Green, dépôts cellule, thermostat isolé et premiers chemins combinés. \\
0209--0225 & Collision SRC CUDA, état particulaire persistant, ponts host/device, workspace cellule persistant et premières validations end-to-end. \\
0227--0244 & Briques CUDA du resampling : garde pauvre/riche, plan de transfert, opérations extraction/insertion et état persistant. Ces briques préparent le futur, mais ne constituent pas encore la fermeture liquide full CUDA. \\
0245--0249b & Streaming/frontières CUDA : périodique, wall-simple, rectangle immergé, piston/mobile wall, inlet/outlet full-face et inlet/outlet segmenté. \\
0251--0259 & État cellule partagé, collision SRC persistante, premières optimisations de transfert et thermostat CUDA validé initialement en périodique. \\
0260--0275 & Passage au SRC classic résident optimisé : périodique, wall-simple, solide rectangle, piston, inlet/outlet full-face, inlet/outlet segmenté, validations consolidées, instrumentation performance et backend CUDA principal. \\
0276--0281 & Extension du thermostat CUDA aux frontières : wall-simple, solide, piston/mobile wall, inlet/outlet full-face et segmenté, puis validation consolidée. \\
0284--0286 & Ajout du solide circulaire CUDA, d'abord en périodique puis avec inlet/outlet full-face pour le cas Von Karman, et consolidation du jalon SRC classic full CUDA. \\
0288--0293 & Stabilisation des ouvertures CUDA : clipping géométrique des segments, modes explicites de sortie \texttt{neumann}, \texttt{equilibrium\_flux} et \texttt{forced\_flux}, commutateur physique unique \texttt{thermostatEnable}, puis pré-extraction outlet avant refill inlet. \\
\bottomrule
\end{tabular}
\end{center}

Le point 0276--0281 est décisif : avant cette étape, le code disposait d'un chemin SRC CUDA résident robuste, mais le thermostat CUDA n'était réellement validé qu'en périodique. Après cette étape, le thermostat appartient bien au chemin SRC classic CUDA pour les familles de conditions limites validées.

\subsection{Thermostat CUDA : séparation collision/thermostat près des frontières}

Le correctif physique principal concerne les particules virtuelles de paroi ou de solide. Pour la collision SRC, les moyennes cellule peuvent inclure des particules virtuelles, car elles servent à imposer le couplage wall/no-slip/thermal attendu. En revanche, le thermostat de fluide ne doit pas utiliser ces particules virtuelles pour reconstruire l'énergie thermique réelle des particules transportées. La règle validée est donc
\begin{align}
  \text{moyenne de collision} & = \text{particules réelles} + \text{particules virtuelles éventuelles}, \\
  \text{moyenne de thermostat} & = \text{particules réelles post-collision uniquement}.
\end{align}

Cette distinction explique l'échec initial du thermostat CUDA hors périodique et le correctif 0276. Dans le chemin fusionné, il ne fallait pas réutiliser aveuglément les moments de collision enrichis par les particules virtuelles ; il fallait reconstruire, après rotation SRC, un dépôt thermostat fondé sur les seules particules réelles. Une fois cette séparation imposée, le même principe s'est étendu aux parois simples et aux solides fixes.

Pour le piston/mobile wall, une autre contrainte apparaît. Le cas validé peut activer Q6, resampling, viriel ou réponse de capacité entre collision et thermostat. Il serait alors incorrect d'appliquer le thermostat immédiatement après la collision CUDA. La validation piston utilise donc le thermostat CUDA en mode post-CPU : les étapes CPU/OpenMP intermédiaires modifient l'état, puis le thermostat CUDA agit sur l'état mis à jour.


\subsection{Conditions limites couvertes}

Les familles de conditions limites validées couvrent maintenant les cas principaux du dépôt. Le tableau suivant utilise des noms de cas abrégés afin de mettre l'accent sur la famille physique plutôt que sur le nom exact des scripts.
\begin{center}
\begin{tabular}{p{0.24\linewidth}p{0.27\linewidth}p{0.39\linewidth}}
\toprule
Famille & Cas de validation & Chemin CUDA validé \\
\midrule
Périodique & Taylor--Green / SRC périodique & Streaming périodique, collision SRC, thermostat CUDA. \\
Wall-simple & Poiseuille wall & Paroi simple, collision SRC, reconstruction thermostat réelle, thermostat CUDA. \\
Solide fixe & Rectangle périodique & Gestion solide/rectangle, collision SRC, thermostat CUDA. \\
Solide circulaire & Cercle périodique & Réflexion cercle CUDA, fraction solide circulaire, collision SRC et thermostat CUDA. \\
Piston/mobile wall & Piston legacy & Mobile wall et SRC CUDA ; thermostat CUDA post-CPU pour préserver Q6/resampling/viriel. \\
Inlet/outlet full-face & Obstacle rectangle ouvert & Inlet/outlet résident full-face, collision SRC, thermostat CUDA fusionné. \\
Inlet/outlet segmenté & U-turn segmenté & Inlet/outlet segmenté résident, état partagé, collision SRC, thermostat CUDA fusionné. \\
Cercle + inlet/outlet & Cylindre/Von Karman & Full-face inlet/outlet, réservoir dur, cercle solide CUDA, collision SRC et thermostat CUDA fusionné. \\
Régimes outlet explicites & 0291--0293 & Sortie passive \texttt{neumann}, équilibrage local \texttt{equilibrium\_flux} et extraction imposée \texttt{forced\_flux}, applicables aux chemins full-face et segmentés résidents. \\
\bottomrule
\end{tabular}
\end{center}

Le correctif 0280c est particulièrement important pour les ouvertures. Les chemins résidents inlet/outlet refusaient initialement \texttt{thermostatEnable=true}. Le garde a été relâché uniquement dans le cas sûr : thermostat CUDA persistant demandé, Q6/resampling/viriel/capacité désactivés, et chemin classic-only. Ainsi, l'activation du thermostat ne désactive plus silencieusement le chemin résident, mais les modules de fermeture liquide restent protégés.

\subsection{Ouvertures CUDA complètes : full-face, segments et régimes de sortie}

Les conditions limites ouvertes CUDA ont été complétées au-delà du simple couple \texttt{inlet}/\texttt{outlet} géométrique. Le code distingue maintenant deux aspects : la \emph{géométrie} de l'ouverture et le \emph{régime physique} appliqué à la sortie. La géométrie peut être une face complète, par exemple \texttt{bcLeft=inlet} et \texttt{bcRight=outlet}, ou bien une liste de segments portés par une face solide. Dans les deux cas, le chemin cible reste résident CUDA lorsque le cas est SRC classic-only : les particules sont advectées, traitées par les frontières ouvertes, éventuellement insérées ou extraites par le réservoir, puis transmises directement à la collision SRC et au thermostat sans retour hôte intermédiaire.

Le traitement segmenté a nécessité un correctif géométrique spécifique. Les premières versions sélectionnaient les cellules de réservoir à partir de leur centre, puis inséraient les particules dans toute la cellule. Lorsqu'une borne de segment coupait une cellule, des particules pouvaient donc être créées légèrement au-delà de la portion réellement ouverte, ce qui produisait une singularité de densité localisée à la lèvre du segment. Le correctif 0288 remplace cette logique par un clipping d'intervalle : une cellule de bord est pondérée par la fraction réellement ouverte, et les positions injectées sont tirées uniquement dans la sous-cellule incluse dans le segment. La suppression de cette singularité confirme que l'inlet segmenté doit être interprété comme une géométrie continue discrétisée, non comme une sélection de centres de cellules.

Les régimes de sortie sont maintenant explicites via
\begin{equation}
  \texttt{openBoundaryOutletMode}
  \in
  \{\texttt{neumann},\texttt{equilibrium\_flux},\texttt{forced\_flux}\}.
\end{equation}
Leur sens est le suivant.
\begin{center}
\begin{tabular}{p{0.22\linewidth}p{0.68\linewidth}}
\toprule
Mode & Interprétation \\
\midrule
\texttt{neumann} & Sortie passive. Le bord laisse partir les particules qui atteignent naturellement l'outlet ; aucune extraction artificielle n'est imposée. Ce mode convient aux canaux ouverts classiques, aux sillages et aux sorties suffisamment aval. \\
\texttt{equilibrium\_flux} & Équilibrage local. Le code tente d'extraire à l'outlet une quantité comparable au gain dû au refill inlet du pas courant. L'extraction reste plafonnée par la masse ou le nombre de particules effectivement disponibles dans la couche outlet ; elle ne garantit donc pas une masse globale constante si la sortie est localement pauvre. \\
\texttt{forced\_flux} & Extraction imposée par l'utilisateur, décorrélée du flux inlet. La quantité extraite est spécifiée par un flux ou par un budget par pas, et le code prélève localement dans la couche outlet tant que des particules y sont disponibles. Ce mode permet de modéliser aspiration, drainage ou pompage local, en laissant à l'utilisateur le choix d'équilibrer ou non les flux entrant et sortant. \\
\bottomrule
\end{tabular}
\end{center}

Les paramètres associés restent volontairement peu nombreux. Un budget d'extraction peut être donné en masse ou en nombre de particules :
\begin{itemize}[nosep]
  \item \texttt{openBoundaryOutletForcedMassFlux} : flux massique demandé ; la masse extraite par pas est proportionnelle à $\Delta t$.
  \item \texttt{openBoundaryOutletForcedMassPerStep} : masse directement demandée par pas, utile pour les essais contrôlés.
  \item \texttt{openBoundaryOutletForcedParticleFlux} : flux en nombre de particules, converti en budget par pas.
  \item \texttt{openBoundaryOutletForcedParticlesPerStep} : nombre de particules à extraire par pas ; prioritaire pour les tests à masse particulaire uniforme.
  \item \texttt{openBoundaryOutletForcedLayerCells} : épaisseur, en cellules, de la couche outlet dans laquelle l'extraction locale est autorisée.
\end{itemize}

Le mode \texttt{equilibrium\_flux} a également conduit à une correction d'ordre d'application. Dans un réservoir hard inlet, le refill consomme des slots \texttt{Inactive}. Si l'extraction outlet équilibrante est effectuée après le refill, le réservoir peut s'épuiser avant que l'outlet n'ait libéré les slots nécessaires. Le correctif 0293 applique donc l'extraction outlet d'équilibre ou forcée avant le refill hard inlet, puis reconstruit le pool disponible. L'ordre pertinent devient
\begin{equation}
  \text{streaming/frontières}
  \longrightarrow
  \text{extraction outlet}
  \longrightarrow
  \text{refill hard inlet}
  \longrightarrow
  \text{collision SRC}
  \longrightarrow
  \text{thermostat}.
\end{equation}
Cette modification ne transforme pas \texttt{equilibrium\_flux} en condition de masse globale parfaite : si la couche outlet ne contient pas assez de particules, l'extraction reste partielle et la masse totale peut augmenter. Elle rend en revanche le recyclage du pool inactif cohérent lorsque l'outlet dispose localement de particules à extraire.

La gestion du thermostat a été clarifiée en parallèle. Le paramètre \texttt{thermostatEnable} est le commutateur physique unique. Lorsqu'il vaut \texttt{false}, le thermostat est désactivé en CPU comme en CUDA, même si des variables d'environnement CUDA sélectionnent des chemins rapides. Lorsqu'il vaut \texttt{true}, les variables d'environnement ne choisissent que le backend et le mode de résidence. Cette règle est importante pour les simulations d'aspiration : \texttt{forced\_flux} avec thermostat actif représente une extraction isotherme, tandis que \texttt{forced\_flux} sans thermostat laisse évoluer l'énergie du gaz restant.

Enfin, les messages d'épuisement du réservoir ont été rendus explicites. Lorsque le pool \texttt{Inactive} ne suffit plus, le code signale maintenant l'étape concernée, le type de réservoir, le fait que l'append GPU dynamique est désactivé, puis les compteurs utiles. Ce choix garde le chemin résident déterministe : les longs runs doivent soit préallouer un pool suffisant, soit régler les flux et l'épaisseur outlet de manière à recycler assez de particules.

\subsection{Validation consolidée 0281}

Le validateur 0281 regroupe les validations thermostat CUDA boundary-aware. Il exécute cinq familles, chacune sur deux tailles de grille, soit dix lignes de synthèse. Toutes les lignes ont passé la comparaison avec zéro métrique échouée sur 76 métriques comparées. La température effective après thermostat est ramenée à la cible \(10^{-3}\) dans tous les cas.


\begin{center}
\small
\begin{tabular}{p{0.27\linewidth}p{0.28\linewidth}p{0.14\linewidth}p{0.10\linewidth}p{0.11\linewidth}}
\toprule
Validateur & Cas & Grille & Statut & Speedup \\
\midrule
wall 0276 & Poiseuille wall & 64x64 s300 & PASS & 0.29 \\
wall 0276 & Poiseuille wall & 128x128 s300 & PASS & 0.42 \\
solid 0277 & Rectangle périodique & 64x64 s300 & PASS & 0.63 \\
solid 0277 & Rectangle périodique & 128x128 s300 & PASS & 0.89 \\
piston 0278 & Piston legacy & 64x64 s300 & PASS & 0.55 \\
piston 0278 & Piston legacy & 128x128 s300 & PASS & 0.72 \\
full-face IO & Obstacle ouvert & 64x64 s300 & PASS & 1.72 \\
full-face IO & Obstacle ouvert & 128x128 s300 & PASS & 3.14 \\
segmenté IO & U-turn segmenté & 64x64 s300 & PASS & 2.55 \\
segmenté IO & U-turn segmenté & 128x128 s300 & PASS & 5.35 \\
\bottomrule
\end{tabular}
\end{center}

Ces speedups doivent être interprétés avec prudence. Les runs wall/solid/piston sont courts, et le coût fixe de lancement, synchronisation, génération d'état et comparaison peut dominer. L'information essentielle est donc d'abord la non-régression physique : mêmes métriques que le chemin CPU de référence et thermostat à la cible. Les gains deviennent nets pour les inlet/outlet résidents, où la résidence GPU et l'état partagé évitent plus de transferts.

\subsection{Politique de garde et conséquences pour les futurs développements}

L'état atteint peut être formulé ainsi : la branche \texttt{SRC\_GPU} fournit maintenant un code \textbf{SRC/MPCD classic full CUDA résident} pour les cas classic-only validés, incluant advection/streaming, décalage de grille, collision/rotation SRC, conditions limites avancées et thermostat. Le chemin OpenMP reste disponible explicitement et sert de référence de validation.

La fermeture liquide reste volontairement séparée. Lorsque \texttt{projectionEnable}, \texttt{resamplingEnable}, la réponse de capacité ou le viriel sont actifs, le code ne doit pas forcer un raccourci CUDA classic-only qui contournerait ces modules. Le bon invariant est :
\begin{itemize}[nosep]
  \item classic-only : chemin résident CUDA fusionné ;
  \item Q6/resampling/viriel : étapes CPU/OpenMP ou hybrides conservées, puis thermostat CUDA post-CPU si activé ;
  \item futur Q6 CUDA : chantier séparé portant sur le solveur elliptique, les opérateurs divergence/gradient, les masques de frontière et la synchronisation avec les autres modules.
\end{itemize}

Ce découplage est une propriété d'architecture, pas une limitation accidentelle. Il permet d'utiliser immédiatement un SRC classic CUDA physique et rapide pour les cas validés, tout en gardant une trajectoire claire vers une fermeture liquide entièrement accélérée.



\section{Resampling CUDA post-SRC : implémentation, paramètres et validation}

\subsection{Principe physique et ordre d'exécution}

Le chantier CUDA du resampling a conduit à préciser un point important d'interprétation. Le resampling n'est pas un processus physique supplémentaire placé avant la collision ; il ne doit pas être lu comme une création ou destruction de particules matérielles. La dynamique physique validée du SRC/MPCD reste celle de l'advection, des conditions limites, du décalage aléatoire de grille et de la collision/rotation SRC. Le resampling CUDA est donc placé \emph{après} le SRC classic, comme dans le développement MATLAB et OpenMP de la méthode pondérée.

L'ordre nominal retenu est :
\begin{equation}
\begin{aligned}
&\text{streaming CUDA} \\
&\quad \to \text{conditions limites CUDA : murs, solides, inlet/outlet, piston} \\
&\quad \to \text{collision SRC CUDA avec shift collisionnel interne} \\
&\quad \to \text{thermostat CUDA si } \texttt{thermostatEnable=true} \\
&\quad \to \text{Q6 CPU/OpenMP éventuel si } \texttt{projectionEnable=true} \\
&\quad \to \text{resampling CUDA post-SRC : survey, reconditionnement, population guard}.
\end{aligned}
\end{equation}

Ce choix garantit que le resampling maintient la validité statistique de l'état fluide produit par SRC sans redéfinir les collisions physiques. Le shift reste attaché à l'opérateur de collision ; le resampling travaille sur la grille physique non shiftée. Pour les domaines ouverts, toute variation physique nette de masse ou d'impulsion doit provenir des conditions inlet/outlet, non du resampling. Dans un domaine fermé, le resampling doit conserver la masse totale, l'impulsion et l'énergie relative locale à la précision numérique.

\subsection{Jalons d'implémentation 0294--0299}

Le développement a été volontairement progressif afin de ne pas casser le chemin SRC classic full CUDA résident validé.

\begin{itemize}[nosep]
  \item \textbf{0294 -- audit.} L'audit a distingué les briques historiques de resampling CUDA déjà présentes, notamment classification poor/rich, plans de transfert, extraction/insertion persistante, rôles \texttt{Fluid}/\texttt{Inactive}/\texttt{Latent}, workspace cellule et état particulaire CUDA partagé. La conclusion d'architecture est que les anciennes briques étaient utiles mais ne formaient pas encore un module post-SRC non destructif.
  \item \textbf{0295 -- survey post-SRC.} Un diagnostic CUDA passif calcule, après SRC, $N_c$, $M_c$, $P_{x,c}$, $P_{y,c}$, $U_{x,c}$, $U_{y,c}$, l'énergie relative et les classes \texttt{empty}/\texttt{poor}/\texttt{target}/\texttt{rich}. Aucun rôle, masse ou vitesse n'est modifié. Les validations strictes TG, Poiseuille, backward step et U-box segmentée sont passées. Von Karman est conservé comme stress test car le thermostat CUDA boundary-aware n'est pas bit-reproductible OFF/OFF dans ce cas, alors que le même cas sans thermostat est reproductible.
  \item \textbf{0296 -- reconditionnement conservatif des masses.} Le module réduit une dispersion excessive des masses particulaires sans changer le support. Dans chaque cellule humide non vide, les masses sont rapprochées de $M_c/N_c$ avec une force \texttt{strength}, puis l'impulsion cellulaire est restaurée. Sur les états à masses uniformes, le module est neutre et produit uniquement son CSV de diagnostic.
  \item \textbf{0297 -- garde local de population.} Le premier module mutateur agit sur les cellules $0<N_c<N_{\min}$ par split de particules représentatives et sur les cellules $N_c>N_{\max}$ par merge local. Le split ne crée pas de fluide : une particule numérique pondérée est remplacée par deux porteurs statistiques dont la masse et l'impulsion totales égalent celles du donneur.
  \item \textbf{0298 -- restauration de l'énergie relative.} Après split/merge, la masse et l'impulsion locales sont conservées, puis les vitesses relatives à $U_c$ peuvent être rescalées pour restaurer $K_{\mathrm{rel},c}=\frac12\sum_i m_i\lVert v_i-U_c\rVert^2$.
  \item \textbf{0299 -- filtrage boundary-aware.} Les cellules candidates sont filtrées près des ouvertures et, sur demande, près des parois ou solides. Le halo ouvert vaut 1 cellule par défaut, afin de ne pas mélanger le contrôle de support interne avec les réservoirs hard inlet/outlet.
\end{itemize}

\subsection{Formulation locale du split, du merge et de la restauration}

Pour une cellule pauvre, le split minimal sélectionne un donneur fluide local de masse $m_d$, position $x_d$ et vitesse $v_d$, puis active un slot \texttt{Inactive} :
\begin{equation}
  m_d' = m_d - \delta m, \qquad
  m_{\mathrm{new}} = \delta m, \qquad
  x_{\mathrm{new}}=x_d, \qquad
  v_{\mathrm{new}}=v_d.
\end{equation}
La masse, l'impulsion et l'énergie portée par la particule donneuse sont alors conservées exactement avant la restauration globale de cellule. Pour une cellule riche, deux particules locales sont fusionnées par moyenne pondérée :
\begin{equation}
  m' = m_1+m_2, \qquad
  v' = \frac{m_1 v_1 + m_2 v_2}{m_1+m_2}.
\end{equation}
Cette opération conserve masse et impulsion mais peut réduire l'énergie relative. C'est la raison du module 0298 : après mutation, le code compare $K_{\mathrm{rel},c}^{\mathrm{before}}$ et $K_{\mathrm{rel},c}^{\mathrm{after}}$, puis applique si nécessaire
\begin{equation}
  v_i \leftarrow U_c + s_c\,(v_i-U_c), \qquad
  s_c = \sqrt{\frac{K_{\mathrm{rel},c}^{\mathrm{target}}}{K_{\mathrm{rel},c}^{\mathrm{current}}}},
\end{equation}
avec un bornage de $s_c$ par \texttt{MPCD\_CUDA\_RESAMPLING\_MOMENT\_RESTORE\_0298\_MAX\_SCALE}. L'opération est donc un remaillage statistique conservatif, non une correction cachée du champ de vitesse.

\subsection{Options utilisateur et variables d'environnement spécifiques 0295--0314}

Les listings de paramètres 0292 fournis avec ce rapport décrivent le socle CUDA 0286/0292 : backend, collisions résidentes, thermostat, inlet/outlet, solides, Q6 préservé et paramètres \texttt{params.kv}. Les options 0295--0303 ci-dessous sont des compléments spécifiques au resampling CUDA post-SRC et aux scripts de démonstration.


\begin{itemize}[nosep]
  \item \path{MPCD_CUDA_RESAMPLING_SUPPORT_SURVEY_0295} : active le survey post-SRC non-mutant ; cadence via \path{MPCD_CUDA_RESAMPLING_SUPPORT_SURVEY_0295_EVERY}.
  \item \path{MPCD_CUDA_RESAMPLING_MASS_RECONDITION_0296} : active le reconditionnement conservatif des masses ; force via \path{MPCD_CUDA_RESAMPLING_MASS_RECONDITION_0296_STRENGTH}.
  \item \path{MPCD_CUDA_RESAMPLING_POPULATION_GUARD_0297} : active le split/merge local post-SRC ; cadence via \path{MPCD_CUDA_RESAMPLING_POPULATION_GUARD_0297_EVERY}.
  \item \path{MPCD_CUDA_RESAMPLING_POPULATION_GUARD_0297_NMIN}, \path{..._NTARGET}, \path{..._NMAX} : triplet de population. Le réglage nominal provisoire est \texttt{12:20:32}.
  \item \path{MPCD_CUDA_RESAMPLING_MOMENT_RESTORE_0298} : restaure l'énergie relative cellulaire ; bornage via \path{MPCD_CUDA_RESAMPLING_MOMENT_RESTORE_0298_MAX_SCALE}.
  \item \path{MPCD_CUDA_RESAMPLING_POPULATION_GUARD_0299_BOUNDARY_AWARE} : active le filtrage boundary-aware ; l'exclusion des ouvertures est contrôlée par \path{MPCD_CUDA_RESAMPLING_POPULATION_GUARD_0299_OPEN_BOUNDARY_HALO_CELLS}.
\end{itemize}


Les scripts de démonstration ajoutés ou mis à jour en 0303 utilisent des alias plus courts afin de faciliter les comparaisons visuelles : \texttt{RESAMPLING\_ENABLE}, \texttt{RESAMPLING\_SURVEY\_ENABLE}, \texttt{GUARD\_NMIN}, \texttt{GUARD\_NTARGET}, \texttt{GUARD\_NMAX}, \texttt{GUARD\_EVERY}, \texttt{RESTORE\_ENABLE}, \texttt{BOUNDARY\_AWARE} et \texttt{OPEN\_BOUNDARY\_HALO\_CELLS}. Ces variables ne sont pas des clés \texttt{params.kv} ; elles configurent les scripts et exportent les variables \texttt{MPCD\_CUDA\_...} appropriées.

\subsection{Sorties et diagnostics produits}

Le module resampling CUDA ajoute plusieurs sorties légères, placées dans le \texttt{outputDir} du run ou dans \texttt{dev\_history/artifacts/} pour les campagnes de validation.

\begin{table}[htbp]
\centering
\footnotesize
\begin{tabular}{p{0.31\linewidth}p{0.59\linewidth}}
\toprule
Fichier & Contenu principal \\
\midrule
\path{cuda_resampling_support_survey_0295.csv} & Population et moments cellulaires post-SRC : \texttt{emptyCells}, \texttt{poorCells}, \texttt{richCells}, \texttt{targetBandCells}, \texttt{meanNActive}, \texttt{stdNActive}, \texttt{minNWet}, \texttt{maxNWet}, \texttt{massRelRmsWet}, $K_{\mathrm{rel}}$ et $k_BT$ pondéré. \\
\path{cuda_resampling_mass_recondition_0296.csv} & Nombre de cellules et particules traitées, cellules sautées ou bornées, erreur de masse et d'impulsion après reconditionnement. \\
\path{cuda_resampling_population_guard_0297.csv} & Compteurs \texttt{splitApplied}, \texttt{mergeApplied}, erreurs maximales de masse, impulsion et $K_{\mathrm{rel}}$, budgets avant/après, exclusions boundary-aware 0299. \\
\path{cuda_resampling_active_sweep_0300_*.csv} & Résumé par run, comparaison au classic et manifest de la campagne active TG/Poiseuille/step/U-box. \\
\path{cuda_resampling_backward_step_long_0301_*.csv} & Séries temporelles et comparaisons longues sur marche arrière pour plusieurs vitesses d'entrée et triplets $N_{\min}/N_{\mathrm{target}}/N_{\max}$. \\
\bottomrule
\end{tabular}
\caption{Diagnostics ajoutés par la validation resampling CUDA 0295--0301.}
\end{table}

\subsection{Scripts de démonstration et post-traitement}

Les quatre scripts de démonstration CUDA historiques ont été rendus compatibles resampling en 0303 : Taylor--Green forcé, Poiseuille périodique forcé, backward step inlet/outlet et U-box segmentée. Ils peuvent être lancés avec \texttt{RESAMPLING\_ENABLE=0} pour obtenir une référence SRC classic avec survey passif, ou avec \texttt{RESAMPLING\_ENABLE=1} pour activer le triplet nominal \texttt{12:20:32}. Les répertoires sont séparés en \texttt{classic\_survey/} et \texttt{resampling\_guard\_nmin12\_nt20\_nmax32/} afin de permettre une comparaison directe par les outils MATLAB/post-traitement existants.

Un script Von Karman séparé, \texttt{run\_demo\_src\_resampling\_cuda\_von\_karman\_cylinder\_0303.sh}, a été ajouté pour ne pas interférer avec le script local \texttt{0285} utilisé pour les essais de développement. Le cas Von Karman avec thermostat est conservé comme stress test physique ; pour les validations bit-à-bit strictes, le même cas est lancé avec thermostat désactivé car le thermostat CUDA boundary-aware présente une non-reproductibilité OFF/OFF mesurable sur ce cas complexe.

\subsection{Validation courte 0300 : quatre cas reproductibles}

La campagne 0300 a comparé \texttt{classic}, \texttt{mass0296} et trois gardes actifs sur quatre cas courts : Taylor--Green périodique, Poiseuille wall, backward step et U-box segmentée. Tous les runs se sont terminés avec \texttt{exitCode=0}. TG et Poiseuille restent essentiellement passifs, ce qui vérifie l'absence d'effet parasite sur des supports déjà bien conditionnés. Le backward step déclenche le guard et fournit le premier cas actif court ; la U-box courte reste peu sollicitante mais confirme la compatibilité avec les ouvertures segmentées.

\begin{table}[htbp]
\centering
\footnotesize
\begin{tabular}{p{0.23\linewidth}rrrrr}
\toprule
Cas 0300 & split & merge & poor & rich & erreur max. $K_{\mathrm{rel}}$ \\
\midrule
TG, \texttt{12:20:32} & 0 & 0 & 1 & 0 & $1.4\times10^{-17}$ \\
Poiseuille, \texttt{12:20:32} & 0 & 0 & 0 & 0 & $1.4\times10^{-17}$ \\
Backward step, \texttt{12:20:32} & 17 & 0 & 17 & 0 & $2.1\times10^{-17}$ \\
U-box segmentée, \texttt{12:20:32} & 0 & 0 & 0 & 0 & $1.4\times10^{-17}$ \\
\bottomrule
\end{tabular}
\caption{Synthèse du sweep actif court 0300 pour le réglage nominal provisoire \texttt{12:20:32}.}
\end{table}

\subsection{Validation longue 0301--0302 : marche arrière}

Le backward step a été choisi comme banc d'essai principal parce que les simulations classic perdent rapidement des particules derrière la marche lorsque la vitesse d'entrée augmente. Sur une grille modérée $96\times48$, avec $3000$ pas et $U_{\mathrm{in}}=0.30$, $0.45$ puis $0.60$, la référence classic développe respectivement $21$, $98$ et $191$ cellules humides vides en fin de run. Le réglage nominal \texttt{12:20:32} supprime ces cellules vides finales dans les trois cas.

\begin{table}[htbp]
\centering
\footnotesize
\begin{tabular}{c r r r r r r}
\toprule
$U_{\mathrm{in}}$ & empty classic & empty guard & poor classic & poor guard & stdN classic & stdN guard \\
\midrule
0.30 & 21 & 0 & 543 & 370 & 8.05 & 5.18 \\
0.45 & 98 & 0 & 845 & 632 & 10.56 & 6.37 \\
0.60 & 191 & 0 & 1074 & 775 & 12.81 & 7.33 \\
\bottomrule
\end{tabular}
\caption{Effet du resampling CUDA nominal \texttt{12:20:32} sur le support particulaire du backward step long 0301.}
\end{table}

Pour $U_{\mathrm{in}}=0.60$, le réglage nominal effectue $10832$ splits et $7759$ merges cumulés. Les erreurs maximales observées restent au bruit numérique : masse $3.6\times10^{-15}$, impulsion $1.8\times10^{-14}$ et énergie relative restaurée $2.8\times10^{-13}$. Le moment longitudinal final diffère du classic d'environ $142$ sur $P_x\simeq1.57\times10^4$, soit moins de un pour cent. Cette différence est acceptable pour un run actif où le support particulaire est volontairement reconditionné et où le classic souffre précisément de cellules vides persistantes.

La conclusion expérimentale de 0301--0302 est donc que \texttt{12:20:32} constitue le meilleur réglage nominal provisoire. Le triplet \texttt{10:20:34} est plus prudent mais moins correcteur ; \texttt{14:20:30} réduit davantage les cellules pauvres mais devient plus intrusif en nombre de mutations et en décalage global du moment. Le réglage nominal retenu est suffisamment actif pour supprimer les vides derrière la marche, tout en gardant les budgets $M/P/K_{\mathrm{rel}}$ propres.


\subsection{Diagnostics 0304--0306 : flags de support, classification géométrique et outliers de vitesse}

Après la validation 0301--0302, les visualisations longues ont montré que le seul comptage des cellules vides ne suffisait pas à qualifier la qualité du support. Certaines cellules, notamment près des parois ou des solides, pouvaient rester peuplées tout en présentant des vitesses cellulaires aberrantes. Cette observation a conduit à distinguer trois défauts différents : la cellule vide ($N_c=0$), la cellule pauvre mais non vide, et la cellule bien peuplée en nombre mais mal représentée du point de vue des masses ou des vitesses.

Le jalon 0304 ajoute un flag passif post-SRC/post-thermostat sur la grille physique non shiftée. Ce flag ne déclenche encore aucune correction ; il compte seulement les cellules humides vides, les cellules sous un seuil $N_c\leq N_{\mathrm{trigger}}$, les extrema de population et les coûts du diagnostic. L'intérêt de cette étape est double. D'une part, elle montre que le déclenchement adaptatif futur peut être fondé sur des informations déjà disponibles après le dépôt cellulaire, sans nouveau survey lourd. D'autre part, elle prépare un mécanisme commun qui pourra plus tard déclencher soit un guard anticipé, soit un refill de cellule vide.

Le jalon 0305 ajoute une classification géométrique des cellules problématiques : bulk, wall-adjacent, solid-adjacent, open-boundary-adjacent et corner-adjacent. Cette classification a permis de clarifier l'interprétation des cellules vides. Dans le cas Taylor--Green périodique avec trou initial, les cellules vides sont bien classées comme bulk et le resampling les résorbe par advection et correction de support. En revanche, les essais de sillage et de marche ont montré que les zones proches paroi ou proches solide exigent un traitement plus prudent que le bulk.

Le jalon 0306 ajoute enfin un diagnostic d'outliers de vitesse par classe de population et par classe géométrique. Les cellules sont réparties en bins $N=1$, $N=2$, $N=3$, $4\leq N\leq 6$, $7\leq N<N_{\min}$ et $N\geq N_{\min}$. Pour chaque bin, le code suit les cellules dont $\lVert U_c\rVert$ dépasse un seuil utilisateur. Un fichier de pires cellules garde, pour chaque pas diagnostiqué, les coordonnées, $N_c$, $M_c$, $U_x$, $U_y$, $\lVert U_c\rVert$, $K_{\mathrm{rel}}$, $k_BT$ et les indicateurs géométriques. Ce diagnostic a montré que les vitesses aberrantes observées avant 0307 n'étaient pas d'abord des excès thermiques : dans les pires cas, $K_{\mathrm{rel}}=0$ et $k_BT=0$, mais $M_c$ est de l'ordre de $10^{-9}$ et $\lVert U_c\rVert$ peut atteindre plusieurs dizaines. Un rethermostatage relatif n'aurait donc pas corrigé cette pathologie, car il conserve la vitesse moyenne cellulaire.

\subsection{Split-safety 0307--0308 : prévention des cascades de masse faible}

Le diagnostic 0306 a conduit à identifier un mécanisme fondamental : le split local conservatif peut être réappliqué plusieurs fois à des particules déjà allégées. Dans une cellule structurellement pauvre, par exemple solid-adjacent ou proche d'une arête de marche, la règle initiale ``si $0<N_c<N_{\min}$, splitter localement'' peut donc produire une cascade géométrique
\begin{equation}
  m=1 \to \frac12 \to \frac14 \to \frac18 \to \cdots,
\end{equation}
qui finit par créer des particules quasi sans masse. Ces particules ne pèsent presque rien dans les bilans globaux, mais elles peuvent porter une vitesse cellulaire énorme par la formule $U_c=P_c/M_c$ lorsque $M_c$ devient dégénéré.

Le jalon 0307 introduit une sécurité de split activable : choix préférentiel du donneur le plus massif, masse minimale du donneur, masse minimale de la nouvelle particule, masse minimale restante après split, et modes de prudence pour les cellules solid-adjacent. Le mode nominal provisoire conserve le split près des solides mais interdit les fragments sous-massifs :
\begin{equation}
  m_{\mathrm{donor}} \geq 0.5, \qquad
  m_{\mathrm{new}} \geq 0.25, \qquad
  m_{\mathrm{donor}}' \geq 0.25.
\end{equation}
L'objectif n'est pas de corriger a posteriori les particules faibles masses, mais d'empêcher leur création.

Les résultats 0307 sont très discriminants. Sur Von Karman, le mode diagnostic sans prévention observe $88863$ splits, dont $87666$ proviennent de donneurs de masse inférieure à $0.5$ et $86758$ de donneurs de masse inférieure à $0.1$ ; la masse minimale de donneur atteint $2.6\times10^{-13}$. Avec le mode \texttt{safe\_floor}, le nombre de splits chute à $1126$ et aucun split ne provient d'un donneur sous $0.5$. Sur backward step, la masse minimale de nouvelle particule passe d'environ $3.3\times10^{-3}$ à $0.2501$. Ces résultats montrent que la très forte activité du guard avant 0307 était en partie une activité pathologique de resplit de fragments, non un vrai contrôle utile du support.

\begin{table}[htbp]
\centering
\footnotesize
\begin{tabular}{p{0.22\linewidth}p{0.20\linewidth}rrrr}
\toprule
Cas & Mode & splits & donneurs $<0.5$ & donneurs $<0.1$ & $m_{\min}^{\mathrm{new}}$ \\
\midrule
Backward step & diagnostic & 3320 & 1605 & 433 & $3.3\times10^{-3}$ \\
Backward step & safe floor & 2008 & 0 & 0 & $2.50\times10^{-1}$ \\
Von Karman & diagnostic & 88863 & 87666 & 86758 & $2.3\times10^{-14}$ \\
Von Karman & safe floor & 1126 & 0 & 0 & $2.50\times10^{-1}$ \\
\bottomrule
\end{tabular}
\caption{Effet du split-safety 0307 sur les cascades de masse faible. Les chiffres sont ceux des runs de validation 0307 ; ils montrent que le mode \texttt{safe\_floor} coupe la cascade de splits sans désactiver globalement le resampling.}
\end{table}

Le jalon 0308 consolide ce comportement comme mode nominal des scripts de resampling CUDA. La validation stricte 0308, avec complétion contrôlée jusqu'à $6000$ pas, donne quatre runs PASS : backward step classic, backward step resampling, Von Karman classic et Von Karman resampling. Les outliers de vitesse disparaissent : \texttt{sumHighUAllBins=0} dans tous les cas, et les pires cellules ont désormais des masses normales. Par exemple, pour backward step resampling, la pire cellule a $N=20$, $M\simeq19.99$ et $\lVert U\rVert\simeq0.67$ ; pour Von Karman resampling, elle a $N=14$, $M\simeq13.25$ et $\lVert U\rVert\simeq0.49$. La cause principale des vitesses folles était donc la dégénérescence de masse par split répété, non le thermostat.

\subsection{Scripts de visualisation 0309 et diagnostic particulaire MATLAB}

Le jalon 0309 ajoute des scripts de démonstration visuelle indépendants pour Poiseuille, backward step et Von Karman. Chaque cas produit deux répertoires distincts : une référence SRC classic et un run \texttt{resampling\_split\_safe}. Les scripts activent par défaut le mode 0307 consolidé et conservent le survey de support afin de permettre une comparaison quantitative et visuelle.

Les outils MATLAB de visualisation ont également été corrigés pour éviter une confusion importante. Dans les états \texttt{.smpcd}, les slots \texttt{Inactive} gardent leurs positions ; il est donc possible qu'une zone ait $N=0$ en particules fluides tout en affichant des points si le visualiseur trace tous les slots. Les fonctions de lecture/affichage distinguent désormais les rôles \texttt{Fluid}, \texttt{Inactive} et \texttt{Latent}, et peuvent colorer les particules par rôle, masse ou vitesse. Des filtres de seuils sur la masse et sur la vitesse permettent de suivre les particules suspectes sans surcharger la figure. Ces outils sont devenus essentiels pour vérifier que le mode split-safe supprime les particules quasi sans masse dans les sillages longs.

\subsection{Coût des slots inactifs, fast path et dumps compacts 0310--0314}

Les runs longs avec hard inlet et resampling ont montré qu'un grand réservoir \texttt{Inactive} est nécessaire pour éviter l'épuisement des slots lorsque l'inlet injecte un flux net. En revanche, il ne faut pas que les kernels physiques ou les dumps paient systématiquement le coût de la capacité totale. L'audit 0310 a d'abord révélé un défaut de script : certains scripts de démonstration forçaient accidentellement un réservoir de cinq millions de slots inactifs. Cette valeur a été supprimée en 0311.

Le jalon 0312 introduit un audit indépendant des scripts de démonstration. Il génère directement l'état initial et le fichier \texttt{params.kv}, puis lance le binaire sans dépendre des scripts modifiés pour la visualisation. Sur backward step $96\times48$, $400$ pas et \texttt{GUARD\_EVERY=20}, l'audit valide que \texttt{inactiveLogged} suit bien \texttt{inactiveSlotsRequested}. Le temps augmente avec la capacité inactive : $3.06$ s pour $8000$ slots, $3.45$ s pour $50000$ slots et $3.80$ s pour $100000$ slots. Le coût est donc réel, même s'il reste modéré dans cette configuration.

Le jalon 0313 ajoute un fast path de pool inactif en queue de tableau. Lorsqu'une grande réserve inactive est préallouée en fin d'état particulaire, il est inutile de scanner toute la capacité pour trouver des slots disponibles. Le code scanne d'abord une fenêtre bornée dans la queue inactive ; si elle suffit, l'ancien scan complet est évité, sinon un fallback exact conserve la robustesse. Ce fast path concerne le hard inlet full-face, l'inlet segmenté et le population guard lors des splits. Il ne modifie ni la collision SRC, ni le thermostat, ni la conservation masse/impulsion/énergie du resampling.

Le jalon 0314 traite l'autre source de coût : les summaries et dumps. Deux nouvelles clés \texttt{params.kv} contrôlent le rôle à écrire ou à résumer :
\begin{equation}
  \texttt{summaryRoleFilter} \in \{\texttt{all},\texttt{fluid}\}, \qquad
  \texttt{dumpRoleFilter} \in \{\texttt{all},\texttt{fluid}\}.
\end{equation}
Avec \texttt{fluid}, le code produit des états compacts destinés à la visualisation et au post-traitement : seuls les slots fluides sont téléchargés et écrits. Ce mode réduit le coût des dumps et évite de charger inutilement des centaines de milliers de slots inactifs dans MATLAB. Il n'est pas restart-compatible au sens strict, car le réservoir \texttt{Inactive} n'est plus présent dans le dump ; les dumps complets restent disponibles avec \texttt{dumpRoleFilter=all}.

\subsection{Paramètres additionnels 0304--0314}

Les inventaires CSV accompagnant cette version du rapport ont été consolidés pour le socle historique puis prolongés jusqu'au cycle 0493x9r pour les paramètres et gates capillaires, avec les principaux alias des runners x9s. Ils incluent les variables d'environnement du code, les clés \texttt{params.kv} de Q6 multi-espèces/Q6-g-f, les sémantiques x8 de profils segmentés et de sortie Neumann, ainsi que les contrôles de courbure, tension superficielle et mouillage x9. Les paramètres les plus importants sont listés ci-dessous.

\begin{itemize}[nosep]
  \item \path{MPCD_CUDA_RESAMPLING_ADAPTIVE_FLAG_0304} et \path{..._EVERY} : diagnostic post-SRC des cellules low-N ou vides sur grille physique non shiftée.
  \item \path{MPCD_CUDA_RESAMPLING_GEOMETRY_DIAG_0305_HIGH_U} : seuil de vitesse cellulaire utilisé par la classification géométrique 0305.
  \item \path{MPCD_CUDA_RESAMPLING_OUTLIER_0306_U_THRESHOLD} : seuil d'outlier de vitesse utilisé dans les worst-cells 0306.
  \item \path{MPCD_CUDA_RESAMPLING_SPLIT_SAFETY_0307} : active la prévention des cascades de split ; le mode nominal combine \path{MPCD_CUDA_RESAMPLING_SPLIT_PREFER_MAX_MASS_DONOR_0307=1}, \path{MPCD_CUDA_RESAMPLING_SPLIT_DONOR_MIN_MASS_0307=0.5} et \path{MPCD_CUDA_RESAMPLING_SPLIT_NEW_PARTICLE_MIN_MASS_0307=0.25}.
  \item \path{MPCD_CUDA_RESAMPLING_SOLID_ADJACENT_SPLIT_MODE_0307} : mode de prudence pour les cellules solid-adjacent. Le défaut nominal reste le mode normal, car le plancher de masse suffit à supprimer la pathologie sans désactiver la correction près des solides.
  \item \path{MPCD_CUDA_INACTIVE_TAIL_POOL_0313} : active le fast path de recherche de slots inactifs en queue de tableau ; les paramètres \path{..._MIN_SCAN_0313}, \path{..._MAX_SCAN_0313} et \path{..._SCAN_MULT_0313} règlent la fenêtre de scan.
  \item \path{summaryRoleFilter} et \path{dumpRoleFilter} : clés \texttt{params.kv} 0314 contrôlant le mode complet ou compact des summaries et dumps.
\end{itemize}

\subsection{Position par rapport à Q6 et aux futurs développements}

Le resampling CUDA et Q6 traitent deux problèmes distincts. Le premier modifie la représentation particulaire afin de maintenir un support statistique local ; le second corrige soit la composante compressive du champ barycentrique dans les modes historiques \texttt{common}/\texttt{weighted}, soit des champs masqués propres à chaque espèce dans le mode \texttt{independent\_masked}, sans changer la population, les masses ou les types. Depuis les jalons 0490--0493w8, ces deux briques disposent d'une extension multi-espèces et peuvent être activées indépendamment. La section suivante remplace donc l'ancien statut, dans lequel Q6 multi-espèces était encore présenté comme un chantier futur.


\section{Traitement multi-espèces}

Les cycles 0490--0493 ont introduit deux extensions orthogonales. Le cycle 0490a--0490p rend le resampling conservatif par espèce et résident CUDA dans son domaine de qualification. Le cycle 0491 conserve les modes Q6 barycentriques \texttt{common} et \texttt{weighted}. Le cycle 0493w5--0493w8 ajoute le mode \texttt{independent\_masked}, dans lequel chaque espèce de force Q6 strictement positive possède son propre support, son propre second membre, son propre solveur masqué et sa propre correction particulaire. Cette séparation préserve les quatre chemins de simulation et rend explicite le choix de l'opérateur Q6 :

\begin{table}[htbp]
\centering
\small
\caption{Chemins multi-espèces et opérateurs Q6 conservés dans l'architecture 0490--0493w8.}
\begin{tabular}{@{}p{0.20\textwidth}ccp{0.53\textwidth}@{}}
\toprule
Chemin & Resampling & Q6 & Contrat principal \\
\midrule
SRC & non & non & Le \texttt{type} est transporté sans conversion ; la collision SRC couple les populations présentes dans une même cellule sans dépendre de l'étiquette de type. \\
SRC+resampling & oui & non & Les opérations de population, les transferts et la fermeture de masse sont conservatifs par espèce. \\
SRC+Q6 & non & oui & \texttt{common}/\texttt{weighted} utilisent une projection barycentrique du mélange ; \texttt{independent\_masked} utilise une projection masquée distincte pour chaque espèce de force positive. \\
SRC+resampling+Q6 & oui & oui & Les deux briques partagent l'état CUDA ; Q6 précède thermostat et resampling, et les dépôts cellule--espèce sont invalidés et reconstruits aux points d'autorité nécessaires. \\
\bottomrule
\end{tabular}
\end{table}

Le multi-espèces considéré ici est un modèle de transport et de couplage cinématique. Il n'introduit pas encore une équation d'état propre à chaque phase, une diffusion Maxwell--Stefan explicite, une réaction chimique, une tension superficielle ou une règle de collision dépendant de la paire d'espèces. Ces compléments pourront réutiliser les dépôts et les diagnostics décrits ci-dessous, mais ils ne doivent pas être confondus avec le contrat actuellement validé.

\subsection{État commun, registre d'espèces et dépôts cellulaires}

\subsubsection{Identité physique et état algorithmique}

Le champ particulaire \texttt{role} décrit l'état algorithmique d'un slot : \texttt{Fluid} pour une particule active, \texttt{Inactive} pour un slot libre du pool et \texttt{Latent} pour une particule préallouée mais encore hors dynamique. Le champ \texttt{type} décrit au contraire l'identité physique transportée. Une opération de resampling peut modifier le \texttt{role}, mais ni le resampling, ni Q6, ni le thermostat ne doivent convertir implicitement une espèce en une autre.

Le registre est activé par
\begin{equation}
  \texttt{speciesRegistryEnable=true},
  \qquad \texttt{speciesCount}=N_s,
\end{equation}
et par des déclarations de la forme
\begin{equation}
\begin{aligned}
  \texttt{speciesK} ={}& \texttt{type name phaseFamily}\\
  &\texttt{q6StrengthDeclared massClosureStrengthDeclared}\\
  &[\texttt{referenceCellMassDeclared}].
\end{aligned}
\end{equation}
La famille vaut actuellement \texttt{liquid} ou \texttt{gas}. Le coefficient \texttt{q6StrengthDeclared} intervient dans la distribution Q6 décrite plus loin ; \texttt{massClosureStrengthDeclared} et la masse cellulaire de référence interviennent dans la fermeture du resampling. La clé \texttt{speciesRequireRegisteredTypes} impose que tout slot physique actif ou latent appartienne au registre. Les slots \texttt{Inactive}, qui ne portent pas de masse fluide active, sont exclus de cette exigence.

\subsubsection{Dépôt cellule--espèce partagé}

Pour une cellule physique non décalée $c$ et une espèce enregistrée $s$, le dépôt reconstruit
\begin{align}
  N_{c,s} &= \sum_{i\in c}\mathbf{1}_{\tau_i=s},\\
  M_{c,s} &= \sum_{i\in c}m_i\mathbf{1}_{\tau_i=s},\\
  \bm{P}_{c,s} &= \sum_{i\in c}m_i\bm{v}_i\mathbf{1}_{\tau_i=s}.
\end{align}
Les grandeurs barycentriques sont
\begin{equation}
  M_c=\sum_sM_{c,s},\qquad
  \bm{P}_c=\sum_s\bm{P}_{c,s},\qquad
  \bm{u}_c=\frac{\bm{P}_c}{M_c},
\end{equation}
et les fractions massiques cellulaires sont
\begin{equation}
  Y_{c,s}=\frac{M_{c,s}}{M_c},
  \qquad \sum_sY_{c,s}=1
\end{equation}
dans toute cellule humide. Pour la politique de resampling, on utilise aussi l'occupation normalisée
\begin{equation}
  \theta_{c,s}=\frac{M_{c,s}}{M_s^{\mathrm{ref}}},
\end{equation}
notation distincte du poids Q6 $w_{c,s}$ introduit plus loin.

Le dépôt CPU déterministe 0490b constitue la référence. Le workspace CUDA 0490h contient notamment les types enregistrés, les coefficients Q6, les masses de référence, les familles de phase, l'activation du resampling par espèce, puis les tableaux denses $N_{c,s}$, $M_{c,s}$, $P_{x,c,s}$, $P_{y,c,s}$, $Y_{c,s}$ et $\theta_{c,s}$. Les allocations sont conservées et réutilisées ; les tableaux denses ne sont téléchargés vers l'hôte que dans les smokes d'équivalence. Cette brique, initialement créée pour le resampling, devient avec 0491 un service partagé accessible même lorsque \texttt{resamplingEnable=false}.

\subsection{Resampling multi-espèces : formulation et implémentation 0490}

\subsubsection{Invariants des opérations discrètes}

Le resampling modifie le nombre de représentants et éventuellement leur masse, mais doit conserver l'identité physique. Ses invariants locaux et globaux sont :
\begin{equation}
  \Delta M_s=0,
  \qquad \Delta\bm{P}_s=\bm{0}
  \qquad \text{pour chaque espèce }s
\end{equation}
dans un domaine fermé, à l'exception des flux explicitement imposés par des ouvertures.

Un split hérite du type du parent. Dans sa forme élémentaire,
\begin{equation}
  (m,\bm{v},s)
  \longrightarrow
  \left(\frac{m}{2},\bm{v},s\right)
  +\left(\frac{m}{2},\bm{v},s\right),
\end{equation}
ce qui conserve exactement masse, quantité de mouvement et énergie cinétique du parent avant déplacement géométrique des filles. Un merge n'est autorisé qu'entre deux particules de même type ; la vitesse du représentant fusionné est choisie barycentriquement,
\begin{equation}
  m'=m_1+m_2,
  \qquad
  \bm{v}'=\frac{m_1\bm{v}_1+m_2\bm{v}_2}{m_1+m_2},
\end{equation}
puis la restauration d'énergie relative existante est appliquée lorsque ce chemin est activé.

Le garde de population conserve la bande globale $N_{\min}/N_{\mathrm{target}}/N_{\max}$, mais répartit la cible entre les espèces présentes. Une méthode déterministe du plus fort reste attribue les représentants, avec au moins un porteur par espèce présente lorsque $N_{\mathrm{target}}$ le permet. Le split est dirigé vers l'espèce la plus déficitaire ; le merge vers une espèce excédentaire possédant au moins deux représentants. Les opérations entre types différents sont interdites, même lorsque leur fusion serait numériquement commode.

\subsubsection{Refill et transferts conservatifs par type}

Le remplissage d'une cellule temporairement vide utilise une mémoire résidente $M^{\mathrm{last}}_{c,s}$. La population recréée est distribuée selon la dernière composition admissible, puis une correction globale par espèce restaure masse et impulsion. Le refill est refusé lorsque la population cible est inférieure au nombre d'espèces à représenter ou lorsqu'une espèce disparue ne possède plus aucun porteur actif permettant une compensation conservative exacte.

Les transferts donneur--receveur sont également planifiés par type. Pour une espèce $s$, une cellule receveuse $r$ ne peut consommer que l'excès disponible chez un donneur $d$ de la même espèce. Le plan contient donc explicitement
\begin{equation}
  (r,d,s,\Delta M_{r\leftarrow d,s}),
\end{equation}
et le matérialiseur rejette tout candidat dont le \texttt{type} ne correspond pas. Une cellule liquide pure ne peut pas être alimentée par un donneur gazeux sous prétexte qu'il est plus proche ou plus riche.

\subsubsection{Fermeture de masse dépendante de la composition}

Pour une cellule non vide, on définit
\begin{equation}
  W_c=\sum_s\frac{M_{c,s}}{M_s^{\mathrm{ref}}},
  \qquad
  M_c^{\star}=\frac{M_c}{W_c}.
\end{equation}
Le coefficient de fermeture est une moyenne d'occupation des coefficients déclarés :
\begin{equation}
  \beta_c=\operatorname{clamp}\!\left[
  \beta_{\mathrm{global}}
  \frac{\sum_s\theta_{c,s}\beta_s}
       {\sum_s\theta_{c,s}},0,1\right],
\end{equation}
et la cible cellulaire devient
\begin{equation}
  M_c^{\mathrm{eff}}
  =M_c+\beta_c\left(M_c^{\star}-M_c\right).
\end{equation}
Une espèce liquide avec $\beta_s\simeq1$ participe fortement à la fermeture ; une espèce gazeuse avec $\beta_s\simeq0$ reste plus libre ; une cellule mixte reçoit un comportement intermédiaire.

La première réalisation 0490d appliquait un facteur commun dans chaque cellule et préservait ainsi exactement la composition locale. Une succession de facteurs différents dans des cellules de compositions variables peut toutefois conserver la masse totale tout en faisant dériver les masses globales par espèce. Le chemin résident ultérieur utilise donc une matrice d'échelles cellule--espèce, équilibrée alternativement sur les contraintes cellulaires et sur les sommes globales par espèce. La dernière normalisation impose $M_s$ comme invariant dur ; un décalage barycentrique commun restaure ensuite la vitesse cellulaire lorsque les facteurs propres aux espèces ont modifié le moment local.

\subsubsection{Chaîne CUDA résidente 0490h--0490p}

Les jalons d'implémentation sont :
\begin{itemize}[nosep]
  \item \textbf{0490a--0490b} : registre d'espèces et dépôt CPU cellule--espèce de référence ;
  \item \textbf{0490c--0490d} : split/merge conservatifs par type et fermeture de masse dépendante de la composition ;
  \item \textbf{0490h} : dépôt $N_{c,s}$, $M_{c,s}$ et $\bm{P}_{c,s}$ résident CUDA ;
  \item \textbf{0490i} : fermeture de masse sur l'état CUDA partagé ;
  \item \textbf{0490j} : sélection d'espèce du garde de population ;
  \item \textbf{0490k} : plan donneur--receveur CUDA ordonné par receveur, espèce puis donneur compatible ;
  \item \textbf{0490l} : validation stricte sans plan ni vecteur d'opérations piloté par le CPU ;
  \item \textbf{0490m} : application directe du plan device, matérialisation et mutation de l'état partagé ;
  \item \textbf{0490n} : dépôts post-opération et pool \texttt{Fluid}/\texttt{Latent}/\texttt{Inactive} résidents ;
  \item \textbf{0490p} : politique wet/poor/rich/target calculée et consommée sur device.
\end{itemize}

La correction 0490n-fix2 distingue les sous-remplissages d'une entrée individuelle, possibles à cause de l'indivisibilité des particules, des sous-remplissages agrégés par couple $(d,s)$. Les premiers sont diagnostiques ; les seconds restent fatals, car un excès provenant d'une autre espèce ou d'un autre donneur ne doit pas masquer un vrai déficit.

Les paramètres structurants sont regroupés par fonction :
\begin{itemize}[nosep]
  \item fermeture : \path{speciesResamplingMassClosureEnable} et son backend CUDA ;
  \item population : \path{speciesResamplingPopulationGuardEnable} et son backend CUDA ;
  \item transferts : \path{speciesResamplingTransferEnable} et son backend CUDA ;
  \item production résidente : fast path, dépôts, pool et maintenance stricte.
\end{itemize}
Les diagnostics principaux sont :
\begin{itemize}[nosep]
  \item \path{species_runtime_0490a.csv} et \path{species_cell_runtime_0490b.csv} ;
  \item \path{species_cell_cuda_equivalence_0490h.csv} ;
  \item \path{cuda_species_mass_closure_0490i.csv} ;
  \item \path{cuda_species_transfer_plan_0490k.csv} ;
  \item \path{cuda_species_resident_fast_path_0490m.csv} ;
  \item \path{cuda_species_resident_maintenance_0490n.csv}.
\end{itemize}

\subsubsection{Statut physique et recommandation d'usage}

Le cycle 0490 établit que le resampling peut être rendu conservatif par espèce et exécuté sans retour complet vers le CPU. Il ne démontre pas qu'il est nécessaire ni bénéfique dans toutes les simulations. Les campagnes ultérieures montrent un coût important lorsque de nombreux couples cellule--espèce sont réparés et une sensibilité possible de l'autodiffusion à la fragmentation et aux transferts, alors que les grandeurs bulk de type Taylor--Green, température et composition restent généralement proches de la référence.

Le resampling multi-espèces ne doit donc pas être le réglage nominal du fluide. Il doit rester une correction \emph{opt-in}, déclenchée pour un défaut de support identifié et comparée à un cas sans resampling. Les critères de décision doivent inclure le nombre de cellules réellement réparées, le coût par pas, la distribution des masses particulaires, la conservation par espèce, l'autodiffusion et l'effet sur les observables physiques ciblées. Les cellules vides, les cellules de coupure géométrique et les ouvertures doivent en outre être traitées selon leur origine physique ; une paroi ou une injection ne peut pas être remplacée par une création locale de masse déguisée.

\subsection{Projection Q6 multi-espèces : modes barycentriques et indépendants}

\subsubsection{Pourquoi le mode \texttt{weighted} ne suffit pas pour liquide incompressible--gaz compressible}

Le cycle 0491 a défini une projection barycentrique unique du mélange. Pour une cellule $c$, Q6 calcule une correction commune $\Delta\bm u^{Q6}_c$, puis le mode \texttt{weighted} la redistribue entre espèces à l'aide de poids $w_{c,s}$ satisfaisant
\begin{equation}
  \sum_s Y_{c,s}w_{c,s}=1,
  \qquad
  \sum_s M_{c,s}\Delta\bm u_{c,s}
  =M_c\Delta\bm u^{Q6}_c.
\end{equation}
Cette identité conserve exactement la correction barycentrique du mélange. Elle constitue une propriété utile de non-régression et le mode \texttt{weighted} est conservé inchangé pour la reproductibilité des campagnes 0491.

Ce contrat n'est toutefois pas adapté lorsque les espèces n'ont pas la même contrainte physique. Dans une cellule pure, la normalisation impose $w=1$ et le repli \texttt{fallback=common} peut encore appliquer la correction commune à une espèce dont la force Q6 déclarée est nulle. Ainsi,
\begin{equation}
  \texttt{q6StrengthDeclared}_s=0
\end{equation}
ne garantissait pas, dans le mode barycentrique, $\Delta\bm u^{Q6}_s=\bm 0$. Cette propriété est rédhibitoire pour le modèle cible liquide incompressible--gaz compressible : le gaz doit pouvoir participer aux collisions SRC et au transport sans recevoir directement la projection d'incompressibilité du liquide.

\subsubsection{Support d'occupation et contrat \texttt{independent\_masked}}

Le cycle 0493w5 introduit
\begin{equation}
  \texttt{speciesQ6Mode=independent\_masked}
\end{equation}
et le seuil runtime
\begin{equation}
  \texttt{speciesQ6MinOccupancyFraction}=\varphi_{\min}.
\end{equation}
Pour l'espèce $s$ dans la cellule $c$, on définit le poids volumique approché
\begin{equation}
  \omega_{c,s}=\frac{M_{c,s}}{M_s^{\mathrm{ref}}},
  \qquad
  \varphi_{c,s}=\frac{\omega_{c,s}}{\sum_r\omega_{c,r}}.
\end{equation}
La fraction massique brute n'est pas utilisée : à grand rapport de densité liquide/gaz, une faible fraction volumique de liquide pourrait dominer artificiellement la masse de la cellule. Le support Q6 de l'espèce est
\begin{equation}
  \mathcal M_{c,s}
  =\mathbf 1_{\alpha_s>0}
   \mathbf 1_{M_{c,s}>0}
   \mathbf 1_{\varphi_{c,s}\geq\varphi_{\min}},
\end{equation}
avec $\alpha_s=\texttt{q6StrengthDeclared}_s$.

Le contrat dur du nouveau mode est
\begin{equation}
  \alpha_s=0
  \quad\Longrightarrow\quad
  \Delta\bm u^{Q6}_s=\bm 0
\end{equation}
au niveau structurel : l'espèce ne construit pas de masque actif, ne lance pas de solveur et n'entre pas dans le kernel d'application. Il n'existe dans ce mode ni redistribution barycentrique, ni \texttt{fallback=common}, ni correction compensatoire imposée à une autre espèce.

\subsubsection{Opérateur masqué par espèce et application atomique}

Pour chaque espèce activée, et toujours depuis le même état particulaire pré-Q6, le code résident :
\begin{enumerate}[nosep]
  \item dépose $M_{c,s}$ et $\bm P_{c,s}$ ;
  \item construit la vitesse propre $\bm u_{c,s}=\bm P_{c,s}/M_{c,s}$ et le masque $\mathcal M_{c,s}$ ;
  \item assemble la divergence de $\bm u_s$ sur ce support ;
  \item résout un système de Poisson masqué distinct ;
  \item stocke la correction de l'espèce sur le device ;
  \item n'applique les corrections aux particules qu'après convergence de tous les solveurs activés.
\end{enumerate}
Cette dernière règle évite l'ordre de passage entre espèces et empêche un état partiellement corrigé si un solveur ultérieur échoue.

Une cellule interne hors support est affectée d'une pression de correction nulle,
\begin{equation}
  \pi_s=0,
\end{equation}
ce qui définit une condition de Dirichlet simple à l'interface du support. Deux gouttes déconnectées forment ainsi deux blocs algébriques indépendants sans recherche explicite de composantes connexes. Cette condition est numériquement bien définie, mais elle ne constitue pas encore une loi physique complète d'interface liquide--gaz : elle ne contient ni pression gazeuse couplée, ni saut capillaire, ni tension superficielle, ni mouillage.

La correction globale uniforme de quantité de mouvement du Q6 historique n'est pas appliquée dans ce mode, car elle modifierait également les espèces déclarées compressibles. L'impulsion directe éventuellement produite par la condition de Dirichlet masquée est donc diagnostiquée explicitement. Pour les comparaisons strictes monoespèces, la qualification 0493w8 désactive la même correction dans le chemin historique par
\begin{equation}
  \texttt{projectionMomentumCorrectionEnable=false}.
\end{equation}

\subsubsection{Conditions limites et résidence CUDA}

Le cycle 0493w7 lève la restriction périodique initiale et fait suivre au solveur d'espèce les topologies déjà acceptées par Q6 CUDA résident :
\begin{itemize}[nosep]
  \item périodique en $x/y$ ;
  \item canal périodique en $x$ avec parois en $y$ ;
  \item entrée/sortie pleine face ;
  \item entrée/sortie segmentée sur des faces autrement imperméables ;
  \item mêmes topologies avec forçage Darcy--Brinkman.
\end{itemize}
Une frontière physique n'est pas assimilée à une cellule voisine où l'espèce serait absente. Les faces périodiques sont raccordées, les parois imposent un flux normal de base nul avec stencil de correction de type Neumann, et les ouvertures utilisent le flux normal prescrit par Q6 résident. Une entrée segmentée typée n'agit que sur l'espèce correspondante. Lorsqu'une seule espèce a une force Q6 positive, elle reçoit le flux ouvert complet ; lorsque plusieurs espèces sont projetées, les flux non typés sont répartis selon l'occupation locale afin de ne pas dupliquer le flux total.

Les exclusions générales du backend résident restent inchangées : domaine mobile ou tronqué, masque de projection de solide immergé non pris en charge, configuration d'ouverture non supportée et réponse de capacité fermée ne sont pas activés implicitement. Darcy reconnaît désormais tous les writers préfixés par \texttt{cuda\_q6\_resident\_0400}, y compris le writer indépendant ; l'état Q6 frais est consommé directement sans ré-upload particulaire.

Le chemin demeure CUDA résident sur toutes les étapes : dépôt cellule--espèce, masque, assemblage, CG, stockage des corrections, application aux particules et redépôt post-application. Seules des réductions scalaires diagnostiques sont copiées vers l'hôte ; aucun état particulaire complet, aucun masque dense et aucun champ cellule--espèce dense n'est matérialisé sur le CPU.

\subsubsection{Diagnostic post-application 0493w6}

Le flux de face auxiliaire annulé par le solveur n'est pas nécessairement identique au champ cellulaire reconstruit après application aux particules. Le cycle 0493w6 distingue donc explicitement
\begin{align}
  D^{\mathrm{face}}_s
  &=\left\|\nabla_h\cdot
    \left(\bm u^{\mathrm{face}}_s+\Delta\bm u^{\mathrm{face}}_s\right)
    \right\|_{2,\mathcal M_s},\\
  D^{\mathrm{cell}}_s
  &=\left\|\nabla_h\cdot\bm u^{\mathrm{redéposée}}_s
    \right\|_{2,\mathcal M_s}.
\end{align}
Après application, masse et quantité de mouvement par espèce sont redéposées sur le GPU avec le masque exact conservé sur le device. Le fichier
\begin{center}
  \path{cuda_species_q6_independent_masked_0493w5.csv}
\end{center}
conserve les colonnes historiques \texttt{divAfterRms} et \texttt{divAfterMaxAbs} pour le flux projeté et ajoute :
\begin{center}
\texttt{divAfterProjectedFaceFluxRms},
\texttt{divAfterProjectedFaceFluxMaxAbs},\\
\texttt{divAfterAppliedCellVelocityRms},
\texttt{divAfterAppliedCellVelocityMaxAbs}.
\end{center}

Cette séparation a révélé une limite précise du stencil actuel sur support interne partiel. Pour deux îlots liquides déconnectés, le flux projeté est annulé à environ $1.2\times10^{-15}$ de la divergence initiale, mais la divergence du champ redéposé reste à environ $0.331$ de sa valeur initiale. Le cas d'ouverture segmentée donne un rapport post-application voisin de $5.37\times10^{-2}$. Le support complet, le canal à parois, l'ouverture pleine face et Darcy atteignent au contraire des rapports de l'ordre de $10^{-12}$ ou moins. Cette limite est localisée à la conversion face--cellule et aux transitions de masque ; elle n'est ni cachée ni confondue avec la convergence du CG.

\subsubsection{Paramètres utilisateur et modes conservés}

Les clés \texttt{params.kv} de la projection multi-espèces sont :
\begin{itemize}[nosep]
  \item \path{speciesQ6Enable} : active la logique sensible aux espèces, sans activer implicitement Q6 ;
  \item \path{speciesQ6Mode} : \texttt{common}, \texttt{weighted} ou \texttt{independent\_masked} ;
  \item \path{speciesQ6Sensitivity} : interpolation $[0,1]$ du mode \texttt{weighted} ;
  \item \path{speciesQ6AlphaEpsilon} et \path{speciesQ6FallbackMode} : robustesse du mode \texttt{weighted} historique ;
  \item \path{speciesQ6ComparisonTolerance} : tolérance des comparaisons de qualification ;
  \item \path{speciesQ6MinOccupancyFraction} : seuil $[0,1]$ du support \texttt{independent\_masked}.
\end{itemize}

Le mode indépendant exige une masse cellulaire de référence finie et positive pour chaque espèce et une force Q6 déclarée dans $[0,1]$. Le seuil $0.5$ est le défaut de test initial, pas une calibration universelle d'atomisation. Pour le benchmark \texttt{dual\_identical}, le seuil est fixé à $0$ avec la condition supplémentaire $M_{c,s}>0$, afin que les deux étiquettes d'un même fluide soient projetées partout où elles sont présentes et qu'un basculement artificiel majoritaire/minoritaire ne soit pas introduit.

\subsubsection{Validation progressive 0493w5--0493w7}

Le smoke périodique 0493w5 couvre un support complet, deux îlots déconnectés, une occupation liquide $0.6$ au-dessus du seuil $0.5$ et une occupation $0.4$ en dessous. Il se termine par
\begin{equation}
  \texttt{PASS},\qquad
  \Delta u_{\max,\mathrm{gaz}}=4.34\times10^{-19},
  \qquad
  \Delta u_{\max,\mathrm{liquide}}=2.17\times10^{-2}.
\end{equation}
Le gaz de force nulle ne reçoit donc aucune correction directe mesurable, tandis que le liquide est effectivement projeté.

La matrice 0493w7 qualifie cinq topologies :
\begin{table}[htbp]
\centering
\small
\caption{Qualification multi-conditions-limites de \texttt{independent\_masked}. Le rapport est la divergence cellulaire post-application divisée par la divergence initiale.}
\begin{tabular}{@{}lccp{0.36\textwidth}@{}}
\toprule
Cas & Statut & $D^{\mathrm{cell}}_{\mathrm{après}}/D_{\mathrm{avant}}$ & Observation \\
\midrule
Périodique & PASS & $5.91\times10^{-16}$ & Support complet. \\
Canal à parois & PASS & $1.60\times10^{-12}$ & Stencil de paroi résident. \\
Entrée/sortie pleine face & PASS & $6.13\times10^{-12}$ & Flux ouvert complet porté par le liquide projeté. \\
Entrée/sortie segmentée & PASS & $5.37\times10^{-2}$ & Résidu localisé aux transitions ouverture/paroi. \\
Canal Darcy & PASS & $1.66\times10^{-12}$ & État Q6 résident consommé directement, \texttt{darcyFresh=1}. \\
\bottomrule
\end{tabular}
\end{table}
Dans les cinq cas, le nombre de particules gazeuses directement corrigées est nul et aucun fallback CPU ni téléchargement complet n'est utilisé.

\subsubsection{Qualification Taylor--Green mono/dual identique 0493w8}

Le banc 0493w8 reprend Taylor--Green sans resampling et compare trois scénarios sur \texttt{src} et \texttt{src-q6} :
\begin{itemize}[nosep]
  \item \texttt{mono\_legacy} : Q6 monoespèce historique sans registre ;
  \item \texttt{mono\_independent} : une espèce enregistrée, support complet et \texttt{independent\_masked} ;
  \item \texttt{dual\_identical} : le même fluide partagé à parts égales entre deux identifiants de type, avec masses, métadonnées physiques, masses cellulaires de référence et forces Q6 identiques.
\end{itemize}
Les positions, vitesses, masses et rôles initiaux des états mono et dual sont identiques ; seul le tableau de types diffère. Les paires de vitesses particulières opposées sont attribuées intégralement à un type, de sorte que chaque sous-population possède exactement la vitesse Taylor--Green demandée au pas initial. Le réglage nominal est $24\times24$, $\gamma=32$, $300$ pas, trois graines 493801--493803 et aucun resampling.

La qualification finale donne \texttt{PASS checks=141 failed=0}. La non-régression monoespèce est stricte : selon les graines, l'écart relatif sur la viscosité historique/indépendante est compris entre $0$ et $3.5\times10^{-16}$, le RMS des courbes TG est inférieur à $7.6\times10^{-17}A_0$, et le RMS des vitesses particulaires est voisin de $1.2\times10^{-15}U_0$. La neutralité du registre et du découpage en deux types sous SRC est également vérifiée à environ $7\times10^{-16}$ sur les états.

Pour le fluide physiquement identique projeté par deux solveurs indépendants, les résultats sont :
\begin{table}[htbp]
\centering
\small
\caption{Qualification Taylor--Green 0493w8, trois graines.}
\begin{tabular}{@{}rrrrrr@{}}
\toprule
Graine & $\nu_{\mathrm{mono\ Q6}}$ & $\nu_{\mathrm{dual\ Q6}}$ & écart relatif & RMS courbe/$A_0$ & slip/$U_0$ \\
\midrule
493801 & 0.11926964 & 0.11734371 & $1.63\,\%$ & $2.84\times10^{-3}$ & $1.43\,\%$ \\
493802 & 0.10701807 & 0.10843055 & $1.31\,\%$ & $4.82\times10^{-3}$ & $1.60\,\%$ \\
493803 & 0.10306556 & 0.10407460 & $0.97\,\%$ & $2.75\times10^{-3}$ & $1.76\,\%$ \\
\bottomrule
\end{tabular}
\end{table}
La moyenne vaut $\bar\nu_{\mathrm{mono}}=0.1097844$ et $\bar\nu_{\mathrm{dual}}=0.1099496$, soit un écart de moyennes d'environ $0.15\,\%$. La divergence résiduelle moyenne représente environ $14.0\,\%$ de la divergence SRC en monoespèce Q6 et $5.2\,\%$ en dual indépendant. L'appariement particule par particule n'est pas un critère du cas dual : deux sous-populations finies projetées séparément développent des trajectoires microscopiques différentes, alors que les observables barycentriques et de transport restent proches.

La dérive moyenne maximale mesurée est $1.685\times10^{-3}U_0$ pour la graine 493803. Elle est identique, à moins de $4\times10^{-17}U_0$, entre Q6 historique et Q6 indépendant. Elle est donc conservée comme diagnostic commun de l'opérateur sans correction uniforme d'impulsion, et non interprétée comme une régression du nouveau mode.

\subsubsection{Lecture physique et recommandations}

Les résultats établissent les points suivants :
\begin{itemize}
  \item conserver \texttt{common} comme référence de non-régression et \texttt{weighted} pour la reproduction des campagnes 0491 ;
  \item utiliser \texttt{independent\_masked} lorsque des espèces ont des contraintes de compressibilité différentes, notamment liquide projeté et gaz non projeté ;
  \item fixer explicitement les forces Q6 et les masses cellulaires de référence de chaque espèce ;
  \item calibrer le seuil d'occupation selon la résolution et la statistique locale, et ne pas interpréter la valeur par défaut $0.5$ comme une loi physique universelle ;
  \item distinguer dans les diagnostics la convergence du flux de face et la divergence du champ réellement appliqué ;
  \item traiter le résidu des supports partiels et des extrémités de segments comme un chantier de stencil face--cellule, sans remettre en cause la sélectivité déjà démontrée ;
  \item conserver Q6 et resampling comme deux axes indépendants et laisser le resampling désactivé pendant la qualification de projection ;
  \item réserver les compléments suivants à la physique interfaciale : pression de phase, tension superficielle, mouillage, traînée et diffusion inter-espèces.
\end{itemize}

Une comparaison exploratoire du runner liquide type 1 injecté dans un gaz type 2 confirme qualitativement cette direction : le mode \texttt{independent\_masked} conserve un jet dense et cohérent, avec un front et des déformations de caractère liquide, alors que le mode \texttt{weighted} produit une filamentation beaucoup plus diffuse. Cette observation LiveVis est cohérente avec le contrat sélectif, mais elle ne remplace pas une qualification quantitative d'interface ; l'échelle de couleur et les conditions de prise de vue doivent rester identiques pour toute comparaison publiée.

\section{Addendum CUDA-TOPO : Darcy--Brinkman et champ \texorpdfstring{$\chi$}{chi}}

\subsection{Objectif du chantier topologique}

Le chantier CUDA-TOPO a été ouvert après la stabilisation du chemin SRC classic CUDA résident et du resampling post-SRC. Son objectif immédiat n'était pas encore de fournir une optimisation topologique complète, mais de vérifier si le solveur pouvait accepter un champ géométrique externe, noté $\chi$, et l'utiliser pour imposer une résistance volumique de type Darcy--Brinkman. L'intérêt attendu est double. D'une part, un automate extérieur d'optimisation peut modifier $\chi$ sans modifier le code source du solveur. D'autre part, le même mécanisme doit pouvoir représenter des géométries non analytiques, telles que des coudes, des canaux diagonaux, des obstacles elliptiques ou des profils NACA.

Dans cette étape, $\chi=1$ désigne le fluide libre et $\chi=0$ la zone pénalisée. La convention retenue est donc opposée à une variable de porosité solide : le champ $\chi$ est un indicateur de liberté fluide. La voie validée à ce stade est une voie SRC classic CUDA-VIZ avec pénalisation Darcy--Brinkman. Elle n'active pas Q6, le resampling ni le viriel dans les scripts de validation topologique, mais l'architecture conserve la possibilité de réintroduire ces modules dans un chantier ultérieur.

\subsection{Formulation Darcy--Brinkman utilisée}

La résistance locale est calculée à partir du champ $\chi$ par une interpolation de type Borrvall--Petersson, dans l'esprit des formulations de pénalisation de fluide utilisées en optimisation topologique et des modèles de Brinkman pour milieux poreux \cite{Brinkman1949,BorrvallPetersson2003} :
\begin{equation}
  \alpha(\chi)=\alpha_{\min}+\left(\alpha_{\max}-\alpha_{\min}\right)\frac{q(1-\chi)}{q+\chi}.
\end{equation}
Cette forme donne $\alpha\simeq \alpha_{\min}$ dans le fluide libre et $\alpha\simeq \alpha_{\max}$ dans la zone solide/poreuse. Le paramètre $q$ règle la raideur de l'interpolation près de l'interface. La pénalisation agit ensuite sur la vitesse moyenne cellulaire, sous la forme d'un amortissement exponentiel vers une vitesse solide prescrite $u_s$ :
\begin{equation}
  u^{\star}=u_s+\exp(-\alpha\Delta t)(u-u_s),
\end{equation}
ou, de manière équivalente,
\begin{equation}
  \Delta u=-\lambda(\alpha)(u-u_s),\qquad
  \lambda(\alpha)=1-\exp(-\alpha\Delta t).
\end{equation}
Pour $\Delta t=5\times10^{-4}$, on obtient par exemple $\lambda\simeq0.393$ pour $\alpha=1000$, $\lambda\simeq0.918$ pour $\alpha=5000$, et $\lambda\simeq0.993$ pour $\alpha=10000$. Ces valeurs expliquent pourquoi des pénalisations modérées peuvent annuler correctement la vitesse dans la zone $\chi\simeq0$ sans reproduire pour autant une frontière solide équivalente.

La force diagnostique associée est une force volumique effective de pénalisation. Elle ne doit pas être lue comme une force aérodynamique pariétale exacte : elle mesure l'impulsion retirée au fluide par le frein Darcy dans les cellules pénalisées. Les observables ajoutées dans le chantier sont donc appelées \emph{proxies} de traînée et de portance.

\subsection{Implémentation CUDA-VIZ et contrôles interactifs}

Le premier jalon du chantier ajoute un chemin CUDA-VIZ pour appliquer la pénalisation Darcy--Brinkman dans le pas SRC classic résident. Les champs livevis disponibles incluent notamment $\chi$, $\alpha$ et la puissance dissipée Darcy. Le fichier \texttt{darcy\_cost\_0343.csv} consigne les grandeurs globales associées à la pénalisation : puissance Darcy, vitesse moyenne, fuite solide et diagnostics de coût.

Un second jalon rend le contrôle livevis persistant par le fichier racine \texttt{livevis\_control.kv}. Ce fichier permet de changer pendant l'exécution le champ visualisé, le gain, les bornes de clipping, le nombre de passes de lissage et la colormap. Ce contrôle s'est révélé important pour distinguer les effets physiques de pénalisation des simples effets d'échelle de couleur.

Le troisième jalon transforme $\chi$ en véritable entrée du solveur. Deux modes coexistent : un mode analytique pour les validations rapides et un mode \texttt{file} où le champ $\chi$ est lu une fois côté CPU, envoyé une fois sur GPU, puis utilisé en résident. Le pré-calcul de $\chi$, $\alpha$ et $\lambda$ évite de recalculer l'interpolation à chaque pas. La comparaison cercle analytique/cercle fichier a montré des résultats quasi identiques, ce qui valide le chemin d'entrée externe.

\subsection{Générateurs de géométrie et cas de validation}

Plusieurs générateurs de champ $\chi$ ont été ajoutés progressivement. Les premiers cas ont porté sur des canaux courbes et un canal diagonal, afin de vérifier que le code n'était plus limité aux formes analytiques simples. Les cas \texttt{bend\_pipe} et \texttt{diagonal\_channel} ont validé la lecture de $\chi$ externe, la visualisation directe de la géométrie et l'effet qualitatif de la pénalisation sur le flux.

La famille suivante de générateurs a introduit des obstacles de canal : cylindre, ellipse et profils NACA à quatre chiffres. Le profil NACA0012 a servi de cas de référence car sa symétrie impose une portance nulle à incidence nulle et une portance antisymétrique lorsque l'incidence change de signe. Les scripts de sweep créent les champs $\chi$, lancent les cas Darcy--Brinkman, appliquent les statistiques de fenêtre finale et construisent des polaires proxy.

\subsection{Observables de force et statistiques de fenêtre}

Les observables de benchmark sont volontairement optionnelles et cadencées pour limiter le coût des atomiques GPU. Le fichier \texttt{topo\_benchmark\_0348.csv} contient les grandeurs de force cellulaire projetées sur une direction de traînée et une direction de portance. Le post-traitement \texttt{0348b} extrait ensuite des moyennes, écarts-types, minima, maxima et dernières valeurs sur une fenêtre finale. Cette séparation évite de confondre le transitoire initial et le régime établi.

Les grandeurs principales sont
\begin{itemize}[leftmargin=*]
  \item \texttt{dragProxy}, projection de la force Darcy sur la direction de l'écoulement ;
  \item \texttt{liftProxy}, projection transverse ;
  \item \texttt{darcyPower}, puissance dissipée par la pénalisation ;
  \item \texttt{solidLeakRms}, fuite de vitesse dans la zone pénalisée ;
  \item les rapports \texttt{absLiftOverAbsDrag} et \texttt{liftOverDragProxy}.
\end{itemize}
Ces grandeurs sont cohérentes entre cas, mais elles restent des indicateurs de pénalisation volumique et non des coefficients $C_L$ ou $C_D$ calibrés.

\subsection{Calibration Poiseuille et échelle de Reynolds effective}

Afin de comparer les sweeps NACA à une échelle hydrodynamique, un script de calibration Poiseuille a été ajouté. Pour la grille cible NACA $N_x=600$, $N_y=160$, $\gamma=10$, $\Delta t=5\times10^{-4}$ et $k_BT=0.01$, le fit Poiseuille donne
\begin{equation}
  \nu_{\mathrm{eff}}\simeq 2.328\times 10^{-3}.
\end{equation}
Le cas $k_BT=0.01$ est retenu comme calibration préférée car le fit présente un coefficient $R^2\simeq0.977$, nettement plus propre que le cas $k_BT=0.1$. Pour une corde $c=0.22$ et une vitesse $U_0=1$, le Reynolds effectif vaut alors
\begin{equation}
  Re_{\mathrm{eff}}=\frac{U_0 c}{\nu_{\mathrm{eff}}}\simeq94.5.
\end{equation}
Cette valeur est faible devant les références expérimentales NACA classiques à haut Reynolds. Les comparaisons doivent donc être interprétées comme des comparaisons de forme normalisée, pas comme des validations quantitatives de coefficients aérodynamiques.

\subsection{Résultats NACA : comparaison normalisée aux polaires de référence}

Le sweep NACA0012 calibré à $Re_{\mathrm{eff}}\simeq94.5$ montre une réponse de portance proxy très structurée. La portance est quasi nulle à incidence nulle, change de signe lorsque l'angle d'attaque change de signe et croît presque linéairement sur l'intervalle étudié. Après normalisation par la valeur à $10^\circ$, la forme de la courbe de portance est proche de la courbe de référence NACA/Ladson utilisée comme repère qualitatif.

La traînée proxy est plus problématique. Elle est bien paire et minimale autour de $0^\circ$, mais sa croissance avec $|\alpha|$ reste trop faible. Dans le sweep de référence, la traînée normalisée augmente d'environ $28\%$ entre $0^\circ$ et $14^\circ$, alors que la référence NACA normalisée utilisée dans la comparaison augmente d'un facteur supérieur à deux. La conséquence est une finesse proxy trop favorable aux grands angles d'attaque, car la portance augmente raisonnablement alors que la traînée ne croît pas assez.

Un sweep de sensibilité en $\alpha_{\max}$ et $U_0$ a ensuite montré que l'augmentation de $\alpha_{\max}$ améliore l'imperméabilité locale et réduit la fuite solide, mais ne rend pas la courbe de traînée suffisamment convexe. L'augmentation de $U_0$ accroît davantage la courbure de traînée, mais au prix d'un bruit plus important et sans rétablir une vraie polar d'obstacle solide. Cette observation suggère que la limite ne provient pas uniquement d'une pénalisation trop faible, mais de la nature même du modèle Darcy pur.

\subsection{Comparaison Von Karman : solide immergé et Darcy pur}

Un cas autonome Von Karman a été ajouté pour comparer proprement le cylindre solide défini dans le code et un cylindre représenté uniquement par le champ Darcy. Les scripts ont été réalignés sur le chemin CUDA périodique/cercle validé du script portable : même chemin \texttt{periodic\_circle}, même thermostat non partagé, même angle de rotation effectif et mêmes paramètres de livevis. Après cet alignement, le mode \texttt{immersedSolid} du script autonome reproduit le comportement du script portable.

Le test critique compare ensuite trois modèles : \texttt{immersed}, \texttt{darcy} et, pour diagnostic seulement, \texttt{both}. Le mode \texttt{immersed} active la frontière solide codée et exclut les particules du cylindre dans l'état initial. Le mode \texttt{darcy} désactive la frontière solide et laisse le cylindre rempli de particules ; seule la pénalisation volumique freine la vitesse moyenne dans la zone $\chi\simeq0$.

Les diagnostics de densité de sillage sur $5000$ pas donnent le tableau synthétique suivant :
\begin{center}
\begin{tabular}{lcccc}
\toprule
Modèle & $\langle N_{\mathrm{wake}}/N_{\mathrm{up}}\rangle$ & Dernier $N_{\mathrm{wake}}/N_{\mathrm{up}}$ & $\langle U_{x,\mathrm{wake}}/U_{x,\mathrm{up}}\rangle$ & Dernier $U_{x,\mathrm{wake}}/U_{x,\mathrm{up}}$\\
\midrule
Solide immergé & 0.959 & 0.969 & 0.766 & 0.807\\
Darcy pur & 1.000 & 1.000 & 0.999 & 0.997\\
\bottomrule
\end{tabular}
\end{center}
Le résultat est très instructif. Le solide immergé crée un sillage réel : légère déplétion de population et déficit de vitesse d'environ $20$ à $25\%$ dans la boîte de sillage proche. Le Darcy pur, malgré une vitesse faible dans la zone pénalisée elle-même, ne crée pratiquement ni déficit de population ni déficit de vitesse en aval. Le fluide traverse donc un frein volumique poreux plus qu'il ne contourne une frontière solide.

Cette observation invalide l'hypothèse initiale selon laquelle les mauvais proxies de traînée seraient dus à un appauvrissement excessif du sillage Darcy. C'est au contraire le solide immergé qui crée le sillage plus marqué. La limite du Darcy pur est plutôt son incapacité à imposer une condition de non-pénétration et une production de vorticité pariétale équivalentes à celles d'une paroi solide.



\subsection{Complément 0418--0426 : de Darcy pur vers une frontière pilotée par \texorpdfstring{$\chi$}{chi}}

Les premiers tests ont montré qu'une pénalisation Darcy pure arrête bien la vitesse moyenne dans la zone $\chi\simeq0$, mais ne crée pas le même sillage qu'une frontière solide codée. La suite du chantier a donc consisté à transformer le champ $\chi$ en support géométrique plus riche, sans revenir à une liste de formes analytiques figées. Trois ajouts sont importants.

Le premier est la désactivation initiale des particules dans la zone solide, contrôlée par \texttt{darcyInitialDeactivateBelowChi}. Lorsque ce seuil vaut par exemple $0.05$, les particules placées initialement dans les cellules très solides sont basculées hors du rôle fluide. Le champ $\chi$ agit donc dès l'état initial comme un masque géométrique, et non seulement comme une force volumique appliquée après coup. Ce point est essentiel pour les géométries de type backward step, bend pipe ou obstacle, car il évite de commencer le calcul avec une masse fluide artificielle dans le solide.

Le deuxième ajout est le mode \texttt{mean} de la pénalisation. Le code dépose masse et quantité de mouvement sur la grille, reconstruit la vitesse moyenne cellulaire, puis applique le kick Darcy--Brinkman sur cette vitesse moyenne avant de redistribuer la correction aux particules. Dans ce mode, la force n'est pas une correction particule par particule indépendante : elle est cohérente avec le moment cellulaire local. Le mode \texttt{mean\_outward\_bath} ajoute ensuite un bain orienté par la normale locale. La normale est reconstruite à partir de $\nabla\chi$ ; avec la convention $\chi=0$ solide et $\chi=1$ fluide, elle pointe du solide vers le fluide. Le bain corrige la composante relative des vitesses de manière à favoriser une réémission hors du solide et à limiter les fuites résiduelles au voisinage de l'interface. Dans les validations backward step, ce bain n'est pas apparu comme un coût dominant ; au contraire, il stabilise localement la dynamique et réduit les corrections parasites.

Le troisième ajout est \texttt{chiVP}, c'est-à-dire une contribution de particules virtuelles de collision dérivée de $\chi$. Il ne s'agit pas de créer, streamer ou dumper des particules virtuelles persistantes. Le mécanisme ajoute seulement une masse et une quantité de mouvement effectives dans la collision SRC des cellules d'interface. Si $M$ et $P$ sont la masse et le moment réels déposés dans une cellule et si $M_{\rm vp}$ désigne la masse virtuelle issue de $\chi$, le centre de masse de collision est remplacé par
\begin{equation}
  u_{\rm cm}^{\rm eff} = \frac{P + M_{\rm vp} u_s}{M+M_{\rm vp}}.
\end{equation}
Dans les scripts, \texttt{darcyChiCollisionVpGamma=-1} signifie que l'occupation virtuelle de référence est héritée de la valeur MPCD nominale, \texttt{darcyChiCollisionVpLayers} fixe l'épaisseur de bande, \texttt{darcyChiCollisionVpThreshold} définit le seuil d'interface et \texttt{darcyChiCollisionVpStrength} règle l'intensité relative. Sur le cas Von Karman filtré, la valeur \texttt{0.25} a donné le meilleur compromis global entre proche sillage, sillage aval et zone proche paroi. Cette valeur a donc été retenue pour le backward step Darcy/chiVP.

\subsection{Correction de performance 0426 : fastflags CUDA résidents}

Une ablation backward step $960\times480$ a isolé un problème d'aiguillage important. Les premiers runs Darcy segmentés semblaient environ trois à quatre fois plus lents que le solide rectangle full-face. Les tests successifs ont écarté les hypothèses naturelles : désactiver \texttt{chiVP}, enlever le bain \texttt{outward}, mettre $\alpha=0$ ou désactiver \texttt{darcyBrinkmanEnable} ne ramenait pas le temps vers le chemin solide. L'activation du profil interne a montré que le temps était absorbé par \texttt{src\_collision}, \texttt{force\_stream} et un bloc géométrique générique, et non par les kernels Darcy proprement dits.

La cause était finalement l'absence, dans les scripts Darcy de démonstration, du même groupe de drapeaux CUDA résidents que dans les scripts solides optimisés. Le bloc 0426 active par défaut, entre autres, la rotation sur device, le diagnostic thermostat rapide, le dépôt streaming fusionné, le contrôle paresseux des kernels, la suppression de synchronisations finales ou de setup inutiles, la suppression du téléchargement du workspace et l'évitement du remplissage hôte des cell-id lorsque le chemin résident le permet. Il est encapsulé dans les scripts par \texttt{export\_src\_cuda\_resident\_fastflags\_0426} et peut être désactivé par \texttt{MPCD\_DARCY\_FASTFLAGS\_ENABLE=0} pour diagnostic.

Le résultat de performance est net. Sur $500$ pas du backward step $960\times480$, le cas \texttt{darcyOff} segmenté passe d'environ $109\,\si{\second}$ profilées à environ $28\,\si{\second}$ après activation des fastflags. Le vrai cas physique \texttt{mean\_outward\_bath+chiVP}, avec $\alpha_{\max}=8\times10^5$, donne ensuite environ $28\,\si{\second}$ profilées et $31\,\si{\second}$ de temps mural pour $500$ pas, soit un coût voisin du \texttt{solid\_fullface} profilé. La conclusion est donc forte : le chemin Darcy/chiVP n'est pas intrinsèquement coûteux ; le ralentissement provenait d'une configuration incomplète du backend résident.

\subsection{Conséquence pour les formes solides arbitraires}

Après cette correction, le champ $\chi$ devient une interface pratique pour imposer des formes solides quelconques dans le chemin SRC classic CUDA résident. Le solveur n'a plus besoin d'un \texttt{caseType} ou d'une forme analytique dédiée pour chaque géométrie : un générateur externe, un automate d'optimisation ou un script de prétraitement peut écrire $\chi$, puis le même pipeline applique désactivation initiale, pénalisation moyenne, bain orienté et collision virtuelle d'interface. Cette souplesse est particulièrement importante pour les coudes, marches, obstacles composites et profils issus d'une boucle d'optimisation.

Il reste néanmoins utile de garder la distinction conceptuelle suivante. Le mode Darcy pur représente un milieu poreux ou une résistance volumique. Le mode augmenté \texttt{mean\_outward\_bath+chiVP} fournit une approximation beaucoup plus solid-aware de la frontière, et devient la voie opérationnelle pour les géométries arbitraires pilotées par $\chi$. Une condition de paroi strictement équivalente à un bounce-back analytique reste à qualifier par des métriques de sillage, de production de vorticité et de fuite normale, mais le verrou principal de performance est levé.

\subsection{Statut validé et limite identifiée}

Le point d'étape Darcy--Brinkman doit maintenant être formulé en distinguant clairement trois niveaux :
\begin{itemize}[leftmargin=*]
  \item le chemin CUDA-VIZ applique correctement une pénalisation Darcy--Brinkman issue d'un champ $\chi$ analytique ou externe ;
  \item les champs $\chi$, $\alpha$, vitesse, vorticité et puissance Darcy sont visualisables en temps réel, avec contrôle persistant par \texttt{livevis\_control.kv} ;
  \item les générateurs de géométrie et les sweeps NACA démontrent que $\chi$ peut piloter des formes complexes sans modifier le solveur ;
  \item les observables de force Darcy restent cohérentes comme proxies de pénalisation volumique ;
  \item Darcy pur reste un modèle poreux et ne reproduit pas, à lui seul, une frontière solide impenetrable ni le même sillage qu'un solide immergé codé ;
  \item la chaîne augmentée \texttt{mean\_outward\_bath+chiVP}, combinée aux fastflags 0426, fournit désormais une voie efficace et flexible pour traiter des formes solides quelconques à partir de $\chi$.
\end{itemize}

Cette distinction est importante pour l'optimisation topologique. Le champ $\chi$ est validé comme entrée géométrique externe et comme champ de pénalisation poreuse. Il n'est pas encore validé comme représentation complète d'un solide équivalent. Pour atteindre cet objectif, un futur chantier devra construire une condition solide dérivée de $\chi$ : détection d'interface, normale locale, réflexion ou bounce-back, éventuellement particules virtuelles de paroi dérivées de $\nabla\chi$. On peut alors distinguer deux usages futurs :
\begin{itemize}[leftmargin=*]
  \item \texttt{CYLINDER\_MODEL=darcy}, pour représenter un milieu poreux ou une résistance volumique continue ;
  \item \texttt{CYLINDER\_MODEL=chi\_solid}, à développer, pour représenter un obstacle solide général piloté par $\chi$.
\end{itemize}

Le chantier Darcy--Brinkman est donc validé comme brique topologique de pénalisation et, depuis 0426, comme interface pratique pour généraliser le traitement des formes solides à partir de $\chi$ sans dépendre de formes prédéfinies. La validation fine des conditions pariétales équivalentes -- fuite normale, production de vorticité, longueur de sillage et efforts -- reste à poursuivre, mais l'obstacle principal de performance du chemin Darcy/chiVP a été levé.

\section{Étalonnage du fluide SRC effectif et référence de co-simulation (0493w1)}
\label{sec:etalonnage-src-effectif}

\subsection{Motivation : caractériser le fluide avant de corriger la simulation}

Les développements consacrés aux cellules pauvres, au resampling, aux parois et aux obstacles ont montré qu'une situation apparemment pathologique peut avoir deux origines très différentes. Elle peut provenir d'un défaut numérique réel de support particulaire, mais elle peut aussi être la conséquence normale d'un fluide SRC dont les coefficients de transport et les échelles acoustiques sont mal adaptés au domaine imposé. Une zone de sillage peu repeuplée, une faible diffusion transverse ou une accumulation à proximité d'une ouverture ne suffisent donc pas, à elles seules, à justifier une correction de population.

La première calibration appliquée au cas Darcy segmenté historique a précisément révélé un régime extrême. Pour
\begin{equation}
  a=\frac{1}{128},\qquad
  \Delta t=10^{-3},\qquad
  k_BT=10^{-3},\qquad
  \gamma=20,
\end{equation}
avec une rotation SRC de $90^\circ$, décalage aléatoire de grille et thermostat \texttt{cell\_relative\_rescale} appliqué à chaque pas, les estimations obtenues étaient
\begin{equation}
  \nu\simeq5.09\times10^{-3},\qquad
  D_{\mathrm{self}}\simeq7.64\times10^{-7},\qquad
  Sc\simeq6.66\times10^{3},\qquad
  \frac{\lambda}{a}\simeq5.1\times10^{-3}.
\end{equation}
Pour $U=0.15$ et $L=0.24$, le Reynolds n'était que $Re\simeq7.1$. L'identification acoustique de ce premier essai était invalide du point de vue de la quantité de mouvement ; seules les bornes idéales isotherme et adiabatique 2D pouvaient être utilisées comme proxies, conduisant à $Ma\simeq3.35$--$4.74$. L'échelle combinée $c_sL/\nu=Re/Ma$ n'était alors que de l'ordre de $1.5$--$2.1$. Le cas réunissait donc faible Reynolds, très forte compressibilité à la vitesse imposée, autodiffusion extrêmement faible et libre parcours presque nul devant la cellule. Une partie des besoins de réparation observés dans ce régime doit par conséquent être requalifiée avant d'être érigée en exigence générale du resampling.

La règle méthodologique désormais retenue est la suivante :
\begin{quote}
Toute configuration SRC destinée à un calcul physique doit être caractérisée comme un fluide effectif avant l'activation de Q6/Q9, du resampling ou d'une fermeture quasi incompressible. Les corrections ultérieures doivent être évaluées relativement à cette référence, et non relativement à des proxies théoriques ou à une vitesse imposée arbitrairement.
\end{quote}

\subsection{Configuration qui définit le fluide effectif}

Le fluide n'est pas défini par la seule grille. La signature minimale de la calibration est
\begin{equation}
\Theta_{\mathrm{SRC}}=
\left(
L_x,L_y,N_x,N_y,\gamma,\Delta t,k_BT,m,\alpha,
\mathcal{S}_{\pm},\mathcal{G},\mathcal{T}
\right),
\end{equation}
où $\mathcal{S}_{\pm}$ décrit le choix aléatoire du signe de rotation, $\mathcal{G}$ le décalage aléatoire de grille et $\mathcal{T}$ la politique de thermostat, y compris sa cadence, sa température cible et sa population minimale. La taille locale de cellule
\begin{equation}
  a_x=\frac{L_x}{N_x},\qquad a_y=\frac{L_y}{N_y}
\end{equation}
entre directement dans les coefficients de transport. À paramètres locaux identiques, une augmentation simultanée de $L$ et de $N$ conserve $a$ et constitue un contrôle de taille de boîte ; à domaine fixé, une augmentation de $N$ modifie $a$ et définit un autre fluide SRC.

Le thermostat appartient explicitement à cette signature. Un SRC brut et un SRC thermostaté à chaque pas ne sont pas deux réalisations d'un même fluide : ils peuvent présenter des viscosités, diffusions et réponses acoustiques différentes. Le calibrateur classe donc la configuration comme \texttt{raw\_src} ou \texttt{thermostatted\_src}. Q6 et le resampling sont volontairement désactivés afin de mesurer le fluide sous-jacent sur lequel ces modules agiront ensuite.

\subsection{Runner canonique de cas unique}

Le runner de référence est
\begin{center}
  \texttt{scripts/run\_0493w1\_src\_fluid\_calibrator.sh}.
\end{center}
Il prend ses paramètres par variables d'environnement, produit les états périodiques nécessaires, lance les expériences et appelle l'analyseur \texttt{analyze\_0493w1\_src\_fluid\_calibrator.py}. Son contrat courant est résumé dans le tableau suivant.

\begin{center}
\small
\begin{tabular}{p{0.34\textwidth}p{0.58\textwidth}}
\toprule
Entrées & Rôle physique ou numérique \\
\midrule
\texttt{Lx, Ly, NX, NY} & Domaine périodique de calibration, grille et taille des cellules de collision. \\
\texttt{GAMMA, DT, KBT, PARTICLE\_MASS} & Occupation moyenne, intervalle de collision, énergie thermique et masse particulaire. \\
\texttt{ROTATION\_ANGLE, RANDOM\_ROTATION\_SIGN, GRID\_SHIFT\_ENABLE} & Règle de collision SRC et restauration de l'invariance galiléenne. \\
\texttt{THERMOSTAT\_*} & Activation, mode, cadence, cible $k_BT$ et population minimale du thermostat de production. \\
\texttt{TG\_MODE\_X/Y, TG\_STEPS, TG\_DUMP\_COUNT} & Longueur d'onde, durée et échantillonnage de la mesure de viscosité. \\
\texttt{SOUND\_MODE\_X, SOUND\_REPLICATES, SOUND\_STEPS, SOUND\_DUMP\_COUNT} & Longueur d'onde et ensemble de réalisations pour la mesure acoustique. \\
\texttt{MSD\_GRID\_MODE, MSD\_STEPS, MSD\_DUMP\_COUNT} & Domaine et durée de la mesure d'autodiffusion. \\
\texttt{CHARACTERISTIC\_U, CHARACTERISTIC\_L} & Échelles de la future simulation pour calculer $Re$, $Ma$, $Pe$ et les vitesses compatibles avec des cibles données. \\
\bottomrule
\end{tabular}
\end{center}

Une calibration standard avec \texttt{SOUND\_REPLICATES=4} effectue six simulations SRC indépendantes : un Taylor--Green, quatre modes longitudinaux et un MSD. Le préflight, activé par \texttt{PREFLIGHT\_ONLY=1}, n'exécute aucune simulation et vérifie notamment la géométrie, la résolution des longueurs d'onde, le nombre d'intervalles de régression, le nombre d'échantillons par période acoustique proxy, le volume estimé des dumps et le nombre de \emph{particle-steps}. Les options \texttt{MAX\_DUMP\_GB} et \texttt{ALLOW\_LARGE\_DUMPS} évitent une campagne volumineuse involontaire. \texttt{ANALYZE\_ONLY=1} permet de reconstruire les résultats à partir de dumps existants.

Le runner caractérise un fluide homogène dans une boîte périodique ; il ne reproduit pas l'obstacle ou la paroi de la simulation finale. Deux usages sont distingués :
\begin{itemize}[leftmargin=*]
  \item le mode \emph{cell-equivalent}, qui conserve $a$, $\Delta t$, $k_BT$ et les règles SRC dans une boîte périodique compacte, afin de mesurer à faible coût les propriétés locales du fluide ;
  \item le mode \emph{domaine complet}, qui conserve le domaine et la grille de la simulation cible pour contrôler les effets de taille finie, de statistique et de longueur d'onde. L'option \texttt{MSD\_GRID\_MODE=full} impose alors le domaine complet pour l'autodiffusion.
\end{itemize}

\subsection{Mesure de la viscosité par décroissance Taylor--Green}

Le générateur initialise le mode transverse
\begin{equation}
\begin{aligned}
  u_x(x,y,0)&=A_0\sin(k_xx)\cos(k_yy),\\
  u_y(x,y,0)&=-A_0\cos(k_xx)\sin(k_yy),
\end{aligned}
\qquad
k_x=\frac{2\pi n_x}{L_x},\quad
k_y=\frac{2\pi n_y}{L_y}.
\end{equation}
L'amplitude est reconstruite par projection massique des vitesses particulaires sur cette base. Dans le régime hydrodynamique linéaire,
\begin{equation}
  A(t)=A_0\exp\left[-\nu(k_x^2+k_y^2)t\right],
\end{equation}
d'où
\begin{equation}
  \nu=-\frac{1}{k_x^2+k_y^2}\frac{\mathrm d}{\mathrm dt}\ln|A(t)|.
\end{equation}

L'analyseur construit plusieurs fenêtres de fit, écarte les points situés sous $2.5\,\%$ de l'amplitude initiale afin de ne pas ajuster le plancher stochastique, impose au moins sept points et une décroissance logarithmique suffisante, puis choisit la fenêtre de meilleure linéarité. Il publie la valeur choisie, $R^2$, le nombre de points et la dispersion des valeurs de $\nu$ sur les fenêtres acceptables. Le statut est \texttt{PASS} pour une valeur positive avec $R^2\geq0.98$, \texttt{REVIEW} pour $0.93\leq R^2<0.98$ et \texttt{INVALID} dans les autres cas.

Une réalisation TG unique peut toutefois rester dominée par le bruit lorsque le mode devient faible. Pour une référence précise, la campagne 0493w1 a montré qu'il faut moyenner les amplitudes signées de plusieurs graines à temps identiques, puis effectuer une seule régression sur le signal d'ensemble. La moyenne des viscosités ajustées séparément est moins robuste, car chaque trajectoire peut sélectionner une fenêtre différente après son entrée dans le bruit.

\subsection{Mesure acoustique par régression hydrodynamique cumulative}

L'état acoustique est initialisé sans vitesse cohérente,
\begin{equation}
  \rho(x,0)=\rho_0\left[1+\varepsilon\cos(kx)\right],
  \qquad u_x(x,0)=0,
\end{equation}
avec $\varepsilon=0.08$ dans la campagne fix3. À chaque dump, l'analyseur calcule les modes complexes normalisés de densité et de vitesse longitudinale. Les quatre réalisations thermiques indépendantes sont synchronisées et moyennées avant le fit.

Avec la convention de Fourier $\exp(-ikx)$, le modèle linéaire utilisé est
\begin{equation}
  \frac{\mathrm d\widehat\rho}{\mathrm dt}=-ik\widehat u,
\end{equation}
\begin{equation}
  \frac{\mathrm d\widehat u}{\mathrm dt}
  =-ikc_s^2\widehat\rho-\nu_Lk^2\widehat u.
\end{equation}
Plutôt que de différencier des données bruitées ou d'exiger une période acoustique complète, l'analyseur intègre ces deux équations depuis l'état initial par la règle du trapèze :
\begin{equation}
  \widehat\rho(t)-\widehat\rho(0)
  =-ik\int_0^t\widehat u(\tau)\,\mathrm d\tau,
\end{equation}
\begin{equation}
  \widehat u(t)-\widehat u(0)
  =-ikc_s^2\int_0^t\widehat\rho(\tau)\,\mathrm d\tau
  -\nu_Lk^2\int_0^t\widehat u(\tau)\,\mathrm d\tau.
\end{equation}
Une régression complexe à deux paramètres avec contraintes de non-négativité fournit $c_s^2$ et $\nu_L$. Les résidus relatifs des balances de continuité et de quantité de mouvement, le conditionnement, les contraintes actives et un jackknife par retrait d'une réalisation quantifient la qualité. Une valeur acoustique \texttt{INVALID} ou \texttt{REVIEW} ne doit pas servir à publier un Mach mesuré ; le runner conserve alors seulement les proxies idéaux
\begin{equation}
  c_{s,\mathrm{iso}}=\sqrt{\frac{k_BT}{m}},
  \qquad
  c_{s,\mathrm{ad},2D}=\sqrt{\frac{2k_BT}{m}}.
\end{equation}
Dans les six fluides de la campagne physique, la vitesse mesurée est proche de la branche isotherme, ce qui est cohérent avec le thermostat cellulaire appliqué à chaque pas.

\paragraph{Limite de cette mesure lorsque Q6 est actif.}
La régression acoustique précédente caractérise le fluide SRC compressible sous-jacent. Elle ne doit pas être appliquée telle quelle à un chemin \texttt{src-q6}. Q6 retire précisément la composante longitudinale divergente dont l'oscillation permet d'identifier $c_s$ ; le système ajusté n'est alors plus fermé par les deux seules équations acoustiques ci-dessus. Pour les variantes Q6, le calibrateur doit publier un statut du type \texttt{soundSpeedStatus=NOT\_APPLICABLE\_Q6} et remplacer le fit de $c_s$ par des diagnostics de réponse incompressible : décroissance du mode longitudinal, réduction de divergence, évolution de la modulation de densité et impulsion retirée par la projection. La viscosité Taylor--Green et l'autodiffusion restent, en revanche, des comparaisons pertinentes entre \texttt{src} et \texttt{src-q6}.


\subsection{Mesure de l'autodiffusion par MSD périodique}

Le troisième état est un fluide homogène à l'équilibre. Un sous-ensemble persistant de particules est suivi entre les dumps. Les incréments périodiques sont déroulés, puis le déplacement du centre de masse de l'échantillon est retiré afin de ne pas confondre diffusion et dérive globale. En deux dimensions,
\begin{equation}
  \left\langle|\Delta\bm r(t)|^2\right\rangle=4D_{\mathrm{self}}t.
\end{equation}
L'analyseur ajuste les portions tardives commençant à $20$, $30$, $40$ ou $50\,\%$ de la durée totale et retient la fenêtre de meilleur $R^2$. Il publie également l'anisotropie $MSD_x/MSD_y$, le plus grand incrément relatif à la boîte et le paramètre non gaussien 2D
\begin{equation}
  \alpha_2=\frac{\langle r^4\rangle}{2\langle r^2\rangle^2}-1.
\end{equation}
Le statut de diffusion est \texttt{PASS} pour $R^2\geq0.995$, \texttt{REVIEW} pour $0.98\leq R^2<0.995$ et \texttt{INVALID} sinon.

\subsection{Grandeurs dérivées et lecture physique}

À partir des trois mesures, le runner calcule notamment
\begin{equation}
  Sc=\frac{\nu}{D_{\mathrm{self}}},
  \qquad
  Re=\frac{UL}{\nu},
  \qquad
  Ma=\frac{U}{c_s},
  \qquad
  Pe_{\mathrm{mass}}=\frac{UL}{D_{\mathrm{self}}}.
\end{equation}
Le libre parcours indiqué est un proxy de déplacement balistique thermique pendant un intervalle de collision,
\begin{equation}
  \lambda_{\mathrm{mean}}=\sqrt{\frac{\pi k_BT}{2m}}\,\Delta t,
  \qquad
  \frac{\lambda}{a}=\frac{\lambda_{\mathrm{mean}}}{\sqrt{a_xa_y}},
\end{equation}
et non un libre parcours moléculaire mesuré. Le runner fournit aussi
\begin{equation}
  \frac{Re}{Ma}=\frac{c_sL}{\nu},
\end{equation}
qui mesure directement la capacité du fluide à atteindre simultanément un Reynolds élevé et un Mach faible, ainsi que la vitesse compatible avec $Ma=0.3$, le Reynolds correspondant, la vitesse imposant $Re=50$ et le Mach résultant.

\subsection{Pré-balayage physique de six configurations}

Le pré-balayage 0493w1 a conservé $\gamma=20$, $m=1$, une rotation de $90^\circ$, le signe aléatoire, le décalage de grille et le thermostat \texttt{cell\_relative\_rescale} à chaque pas. Il a varié $a$, $\Delta t$ et $k_BT$. Les valeurs du tableau suivant proviennent d'une calibration complète par cas et ont toutes obtenu les trois statuts \texttt{PASS}. Elles constituent un balayage de sélection ; la viscosité du candidat final est ensuite affinée par un ensemble multi-graines.

\begin{center}
\scriptsize
\begin{tabular}{lrrrrrrrr}
\toprule
Cas & $a$ & $\Delta t$ & $k_BT$ & $\lambda/a$ & $\nu$ & $c_s$ & $D_{\mathrm{self}}$ & $Sc$ \\
\midrule
\texttt{a192\_dt002\_k030} & $1/192$ & .002 & .030 & .083 & $1.025\,10^{-3}$ & .17054 & $4.226\,10^{-5}$ & 24.25 \\
\texttt{a192\_dt002\_k125} & $1/192$ & .002 & .125 & .170 & $1.252\,10^{-3}$ & .35058 & $1.650\,10^{-4}$ & 7.59 \\
\texttt{a256\_dt002\_k030} & $1/256$ & .002 & .030 & .111 & $5.385\,10^{-4}$ & .17229 & $3.892\,10^{-5}$ & 13.84 \\
\texttt{a256\_dt002\_k125} & $1/256$ & .002 & .125 & .227 & $4.891\,10^{-4}$ & .35460 & $1.626\,10^{-4}$ & 3.01 \\
\texttt{a256\_dt004\_k125} & $1/256$ & .004 & .125 & .454 & $3.113\,10^{-4}$ & .35474 & $3.061\,10^{-4}$ & 1.02 \\
\texttt{a384\_dt002\_k125} & $1/384$ & .002 & .125 & .340 & $2.974\,10^{-4}$ & .35425 & $1.616\,10^{-4}$ & 1.84 \\
\bottomrule
\end{tabular}
\end{center}

Pour $U=0.15$ et $L=0.24$, les indicateurs de sélection étaient :

\begin{center}
\small
\begin{tabular}{lrrr}
\toprule
Cas & $Re(U=0.15)$ & $Ma(U=0.15)$ & $Re(Ma=0.3)$ \\
\midrule
\texttt{a192\_dt002\_k030} & 35.13 & .880 & 11.98 \\
\texttt{a192\_dt002\_k125} & 28.75 & .428 & 20.16 \\
\texttt{a256\_dt002\_k030} & 66.85 & .871 & 23.03 \\
\texttt{a256\_dt002\_k125} & 73.60 & .423 & 52.20 \\
\texttt{a256\_dt004\_k125} & 115.64 & .423 & 82.05 \\
\texttt{a384\_dt002\_k125} & 121.05 & .423 & 85.76 \\
\bottomrule
\end{tabular}
\end{center}

Les deux fluides à $k_BT=0.03$ restent trop compressibles à $U=0.15$. Les variantes \texttt{a256\_dt004\_k125} et \texttt{a384\_dt002\_k125} permettent un Reynolds plus élevé à Mach donné, mais leur $Sc$ proche de l'unité ou de deux et leur $\lambda/a$ élevé les rapprochent d'un régime cinétique/gazeux. Le cas \texttt{a256\_dt002\_k125} a été retenu comme compromis : libre parcours appréciable sans être dominant, $Sc$ de quelques unités, vitesse du son élevée et résolution de l'obstacle améliorée. Les valeurs du balayage ne doivent toutefois pas être prises comme valeurs finales de viscosité, car elles ne reposent que sur une réalisation TG.

\subsection{Contrôle de taille de boîte et ensemble multi-graines}
\label{sec:w1-multiseed-viscosity}

Le candidat a été recalibré avec la même taille de cellule $a=1/256$ dans deux boîtes :
\begin{equation}
\begin{aligned}
  64^2 &: L_x=L_y=0.25,\quad (n_x,n_y)_{TG}=(1,1),\quad n_{s}=2,\\
  128^2 &: L_x=L_y=0.50,\quad (n_x,n_y)_{TG}=(2,2),\quad n_{s}=4.
\end{aligned}
\end{equation}
Les longueurs d'onde physiques sont ainsi identiques : $64$ cellules pour Taylor--Green et $32$ cellules pour le son. Les durées et cadences ont été appariées : $2000$ pas et $50$ dumps demandés pour TG, $1200$ pas et $120$ dumps demandés pour chacune des quatre réalisations acoustiques, $3000$ pas et $60$ dumps demandés pour le MSD complet. Quatre graines de base ont été utilisées : $493201$, $593201$, $693201$ et $793201$.

Les viscosités ajustées séparément étaient beaucoup plus dispersées en $64^2$ qu'en $128^2$ :
\begin{equation}
  \nu_{64,\mathrm{ind}}=(5.880\pm1.186)\times10^{-4},
  \qquad
  \nu_{128,\mathrm{ind}}=(6.051\pm0.320)\times10^{-4},
\end{equation}
où les incertitudes sont les écarts-types entre graines. Cette différence de dispersion provient du meilleur rapport signal/bruit de la projection TG dans la grande boîte, qui contient quatre fois plus de particules.

La méthode de référence consiste à moyenner les amplitudes TG signées avant la régression. Sur la fenêtre commune $0.08\leq t\leq1.36$, on obtient
\begin{equation}
  \nu_{64,\mathrm{ens}}=5.8526\times10^{-4},
  \qquad R^2=0.99870,
\end{equation}
avec une sensibilité jackknife $3.71\times10^{-5}$, et
\begin{equation}
  \nu_{128,\mathrm{ens}}=5.9751\times10^{-4},
  \qquad R^2=0.99972,
\end{equation}
avec une sensibilité jackknife $1.29\times10^{-5}$. L'écart entre les deux boîtes n'est que $2.09\,\%$ et reste inférieur à l'incertitude de l'ensemble $64^2$. Aucun effet significatif de taille de boîte n'est donc détecté à taille de cellule et nombre d'onde physique fixés. La boîte $128^2$ est néanmoins recommandée pour une référence précise, car elle réduit fortement la variance de l'estimateur TG.

Sur les quatre calibrations $128^2$, les autres propriétés donnent
\begin{equation}
  c_s=0.35459\pm0.00050,
  \qquad
  D_{\mathrm{self}}=(1.6588\pm0.0268)\times10^{-4},
\end{equation}
les incertitudes étant les écarts-types inter-graines. En combinant $D_{\mathrm{self}}$ avec la viscosité du signal d'ensemble,
\begin{equation}
  \boxed{\nu=5.98\times10^{-4}},\qquad
  \boxed{c_s=0.35459},\qquad
  \boxed{D_{\mathrm{self}}=1.659\times10^{-4}},\qquad
  \boxed{Sc\simeq3.60}.
\end{equation}

Pour $L=0.24$ et $U=0.1064$,
\begin{equation}
  Ma=0.3001,\qquad Re=42.74,
\end{equation}
et
\begin{equation}
  \frac{c_sL}{\nu}=142.4.
\end{equation}
Atteindre $Re=50$ avec ce fluide demande
\begin{equation}
  U\simeq0.1245,\qquad Ma\simeq0.351.
\end{equation}
Le compromis physiquement documenté est donc $Re\simeq43$ à $Ma\simeq0.30$, ou $Re=50$ au prix d'un Mach voisin de $0.35$. Ces valeurs remplacent les premières estimations fondées sur une seule trajectoire TG.

\paragraph{Portée de la viscosité de référence.}
Les valeurs précédentes caractérisent le \emph{SRC thermostaté sous-jacent}, Q6 et le resampling étant désactivés dans cette campagne. Elles ne doivent donc pas être transplantées sans contrôle à un calcul dont le chemin numérique complet contient Q6-g-f. Pour un benchmark dont le Reynolds doit être connu à quelques pourcents, la viscosité transverse doit être requalifiée sous le chemin effectivement utilisé et son caractère constitutif doit être vérifié : reproductibilité multi-graines, indépendance raisonnable vis-à-vis de l'amplitude du mode transverse et contrôle de longueur d'onde avant d'interpréter $Re=UL/\nu$. Cette distinction devient déterminante pour la comparaison Zovatto--Pedrizzetti du cycle 0493x8, voir section~\ref{sec:x8-vk-validation}.

\subsection{Recommandations pour les domaines et modules à traiter}

\paragraph{Écoulements périodiques et benchmarks bulk.}
La calibration du SRC sous-jacent doit être effectuée avant toute comparaison Q6/Q9. Taylor--Green, onde longitudinale et MSD fournissent la référence à laquelle les variantes projetées doivent être comparées. Le resampling doit rester désactivé lorsqu'aucun défaut de support n'est observé. Une boîte $128^2$ et plusieurs graines sont préférables lorsque le Reynolds doit être fixé avec précision.

\paragraph{Obstacles, sillages et pénalisation Darcy/$\chi$.}
La vitesse d'entrée doit d'abord être choisie à partir de $c_s$ et de $\nu$, et non imposée indépendamment. Les sillages du cas Darcy historique doivent être rejoués avec le fluide calibré \texttt{a256\_dt002\_k125} avant de conclure à un défaut de repeuplement. Il faut distinguer la pénalisation poreuse volumique d'une frontière solide et vérifier séparément fuite normale, production de vorticité, déficit de vitesse et bilan de masse. Une absence de sillage à $Re\simeq7$ et $Ma>3$ ne constitue pas un test pertinent d'un mécanisme de resampling destiné à un écoulement faiblement compressible.

\paragraph{Parois et cellules coupées.}
Les particules virtuelles peuvent compléter les moments de collision dans les fractions solides, mais elles ne doivent pas être comptées comme masse physique. Dans une cellule coupée, la cible de masse doit être liée à la fraction accessible et aux conditions de flux, non à une occupation bulk uniforme. Les cellules totalement solides ne doivent pas conserver une population fluide persistante. Toute réparation de support doit préserver masse, quantité de mouvement et énergie thermique relatives dans la partie fluide.

\paragraph{Ouvertures et domaines segmentés.}
Les distributions injectées doivent être compatibles avec le $k_BT$, la masse, le débit et le Mach calibrés. Avant tout split ou transfert, il faut auditer la suppression aux sorties, la réinjection, le bilan de masse par face et la capacité du transport thermique à alimenter les zones transverses. Un dépeuplement créé par une condition limite incohérente ne doit pas être compensé silencieusement par le resampling.

\paragraph{Resampling.}
Le resampling est un contrôle de représentation, pas une équation d'état ni une fermeture incompressible. Il ne doit être activé qu'après démonstration qu'un support statistique insuffisant subsiste dans un régime physique correctement calibré. Les critères doivent être géométriques et statistiques, dépendre de la population effective et de la fraction fluide, rester paramétrables, et conserver exactement masse, moment, énergie thermique et espèces. Le split seul peut supprimer les cellules pauvres mais provoquer une cascade de masses faibles et un coût élevé ; les résultats 0493o doivent donc être requalifiés sur les nouveaux régimes avant promotion en chemin de production.

\paragraph{Fermeture quasi incompressible Q6/Q9 et viriel.}
Une projection ne doit pas servir à masquer un Mach intrinsèquement trop élevé. La première action consiste à choisir $U$, $a$, $\Delta t$ et $k_BT$ afin d'obtenir un domaine $Re$--$Ma$ accessible. Q6/Q9 doit ensuite être évalué comme correction du fluide calibré : réduction des modes compressifs, conservation du transport transverse et absence de modification indue de $\nu$, $c_s$ ou $D_{\mathrm{self}}$. La fermeture virielle doit rester un modèle EOS explicite, distinct du maintien du support particulaire.

\paragraph{Mélanges multi-espèces.}
La calibration doit être répétée pour toute modification de masses, températures, occupations, règles de collision ou thermostat par espèce. Les grandeurs barycentriques ne suffisent pas nécessairement : les diffusivités propres et croisées, le Schmidt de chaque composant et la réponse acoustique du mélange peuvent devenir déterminants. Les modes Q6 barycentriques et le mode \texttt{independent\_masked}, ainsi que le resampling conservatif par type, doivent être qualifiés relativement à cette référence multi-espèces. Pour Q6, la mesure de $c_s$ devient non applicable ; la comparaison porte alors sur le transport transverse, la divergence, les modes longitudinaux résiduels, les glissements entre espèces et les corrections directes par type.

\subsection{Contrat cible pour une référence de co-simulation}

Le runner 0493w1 est actuellement un outil de cas unique fonctionnel. À terme, il doit devenir une interface stable du code : l'utilisateur fournit la configuration SRC et les échelles de son domaine, puis reçoit un artefact de référence utilisable par le solveur, un outil externe ou une boucle de co-simulation. Le contrat cible devrait comprendre :
\begin{itemize}[leftmargin=*]
  \item un manifeste exhaustif des paramètres qui définissent le fluide, du domaine de calibration, des modes, cadences et graines ;
  \item l'identifiant du commit, le hash du binaire et un hash canonique de la configuration ;
  \item les statuts indépendants \texttt{viscosityStatus}, \texttt{soundStatus} et \texttt{diffusionStatus}, accompagnés des résidus, $R^2$ et incertitudes ;
  \item un mode ensemble qui moyenne automatiquement les signaux TG et acoustiques avant régression ;
  \item deux profils explicites, \texttt{cell\_equivalent} et \texttt{full\_domain}, sans ambiguïté sur la taille de cellule et les longueurs d'onde ;
  \item des recommandations machine-readable : $U(Ma_{\mathrm{cible}})$, $Re(Ma_{\mathrm{cible}})$, $U(Re_{\mathrm{cible}})$, $Ma(Re_{\mathrm{cible}})$ et avertissements de régime cinétique ;
  \item une politique d'archivage non destructive empêchant l'écrasement silencieux d'une référence existante.
\end{itemize}

Les sorties canoniques actuelles sont
\begin{verbatim}
analysis/fluid_calibration_0493w1.json
analysis/fluid_calibration_0493w1.csv
analysis/README_0493w1_RESULTS.md
analysis/tg_decay_0493w1.csv
analysis/sound_mode_0493w1.csv
analysis/msd_0493w1.csv
\end{verbatim}
Le JSON constitue l'interface principale pour la co-simulation ; les trois séries CSV permettent de recalculer les fits, de construire un ensemble multi-graines et d'auditer les diagnostics sans conserver les dumps particulaires volumineux.

Un appel représentatif du candidat de référence s'écrit :
\begin{verbatim}
LIVE_PROGRESS=1 ALLOW_LARGE_DUMPS=1 \
RUN_ROOT=runs/calibrations/a256_dt002_k125_size128 \
Lx=0.5 Ly=0.5 NX=128 NY=128 GAMMA=20 \
DT=0.002 KBT=0.125 PARTICLE_MASS=1 \
ROTATION_ANGLE=1.5707963267948966 \
RANDOM_ROTATION_SIGN=true GRID_SHIFT_ENABLE=true \
THERMOSTAT_ENABLE=true THERMOSTAT_MODE=cell_relative_rescale \
THERMOSTAT_EVERY=1 THERMOSTAT_TARGET_KBT=0.125 \
THERMOSTAT_MIN_PARTICLES=3 \
TG_MODE_X=2 TG_MODE_Y=2 TG_STEPS=2000 TG_DUMP_COUNT=50 \
SOUND_MODE_X=4 SOUND_REPLICATES=4 \
SOUND_STEPS=1200 SOUND_DUMP_COUNT=120 \
MSD_GRID_MODE=full MSD_STEPS=3000 MSD_DUMP_COUNT=60 \
CHARACTERISTIC_U=0.1064 CHARACTERISTIC_L=0.24 \
bash scripts/run_0493w1_src_fluid_calibrator.sh
\end{verbatim}

Cette étape d'étalonnage ne rend pas inutiles les développements de support, de projection ou de fermeture. Elle fournit le référentiel qui permet enfin de décider lesquels sont nécessaires, dans quelles zones, et avec quel rapport bénéfice/coût. Les validations futures doivent donc être organisées en deux niveaux : d'abord le fluide SRC calibré et ses conditions limites, puis l'effet incrémental de Q6/Q9, du resampling, du viriel ou des traitements $\chi$/paroi.



\section{Cycle 0493x : Q6-g, Q6-g-f, dam-break et condition de pression liquide--gaz}
\label{sec:q6g-dambreak}

\subsection{Le dam-break comme cas discriminant historique}

Le cas \emph{dam-break} occupe une place particulière dans le développement de la méthode. Il combine simultanément une forte accélération volumique, un domaine liquide partiel, une interface fortement déformable, des impacts sur paroi, des zones de faible population et, lorsque le gaz est explicite, un contraste de masse important. Dans les générations historiques, il a donc échoué pour plusieurs raisons différentes. La redistribution dure et la fermeture virielle permettaient d'explorer la réponse à la compression, mais leur coût et leur caractère intrusif rendaient difficile une simulation suffisamment résolue et longue pour distinguer défaut de support, défaut de fermeture et véritable dynamique. Les premiers chemins Q6/Q9 supprimaient mieux les modes compressifs, mais ne garantissaient pas à eux seuls un support particulaire local. Le chantier resampling a ensuite traité ce problème de représentation, au prix d'une architecture complexe dont le bénéfice doit être évalué séparément de la projection.

Le passage au SRC CUDA résident change ce contexte expérimental. Il devient possible d'utiliser le dam-break non plus seulement comme démonstrateur qualitatif, mais comme test de régression long, avec plusieurs centaines de milliers de particules, dumps réguliers et diagnostics de géométrie/interfacialité. Le jalon 0493x est donc important au regard des chantiers resampling et viriel : pour la première fois, la chaîne de projection elle-même peut être isolée, instrumentée puis qualifiée sur ce cas sans que le coût de calcul ne masque immédiatement les mécanismes en jeu.

\subsection{Découverte du défaut de séquencement force--advection--projection}

Le défaut décisif a été mis en évidence sur un cas plus simple que le dam-break : une boîte fermée entièrement remplie de liquide, sous gravité uniforme, avec Q6 actif. La séquence historique applique le kick de force dans le transport, advecte les particules, effectue la collision, puis projette le champ de vitesse. Si
\begin{equation}
  v_y^{\star}=v_y^n+g\Delta t,
\end{equation}
le streaming utilise déjà
\begin{equation}
  y^{n+1}=y^n+v_y^{\star}\Delta t
  =y^n+v_y^n\Delta t+g\Delta t^2.
\end{equation}
Une projection Q6 appliquée après ce déplacement peut corriger la vitesse, mais ne peut pas annuler le déplacement déterministe $g\Delta t^2$ déjà accumulé. Pour $g=-0.5$ et $\Delta t=0.005$, ce terme vaut $-1.25\times10^{-5}$ par pas, valeur du même ordre que la lente sédimentation observée dans le liquide pourtant quasi incompressible.

Cette observation modifie l'interprétation du rôle de Q6. Le point essentiel n'est pas seulement d'obtenir une vitesse post-collision à faible divergence ; c'est que la \emph{vitesse de transport} utilisée par le streaming satisfasse la contrainte incompressible après prise en compte de la force. La branche 0493x3 introduit donc le terme Q6-g, qui désigne ici un Q6 \emph{force-aware} et non un nouvel opérateur elliptique : la force est incorporée à la vitesse tentative, cette vitesse est projetée, puis seulement utilisée pour déplacer les particules.

\subsection{De la preuve de concept à un Q6-g efficace}

La première variante, \texttt{q6ForceProjectionMode=prestream}, conservait volontairement deux projections par pas forcé : une projection avant streaming pour tester l'hypothèse, puis la projection post-collision historique. Les tests Taylor--Green à force nulle ont vérifié la neutralité du nouveau chemin ; sous forçage solénoïdal, l'écoulement restait cohérent ; dans la boîte liquide sous gravité, la sédimentation systématique disparaissait. Cette version démontrait la cause mais doublait pratiquement le coût de Q6.

Le mode \texttt{prestream\_single} 0493x4a repose sur l'observation que collision SRC et thermostat relatif ne déplacent pas les particules. Une seule projection avant le transport suffit donc : la vitesse issue de la collision est projetée au début du pas suivant avant de pouvoir servir au streaming. La séquence forcée devient
\begin{equation}
  \text{force}\;\longrightarrow\;Q6\;\longrightarrow\;\text{streaming}
  \;\longrightarrow\;\text{collision SRC}\;\longrightarrow\;\text{thermostat}.
\end{equation}
Sur le liquide fermé, le coût mesuré sur $1000$ pas passe d'environ $142.9\,\mathrm{s}$ pour le double solve à $95.5\,\mathrm{s}$ pour le solve unique, sans réintroduire la dérive gravitaire.

Le mode \texttt{prestream\_single\_fused} 0493x4b retire ensuite le dernier passage particulaire dédié au kick. Le dépôt CUDA construit directement le moment tentative $m(\bm v+\bm a\Delta t)$, le solveur Q6 calcule la correction, puis un même passage applique successivement la force physique et la correction Q6. La quantité de mouvement apportée par la force n'est pas incluse dans la correction de moment de Q6. Sur le même contrôle liquide, le temps tombe à environ $91.9\,\mathrm{s}$ pour $1000$ pas, et un run $5000$ pas reste stable. Ce mode devient donc le support temporel retenu pour la suite du chantier 0493x.

\subsection{Du liquide plein à la surface libre : \texttt{free\_surface\_masked}}

Le passage à un liquide partiellement rempli nécessite de distinguer l'incompressibilité du liquide et le vide ou le gaz extérieur. Le mode \texttt{free\_surface\_masked}, ajouté à \texttt{speciesQ6Mode}, construit pour la phase liquide un proxy de remplissage absolu
\begin{equation}
  f_c=\frac{M_{l,c}}{M_{l,\mathrm{ref}}},
\end{equation}
et utilise \texttt{speciesQ6MinOccupancyFraction} comme seuil de support. Le défaut global de ce paramètre reste $0.5$ ; les runners 0493x utilisent typiquement $0.25$, soit trois particules sur une occupation nominale $\gamma=10$ pour rester dans le support brut. Cette valeur est un réglage de qualification et non une constante physique.

Dans 0493x5a, une cellule liquide active adjacente à une cellule extérieure impose encore une pression de jauge nulle au milieu de la face. Pour une interface placée à une demi-maille, le terme Dirichlet contribue avec un facteur $2/h^2$ dans l'opérateur et $2/h$ dans la correction de vitesse. Un liquide demi-boîte $200\times200$ sous gravité reste alors sans sédimentation systématique sur $2000$ pas, avec une dérive verticale du centre de masse de l'ordre de $5\times10^{-5}$. Le dam-break vide 0493x5a2 est robuste avant impact sur la paroi opposée, mais la fragmentation post-impact révèle que le bord du support numérique ne peut pas être assimilé à l'interface physique.

L'extension 0493x5b ajoute un gaz explicite compressible, enregistré par \texttt{phaseFamily=gas} et \texttt{q6StrengthDeclared=0}. Le gaz participe aux collisions SRC et reçoit la force volumique mais n'est pas directement projeté par Q6. Cette séparation est indispensable : un Q6 barycentrique commun rendrait également le gaz incompressible. Le cas de qualification utilise volontairement un rapport de masse liquide/gaz de $1000$. Avec $m_g=1$, $k_BT=0.05$ et $|g|=0.5$, l'échelle de hauteur idéale $H_g=k_BT/(m_g|g|)$ vaut seulement $0.1$ pour une boîte de hauteur unité ; le gaz initialement uniforme se stratifie donc très fortement. Ce régime est un stress-test utile, mais ne doit pas être interprété comme un modèle quantitatif de l'air au-dessus de l'eau.

\subsection{Séparer support numérique et interface physique}

Les étapes 0493x6a--x6f ont transformé ce constat en architecture. Le diagnostic 0493x6a a d'abord reconstruit une pression gazeuse EOS sans modifier la physique. Le diagnostic 0493x6b a ensuite comparé le seuil de support Q6 à l'isovaleur physique $\alpha=0.5$, avec
\begin{equation}
  \alpha_{\rm raw}(c)=
  \frac{\sum_{s\in\mathrm{Liquid}} M_s(c)}
       {\sum_{s\in\mathrm{Liquid}} M_{\mathrm{ref},s}}.
\end{equation}
Le cycle 0493x6c rend ce champ résident CUDA et lui applique un lissage conservatif à cinq points, avec le coefficient expérimental fixe $\lambda=0.125$. Le coût mesuré est de l'ordre de $0.1\,\mathrm{ms}$ par projection sur $300\times150$, pour environ $0.69\,\mathrm{MiB}$ de stockage, donc négligeable devant le solveur Q6.

Le test 0493x6d utilisait directement la position locale de l'isovaleur pour remplacer le facteur demi-maille par $1/\theta$. Il est stable mais augmente sensiblement le conditionnement du CG et conserve un défaut conceptuel : il n'agit que lorsque l'interface coïncide avec le bord du support. Le diagnostic topologique 0493x6e démontre au contraire que, lorsque le dam-break se déforme, une part importante des traversées $\alpha=0.5$ se trouve entre deux cellules encore actives du carrier. Les catégories active--active, active--inactive et inactive--inactive sont donc utiles pour l'audit, mais ne doivent pas définir la physique de l'interface.

Le stencil 0493x6f sépare alors explicitement les deux objets :
\begin{equation}
  \text{pressureMask}=\text{carrierMask}\cap\{\alpha\ge 0.5\}.
\end{equation}
Une passe CUDA $O(N_{\rm cells})$ prépare une fois par solve les coefficients de faces. Une traversée physique $\alpha=0.5$ utilise $1/\theta$ ; les très petites coupures $\theta<0.10$ conservent provisoirement le facteur stabilisé $2$ ; une fin du carrier sans traversée physique n'est plus convertie artificiellement en surface à pression nulle. À $200$ pas, la couverture de l'interface est comprise entre $99.1\%$ et $100\%$ ; sur le dam-break $1000$ pas, elle reste entre environ $97.9\%$ et $100\%$, malgré la fragmentation post-impact. Le coût moyen de préparation est d'environ $50\,\mu\mathrm{s}$ par solve sur $300\times150$.

\subsection{Condition interfaciale \texorpdfstring{$p_l=p_g$}{pl=pg} et validation de jauge}

Le backend résident 0493x utilise $\phi$ comme potentiel de correction de vitesse. Pour un liquide de masse volumique de référence constante, la conversion pression--potentiel s'écrit
\begin{equation}
  \phi=\frac{\Delta t}{\rho_{l,\mathrm{ref}}}p
\end{equation}
à une constante de jauge près. Le cycle 0493x6g est le premier consommateur physique du stencil x6f et impose
\begin{equation}
  p_l|_\Gamma=p_g,
  \qquad
  \phi_{\Gamma,f}=\frac{\Delta t}{\rho_{l,\mathrm{ref}}}
  \left(p_{g,f}-p_{\mathrm{ref}}\right).
\end{equation}
La pression gazeuse de qualification est l'EOS scalaire
\begin{equation}
  p_g(c)=\frac{N_g(c)k_BT}{A_c},
\end{equation}
où toutes les espèces de \texttt{phaseFamily=gas} contribuent à $N_g$. La trace de pression est prise du côté $\alpha<0.5$ afin d'éviter une dilution artificielle par une cellule côté liquide ne contenant pas de gaz. La valeur $p_{\mathrm{ref}}$ est une jauge ; le runner EOS choisit par défaut la pression du gaz uniforme initial.

L'architecture est volontairement construite pour que la matrice CG de x6f ne change pas. Pour une face Dirichlet préparée de coefficient $c_f$,
\begin{equation}
  A_f(\phi_c)=\frac{c_f}{h^2}\phi_c,
  \qquad
  b_c\leftarrow b_c+\frac{c_f}{h^2}\phi_{\Gamma,f},
\end{equation}
et la correction de vitesse utilise la même distance effective. La pression gazeuse n'est donc jamais recalculée pendant les itérations CG. Son calcul EOS est fusionné avec la passe de construction géométrique et les valeurs de face sont écrites pendant la préparation du stencil.

Trois niveaux de validation ont été utilisés. D'abord, \texttt{scale=0} est un chemin nul strict : après un correctif, il n'alloue plus les buffers de pression et retombe sur le chemin x6f. Les petites différences particulaires observées entre répétitions sont du même ordre que le non-déterminisme flottant CUDA mesuré par des runs x6f/x6f. Ensuite, une pression interfaciale spatialement constante vérifie l'invariance de jauge. À \texttt{projectionTolerance}=$10^{-5}$, la différence x6f/pression constante reste de l'ordre de $10^{-12}$ sur les positions RMS et $10^{-10}$ sur les vitesses RMS ; à $10^{-7}$ elle descend vers le bruit de répétabilité, et à $10^{-9}$ elle atteint environ $1.8\times10^{-16}$ en position RMS et $1.3\times10^{-14}$ en vitesse RMS. Le défaut résiduel à la tolérance de production $10^{-5}$ est donc bien une erreur contrôlée d'arrêt du CG, non une violation de la jauge. Enfin, la relation entre $\phi_\Gamma$ et la différence de pression est auditée directement au niveau du roundoff.

\subsection{Dam-break liquide--gaz final 0493x6g}

La qualification finale compare deux runs $300\times150$, $\gamma=10$, $90000$ particules liquides et $360000$ particules gazeuses, soit $450000$ particules fluides, sur $1000$ pas : une référence x6f à $p_\Gamma=0$ et un run x6g avec $p_\Gamma=p_g$ EOS. Les deux runs utilisent le même générateur d'état, les mêmes paramètres physiques et la même tolérance de projection $10^{-5}$.

La référence termine en environ $131.1\,\mathrm{s}$ de temps solveur et le cas EOS en environ $122.8\,\mathrm{s}$. Cette différence ne doit pas être interprétée comme une accélération due à $p_g$ ; elle montre simplement qu'aucun surcoût systématique de x6g n'est visible à l'échelle du run. Le nombre maximal d'itérations Q6 passe de $435$ à $450$. Le front liquide diffère de seulement $-9.8\times10^{-5}$, le déplacement du centre de masse de $(-5.3\times10^{-5},-5.9\times10^{-4})$, et le support final reste voisin ($8296$ cellules contre $8321$). Le run EOS reste donc stable et très proche de la référence sur les observables macroscopiques.

Ce résultat est d'autant plus sévère que le gaz du stress-test devient fortement hétérogène : l'écart-type maximal audité de $p_g-p_{\mathrm{ref}}$ atteint environ $3.3\times10^5$ dans les unités du modèle. La faible différence sur le liquide confirme que, pour un fort rapport de densité, les variations de pression gazeuse ne dominent pas nécessairement la dynamique liquide. Elle ne valide pas pour autant l'EOS idéale comme modèle réaliste du gaz dans ce régime fortement stratifié ; elle valide le \emph{couplage numérique} d'une pression gazeuse non uniforme à l'interface.


\subsection{Projection Q6-g-f : restauration de densité intégrée, application face--particule et validation longue}
\label{subsec:q6-g-f}

\subsubsection{Définition et objectif}

Dans la suite du rapport, \textbf{Q6-g-f} désigne la génération de projection actuellement retenue pour le liquide à surface libre. Le nom permet de la distinguer rapidement du Q6 historique, du Q6 multi-espèces \texttt{independent\_masked} et du premier Q6-g uniquement \emph{force-aware}. Q6-g-f n'est pas un second solveur elliptique : il conserve le solveur FV/CG résident et compose, dans un même pas et un même second membre, quatre éléments désormais qualifiés : (i) la vitesse tentative incluant le forçage avant streaming, (ii) le domaine de pression défini par la géométrie $\alpha=0.5$, (iii) l'application particulaire compatible avec les corrections de faces et (iv) un retour lent de la densité liquide vers sa valeur nominale.

L'objectif physique n'est pas d'imposer artificiellement $\nabla\!\cdot\bm u=0$ jusque dans une cellule dont la densité particulaire a déjà dérivé. Le champ de remplissage brut est
\begin{equation}
  f_c=\frac{M_{l,c}}{M_{l,\mathrm{ref}}}.
\end{equation}
Dans le bulk liquide, Q6-g-f impose alors la contrainte relaxée
\begin{equation}
  \left.\nabla\!\cdot\bm u_{\mathrm{proj}}\right|_c
  =\frac{f_c-1}{\tau_\rho},
  \label{eq:q6gf-density-target}
\end{equation}
soit, dans l'ancienne paramétrisation par pas,
\begin{equation}
  \left.\nabla\!\cdot\bm u_{\mathrm{proj}}\right|_c
  =\frac{\beta_\rho}{\Delta t}(f_c-1),
  \qquad
  \beta_\rho=\frac{\Delta t}{\tau_\rho}.
\end{equation}
Une cellule sur-remplie reçoit ainsi une petite divergence positive et se détend ; une cellule sous-remplie reçoit une divergence négative. Lorsque $f_c=1$, la contrainte redevient la contrainte incompressible usuelle. La source est volontairement limitée au bulk du domaine de pression : la bande interfaciale reste régie par le stencil x6f et par sa condition de pression, de façon à ne pas confondre restauration de densité et loi de contrainte à l'interface.

La valeur qualifiée est $\tau_\rho=0.25$. À $\Delta t=0.005$, elle correspond à $\beta_\rho=0.02$ par pas ; à $\Delta t=0.0025$, elle donne automatiquement $\beta_\rho=0.01$. Cette formulation temporelle est essentielle pour pouvoir changer le pas de temps sans changer implicitement la constante physique de restauration.

\subsubsection{Assemblage du second membre et coexistence avec la pression gazeuse}

Le potentiel de projection reste celui de Q6-g. Pour la phase liquide, le second membre peut être résumé par
\begin{equation}
  b=-\nabla\!\cdot\bm u^{\star}
    +b_{\Gamma}(p_g-p_{\mathrm{ref}})
    +\frac{f-1}{\tau_\rho},
  \label{eq:q6gf-rhs}
\end{equation}
où $\bm u^{\star}$ est la vitesse tentative après prise en compte du forçage et $b_{\Gamma}$ désigne la contribution Dirichlet préparée sur les faces qui traversent $\alpha=0.5$. Le terme de pression x6g et le terme de restauration x7d sont donc additifs et consommés par le \emph{même} solveur CG ; aucune seconde projection n'est ajoutée. Le viriel explicite de densité x7a/x7b est désactivé et déclaré mutuellement exclusif avec cette fermeture, car son kick post-projection réintroduisait un ressort conservatif susceptible d'accumuler de l'énergie cohérente.

La géométrie interfaciale conserve la séparation établie par x6f. Le remplissage brut $f$ reste non borné pour le diagnostic et pour la source de restauration, tandis que la géométrie filtrée part de
\begin{equation}
  \alpha_0=\operatorname{clamp}(f,0,1)
\end{equation}
et applique le filtre à cinq points de coefficient $\lambda=0.125$. Le domaine de pression est
\begin{equation}
  \text{pressureMask}=\text{carrierMask}\cap\{\alpha\ge0.5\}.
\end{equation}
Une traversée physique de l'isovaleur utilise le coefficient $1/\theta$ ; les petites coupures $\theta<0.10$ gardent provisoirement le coefficient stabilisé $2$. Un correctif tardif du stencil a remplacé le garde numérique $\alpha_{\rm high}-\alpha_{\rm low}>10^{-14}$ par la condition strictement topologique $\alpha_{\rm high}-\alpha_{\rm low}>0$. Le prédicat de traversée garantit déjà $\alpha_{\rm high}>\alpha_{\rm low}$ ; le changement évite qu'une traversée réelle quasi dégénérée soit comptée par l'audit topologique mais omise par le stencil. Le défaut avait été isolé sur une seule face parmi $2758$ au pas $2700$ d'un run fin long.

\subsubsection{Reconstruction face--particule}

L'analyse du champ réellement redéposé après projection a montré que la convergence du flux de faces n'impliquait pas à elle seule une divergence également faible après application aux particules. Deux corrections ont été conservées dans Q6-g-f. D'abord, le correctif x6h-A reconstruit les corrections des faces ouest et sud sur les frontières physiques basses, qui n'ont pas de cellule propriétaire voisine dans le stockage east/north. Il utilise la même convention cible--avant que les faces hautes et n'applique cette reconstruction qu'aux cellules du \texttt{pressureMask}.

Ensuite, le chemin x6h-B1, activé par \texttt{MPCD\_Q6\_FACE\_TO\_PARTICLE\_RT0\_0493X6H\_B1}, remplace l'incrément constant par cellule par une reconstruction affine entre les corrections de faces. Si $\delta u_W$ et $\delta u_E$ sont les corrections normales aux faces verticales d'une cellule et $\xi\in[0,1]$ la coordonnée locale, l'application suit le schéma
\begin{equation}
  \delta u_x(\xi)=\delta u_W+\xi(\delta u_E-\delta u_W),
\end{equation}
avec une expression analogue entre les faces sud et nord. La reconstruction a la même divergence discrète que le champ FV qui a été projeté. Elle réutilise les buffers de faces déjà résidents, n'ajoute ni stockage persistant par espèce ni passage particulaire supplémentaire et reste limitée au chemin surface libre, espèce liquide projetée unique et application force+Q6 fusionnée actuellement qualifié.

Le diagnostic x6h-B0, normalement désactivé, a servi à localiser le défaut post-application en séparant bulk, interface, paroi et zones mixtes. Il n'appartient pas au chemin de production ; son allocation et son kernel sont absents lorsqu'il est désactivé.

\subsubsection{Implémentation CUDA et diagnostics}

La séquence nominale Q6-g-f est maintenant
\begin{equation}
  \text{dépôt de }\bm u^{\star}
  \;\longrightarrow\;
  \text{construction }\alpha\text{ et stencil}
  \;\longrightarrow\;
  \text{assemblage RHS Q6-g-f}
  \;\longrightarrow\;
  \text{CG}
  \;\longrightarrow\;
  \text{application B1 force+Q6}
  \;\longrightarrow\;
  \text{streaming}
  \;\longrightarrow\;
  \text{collision SRC}
  \;\longrightarrow\;
  \text{thermostat}.
\end{equation}
Le paramètre \texttt{q6DensityRelaxationTime} est l'entrée physique préférée ; \texttt{q6DensityRelaxationBeta} reste disponible pour compatibilité. Les deux sont mutuellement exclusifs lorsqu'ils sont positifs et l'implémentation vérifie que le coefficient effectif $\Delta t/\tau_\rho$ reste dans $[0,1]$. Une valeur nulle est un no-op exact. Le chemin requiert une projection Q6 CUDA active et la qualification actuelle conserve \texttt{virialDensityKickEnable=false}.

L'audit \texttt{cuda\_species\_q6\_independent\_masked\_0493w5.csv} publie désormais le coefficient effectif \texttt{q6DensityRelaxationBeta}, la constante de temps \texttt{q6DensityRelaxationTime} et \texttt{densityRelaxationTargetDivRms}. La colonne \texttt{divAfterProjectedFaceFluxRms} mesure le résidu par rapport à la contrainte active, qui n'est donc plus nécessairement nulle ; \texttt{divAfterAppliedCellVelocityRms} reste le diagnostic du champ cellulaire effectivement obtenu après l'application aux particules. Les audits x6c/x6f/x6g restent inchangés et permettent de vérifier séparément conservation géométrique, couverture du stencil et relation pression--potentiel.

\subsubsection{Validation de fermeture volumique et robustesse au raffinement}

La première validation longue de la restauration de densité a été réalisée sans pression gazeuse interfaciale, afin d'isoler l'opérateur. Sur le dam-break $300\times150$, $\gamma=10$, $\Delta t=0.005$, $5000$ pas et $\tau_\rho=0.25$, la somme brute de remplissage liquide reste $9000$ cellules nominales et la somme géométrique bornée atteint $8554.2$, soit un excès de seulement $445.8$ cellules ou $4.95\,\%$. Le même cas avant restauration de densité, avec les correctifs A+B1 seulement, conservait un excès d'environ $1967.4$ cellules ; Q6-g-f réduit donc d'environ $77\,\%$ ce défaut de compression sans recourir au resampling ni à un kick viriel post-projection.

La qualification x7e réactive ensuite simultanément la pression gazeuse x6g et la restauration x7d. À $t=2.5$, le couple $300\times150$, $\Delta t=0.005$, $500$ pas et $600\times300$, $\Delta t=0.0025$, $1000$ pas donne des excès géométriques normalisés de $7.37\,\%$ et $7.43\,\%$, soit un rapport fine/coarse de $1.009$. Les deux contributions sont présentes dans le même RHS et les audits x6g/x7d correspondent sur tous les pas contrôlés.

Le test décisif est le temps long $t=25$. Avec $p_l|_\Gamma=p_g$ actif, le coarse $300\times150$, $5000$ pas termine avec
\begin{equation}
  E_{V,c}=5.016\,\%,
\end{equation}
et le fine $600\times300$, $10000$ pas avec
\begin{equation}
  E_{V,f}=5.864\,\%.
\end{equation}
Les résidus de la contrainte projetée restent de l'ordre de quelques $10^{-5}$ avec \texttt{projectionTolerance}=$10^{-5}$, les audits x6g et x7d sont appariés sur $201/201$ pas, et aucun retournement tardif de volume analogue à celui du kick viriel explicite n'apparaît. Sans pression gazeuse, les deux défauts correspondants sont $4.953\,\%$ et $5.669\,\%$ ; l'essentiel de l'écart de maille existe donc déjà dans l'opérateur/liquide sous-jacent et n'est pas créé par x6g.

Cette comparaison ne doit pas être surinterprétée comme une convergence hydrodynamique classique. Une calibration SRC homogène donne $\nu_c\simeq6.97\times10^{-4}$ pour la maille coarse et $\nu_f\simeq3.03\times10^{-4}$ pour la maille fine ; pour l'échelle gravitaire du dam-break, le Reynolds passe ainsi d'environ $726$ à $1670$. Halver simultanément la maille et $\Delta t$ à $\gamma$ fixé ne conserve donc pas les coefficients de transport MPCD. En revanche, au regard de l'objectif de Q6-g-f, le résultat est précisément favorable : malgré un changement important du régime visqueux et un facteur deux de résolution, l'erreur volumique reste bornée à quelques pourcents pendant $t=25$ et ne présente pas de dérive séculaire. Q6-g-f est donc considéré comme \emph{qualifié pour la fermeture volumique/incompressible} dans ce périmètre, sans que cela constitue encore une preuve de convergence complète de la solution du dam-break.

\subsection{Positionnement par rapport au resampling, au viriel et à la tension superficielle}

Le cycle 0493x clarifie finalement les rôles respectifs des trois grands chantiers historiques. Le viriel est une fermeture d'équation d'état et une source de contrainte bulk/paroi ; il ne doit pas compenser un défaut de séquencement temporel de la projection. Le resampling est un mécanisme de qualité du support particulaire ; il doit intervenir lorsque la représentation locale devient insuffisante, pas pour imposer l'incompressibilité. Q6-g agit sur la vitesse de transport : il garantit que la force est incluse dans la vitesse tentative avant la projection qui précède le streaming. Q6-g-f complète cette fonction par une restauration de densité intégrée à la contrainte de divergence et par une application face--particule compatible avec le champ FV projeté. Enfin, le stencil x6f/x6g traite la condition de contrainte à l'interface sans identifier celle-ci au bord du support.

Cette séparation a servi directement de point d'entrée au cycle 0493x9. Avec la convention de signe retenue pour la courbure, la condition normale est
\begin{equation}
  p_l|_\Gamma=p_g+\sigma\kappa,
\end{equation}
et donc
\begin{equation}
  \phi_{\Gamma,f}=
  \frac{\Delta t}{\rho_{l,\mathrm{ref}}}
  \left(p_{g,f}+\sigma\kappa_f-p_{\mathrm{ref}}\right).
\end{equation}
Cette extension est maintenant implémentée sans modifier la matrice Q6 : la courbure est préparée sur un champ auxiliaire p3/Scharr, interpolée au crossing physique $\alpha=0.5$, puis ajoutée au potentiel de Dirichlet x6g avant le solve CG et l'application B1. Les qualifications x9 montrent que cette voie produit une tension superficielle bulk effective ; les détails, la régularisation des rayons sous-résolus et les limites du mouillage sont donnés dans la section 0493x9 ci-dessous.


\subsection{Paramètres, flags expérimentaux et runners du cycle 0493x}

Les choix qui affectent la physique ou l'ordre d'intégration sont désormais séparés des gates d'exécution et de diagnostic. Le paramètre \texttt{q6ForceProjectionMode} accepte \texttt{legacy}, \texttt{prestream}, \texttt{prestream\_single} et \texttt{prestream\_single\_fused}; le dernier est le chemin qualifié pour Q6-g-f. Le mode \texttt{speciesQ6Mode=free\_surface\_masked} réutilise \texttt{speciesQ6MinOccupancyFraction} comme seuil de remplissage absolu du liquide. Les qualifications finales utilisent $0.10$ afin de conserver une bande numérique de carrier suffisamment large sans confondre ce seuil de support avec l'interface physique $\alpha=0.5$.

La restauration de densité dispose toujours des entrées historiques \texttt{q6DensityRelaxationBeta} et \texttt{q6DensityRelaxationTime}; la seconde est l'entrée physique recommandée et impose $\beta_{\rm step}=\Delta t/\tau_\rho$. Le jeu signé ajoute quatre paramètres :
\begin{itemize}[leftmargin=2em]
  \item \texttt{q6DensityRelaxationCompressionGateEnable}, booléen, qui active l'admission cohérente de la branche de compression ;
  \item \texttt{q6DensityRelaxationCompressionThresholdFill}, seuil positif sur le défaut de remplissage ;
  \item \texttt{q6DensityRelaxationTractionThresholdFill}, seuil positif pour la branche de déplétion ;
  \item \texttt{q6DensityRelaxationTractionGain}, gain de cette branche négative.
\end{itemize}
Le gate est un classifieur et non une zone morte : après admission, le RHS conserve le défaut complet. Le jeu de production actuellement qualifié est $\tau_\rho=0.25$, gate de compression actif, seuils exprimés dans les runners comme $+3$ et $-6$ particules relativement à $\gamma$, et gain de traction égal à $1$. Ainsi, pour Poiseuille à $\gamma=20$, les seuils écrits dans \texttt{params.kv} valent respectivement $0.15$ et $0.30$ ; pour le dam-break à $\gamma=10$, ils valent $0.30$ et $0.60$. La variante brute signée sans gate reste utile comme ablation, mais n'est plus le profil de production.

Le paramètre \texttt{projectionTolerance} garde son défaut C++ $10^{-10}$ ; la campagne de qualification finale impose toutefois \texttt{PROJECTION\_TOLERANCE=1e-5} et, par défaut, \texttt{projectionMaxIterations=800}. Le gate \texttt{MPCD\_Q6\_FACE\_TO\_PARTICLE\_RT0\_0493X6H\_B1} reste activé dans Q6-g-f. Le nouveau gate \texttt{MPCD\_Q6\_G\_F\_RESIDENT\_CG\_0493X7J} est activé par défaut dans le backend : il change seulement l'exécution du CG, pas l'opérateur ni les paramètres physiques. Le runner x7i force le chemin x7j multi-bloc pour Q6-g-f et conserve la politique 0407 rapide du Q6 historique sur les petits domaines, ce qui évite de ralentir artificiellement la référence \texttt{src-q6}.

Les alias génériques de \texttt{scripts/src\_mpcd\_run\_common\_0434.sh} sont désormais \texttt{Q6\_GF\_DENSITY\_RELAXATION\_TIME}, \texttt{Q6\_GF\_MIN\_FILL\_FRACTION}, \texttt{Q6\_GF\_DENSITY\_COMPRESSION\_GATE\_ENABLE}, \texttt{Q6\_GF\_DENSITY\_COMPRESSION\_THRESHOLD\_PARTICLES}, \texttt{Q6\_GF\_DENSITY\_TRACTION\_THRESHOLD\_PARTICLES} et \texttt{Q6\_GF\_DENSITY\_TRACTION\_GAIN}. Les variantes monophasées peuvent synthétiser un registre d'une espèce, ou utiliser un registre externe via \texttt{Q6\_GF\_EXTERNAL\_SPECIES}; \texttt{Q6\_GF\_HAS\_GAS\_PHASE} sélectionne explicitement la présence d'une phase gazeuse pour le chemin x6g/x7m. Le runner de qualification \texttt{run\_0493x7i\_q6\_g\_f\_physical\_qualification.sh} encapsule ces choix sous le préfixe \texttt{X7I\_*}, avec notamment $\tau_\rho=0.25$, seuils $3/6$, gain $1$, remplissage minimal $0.10$, tolérance $10^{-5}$ et profil Poiseuille low-Mach $(\gamma,N_x,N_y,\Delta t,a_x)=(20,192,96,10^{-3},4\times10^{-4})$.

Les gates x6 de géométrie et de pression gazeuse restent recensés dans l'inventaire CSV. x7k et x7l ne créent aucun interrupteur supplémentaire : les diagnostics Q6-g-f et les réductions hôte du thermostat résident sont simplement limités au premier pas et à la cadence \texttt{summaryEvery}, tandis que l'opération physique du thermostat reste exécutée à chaque pas. x7m, x7o, x7p et x7q sont également automatiques une fois leur chemin applicable sélectionné. En particulier, \textbf{x7q n'a pas de flag utilisateur} : la fermeture exacte du mode $k=0$ est activée uniquement lorsque B1 est actif, que l'espèce projetée occupe le domaine de pression complet, qu'au moins une direction est périodique et que le kick viriel de densité n'est pas demandé.

\subsection{Élargissement topologique et intégration des runners : 0493x7f--x7i}

Le jalon x7f transforme Q6-g-f d'un chemin essentiellement fermé de dam-break en une brique de projection utilisable sur les familles statiques déjà représentées par le backend Q6 CUDA résident. Le contrat couvre la boîte périodique, le canal périodique en $x$ avec parois statiques en $y$, la boîte fermée statique, les paires inlet/outlet pleine face et les ouvertures segmentées qualifiées. Il ne généralise pas silencieusement les pistons, les domaines tronqués mobiles, les solides immergés ou le resampling : ces briques conservent leurs garde-fous propres. Le correctif x7f-fix2 relaxe seulement le garde historique du chemin wall-simple : une face \texttt{specular} de réflexion pure est désormais admise sans imposer artificiellement la création de particules virtuelles ; l'arithmétique de collision n'est pas modifiée.

x7g déplace ensuite le Darcy--Brinkman dans la vitesse tentative de Q6-g-f. Les modes basés sur une relaxation moyenne --- \texttt{mean}, \texttt{classic}, \texttt{cell\_mean} et les variantes combinées \texttt{mean\_outward\_bath}, \texttt{mean\_oriented\_bath}, \texttt{brinkman\_outward\_bath} --- sont appliqués sur l'état CUDA avant la projection, puis ne sont plus rejoués après collision. Le champ $\chi$ reste alors un champ de domaine fictif : \texttt{darcyInitialDeactivateBelowChi} doit être négatif afin qu'un obstacle Darcy ne soit pas interprété comme un trou de support liquide par \texttt{free\_surface\_masked}. Les contributions de particules virtuelles $\chi$ associées à la collision restent dans l'étage SRC auquel elles appartiennent.

x7h factorise cette configuration dans les runners \texttt{run\_ok}. Le mode \texttt{src-q6-g-f} sélectionne \texttt{prestream\_single\_fused}, B1, le stencil x6f, le registre d'espèces et la restauration de densité sans modifier les paramètres physiques des modes de comparaison \texttt{src} et \texttt{src-q6}. x7i construit enfin une campagne physique commune sur quatre cas --- Taylor--Green forcé, Poiseuille, bend-pipe Darcy et same-face IO --- et compare les trois modes à partir des mêmes états et observables. Cette campagne devient le test multi-CL de référence du chemin de production.

\subsection{Restauration de densité signée et filtrage cohérent du bruit d'occupation}

L'expérience Poiseuille a montré qu'une restauration locale fondée directement sur \texttt{rawFill-1} peut injecter dans le RHS un signal du même ordre que les fluctuations statistiques MPCD. Le problème est différent de la fuite de moment corrigée plus tard par x7q : même avec une conservation globale exacte, un RHS brut peut produire des corrections locales inutilement fortes et un profil spatial irrégulier.

La variante signée retenue applique donc deux seuils cohérents. Pour la compression, une cellule est admise seulement si son défaut positif dépasse \texttt{q6DensityRelaxationCompressionThresholdFill} et si au moins un voisin de face dépasse le même seuil. Pour la déplétion, le critère est analogue sous le seuil négatif \texttt{q6DensityRelaxationTractionThresholdFill}, puis la cible est multipliée par \texttt{q6DensityRelaxationTractionGain}. Ce test spatial élimine les excursions isolées sans changer la loi de relaxation dans les structures cohérentes : après admission, le défaut complet est utilisé. Le jeu $(\tau_\rho,n_+,n_-,g_-)=(0.25,3,6,1)$ a été conservé pour la qualification finale.

Cette distinction est importante pour l'interprétation : le gate ne constitue ni un resampling caché, ni une loi de pression additionnelle. Il ne modifie ni masse ni population ; il sélectionne seulement la partie spatialement cohérente de la cible de divergence associée à la restauration Q6-g-f. Le resampling reste une correction distincte du support particulaire, et le viriel une fermeture distincte de pression/capacité.

\subsection{Résidence CUDA et réduction du coût : 0493x7j--x7l}

L'audit x7i a mis en évidence un défaut d'architecture plutôt que de physique : le CG masqué Q6-g-f réalisait des réductions et synchronisations hôte à chaque itération. À tolérance identique $10^{-5}$ sur Taylor--Green $128\times128$, Q6 et Q6-g-f demandaient un nombre d'itérations comparable, mais le solveur masqué Q6-g-f pouvait consacrer à lui seul environ $44\,\mathrm{ms}$ par pas aux itérations et transferts associés.

x7j remplace ce chemin par un CG coopératif entièrement résident. Un unique kernel CUDA porte la récurrence de Krylov complète, y compris $Ap$, $p^TAp$, la mise à jour de $\phi$ et $r$, le test $r^Tr$, la mise à jour de $p$ et la suppression périodique du mode constant pour les domaines complets. Les réductions inter-blocs utilisent \texttt{cooperative\_groups}; aucune donnée GPU--CPU n'est transférée au milieu d'une itération et le host ne récupère qu'un petit état final du solveur. Le domaine complet utilise le Laplacien volumes finis standard ; les domaines partiels conservent le masque et les coefficients de faces préparés par x6f. Un fallback vers l'ancien CG reste possible si le lancement coopératif n'est pas disponible.

x7k et x7l suppriment ensuite les coûts de télémétrie du chemin chaud. Les audits détaillés Q6-g-f ne sont calculés qu'au premier pas et aux pas de résumé ; le thermostat cellulaire résident reste physiquement exécuté à chaque pas, mais ses réductions hôte de diagnostic suivent la même cadence. Ces changements ne modifient ni la projection ni la loi thermique ; ils rapprochent le coût mesuré du coût intrinsèque de la chaîne CUDA résidente.

\subsection{Topologie monophase et symétrie discrète : 0493x7m, x7o et x7p}

x7m corrige une ambiguïté de topologie apparue lorsque \texttt{free\_surface\_masked} est réutilisé sur un fluide monophase rempli. Le seuil $\alpha=0.5$ ne représente une interface physique liquide--gaz que si le registre contient effectivement une espèce gazeuse. En l'absence de gaz enregistré, le domaine de pression reste monophase et complet ; une fluctuation d'occupation autour de $\alpha=0.5$ ne crée donc plus une frontière interne artificielle de Dirichlet. Cette distinction est essentielle pour Poiseuille et Taylor--Green, où le masque numérique de l'espèce ne doit pas fabriquer une surface libre inexistante.

x7o supprime ensuite un biais d'orientation de la reconstruction indépendante. Dans le domaine complet, les vitesses de face sont formulées de manière équivariante par réflexion et la correction centrée d'une cellule est reconstruite par la moyenne de ses deux faces opposées, au lieu de privilégier implicitement le propriétaire est/nord. x7p applique la même convention au Q6 commun. Ces changements ne touchent pas la condition interfaciale x6f des domaines partiels ; ils garantissent que la discrétisation monophase ne dépend pas du choix arbitraire d'une orientation de stockage des faces.

\subsection{Fermeture exacte du moment périodique au niveau particulaire : 0493x7q}

La dernière anomalie majeure a été isolée par un Poiseuille low-Mach. Le correctif x7d-v2-fix2 neutralisait le moment de la correction B1 à partir d'une estimation centrée cellule. Or B1 applique à une particule $i$ de la cellule $c$ une reconstruction affine de type RT0. Dans la direction $x$,
\begin{equation}
 \delta u_{x,i}
 = c_{x,c} + (2\xi_i-1)\left(\delta u_{e,c}-c_{x,c}\right),
\end{equation}
où $c_{x,c}$ est la correction centrée, $\delta u_{e,c}$ la correction de la face est et $\xi_i\in[0,1]$ la coordonnée locale de la particule. Le moment réellement appliqué dans la cellule vaut donc
\begin{equation}
 \sum_{i\in c}m_i\delta u_{x,i}
 = M_c c_{x,c}
 + \left(\delta u_{e,c}-c_{x,c}\right)
   \sum_{i\in c}m_i(2\xi_i-1).
\end{equation}
Le second terme ne s'annule que si le barycentre massique instantané des particules coïncide exactement avec le centre de la cellule. Ce n'est pas vrai pour un échantillon MPCD fini. Une fermeture basée sur $\sum_c M_cc_{x,c}$ laisse donc le résidu
\begin{equation}
 \Delta P_x^{\rm res}
 = \sum_c \left(\delta u_{e,c}-c_{x,c}\right)
   \sum_{i\in c}m_i(2\xi_i-1),
\end{equation}
avec une expression analogue en $y$. Ce résidu est petit localement mais s'accumule à chaque pas ; dans une direction périodique, il contamine directement le mode uniforme $k=0$ que le gradient d'une pression interne ne doit pas modifier.

Le probe discriminant utilise un domaine entièrement périodique, $N=368640$, $a_x=4\times10^{-4}$, $\Delta t=10^{-3}$ et $t=2$. La dynamique imposée exige $\bar u_x=a_xt=8\times10^{-4}$ et $P_x=294.912$. Avant x7q, le cas Q6-g-f signé ne donne plus que $\bar u_x=5.8637\times10^{-5}$ et $P_x=21.616$, avec en outre $P_y\simeq29.47$, alors que le résidu du solveur de projection reste de l'ordre de $10^{-5}$. L'échec est donc dans l'application face--particule, non dans le Poisson/CG.

x7q conserve le pré-correctif centré x7d-v2 mais ajoute, seulement pour B1, \texttt{fullDomain} et les directions périodiques, une mesure exacte du moment RT0 réellement échantillonné aux positions particulaires. Un premier kernel spécialisé applique force et B1 et réduit sur GPU le moment brut et la masse corrigée ; un petit kernel de finalisation calcule la vitesse uniforme résiduelle ; un second passage particulaire soustrait cette vitesse dans les seules directions périodiques. Il n'y a pas de téléchargement ni de synchronisation hôte dans ce chemin de production. Les composantes non périodiques restent libres de transmettre une réaction aux parois. Le chemin n'est pas activé avec le kick viriel de densité.

La portée est volontairement étroite : un dam-break partiel a \texttt{fullDomain=0} et continue d'utiliser le kernel B1 historique, sans passage particulaire supplémentaire. Sur le même probe périodique après x7q, $\bar u_x=7.999999999989\times10^{-4}$, $P_x=294.9119999996$, $\bar u_y\simeq3.9\times10^{-17}$ et $P_y\simeq1.4\times10^{-11}$. La fermeture du mode $k=0$ est ainsi vérifiée au niveau du roundoff sans modifier la condition de surface libre.

\subsection{Qualification physique finale de Q6-g-f signé + x7j + x7q}

Le dam-break $300\times150$, $\gamma=10$, $5000$ pas, reste qualitativement identique aux runs qualifiés précédents après x7q. Le temps de simulation reporté par le solveur est d'environ $113.5\,\mathrm{s}$ pour $t=25$ ; l'absence de régression est cohérente avec le fait que x7q n'entre jamais dans le chemin \texttt{fullDomain=0}.

Poiseuille est le test le plus discriminant de la réparation. Avec le jeu signé $\tau_\rho=0.25$, seuils $3/6$ particules et gain $1$, le run $60000$ pas donne sur la fenêtre $55000$--$60000$ : $R^2=0.95939$, $\nu_{\rm eff}=0.0082273$, $\bar u=0.0040751$ et $u_c=0.0065069$, à comparer à la calibration Taylor--Green $\nu\simeq0.008424$. L'analyse $40000$--$60000$ donne $\nu_{\rm eff}=0.008276\pm0.000181$. Le résidu spatial moyen représente environ $7.49\,\%$ de l'amplitude du profil et la fluctuation temporelle environ $10.54\,\%$, mais les résidus moyens des fenêtres $40000$--$50000$ et $50000$--$60000$ ne sont corrélés qu'à $0.158$. L'amplitude des fluctuations persiste donc statistiquement, sans mise en évidence d'une déformation spatiale stationnaire reproductible.

La campagne x7i finale compare ensuite les trois modes sur quatre topologies avec les mêmes conventions de post-traitement :
\begin{center}
\begin{tabular}{@{}lccc@{}}
\toprule
Observable & SRC & Q6 & Q6-g-f \\
\midrule
TG : déficit relatif aux c\oe urs & $-0.99972$ & $-0.05748$ & $-0.00890$ \\
TG : $\mathrm{std}(N)/\langle N\rangle$ & $2.971$ & $0.295$ & $0.205$ \\
Poiseuille : $R^2$ & $0.99475$ & $0.99484$ & $0.99395$ \\
Poiseuille : $\nu_{\rm eff}$ & $0.009867$ & $0.010171$ & $0.010314$ \\
Poiseuille : $\bar u$ & $0.003561$ & $0.003496$ & $0.003439$ \\
Bend : $u_{\rm coh,tail}/u_{\rm ref}$ & $0.400$ & $1.044$ & $1.115$ \\
Bend : $u_{\rm far,tail}/u_{\rm ref}$ & $0.057$ & $1.049$ & $1.097$ \\
Same-face IO : $u_{\rm coh,tail}/u_{\rm ref}$ & $0.189$ & $0.360$ & $0.345$ \\
Same-face IO : $u_{\rm far,tail}/u_{\rm ref}$ & $0.046$ & $0.068$ & $0.065$ \\
\bottomrule
\end{tabular}
\end{center}

Taylor--Green constitue un PASS fort : Q6-g-f supprime presque entièrement le déficit de population des c\oe urs et réduit encore l'hétérogénéité globale par rapport à Q6. Poiseuille est un PASS : le profil et le transport restent proches de SRC/Q6 dans la campagne commune et le run long confirme la viscosité tardive. Bend-pipe est un PASS fort : la transmission cohérente et le far-field passent d'une réponse faible en SRC à une réponse de l'ordre de $u_{\rm ref}$ avec les deux projections, Q6-g-f étant légèrement supérieur au Q6 précédent. Same-face IO valide la robustesse et l'amplitude de réponse de Q6-g-f ; les temps $T_{50}$ far-field normalisés au plateau propre de chaque mode ne doivent toutefois pas être comparés directement lorsque les plateaux diffèrent, et le seuil physique commun $0.1u_{\rm ref}$ n'est pas atteint au far-field sur la fenêtre disponible.

Le statut combiné est donc : Q6-g-f signé + x7j + x7q est qualifié sur boîte périodique, canal à parois, inlet/outlet segmenté, Darcy/bend-pipe et dam-break à surface libre, sans retour de la pathologie de moment qui détruisait Poiseuille. La prochaine extension physique peut ainsi porter sur la géométrie capillaire et le saut de pression plutôt que sur une nouvelle refonte de la fermeture incompressible de base.


\section{Cycle 0493x8 : conditions limites segmentées Poiseuille--Neumann et validation bibliographique de von Kármán}
\label{sec:x8-vk-validation}

\subsection{Objectif et recentrage du chantier}

Après la qualification multi-topologie 0493x7q, la priorité du cycle 0493x8 n'est plus d'élargir simultanément tous les modes SRC/Q6 historiques, mais de vérifier que le chemin de production \texttt{src-q6-g-f} reproduit un écoulement externe instationnaire de référence avec des conditions limites ouvertes physiquement contrôlées. Le cas choisi est le cylindre circulaire centré entre deux parois parallèles étudié par Zovatto et Pedrizzetti \cite{ZovattoPedrizzetti2001} et, avec une autre convention de nondimensionnalisation, par Sahin et Owens \cite{SahinOwens2004}. Cette validation est particulièrement discriminante : elle sollicite simultanément un profil d'entrée cisaillé, des parois no-slip, un obstacle pénalisé, une sortie ouverte, la projection Q6-g-f et une dynamique de brisure de symétrie à longue échelle.

Le mode SRC sans resampling n'est pas retenu comme cible de validation dans ce cycle. Le benchmark 0493x8 vise explicitement la chaîne Q6-g-f. Le premier calcul de production avait volontairement utilisé un domaine axial réduit afin de qualifier rapidement l'ensemble inlet/outlet + cylindre ; il avait donné une période à environ $2\,\%$ du point tabulé de Zovatto--Pedrizzetti. Ce résultat reste utile comme jalon de développement, mais il est désormais \emph{supplanté pour la comparaison quantitative} par le calcul cluster en domaine bibliographique complet décrit ci-dessous.

\subsection{Inlet segmenté à profil de Poiseuille local : 0493x8k}

La première correction porte sur la sémantique de l'inlet segmenté. Un segment occupant la coordonnée tangentielle $s\in[s_{\min},s_{\max}]$ définit désormais sa coordonnée locale
\begin{equation}
  \eta=\frac{s-s_{\min}}{s_{\max}-s_{\min}},
  \qquad 0\leq \eta\leq 1,
\end{equation}
et le mode \texttt{poiseuille\_y\_max} impose
\begin{equation}
  u_n(\eta)=4U_{\max}\eta(1-\eta).
  \label{eq:x8-local-poiseuille}
\end{equation}
La vitesse normale est donc exactement nulle aux deux extrémités du segment et atteint $U_{\max}$ au milieu. Pour un segment rectiligne, sa moyenne vaut
\begin{equation}
  U_b=\int_0^1 u_n(\eta)\,\mathrm d\eta=\frac{2}{3}U_{\max}.
  \label{eq:x8-poiseuille-mean}
\end{equation}

Le point important n'est pas seulement la forme analytique : la coordonnée est \emph{locale au segment}, et non calculée à partir de la hauteur totale du domaine. Le runner 0493x8k utilise volontairement un inlet partiel $[0.20,0.80]$ pour rendre cette propriété observable. La cible de population du réservoir \texttt{hard\_cell\_density} reste uniforme ; seule la moyenne de vitesse imposée suit la parabole. La même loi est utilisée dans le chemin particulaire SRC d'injection et dans la cible de flux/face consommée par Q6-g-f, ce qui évite de demander simultanément deux profils incompatibles aux deux niveaux de la méthode.

Pour le benchmark final, l'ouverture couvre toute la hauteur et
\begin{equation}
  U_{b,\mathrm{target}}=0.120839323357,\qquad
  U_{\max}=1.5\,U_b=0.181258985036.
\end{equation}
Le champ initial est généré avec le même profil de Poiseuille sans correction ultérieure de dérive moyenne, de sorte que la relation \eqref{eq:x8-poiseuille-mean} est vraie dès l'état initial.

\subsection{Outlet \texttt{neumann} passif : 0493x8q--0493x8t}

La seconde correction traite le fait qu'une sortie dite Neumann ne peut pas être une simple vitesse prescrite renommée. Le cycle x8q--x8t assemble progressivement une fermeture cohérente aux niveaux cinétique et elliptique.

\paragraph{Continuation cinétique x8q.}
Les particules qui traversent naturellement la frontière vers l'extérieur sont supprimées, mais la demi-distribution entrante nécessaire à une frontière ouverte n'est plus obtenue par une copie particule--particule auto-excitante. Le chemin final reconstruit un bain maxwellien local à partir des moments des deux couches de cellules intérieures adjacentes à l'outlet. Ces moments sont accumulés aux positions pré-streaming, avant suppression des particules sortantes ; le bain reconstruit donc densité, vitesse moyenne et température sans rétroaction directe d'un échantillon entrant précédent. L'échantillonnage du demi-espace entrant est pondéré par le flux normal.

\paragraph{Sortie de pression x8r.}
La vitesse de base à la face de sortie conserve l'extrapolation de gradient normal nul,
\begin{equation}
  u^*_{n,\mathrm{out}}=u^*_{n,\mathrm{cell}},
\end{equation}
mais elle n'est plus réutilisée comme une cible Dirichlet au pas suivant. La projection emploie au contraire une référence de correction de pression
\begin{equation}
  \phi_{\mathrm{out}}=0
\end{equation}
à la face physique. L'inlet reste prescrit ; l'outlet devient libre d'acquérir la correction normale exigée par la continuité. La valeur nominale \texttt{UOUT\_NOMINAL} du fichier de paramètres n'est donc plus une cible physique de Q6-g-f dans le mode \texttt{neumann} et n'est conservée que pour compatibilité de métadonnées et de diagnostics.

\paragraph{Conditionnement x8s.}
La condition mixte ainsi obtenue est plus mal conditionnée dans un canal long. Pour le cas pleine hauteur, non périodique en $x$, les trois modes longitudinaux les plus lents de l'opérateur sont connus analytiquement. x8s les résout exactement à l'initialisation du CG et démarre les itérations sur un résidu orthogonal à ces modes. Cette déflation ne modifie ni l'équation, ni la tolérance, ni l'intensité de projection ; elle évite seulement de gaspiller des itérations CG sur les composantes les plus lentes.

\paragraph{Cible de densité sans mode moyen x8t.}
Avec une vraie sortie de pression, le mode constant du second membre devient solvable. Or la branche signée de relaxation de densité x7d peut admettre une cible locale dont la moyenne spatiale n'est pas exactement nulle dans un fluide thermiquement fluctuant. x8t retire donc uniquement cette moyenne,
\begin{equation}
  d_{\rho,c}^{\,\mathrm{centred}}
    =d_{\rho,c}-\langle d_\rho\rangle,
\end{equation}
lorsque la relaxation de densité, la sortie de pression x8r et le domaine de pression complet sont simultanément actifs. La redistribution locale signée est préservée ; on supprime seulement une source volumique globale non intentionnelle.

Le réglage final de l'inlet utilise en outre \texttt{inletThermalNoise=1.0}. Les smokes de fermeture ont montré que ce réglage évite de transformer le réservoir \texttt{hard\_cell\_density} en injection cinétiquement trop déterministe tout en conservant exactement la vitesse moyenne imposée.

\subsection{Du domaine réduit au run cluster complet}

La géométrie physique reste
\begin{equation}
  H=1.5625,\qquad D=0.3125,\qquad \frac{D}{H}=0.2,
\end{equation}
avec $a=1/256$ et $80$ cellules par diamètre. Le calcul définitif abandonne cependant le domaine réduit $4D+11D$ du jalon initial et reprend exactement l'extension axiale de Zovatto--Pedrizzetti :
\begin{equation}
  L_x=17.1875=55D,\qquad
  x_c=4.6875=15D,\qquad
  L_x-x_c=12.5=40D,
\end{equation}
sur une grille
\begin{equation}
  N_x\times N_y=4400\times400.
\end{equation}
Le domaine contient ainsi $1.76\times10^6$ cellules de collision ; à $\gamma=20$, l'ordre de grandeur du support fluide est $3.52\times10^7$ particules, hors slots de réserve. Le cylindre reste centré transversalement, $y_c=0.78125$ et $R=0.15625$.

Cette production a été exécutée au CCUB avec un binaire CUDA natif \texttt{sm\_80}. La projection utilise le backend CUDA résident, l'opérateur \texttt{auto\_fv\_cg}, une tolérance $10^{-5}$ et \texttt{PROJECTION\_MAX\_ITERATIONS=6000}. Les deux premiers segments, $0\ldots25000$ et $25000\ldots50000$, ont été calculés sur l'A100 80~GB de \texttt{webern49}. Le troisième, $50000\ldots75000$, a été poursuivi sur l'A100 PCIe 40~GB de \texttt{webern60}, de même compute capability 8.0, avec exactement le même exécutable \texttt{sm\_80}. Chaque segment redémarre sur le dump final exact du précédent. Les paramètres physiques, géométriques et les contrôles numériques sont conservés ; seule la graine stochastique de segment évolue comme prévu par le runner, et le troisième segment change de matériel d'exécution. Le passage de l'A100 80~GB à l'A100 40~GB ne modifie donc pas la définition du problème ni le binaire CUDA, même si une identité bit-à-bit entre matériels distincts n'est pas revendiquée.

Le segment \texttt{segment\_000} a essentiellement joué le rôle de mise en régime. Le segment \texttt{segment\_001}, entre $25000$ et $50000$ pas globaux, établit déjà une fréquence cohérente mais l'inspection des champs montre encore une allée spatialement trop courte dans les $40D$ aval pour figer la comparaison bibliographique. Un troisième segment de $25000$ pas a donc été lancé spécifiquement pour laisser le sillage peupler la longue zone aval. \texttt{segment\_002} est enregistré dès son début avec
\begin{equation}
  \texttt{liveGrid}=2200\times200,\qquad
  \texttt{recordEvery}=\texttt{filterSampleEvery}=100,
\end{equation}
pour les champs \texttt{rho,ux,uy}. Il fournit $250$ frames entre les pas globaux $50001$ et $74901$. Après inspection visuelle de la croissance et de l'extension spatiale de l'allée, l'analyse définitive retient la fenêtre locale $9001\ldots24901$ de ce segment, soit la fenêtre globale $59001\ldots74901$ et $160$ frames. Cette sélection est fondée sur l'établissement spatial du sillage, et non sur l'optimisation de l'accord fréquentiel avec la littérature. La cadence et la résolution du recorder sont constantes sur toute cette fenêtre quantitative.

\subsection{Viscosité Q6-g-f qualifiée et Reynolds effectif}
\label{sec:x8-viscosity-caveat}

La signature microscopique du calcul est
\begin{equation}
 \gamma=20,\qquad \Delta t=0.002,\qquad k_BT=0.125,\qquad
 m=1,\qquad h=1/256,
\end{equation}
avec rotation SRC de $90^\circ$, signe aléatoire, décalage de grille et thermostat \texttt{cell\_relative\_rescale} à chaque pas. Le benchmark avait été dimensionné avec
\begin{equation}
  \nu_{\rm design}=6.743265812\times10^{-4},
  \label{eq:x8-nu-design}
\end{equation}
ce qui fixe par construction
\begin{equation}
 Re_{H,\rm target}=\frac{U_{b,\rm target}H}{\nu_{\rm design}}=280,
 \qquad
 Re_{D,\rm target}^{(\rm SO)}=\frac{U_{\max}D}{\nu_{\rm design}}=84.
\end{equation}

La qualification constitutive a depuis été reprise avec le calibrateur standalone général décrit en section~\ref{sec:x13-standalone-calibrator}, en sélectionnant directement \texttt{CALIBRATION\_PATH=src-q6-g-f}. Le protocole Taylor--Green est celui du calibrateur historique 0493w1, mais le chemin numérique complet est celui de production Q6-g-f : projection CUDA \texttt{auto\_fv\_cg}, forçage \texttt{prestream\_single\_fused}, tolérance $10^{-5}$, $\tau_\rho=0.25$, seuils de compression/traction $3/6$ particules, gain de traction unité, \texttt{minFill=0.10}, CG x7j résident et resampling désactivé. Huit graines indépendantes ont été exécutées sur le banc périodique $128^2$, mode TG $(2,2)$, soit une longueur d'onde de $64h$.

Le résultat consolidé est
\begin{equation}
  \boxed{
  \overline\nu_{\rm Q6GF}
  =6.651267455\times10^{-4},
  \qquad
  \sigma_\nu=3.258\times10^{-5},
  \qquad
  CV_\nu=4.90\,\%.
  }
  \label{eq:x8-nu-q6gf-qualified}
\end{equation}
Le statut du calibrateur est \texttt{PASS}. L'erreur standard sur la moyenne vaut
\begin{equation}
  SEM_\nu=\frac{\sigma_\nu}{\sqrt 8}
  \simeq1.152\times10^{-5},
\end{equation}
soit environ $1.73\,\%$ de la moyenne. Cette statistique doit être distinguée de la variabilité inter-réalisation de $4.90\,\%$, qui est la bonne échelle pour discuter la dispersion d'un run individuel.

Le réaudit SRC-only du même point microscopique, réalisé avec le Taylor--Green historique à trois graines, donne
\begin{equation}
  \overline\nu_{\rm SRC}=6.68103\times10^{-4},
  \qquad CV\simeq9.25\,\%.
\end{equation}
La différence centrale Q6-g-f/SRC n'est que de $-0.45\,\%$, très inférieure à la dispersion statistique. À la précision de la métrologie disponible, la projection Q6-g-f ne modifie donc pas la viscosité transverse constitutive de ce fluide. La valeur de dimensionnement historique~\eqref{eq:x8-nu-design} est elle-même seulement $1.36\,\%$ au-dessus de la moyenne Q6-g-f désormais qualifiée.

Avec la vitesse bulk effectivement mesurée à $-3D$ sur la fenêtre pleinement développée,
\begin{equation}
  U_{b,\rm meas}=0.120095201,
\end{equation}
la meilleure estimation actuelle du Reynolds du run cluster est
\begin{equation}
  \boxed{
  Re_{H,\rm meas}^{\rm Q6GF}
  =\frac{U_{b,\rm meas}H}{\overline\nu_{\rm Q6GF}}
  =282.125.
  }
  \label{eq:x8-Re-qualified}
\end{equation}
Une propagation exacte de $\overline\nu\pm\sigma_\nu$ donne
\begin{equation}
  Re_H\simeq269.0\;\text{à}\;296.7
  \qquad (\pm1\sigma\ \text{sur }\nu),
  \label{eq:x8-Re-one-sigma}
\end{equation}
tandis que l'erreur standard sur la moyenne correspond à une incertitude d'environ $\pm1.7\,\%$ sur le positionnement central en Reynolds. Le benchmark doit donc être décrit comme un cas $Re_H\simeq282$ avec une variabilité constitutive empirique d'environ $5\,\%$, et non comme un cas dont le Reynolds serait connu à l'unité près.

Cette qualification ferme la réserve principale du rapport précédent. Elle ne modifie ni la fréquence mesurée ni
\begin{equation}
  T^*=\frac{U_b}{fH},
\end{equation}
qui ne contiennent pas $\nu$. Elle fixe en revanche correctement l'abscisse de comparaison : la valeur centrale $Re_H=282.125$ est à seulement $+0.759\,\%$ du point $Re_H=280$ de Zovatto--Pedrizzetti. Une mauvaise estimation de la viscosité ne peut donc plus expliquer à elle seule l'écart de période observé. La dispersion inter-graines de $\nu$ doit toutefois être conservée lorsqu'on compare à une courbe fréquentielle dépendant de $Re$, car elle élargit l'abscisse physique pertinente même lorsque la fréquence elle-même est mesurée avec une précision supérieure.

\subsection{Conservation du débit et de la densité dans la longue zone amont}

L'analyse 0493x8n reconstruit, pour chaque section,
\begin{align}
 U_b(x,t) &= \frac{1}{H}\int_0^H u_x\,\mathrm dy,\\
 Q_v(x,t) &= \int_0^H u_x\,\mathrm dy,\\
 \bar\rho(x,t) &= \frac{1}{H}\int_0^H \rho\,\mathrm dy,\\
 J_\rho(x,t) &= \int_0^H \rho u_x\,\mathrm dy.
\end{align}
Le diagnostic est construit à partir des champs coarse-grainés du recorder ; $J_\rho$ doit donc être interprété comme un flux macroscopique reconstruit, et non comme un audit particulaire microscopique. Ses rapports entre sections restent néanmoins un test direct de la cohérence spatiale du transport.

Sur la fenêtre pleinement développée $59001\ldots74901$, la section de référence est située à $(x-x_c)/D=-3$, donc à $12D$ de l'inlet et $3D$ en amont du cylindre. On obtient
\begin{equation}
 U_{b,\mathrm{meas}}=0.120095201
       =0.993842\,U_{b,\mathrm{target}},
\end{equation}
soit un écart de $-0.6158\,\%$ à la cible. Dans toute la partie amont située au-delà de $1D$ de l'inlet, les dispersions relatives entre sections valent
\begin{equation}
 \mathrm{spread}(Q_v)=0.00223802,\qquad
 \mathrm{spread}(\bar\rho)=0.000222555,\qquad
 \mathrm{spread}(J_\rho)=0.00216157,
 \label{eq:x8-flux-spreads}
\end{equation}
soit respectivement $0.2238\,\%$, $0.0223\,\%$ et $0.2162\,\%$. Le rapport $J_\rho/(\bar\rho Q_v)$ reste voisin de l'unité sur ces sections ; à $-3D$, il vaut environ $0.999881$. La longue zone amont ne présente donc ni dérive massive de densité ni compensation transverse cachée suffisante pour expliquer l'écart fréquentiel final. Par rapport à la fenêtre du segment précédent, le conditionnement est même légèrement meilleur, ce qui renforce l'interprétation du résidu comme un écart de précision du sillage plutôt qu'un défaut de débit entrant.

La couche d'accommodation de l'inlet reste néanmoins visible. Entre $0.1D$ après l'entrée et la section de référence $-3D$, les rapports sont $Q_v=1.00265$, $\bar\rho=0.988831$ et $J_\rho=0.990512$. Il est donc plus exact de parler d'une frontière bien conditionnée après accommodation que d'un inlet mathématiquement transparent. Cette précision est conservée dans l'interprétation du Reynolds mesuré.

\subsection{Allée établie : fréquence, POD et convection des vortex}

Le segment \texttt{segment\_001} avait déjà établi une fréquence proche de $0.099$ mais l'allée restait visuellement trop courte dans le domaine complet. Le segment \texttt{segment\_002} permet de distinguer établissement temporel de la fréquence et établissement spatial du sillage. Sur l'ensemble du segment, les estimateurs restent très proches, $f_{\rm POD}=0.0979188$ et $f_{\rm probe}=0.0979549$, avec seulement $0.0368\,\%$ de désaccord ; cette stabilité montre que le fondamental est déjà identifié. L'inspection des cartes de champs conduit néanmoins à ne retenir pour la comparaison finale que la partie où l'allée est pleinement développée sur la zone aval, soit $59001\ldots74901$.

Sur cette fenêtre finale de $160$ frames,
\begin{equation}
 f_{\mathrm{POD}}=0.096933472,\qquad
 f_{\mathrm{probe}}=0.097401039,
\end{equation}
soit un désaccord relatif de seulement $0.4812\,\%$. La fréquence primaire utilisée pour la nondimensionnalisation est leur moyenne,
\begin{equation}
 f=0.0971672553.
 \label{eq:x8-full-fprimary}
\end{equation}
La phase POD est quasi parfaitement linéaire,
\begin{equation}
 R^2_{\rm phase}=0.999923,
\end{equation}
sur $3.085$ cycles ; la paire POD capte $73.10\,\%$ du champ de brisure de symétrie et l'ajustement commun des sondes donne $R^2_{\rm probe}=0.94160$. La progression spatiale de phase donne $R^2=0.99783$ et
\begin{equation}
 \lambda/D=4.02606.
\end{equation}
La vitesse de convection reconstruite vaut
\begin{equation}
 c_v/U_{b,\mathrm{target}}=1.01411,
 \qquad
 c_v/U_{b,\mathrm{meas}}=1.02040.
\end{equation}
L'accord temporel POD/sondes, la linéarité de phase, la capture POD accrue et la cohérence de propagation spatiale rendent l'extraction fréquentielle concluante sur le régime pleinement développé.

Une analyse intermédiaire de cette même fenêtre avec la plage de recherche probe historique très large avait accroché une harmonique parasite et produit un $R^2_{\rm probe}$ quasi nul ; ce résultat est rejeté. L'analyse finale borne la recherche du fondamental à $St\in[0.18,0.28]$, intervalle déjà identifié indépendamment sur l'ensemble du segment \texttt{segment\_002}. Le fit retenu retrouve alors le même fondamental que le POD avec $R^2_{\rm probe}=0.9416$. Cette correction méthodologique est importante car l'analyseur x8j moyenne POD et probe lorsque les deux fréquences sont positives : une fréquence de sonde de mauvaise qualité ne doit donc pas être propagée dans la comparaison bibliographique.

La réserve quantitative principale devient la longueur statistique de cette fenêtre stricte : $3.085$ cycles suffisent pour démontrer un régime périodique cohérent et estimer $f$ à quelques pourcents, mais pas pour revendiquer une précision sub-pourcent sur la valeur bibliographique elle-même. Les chiffres sont donc rapportés avec leurs décimales de calcul, tandis que l'interprétation physique retient $f\simeq0.0972$ et $T^*\simeq0.79$.

\subsection{Comparaison à Zovatto--Pedrizzetti}

Zovatto et Pedrizzetti utilisent la hauteur du canal et la vitesse moyenne comme échelles \cite{ZovattoPedrizzetti2001} :
\begin{equation}
 Re_{\mathrm{ZP}}=\frac{U_bH}{\nu},
 \qquad
 T^*=\frac{TU_b}{H}=\frac{U_b}{fH}.
 \label{eq:x8-zp-def}
\end{equation}
Pour $D/H=0.2$, cylindre centré ($\gamma_{\rm gap}=2$), leur tableau de périodes donne
\begin{equation}
 Re_{\mathrm{ZP}}=280,\qquad T^*_{\mathrm{ZP}}=0.83.
\end{equation}

La valeur de dimensionnement historique~\eqref{eq:x8-nu-design} fixait $Re_H=280$ par construction avec la vitesse cible. La comparaison définitive utilise désormais la viscosité Q6-g-f qualifiée de~\eqref{eq:x8-nu-q6gf-qualified}. Avec la vitesse bulk cible, on obtient
\begin{equation}
 Re_{\mathrm{ZP,target}}=283.873,
 \qquad
 T^*_{\mathrm{target\,Ub}}=0.79591799,
\end{equation}
soit respectivement $+1.3831\,\%$ et $-4.1063\,\%$ par rapport au point $Re_H=280$, $T^*=0.83$. Avec la vitesse bulk effectivement mesurée à la section $-3D$,
\begin{equation}
 \boxed{Re_{\mathrm{ZP,meas}}=282.125,}
 \qquad
 \boxed{T^*_{\mathrm{meas}}=0.79101677,}
\end{equation}
soit $+0.7588\,\%$ en Reynolds et $-4.6968\,\%$ en période. Les deux estimateurs de fréquence pris séparément donnent, avec la vitesse cible,
\begin{equation}
 T^*_{\mathrm{POD}}=0.79783759\;(-3.8750\,\%),
 \qquad
 T^*_{\mathrm{probe}}=0.79400762\;(-4.3364\,\%).
\end{equation}
L'accord entre ces deux valeurs confirme que l'écart bibliographique n'est pas produit par le choix de l'estimateur de fréquence.

Le troisième segment modifie donc favorablement mais modérément la conclusion quantitative. L'ancien domaine réduit avait fortuitement donné un écart proche de $2\,\%$ ; le domaine bibliographique complet à $50000$ pas donnait encore un résidu de l'ordre de $6\,\%$. Après maturation spatiale du sillage jusqu'à $75000$ pas, la fenêtre pleinement développée ramène cet écart à environ $4$--$5\,\%$. La nouvelle calibration montre simultanément que le point central du calcul est pratiquement au Reynolds bibliographique. L'hypothèse d'une viscosité Q6-g-f fortement erronée ou d'un débit amont mal réglé n'est donc plus plausible comme explication principale. La variabilité constitutive reste finie : $CV_\nu=4.90\,\%$ correspond à l'enveloppe~\eqref{eq:x8-Re-one-sigma} et doit être propagée si l'on interpole une courbe $T^*(Re)$. Le résidu central d'environ $4$--$5\,\%$ doit toutefois être discuté en priorité comme une erreur combinée de discrétisation, de représentation du cylindre/obstacle, de fermeture Q6-g-f et de durée statistique de la fenêtre pleinement développée.

\subsection{Conversion dans la convention de Sahin--Owens}

Sahin et Owens utilisent le maximum du profil d'entrée et le diamètre du cylindre \cite{SahinOwens2004} :
\begin{equation}
 Re_{\mathrm{SO}}=\frac{U_{\max}D}{\nu},
 \qquad
 St_{\mathrm{SO}}=\frac{fD}{U_{\max}}.
 \label{eq:x8-so-def}
\end{equation}
Avec la viscosité Q6-g-f qualifiée et la fréquence finale,
\begin{equation}
 Re_{\mathrm{SO}}=85.1618,\qquad
 St_{\mathrm{SO}}=0.16752145.
\end{equation}
Le point Zovatto $Re_H=280$, $T^*=0.83$, converti avec le même profil de Poiseuille $U_{\max}=1.5U_b$ et $D/H=0.2$, vaut
\begin{equation}
 Re_{\mathrm{SO,ZP}}=84,\qquad
 St_{\mathrm{SO,ZP}}=
 \frac{D/H}{(U_{\max}/U_b)T^*}
 =0.16064257.
\end{equation}
L'écart de Q6-g-f par rapport au point Zovatto ainsi exprimé est désormais $+4.2821\,\%$, cohérent avec la comparaison directe des périodes.

Sahin--Owens donnent par ailleurs, pour $\beta=D/H=0.2$, le seuil convergé de la première bifurcation de Hopf
\begin{equation}
 Re_{\mathrm{crit}}=69.34,\qquad St_{\mathrm{crit}}=0.1567.
\end{equation}
Le rapport $85.1618/69.34=1.22818$ n'est ici qu'un contrôle de régime sous la viscosité Q6-g-f qualifiée. Le point critique de Sahin--Owens n'est pas une référence de fréquence au même Reynolds que le run et ne doit pas être utilisé pour transformer l'écart de $St$ en erreur directe par rapport à leur bifurcation.

\begin{table}[htbp]
\centering
\caption{Synthèse du run cluster VK 0493x8 en domaine Zovatto--Pedrizzetti complet. La fenêtre finale est choisie après inspection des champs, lorsque l'allée est pleinement développée ; la viscosité Q6-g-f est qualifiée par huit Taylor--Green indépendants.}
\label{tab:x8-vk-summary}
\begin{tabular}{@{}lll@{}}
\toprule
Grandeur & Q6-g-f cluster & Référence / écart \\
\midrule
Domaine axial & $15D+40D$ & ZP : $15D+40D$ \\
Grille solveur & $4400\times400$ & $80$ cellules/$D$ \\
Horizon total & $75000$ pas & trois segments restartés \\
Fenêtre finale & pas $59001\ldots74901$ & $160$ frames, $3.085$ cycles \\
$U_b$ à $-3D$ & $0.120095201$ & $-0.6158\,\%$ vs cible \\
spread $Q_v$ amont & $0.2238\,\%$ & sections au-delà de $1D$ inlet \\
spread $\bar\rho$ amont & $0.0223\,\%$ & idem \\
spread $J_\rho$ amont & $0.2162\,\%$ & idem \\
$f_{\rm POD}$ & $0.0969335$ & $R^2_{\rm phase}=0.999923$ \\
$f_{\rm probe}$ & $0.0974010$ & $R^2_{\rm probe}=0.9416$ \\
Accord POD/probe & $0.4812\,\%$ & même fondamental \\
$f_{\rm primaire}$ & $0.0971673$ & moyenne POD/probe \\
Capture paire POD & $73.10\,\%$ & brisure de symétrie \\
$\lambda/D$ & $4.02606$ & $R^2_{\rm phase,x}=0.99783$ \\
$T^*_{\rm ZP}$, $U_b$ cible & $0.795918$ & $0.83$, écart $-4.106\,\%$ \\
$T^*_{\rm ZP}$, $U_b$ mesuré & $0.791017$ & $0.83$, écart $-4.697\,\%$ \\
$\nu_{\rm Q6GF}$, TG8 & $6.65127\times10^{-4}$ & $CV=4.90\,\%$ \\
$Re_{H,\rm meas}$ qualifié & $282.125$ & $269.0$--$296.7$ à $\pm1\sigma_\nu$ \\
$Re_{\rm SO}$ qualifié & $85.1618$ & $Re_{\rm crit}=69.34$ \\
$St_{\rm SO}$ & $0.167521$ & ZP converti $0.160643$, $+4.282\,\%$ \\
\bottomrule
\end{tabular}
\end{table}

\subsection{Statut de validation et réserves restantes}

Le cycle 0493x8 permet désormais de séparer clairement ce qui est acquis de ce qui reste à fermer.
\begin{enumerate}[leftmargin=2em]
  \item \textbf{Conditions limites : acquis.} L'inlet de Poiseuille est local au segment et cohérent entre injection particulaire et projection ; l'outlet est passif aux niveaux cinétique et pression, sans réimposition cachée d'une vitesse nominale.
  \item \textbf{Capacité VK et établissement spatial : acquis.} Dans le domaine exact $15D+40D$, le troisième segment montre enfin une allée développée sur la longue zone aval. Sur la fenêtre retenue, POD et sondes diffèrent de seulement $0.481\,\%$, avec $R^2_{\rm phase}=0.999923$, $R^2_{\rm probe}=0.9416$ et une capture POD de $73.1\,\%$.
  \item \textbf{Conditionnement amont : satisfaisant.} La vitesse bulk à $-3D$ reste à $-0.616\,\%$ de la cible ; au-delà de la couche d'inlet, les dispersions spatiales de $Q_v$, $\bar\rho$ et $J_\rho$ sont respectivement $0.224\,\%$, $0.022\,\%$ et $0.216\,\%$.
  \item \textbf{Reynolds : qualification constitutive fermée au premier ordre.} Huit Taylor--Green Q6-g-f donnent $\nu=6.65127\times10^{-4}$, $CV=4.90\,\%$ et $Re_{H,\rm meas}=282.125$. Le point central est donc pratiquement celui de Zovatto--Pedrizzetti.
  \item \textbf{Comparaison bibliographique : validation dynamique forte, précision encore partielle.} La fenêtre pleinement développée donne $T^*_{\rm meas}=0.791017$ contre $0.83$, soit $-4.70\,\%$. Cet écart ne peut plus être attribué à une mauvaise viscosité centrale, à la longueur du domaine, au débit amont ou à une allée encore en croissance. Il devient un résidu de précision numérique/fermeture. La réserve statistique principale est que la fenêtre la plus stricte ne contient que $3.085$ cycles.
\end{enumerate}

La réserve de longueur de domaine du rapport précédent est donc levée pour Zovatto--Pedrizzetti : le run cluster utilise exactement leurs $15D$ amont et $40D$ aval. Les réserves restantes sont d'une autre nature : une seule résolution principale ($D/h=80$), un seul $\Delta t$ et un seul $\gamma$, l'utilisation d'un cylindre représenté par la chaîne pénalisation/condition solide du code plutôt qu'un solveur de frontière conforme à celui de la référence, la variabilité constitutive de $\nu$ et la durée statistique limitée du régime strictement pleinement développé. Ces éléments justifient une future étude de précision, mais ils ne remettent pas en cause l'observation d'un shedding établi ni la qualification fonctionnelle des conditions limites.

Le statut le plus exact est donc le suivant : \emph{la capacité von Kármán de Q6-g-f, les conditions limites Poiseuille--Neumann, l'établissement spatial du sillage et le positionnement constitutif du run autour de $Re_H=280$ sont validés dans la géométrie bibliographique complète}. Le résidu de période est désormais de l'ordre de $4$--$5\,\%$, et non plus $6\,\%$. Le run cluster prolongé à $75000$ pas doit être conservé comme donnée de référence. Un éventuel segment supplémentaire aurait pour objectif d'augmenter le nombre de cycles pleinement développés et de réduire l'incertitude statistique, non de corriger le Reynolds ou de faire apparaître l'allée.


\section{Cycle 0493x9 : courbure, saut de Laplace, tension superficielle et mouillage}
\label{sec:x9_surface_tension}

\subsection{Positionnement physique et contrat de la fermeture capillaire}

Le cycle 0493x9 ne cherche pas à introduire une attraction moléculaire entre particules SRC, ni à transformer le resampling en mécanisme de cohésion. La fermeture retenue est macroscopique : l'interface géométrique fournit une courbure $\kappa$ et Q6-g-f impose le saut de contrainte normale de Laplace. Pour deux phases $A$ et $B$, avec $A$ du côté $\alpha\geq 0.5$ et $B$ du côté extérieur, la convention est
\begin{equation}
  p_A-p_B=\sigma\kappa,
  \label{eq:x9_laplace_jump}
\end{equation}
où $\sigma\geq0$ est le paramètre \texttt{surfaceTensionSigma}. En deux dimensions, une interface circulaire de rayon $R$ avec normale orientée de $A$ vers $B$ vérifie $\kappa=1/R$ ; il ne faut donc pas transposer directement le facteur $2/R$ d'une sphère tridimensionnelle.

Cette fermeture est volontairement intégrée au solveur de projection. Elle n'ajoute ni force CSF volumique, ni potentiel pair-à-pair, ni kick capillaire direct sur une particule. La capillarité modifie la condition de pression de l'interface ; le champ de correction Q6 résultant est ensuite transmis aux particules par la reconstruction face--particule B1. Cette organisation conserve la séparation construite dans les cycles x6/x7 : transport particulaire, qualité de support, géométrie d'interface et contrainte normale restent des briques distinctes.

\subsection{Géométrie physique $\alpha$ et champ auxiliaire de courbure}

La géométrie physique de surface libre reste celle de x6c. Pour la phase $A$, le remplissage brut de cellule est fondé sur la masse déposée et la masse de référence :
\begin{equation}
  f_c=\frac{M_{A,c}}{M_{A,\mathrm{ref}}},
  \qquad
  \alpha^{(0)}_c=\min(1,\max(0,f_c)).
\end{equation}
Le champ physique $\alpha$ obtenu par le lissage x6c définit l'interface $\alpha=0.5$, le domaine de pression et les crossings sous-maille de x6f. Le cycle x9 ne déplace pas ce crossing pour calculer une courbure plus lisse : il construit un second champ $\alpha_K$ réservé à la géométrie différentielle. Cette distinction évite qu'un lissage nécessaire à $\kappa$ modifie simultanément la quantité de liquide ou la position de l'interface vue par Q6.

La campagne x9a--x9c a comparé plusieurs constructions. La production retient trois passes binomiales $3\times3$ sur le champ de courbure :
\begin{equation}
  \alpha_K = \mathcal{B}^3(\alpha),
  \qquad
  \mathcal{B}=\frac1{16}
  \begin{bmatrix}
  1&2&1\\
  2&4&2\\
  1&2&1
  \end{bmatrix}.
\end{equation}
Le gradient est ensuite évalué par un opérateur de Scharr. Pour une quantité scalaire $q$ et une maille cartésienne, une écriture du gradient horizontal utilisé est
\begin{equation}
 D_x q_{i,j}=\frac{
 3(q_{i+1,j-1}-q_{i-1,j-1})+
 10(q_{i+1,j}-q_{i-1,j})+
 3(q_{i+1,j+1}-q_{i-1,j+1})}
 {32\,\Delta x},
\end{equation}
avec une expression symétrique pour $D_y$. La normale orientée de $A$ vers $B$ est
\begin{equation}
  \bm n_{AB}=-\frac{\nabla_h\alpha_K}
  {\lVert\nabla_h\alpha_K\rVert+\varepsilon_n},
  \label{eq:x9_normal}
\end{equation}
et la courbure est calculée par une seconde divergence de Scharr,
\begin{equation}
  \kappa_h=D_x n_x+D_y n_y.
  \label{eq:x9_curvature}
\end{equation}
Le support effectif de trois passes binomiales suivi du Scharr s'étend sur environ quatre cellules autour du point où la normale est consommée. Cette taille de stencil joue un rôle important pour le mouillage et pour la définition d'une courbure minimale résolue.

Les champs LiveVis \texttt{kappa} et \texttt{curvature} affichent le champ p3 de production ; les alias p1 historiques restent disponibles pour comparaison. Le champ p3 affiché reste un diagnostic géométrique brut : la régularisation x9r décrite plus loin ne modifie pas sa valeur.

\subsection{Injection du saut de Laplace dans Q6-g-f}

Au crossing physique d'une face, x6f interpole la courbure p3 avec la même position sous-maille $\theta$ que celle utilisée pour le stencil de pression. La contribution capillaire au potentiel de Dirichlet est
\begin{equation}
  \phi_{\Gamma,\sigma}
  =\frac{\Delta t}{\rho_{A,\mathrm{ref}}}\,\sigma\,\kappa_{\Gamma}.
  \label{eq:x9_phi_cap}
\end{equation}
Lorsque la phase extérieure fournit une pression $p_B$, la condition complète devient
\begin{equation}
  \phi_{\Gamma}
  =\frac{\Delta t}{\rho_{A,\mathrm{ref}}}
   \left(p_B+\sigma\kappa_{\Gamma}-p_{\mathrm{ref}}\right).
  \label{eq:x9_phi_full}
\end{equation}
Ainsi, x9 compose directement la fermeture capillaire avec x6g. Le système linéaire conserve le même opérateur FV/CG ; seul le terme de frontière de Dirichlet est modifié. Le CG résident x7j résout ensuite le potentiel, les corrections de flux de face sont formées sur CUDA et B1 reconstruit la correction particulaire. La chaîne de production est donc
\begin{equation}
\begin{split}
 &\texttt{prestream\_single\_fused}
 \rightarrow \texttt{free\_surface\_masked}
 \rightarrow \alpha/\alpha_K/\kappa\\
 &\rightarrow \text{stencil x6f/x6g avec }p_B+\sigma\kappa
 \rightarrow \text{CG x7j}
 \rightarrow \text{B1}
 \rightarrow \text{streaming/collision/thermostat}.
\end{split}
\end{equation}
La restauration signée de densité du cycle x7 reste dans le RHS de Q6-g-f et n'est pas remplacée par la capillarité. De même, resampling et kick viriel restent des mécanismes opt-in indépendants.

\subsection{Sélection des phases et cas vide/gaz/paroi}

Le cycle x9g généralise la définition des deux côtés de l'interface. Les paramètres
\begin{verbatim}
phaseInterfaceASelector = family:liquid
phaseInterfaceBSelector = family:gas
\end{verbatim}
acceptent les sélecteurs \texttt{family:liquid}, \texttt{family:gas}, \texttt{family:dispersed}, \texttt{family:unspecified}, \texttt{type:N}, \texttt{vacuum} et \texttt{wall}, avec les restrictions du backend. La phase $A$ doit correspondre à une ou plusieurs espèces enregistrées et fournit la masse de référence, $\alpha$, l'orientation de courbure et le côté projeté. La phase $B$ est le côté extérieur. Lorsque $B$ est une phase gazeuse et que le provider x6g est actif, sa pression EOS alimente $p_B$ dans l'équation~\eqref{eq:x9_phi_full}. Lorsque $B=\texttt{vacuum}$, aucune population gazeuse n'est nécessaire ; cette configuration est particulièrement utile pour qualifier la capillarité du liquide sans payer le coût de particules gazeuses qui n'apportent pas de dynamique utile au cas considéré.

Le mot-clé \texttt{wall} a un sens différent. x9h en fait un fournisseur de géométrie solide, pas un second fluide auquel on appliquerait directement $\sigma\kappa$. Le champ solide combine les parois statiques du domaine et, seulement lorsqu'il est explicitement déclaré comme géométrie de paroi via le mécanisme chi/wallVP existant, un champ $\chi$ avec $S=1-\chi$. La normale de paroi est
\begin{equation}
  \bm n_w=\frac{\nabla S}{\lVert\nabla S\rVert},
\end{equation}
orientée du fluide vers le solide. Pour le mouillage, les deux phases capillaires restent donc $A/B$ et la paroi constitue un troisième objet géométrique $W$.

Cette abstraction ne constitue pas encore un modèle général de deux liquides immiscibles. En particulier, aucune répulsion multi-couleur ou énergie de mélange ne contraint deux espèces particulaires à rester séparées, et le couplage de pression à deux liquides incompressibles n'est pas fermé comme un problème elliptique bilatéral. La capacité actuellement la plus robuste est la surface libre liquide--vide, puis liquide--gaz lorsque la pression gazeuse x6g est pertinente.

\subsection{Validation de la capillarité bulk}

Les qualifications ont volontairement séparé la géométrie passive de l'activation de $\sigma$. Les plans horizontaux ou obliques ont servi à vérifier la neutralité de la courbure loin des parois, tandis que des cercles de plusieurs rayons et plusieurs occupations ont servi à mesurer biais, bruit et sensibilité de phase de grille. Le sweep p1/p2/p3 a conduit à retenir p3 : il réduit fortement le bruit de Scharr sans déplacer l'interface physique $\alpha=0.5$ et s'est montré robuste jusqu'aux petits cercles testés.

L'activation x9d a ensuite été contrôlée par des gouttes statiques appariées. La différence de pression projetée répond linéairement à la cible $\sigma/R$ sur la plage initiale $\sigma=128,256,512$. Une régression contrainte à l'origine donne une pente d'environ $0.9575$ ; la dispersion entre graines place la valeur unité dans l'intervalle d'incertitude utilisé pour cette qualification. La décomposition du transfert montre que l'essentiel du petit déficit vient de la géométrie/interpolation de courbure, alors que le transfert Q6 de la condition de Dirichlet est compatible avec un gain unité. Il n'a donc pas été introduit de renormalisation empirique de $\sigma$.

La validation dynamique est plus discriminante. Une ellipse libre se circularise d'autant plus vite que $\sigma$ augmente, tout en conservant approximativement son aire et son centre de masse. À forte tension, le quadrupôle de forme change de signe : ce comportement n'est pas une simple relaxation monotone mais une oscillation capillaire amortie. Par exemple, les ajustements de deux cas sous-amortis ont donné $\omega_d\simeq1.52$, $\zeta\simeq0.24$ pour $\sigma=10240$, puis $\omega_d\simeq2.71$, $\zeta\simeq0.49$ pour $\sigma=51200$, avec des coefficients de détermination proches de l'unité. Ces nombres ne constituent pas encore une loi de dispersion capillaire convergée, mais ils démontrent que la fermeture transmet une raideur de surface et une dynamique d'oscillation réelles au bulk liquide.

Enfin, le jet gravitaire fournit un test ouvert et topologique. Pour une valeur de démonstration voisine de $\sigma=9450$, le chemin liquide--vide produit de façon reproductible une goutte pendante, un col, une goutte libre cohérente, l'impact sur la paroi inférieure, l'étalement en flaque puis l'impact d'une goutte suivante sur cette flaque. Cette séquence est plus exigeante qu'une goutte isolée parce qu'elle combine création d'interface, changement de topologie, chute, collision avec un solide et interaction liquide--liquide. Les runners x9s étendent ce principe à une goutte initiale paramétrable impactant soit une paroi sèche, soit une flaque initiale ; ils sont destinés à explorer le splash et ne constituent pas encore une qualification quantitative d'un nombre de Weber critique.

\subsection{Régularisation de résolution x9r}

Le jet a également mis en évidence une limite numérique claire. Lorsque la gravité étire un filament jusqu'à l'échelle de la maille, le champ de courbure peut produire $|\kappa|$ correspondant à un rayon inférieur à une cellule. Appliquer alors le saut $\sigma\kappa$ complet transforme une géométrie qui n'est plus résolue en une correction de pression réelle et peut éjecter de petits paquets de particules à grande vitesse. Le cycle x9r introduit donc le paramètre
\begin{verbatim}
surfaceTensionMinRadiusCells = 0
\end{verbatim}
avec $0$ comme no-op exact. Pour $N_R>0$,
\begin{equation}
  \kappa_{\max}=\frac{1}{N_R\min(\Delta x,\Delta y)},
  \qquad
  \kappa_{\rm eff}
  =\operatorname{clip}(\kappa_{\Gamma,\rm raw},-\kappa_{\max},\kappa_{\max}).
  \label{eq:x9r_kappa_clip}
\end{equation}
Seul le terme $\sigma\kappa$ utilise $\kappa_{\rm eff}$ : le champ p3 résident, LiveVis et les diagnostics de courbure restent bruts. Le limiter est appliqué après interpolation au crossing physique, donc il ne modifie ni $\alpha$, ni la position de l'interface, ni les cellules du masque.

Le choix actuellement recommandé est $N_R=3$. Sur le jet sévère $h=0.005$, $\sigma=75600$ et $\Delta t=0.001$, cela donne $\kappa_{\max}=66.67$. Le run de diagnostic contient $73675$ évaluations de faces capillaires, dont seulement $967$, soit $1.31\,\%$, sont limitées. La fraction instantanée maximale est $37.5\,\%$ au tout premier pas, quand l'interface injectée n'est pas encore résolue ; le maximum brut atteint $351.5$ à $t=2.46$ mais ne concerne alors que $4/57$ faces. Le dernier clipping est observé à $t=5.38$ sur $1/80$ faces, puis le limiter reste totalement inactif jusqu'à $t=16$, avec des maxima bruts typiques $59$--$63$. Ce comportement est précisément celui recherché : x9r coupe des singularités discrètes locales puis devient transparent pour le bulk résolu.

Le seuil de trois cellules n'est toutefois pas une constante physique universelle. Il est cohérent avec le support p3+Scharr et avec les tests actuels ; un cutoff à quatre cellules commencerait, dans ce jet, à limiter une courbure macroscopique finale encore résolue. La raideur temporelle reste par ailleurs indépendante de ce cutoff : à $\sigma$ très élevé, réduire $\Delta t$ a encore calmé les éjections résiduelles. Une condition de stabilité capillaire formelle en fonction de $\sigma$, $\rho$, $h$ et $\Delta t$ reste donc à établir avant de considérer la plage de paramètres comme universellement qualifiée.

\subsection{Mouillage et angle de contact}

Le mouillage est le principal volet encore incomplet. Le paramètre
\begin{verbatim}
phaseInterfaceContactAngleDegrees = -1
\end{verbatim}
désactive la fermeture lorsque sa valeur vaut $-1$. Une valeur dans $[0,180]$ représente un angle mesuré à travers la phase $A$. Avec les conventions de normales précédentes, la géométrie de Young s'écrit
\begin{equation}
  \bm n_{AB}\cdot\bm n_w=-\cos\theta_A.
  \label{eq:x9_young}
\end{equation}
Les expériences x9i--x9l ont montré qu'imposer exactement cet angle local ne suffit pas : un remplacement brutal de normale rend $\theta$ exact mais crée une mauvaise divergence, tandis que les extensions de $\alpha$ en cellules fantômes ou les reconstructions au mur peuvent conserver l'angle tout en biaisant fortement $\kappa$.

La fermeture statique la plus convaincante est x9m. Les trois passes binomiales et le Scharr donnent un rayon de support de quatre cellules ; x9m prélève donc la première normale p3 indépendante de la paroi à la couche $j=4$, centre $4.5h$. La normale au mur est fixée par l'équation~\eqref{eq:x9_young}; une seconde normale est mesurée sur l'interface hors support pariétal. Pour un arc circulaire, les deux orientations séparées par une corde de longueur $L$ donnent
\begin{equation}
  \kappa_{\rm contact}
  =\frac{2\sin(\Delta\Phi/2)}{L}.
  \label{eq:x9m_secant_curvature}
\end{equation}
Cette fermeture n'écrase ni $\alpha$ ni le champ normal p3. Sur les calottes circulaires statiques qualifiées, les angles $30^\circ$ à $150^\circ$ donnent des biais de courbure de l'ordre de quelques pourcents, tous inférieurs à environ $8\,\%$ dans le sweep principal ; les tests de rayon, d'ellipse et de phase de grille confirment une bonne robustesse statique.

La qualification dynamique reste en revanche insuffisante. Les runs de goutte sessile montrent le bon sens de mouillage/démouillage lorsque la cible passe de $90^\circ$ vers $60^\circ$ ou $120^\circ$, mais l'équilibre effectif est comprimé vers $90^\circ$ et la courbure de contact devient très bruitée pour les angles éloignés de $90^\circ$. x9m est en outre limité aux parois statiques du domaine dans sa formulation actuelle ; le couplage à une géométrie $\chi$/wallVP et les lignes de contact sur solides mobiles ne sont pas qualifiés. Les extrêmes de mouillage complet $0^\circ$ et $180^\circ$ ne font pas partie du domaine retenu pour x9m. Enfin, la garde actuelle du parseur exige encore $\sigma>0$ lorsqu'un angle de contact est actif ; une ablation strictement $\sigma=0$ doit donc désactiver l'angle ou utiliser une valeur de $\sigma$ numériquement négligeable.

\subsection{Possibilités actuelles, diagnostics et limites}

Le statut pratique de la tension superficielle peut être résumé comme suit :
\begin{itemize}[leftmargin=2em]
  \item \textbf{surface libre bulk :} opérationnelle sur le chemin CUDA Q6-g-f, en liquide--vide et en liquide--gaz lorsque le provider de pression x6g est pertinent ;
  \item \textbf{courbure :} p3/Scharr résident, LiveVis direct, géométrie physique $\alpha=0.5$ inchangée ;
  \item \textbf{saut capillaire :} $p_A-p_B=\sigma\kappa$ injecté comme Dirichlet Q6, sans force CSF ni kick particulaire additionnel ;
  \item \textbf{gouttes libres :} loi de Laplace, circularisation, oscillations amorties et conservation macroscopique démontrées dans les cas qualifiés ;
  \item \textbf{changements de topologie :} dripping, pincement, impact solide, flaque et impact sur flaque observés de manière convaincante dans les cas de démonstration ;
  \item \textbf{petites échelles :} le cutoff x9r borne la courbure utilisée dans $\sigma\kappa$ sans modifier le champ brut. Il s'agit d'un paramètre de résolution et non d'une constante physique universelle : les campagnes x12/JFM et x12yl utilisent actuellement $R_{\min}=4h$, alors que x12cal conserve $3h$ dans un régime de faible courbure ;
  \item \textbf{mouillage :} x9m est une bonne fermeture géométrique statique, mais l'angle de contact dynamique et la courbure de ligne triple restent à perfectionner ;
  \item \textbf{deux fluides :} les sélecteurs A/B sont génériques, mais il n'existe pas encore de modèle complet d'immiscibilité liquide--liquide, de tension interfaciale issue d'une énergie de mélange, ni de pression elliptique bilatérale pour deux liquides incompressibles ;
  \item \textbf{gaz :} un gaz particulaire peut fournir une pression et des collisions SRC, mais les phénomènes aérodynamiques de breakup, la traînée gaz--goutte et les rapports de densité extrêmes ne sont pas encore qualifiés ;
  \item \textbf{stabilité :} la régularisation spatiale ne remplace pas une condition de pas de temps capillaire ; les régimes de $\sigma$ très élevés doivent encore être testés avec raffinement de $\Delta t$ ;
  \item \textbf{calibration :} les résultats actuels valident le mécanisme et plusieurs réponses physiques, mais ne constituent pas encore une étude complète de convergence en $h$, $\Delta t$, $\gamma$, viscosité et nombre de Weber/Ohnesorge pour chaque benchmark.
\end{itemize}

Les sorties principales sont \texttt{cuda\_surface\_tension\_0493x9d.csv} pour le saut capillaire, \texttt{cuda\_surface\_tension\_limiter\_0493x9r.csv} pour le cutoff, \texttt{cuda\_contact\_angle\_offsupport\_0493x9m.csv} pour le mouillage statique et les diagnostics de phase/ellipse existants pour la dynamique de goutte. LiveVis expose directement \texttt{kappa}/\texttt{curvature} pour le contrôle géométrique. Le recorder filtré peut en parallèle conserver les champs supportés tels que \texttt{mass}, \texttt{ux} et \texttt{uy}, qui peuvent ensuite être rejoués pour documenter les séquences de dripping ou de splash sans relancer la simulation.

Le bilan du cycle x9 est donc différent de celui des premières tentatives de ``surface tension'' historiques. La méthode actuelle n'est pas une cohésion cinématique ad hoc : elle réalise explicitement un saut de Laplace dans la fermeture de pression de Q6-g-f. Les cycles x10--x12 ont ensuite montré que la conservation du support particulaire constitue un problème distinct et plus difficile : une interface cinétique trop réfléchissante peut ralentir ou figer les modes physiques, tandis que l'absence de fermeture laisse fuir la masse. La chaîne courante combine donc une géométrie cinétique plus régulière, une relocalisation conservative de support et un traitement thermique local, tout en conservant comme objet séparé la tension superficielle mécanique. Les tentatives abandonnées, la chaîne active et les qualifications splash/JFM sont détaillées dans la section suivante.




\section{Cycles 0493x10--0493x12 : fermeture cinétique, géométrie subcellulaire et qualification splash/JFM}
\label{sec:x10_x12_capillary_kinetic}

\subsection{Séparer capillarité, support particulaire et cinétique d'interface}

Le saut de Laplace x9 impose une contrainte macroscopique correcte au solveur Q6-g-f mais ne suffit pas à définir le devenir individuel des particules thermiques au voisinage de $\alpha=0.5$. À $k_BT=0.125$, une particule peut franchir la géométrie d'interface alors que le flux hydrodynamique moyen reste compatible avec une surface matérielle. Les premiers développements x10 ont montré que deux erreurs opposées sont possibles :
\begin{itemize}[leftmargin=2em]
  \item sans fermeture cinétique, la dynamique normale est peu perturbée mais des particules quittent le support liquide et la masse se disperse ;
  \item avec une fermeture trop proche d'une paroi matérielle, le support est conservé mais la réflexion ou le mélange cinétique ajoutent une force/viscosité parasite et ralentissent les modes interfacials.
\end{itemize}
Le problème n'est donc plus formulé comme la recherche d'une ``force de cohésion'' supplémentaire. La tension superficielle reste portée par $\sigma\kappa$ ; la fermeture cinétique doit seulement maintenir un support particulaire cohérent en perturbant le moins possible les degrés de liberté hydrodynamiques.

\subsection{Inventaire bref des tentatives x10 : ce qui a été appris et ce qui a été abandonné}

Les variantes x10 ont été conservées dans le dépôt lorsqu'elles restent utiles pour les ablations, mais elles ne doivent pas être lues comme une pile cumulative. Le tableau~\ref{tab:x10_attempts_status} résume leur statut.

\begin{table}[htbp]
\centering
\small
\begin{tabular}{>{\raggedright\arraybackslash}p{2.0cm}
                >{\raggedright\arraybackslash}p{7.2cm}
                >{\raggedright\arraybackslash}p{5.0cm}}
\toprule
Jalon & Idée testée & Statut dans le snapshot courant \\
\midrule
x10a--x10e &
Premières géométries de crossing, barrières/seals et réactions locales destinées à empêcher les sorties de support. &
Historique/supplanté ; pas de chaîne de production distincte. x10c/x10e ont notamment été abandonnés comme barrières universelles. \\
x10f--x10g &
Réaction conservative globale et réduction des grandeurs de réaction afin d'éviter les receivers cellule par cellule. &
Définitions CUDA conservées sans call-site actif dans l'orchestration courante. \\
x10h--x10i &
Critère relatif $(\bm v-\bm u_\Gamma)\cdot\bm n$ et réaction mésoscopique par blocs. &
x10h subsiste dans le chemin legacy ; x10i reste le fallback hard-$r=1$ lorsque les branches continues/spéculaires plus récentes sont inactives, mais il est bypassé par x10o dans les runners x12. \\
x10j &
Réflexion spéculaire dans le repère laboratoire. &
Ablation rejetée : conservation de la norme de vitesse mais blocage marqué du dripping. \\
x10k &
Réflexion spéculaire dans le repère local de l'interface. &
Ablation rejetée : le changement de repère seul ne libère pas la dynamique du ligament. \\
x10l &
Diagnostic juste avant le mur cinétique. &
Diagnostic passif conservé. Il a établi que Q6/B1 fournit encore un mouvement normal sortant avant la fermeture cinétique ; la perte de mobilité apparaît donc ensuite. \\
x10m--x10n &
Paroi implicite mobile puis polyligne continue de type marching-squares. &
Étapes architecturales conservées comme ablations ; x10n fournit encore des primitives réutilisées par x10o/Q2/x12a. \\
x10o &
Interface continue Q6 pilotée par une enveloppe thermique locale. &
Socle actif de la chaîne x12. \\
x10p--x10q &
Résolution des recouvrements initiaux ; recherche $3\times3$ normale et fallback $7\times7$ rare lorsqu'aucun segment pertinent n'est trouvé. &
Actif. x10q est intégré structurellement et ne possède pas de flag indépendant. \\
x10r &
Vitesses vectorielles complètes aux endpoints pour améliorer la covariance galiléenne. &
Rejeté en production : excellent test de translation, mais dripping fortement dégradé. Code conservé pour ablation. \\
x10s--x10t &
Projection normale au segment puis cinématique tangentielle rigide. &
Ablations rejetées ; x10t aggrave notamment l'impulsion parasite. Les deux restent OFF dans x12. \\
x10cic/x10biq &
Champ $\alpha$ cinétique par CIC, séparé de x6c, puis vraie reconstruction Q2 tensorielle $3\times3$. &
Actif dans la chaîne x12. \\
x10u &
Relocalisation one-for-one de la même particule, sans changement direct de masse ni vitesse. &
Actif. Supprime l'impulsion spéculaire directe mais ne suffit pas seul à stabiliser le support cinétique. \\
x10v &
Swap local du vecteur vitesse complet avec un partenaire liquide intérieur de même masse. &
Actif. Conserve exactement $P$ et $K$ pour chaque paire de masses égales ; une sous-dispersion capillaire systématique reste cependant observée. \\
x10w &
Limiter thermique local et variante pairwise conservative $P/K$. &
Implémenté mais désactivé. La variante n'a pas fourni une amélioration suffisamment robuste et est mutuellement exclusive avec x12a. \\
\bottomrule
\end{tabular}
\caption{Statut synthétique des principales tentatives x10. Les branches ``rejetées'' sont conservées comme outils d'ablation lorsqu'elles possèdent encore un flag runtime.}
\label{tab:x10_attempts_status}
\end{table}

\subsection{Chaîne cinétique actuellement utilisée : x10o + CIC + Q2 + x10u + x10v}

Dans les runners x12 de production et de calibration, la configuration de référence est explicitement verrouillée par
\begin{verbatim}
MPCD_X10O_Q6_THERMAL_INTERFACE_WALL=1
MPCD_X10_KINETIC_INTERFACE_CIC=1
MPCD_X10_KINETIC_INTERFACE_QUADRATIC=1
MPCD_X10_KINETIC_INTERFACE_ONE_FOR_ONE=1
MPCD_X10_KINETIC_INTERFACE_ONE_FOR_ONE_SWAP=1
MPCD_X10P_INITIAL_OVERLAP_RESOLUTION=1
MPCD_X10R_Q6_THERMAL_FULL_VECTOR_ENDPOINT_VELOCITY=0
MPCD_X10S_Q6_THERMAL_SEGMENT_NORMAL_KINEMATICS=0
MPCD_X10T_Q6_THERMAL_RIGID_TANGENTIAL_KINEMATICS=0
MPCD_X10_KINETIC_INTERFACE_THERMAL_PHASE_LIMITER=0
\end{verbatim}
La fraction cinétique est fixée à
\begin{verbatim}
phaseInterfaceKineticReflectionFraction = 1
phaseInterfaceEvaporationTargetType = -1
\end{verbatim}
dans ces cas. Le nom historique ``reflection fraction'' ne doit toutefois pas être interprété littéralement lorsque x10u est actif : $r=1$ signifie ici que toutes les tentatives de sortie relatives sont prises en charge par la fermeture cinétique, puis x10u substitue à la réflexion matérielle une correction de support purement positionnelle. La chaîne x12 courante est donc qualifiée uniquement dans ce régime hard-$r=1$. Le chemin $0<r<1$ reste un chemin historique de transmission/évaporation préparée et ne bénéficie pas de la totalité de la chaîne x10o+CIC+Q2+x10u+x10v ; aucune loi thermodynamique de changement de phase n'est encore qualifiée.

\subsubsection{Champ CIC dédié à l'interface cinétique}

Le champ x10cic est volontairement distinct de l'alpha x6c de production. Q6, la pression, la restauration de densité, le resampling et la courbure capillaire continuent d'utiliser leurs buffers historiques ; seul le traitement cinétique peut consommer
\texttt{kineticPhaseAlphaCIC0493x10cic}. Le dépôt CIC a été conçu sans ajouter une nouvelle passe particulaire complète : il est fusionné avec le dépôt des moments total-A déjà nécessaire. Le coût additionnel principal est donc un filtre $O(N_{\rm cell})$ et un champ résident de type \texttt{double} par cellule.

Cette séparation est physique autant qu'algorithmique : elle permet de raffiner la géométrie de crossing sans modifier simultanément la géométrie qui entre dans le solveur Q6 ou dans $\sigma\kappa$. Elle évite ainsi qu'une correction conçue pour la covariance subcellulaire de la cinétique modifie indirectement la fermeture de pression.

Une limite d'architecture reste ouverte : malgré le qualificatif \emph{kinetic-interface-only}, l'appel du CIC est encore orchestré depuis \texttt{apply\_independent\_masked\_species\_q6\_0493w5}. Le snapshot courant ne contient donc pas de chemin CIC interface autonome pour un run SRC-only sans Q6.

\subsubsection{Reconstruction biquadratique tensorielle Q2}

Lorsque
\texttt{MPCD\_X10\_KINETIC\_INTERFACE\_QUADRATIC=1},
les neuf échantillons CIC du patch $3\times3$ définissent
\begin{equation}
  P(\xi,\eta)
  =
  \sum_{i=-1}^{1}\sum_{j=-1}^{1}
  a_{ij}L_i(\xi)L_j(\eta),
  \label{eq:x10_q2}
\end{equation}
où $L_i$ sont les polynômes de Lagrange quadratiques sur $\{-1,0,1\}$. Marching-squares ne sert plus qu'à identifier la topologie et le support de branche ; le crossing et sa normale proviennent du champ Q2. L'implémentation charge exactement neuf valeurs alpha par patch actif, les conserve en registres et n'ajoute ni grille de coefficients, ni buffer, ni kernel global supplémentaire. La recherche du crossing est bornée à huit itérations de type sécante/bisection sur les cellules déjà identifiées comme interfaciales.

\subsubsection{Enveloppe thermique x10o et résolution des recouvrements}

x10o conserve l'enveloppe
\begin{equation}
  \delta_{\rm th}
  =
  \min\!\left(
      C_T\Delta t\sqrt{\frac{k_BT}{m}},
      \delta_{\max}
  \right),
  \qquad
  C_T=3,\quad \delta_{\max}=0.75h
  \label{eq:x10o_thermal_envelope_current}
\end{equation}
dans les runners qualifiés. x10p traite les particules qui se trouvent déjà du mauvais côté au début de la résolution ; x10q n'élargit la recherche à $7\times7$ que dans ce cas rare si le voisinage $3\times3$ n'a trouvé aucun segment. Il ne s'agit donc pas d'une recherche large systématique.

\subsubsection{Relocalisation x10u et swap conservatif x10v}

Pour un crossing sélectionné, x10u garde exactement le même slot, la même masse et le même vecteur vitesse. La correction porte seulement sur la position afin de mirrorer l'endpoint qui aurait quitté le support liquide. Il n'y a ni split, ni merge, ni impulsion de paroi directe :
\begin{equation}
  m_i'=m_i,\qquad
  \bm v_i'=\bm v_i
  \quad\text{au niveau de x10u.}
\end{equation}
Cette ablation a démontré qu'une correction de support peut être formulée sans kick direct, mais aussi qu'une vitesse thermique sortante inchangée tend à provoquer des sorties répétées et une forte sensibilité microscopique.

x10v ajoute donc un second étage. Chaque particule relocalisée cherche, dans un voisinage $3\times3$, un partenaire liquide intérieur non relocalisé, de même masse, avec préférence pour le support de plus grand alpha. Les vecteurs vitesse sont échangés littéralement :
\begin{equation}
  (\bm v_i,\bm v_j)
  \longleftarrow
  (\bm v_j,\bm v_i).
\end{equation}
Pour $m_i=m_j$, ce swap conserve exactement
\begin{equation}
  \bm P_{ij}=m(\bm v_i+\bm v_j),
  \qquad
  K_{ij}=\frac{m}{2}\left(\|\bm v_i\|^2+\|\bm v_j\|^2\right).
\end{equation}
L'implémentation réutilise les deux buffers cellule x10m pour les candidats et n'ajoute qu'un marqueur d'un octet par particule. Deux kernels particulaires supplémentaires construisent puis appliquent les swaps.

Cette fermeture est actuellement la meilleure solution pratique obtenue pour éviter la fuite sans inventer de vitesse ou d'énergie, mais elle n'est pas considérée comme physiquement parfaite. Le swap du vecteur complet mélange aussi les composantes tangentielles et peut se comporter comme une dissipation cinétique supplémentaire des modes de surface.

\subsection{x12a : refroidissement thermique local des petites structures}

La branche x10w avait tenté de réduire la rétention dans les ligaments au moyen d'un limiter local fondé sur
\[
  \chi=\delta_{\rm th}/L_{\rm loc},
\]
puis d'une redistribution pairwise conservative. Les seuils
\[
  \chi_{\rm on}=0.022,\qquad
  \chi_{\rm full}=0.028,\qquad
  \eta_{\rm cap}=0.05724334
\]
restent implémentés pour les ablations, mais x10w est désactivé dans la chaîne courante.

x12a adopte une fermeture plus simple : réduire localement l'agitation thermique et l'épaisseur d'interface lorsque l'objet liquide devient mince. Pour chaque segment interfacial, la normale Q2 intérieure est suivie jusqu'à l'interface opposée ; si $d_{\rm opp}$ est la distance complète,
\begin{equation}
  L_{\rm loc}=\frac{d_{\rm opp}}{2}.
\end{equation}
La loi est
\begin{equation}
  f_T
  =
  \frac{k_BT_{\rm eff}}{k_BT}
  =
  \min\!\left[
      1,
      \left(\frac{L_{\rm loc}}{R_c}\right)^2
  \right],
  \qquad
  \frac{R_c}{h}=8\sqrt{10}
  \simeq25.2982.
  \label{eq:x12a_local_cooling}
\end{equation}
La courbure n'intervient pas dans cette loi. L'enveloppe x10o est réduite par $\sqrt{f_T}$ et la même valeur $f_T$ modifie la cible thermique locale du thermostat.

La conception de coût est explicite dans le code : seuls les owners interfacials lancent une recherche bornée le long de la normale, un champ persistant de type \texttt{float} par cellule est ajouté, et le mapping vers le thermostat est fusionné dans son dépôt de moments existant. x12a requiert x10o+CIC+Q2+x10v et est mutuellement exclusif avec x10w.

\subsection{Ce que les ablations d'ondes capillaires ont appris sur la fermeture cinétique}

La campagne $n=4$ a permis de séparer le défaut de support du défaut dynamique. En utilisant le premier passage par zéro comme observable robuste, les essais multi-graines donnent approximativement
\begin{center}
\begin{tabular}{lcc}
\toprule
Fermeture & Fréquence relative au premier zéro & Lecture \\
\midrule
x10u OFF, x10v OFF & $0.939\pm0.035$ & dynamique proche de la théorie, mais fuite de support \\
x10u ON, x10v OFF & $0.921\pm0.181$ & moyenne proche, forte variabilité/chattering \\
x10u ON, x10v ON & $0.811\pm0.129$ & fermeture plus robuste, ralentissement systématique \\
\bottomrule
\end{tabular}
\end{center}
Cette campagne reste importante pour comprendre la généalogie de la fermeture : x10v résout un vrai problème de conservation du support et, sur ce calibrateur historique $n=4$, son activation avait un coût dynamique mesurable. Il ne faut toutefois plus interpréter ce résultat comme un facteur de ralentissement universel de la chaîne courante. La requalification x13h décrite plus loin, effectuée avec le fluide de référence économique et la chaîne complète actuelle, retrouve des fréquences proches de la théorie pour les gouttes oscillantes $n=2,3,4$. Le résultat historique x10v demeure donc une alerte de sensibilité de la fermeture cinétique, pas une loi constitutive à appliquer à toute dynamique capillaire.

\subsection{Calibration x12 : distinguer tension superficielle mécanique et dispersion dynamique}

\subsubsection{x12cal : propriété dynamique des ondes}

Le calibrateur x12cal reprend une interface presque plane perturbée par un mode sinusoïdal et ajuste
\begin{equation}
  \omega^2
  =
  \frac{\sigma_{\rm dyn}}{\rho_{\rm ref}}
  k^3\tanh(kH).
  \label{eq:x12cal_dispersion}
\end{equation}
Il moyenne d'abord les réalisations thermiques d'un même mode en quadratures de Fourier signées, puis ajuste l'ensemble. Sa politique de qualité refuse de publier \texttt{surfaceTensionEffective} si tous les modes demandés ne sont pas cohérents : le champ brut reste disponible, mais un statut \texttt{INVALID} ou \texttt{REVIEW} interdit de le traiter comme une propriété qualifiée.

Les campagnes finales ont précisément montré pourquoi cette séparation est nécessaire. Les modes $n=2$ et $n=3$ donnent des gains bruts de l'ordre de
\[
  G_2=0.780867,\qquad
  G_3=0.782340,
\]
alors que $n=4$ devient nettement plus dispersif ; son fit global tend vers une valeur brute beaucoup plus basse et le premier zéro est un observable plus robuste. La dépendance en $kh$ montre qu'il ne s'agit pas d'un simple facteur multiplicatif de la loi de Laplace. x12cal est donc actuellement un \emph{qualificateur de dispersion dynamique} de la chaîne complète, pas la définition retenue de la tension superficielle mécanique.

\subsubsection{x12yl : définition mécanique/statique de \texorpdfstring{$\sigma_{\rm eff}$}{sigmaeff}}

Le calibrateur x12yl utilise des gouttes statiques à rayons résolus et des paires strictes $(R/h,\mathrm{seed})$ avec et sans tension superficielle. L'analyse moyenne la partie finale de chaque run puis forme
\begin{equation}
  \Delta p_{\rm cap}
  =
  p(\sigma,R,s)-p(0,R,s)
\end{equation}
et régresse
\begin{equation}
  \Delta p_{\rm cap}
  =
  G_\sigma\,\sigma_{\rm declared}
  \langle\kappa\rangle_{\rm active}.
  \label{eq:x12yl_mechanical}
\end{equation}
La courbure est celle mesurée dans le run actif ; elle n'est pas remplacée silencieusement par $1/R$. Les rayons de production par défaut sont $R/h\in\{32,40,48\}$ avec trois graines, \texttt{surfaceTensionMinRadiusCells}=4 et un garde qui évite d'utiliser des rayons proches du cutoff x12a. Le flag x12a reste activé pour reproduire exactement la chaîne de production, mais le préflight exige que le refroidissement local soit inactif sur les gouttes résolues utilisées pour la calibration.

Pour la campagne finale à
\[
  \sigma_{\rm declared}=945,
\]
le résultat rapporté est
\begin{equation}
  \boxed{
  \sigma_{\rm eff}=940.965,\qquad
  G_\sigma=0.99573,\qquad
  R^2=0.999978.
  }
  \label{eq:x12yl_final_gain}
\end{equation}
À l'échelle mécanique/statique résolue, la loi de Laplace est donc reproduite à moins d'un demi-pourcent. Ce résultat remplace l'interprétation ancienne des gains x11 ou x12cal comme renormalisations uniques de $\sigma$. La bonne lecture est désormais :
\begin{equation}
  \sigma_{\rm mech}
  \simeq \sigma_{\rm declared},
  \qquad
  \omega_{\rm num}(k)
  =G_\omega(kh)\,\omega_{\rm th}(k),
\end{equation}
avec une erreur dynamique/dispersive qui dépend du mode et de la fermeture cinétique.

\subsection{Requalification x13h : gouttes oscillantes 2-D et benchmark Taylor--Culick}
\label{sec:x13h_capillary_requalification}

Les campagnes x12cal précédentes étaient volontairement conservatrices : elles montraient qu'un gain dynamique mesuré à longueur d'onde finie ne devait pas être transformé en une renormalisation unique de $\sigma$. Une requalification a ensuite été menée sur le fluide économique x13h, avec la microphysique et le chemin numérique effectivement utilisés par les grandes applications :
\begin{equation}
  \gamma=8,\qquad
  \alpha_{\rm rot}=120^\circ,\qquad
  \lambda/h=0.72,\qquad
  k_BT=0.125,\qquad
  h=1/256,
\end{equation}
avec thermostat \texttt{cell\_relative\_rescale}, grid shift et signe de rotation aléatoire, et
\begin{equation}
  \Delta t=0.0063471328149122585.
\end{equation}
Cette campagne utilise la chaîne complète Q6-g-f + x9 + x10/x12 de production ; elle ne constitue donc pas une relecture du calibrateur historique sur une physique différente.

\subsubsection{Gouttes oscillantes : fréquence dynamique proche de la théorie}

Des gouttes bidimensionnelles faiblement déformées ont été initialisées sur les modes azimutaux $n=2,3,4$ et suivies assez longtemps pour séparer le transitoire initial du régime d'oscillation amortie. Le rapport entre fréquence numérique et fréquence théorique du modèle 2-D correspondant est
\begin{center}
\begin{tabular}{ccc}
\toprule
Mode & $G_\omega=\omega_{\rm num}/\omega_{\rm th}$ & Statut \\
\midrule
$n=2$ & $\simeq0.992$ & validé \\
$n=3$ & $\simeq0.999$ & validé \\
$n=4$ & $\simeq0.987$ & validé \\
\bottomrule
\end{tabular}
\end{center}
Les amortissements ajustés sont positifs et compatibles avec les références utilisées pour cette qualification. Le point essentiel est que les trois modes résolus ne présentent plus le ralentissement universel de l'ordre de $20\,\%$ que pouvait suggérer la lecture brute de certains calibrateurs historiques. La dynamique oscillatoire résolue du fluide x13h est donc qualifiée sur $n=2,3,4$.

Ce résultat fournit aussi un test de cohérence temporelle. Les fréquences de goutte et la vitesse de Taylor--Culick dépendent toutes deux de l'échelle $\sqrt{\sigma/\rho}$ et du même temps numérique $t=N\Delta t$. Si un facteur global $a$ reliait mal le temps du code au temps utilisé dans les références, il affecterait simultanément $G_\omega$ et les vitesses mesurées. L'observation $G_\omega\simeq1$ sur les gouttes à $\sigma=10000$, alors que le benchmark Taylor--Culick du même fluide donne $G_{\rm TC}\simeq0.81$, exclut donc un simple décalage multiplicatif global de l'horloge ou de $\sigma/\rho$.

\subsubsection{Taylor--Culick : définition de la référence}

Pour une nappe liquide bidimensionnelle de hauteur totale $H$, deux surfaces libres exercent une traction totale $2\sigma$ par unité de profondeur. Le bilan classique stationnaire donne
\begin{equation}
  \rho H U_{\rm TC}^2=2\sigma_{\rm mech},
  \qquad
  \boxed{
  U_{\rm TC}=\sqrt{\frac{2\sigma_{\rm mech}}{\rho H}}.
  }
  \label{eq:x13n_tc_velocity}
\end{equation}
La grandeur qui doit entrer dans cette expression est la tension superficielle mécanique du modèle. Le calibrateur x12yl donne
\begin{equation}
  \sigma_{\rm mech}
  =G_\sigma\sigma_{\rm declared},
  \qquad
  G_\sigma=0.99573
\end{equation}
pour la campagne qualifiée à $\sigma_{\rm declared}=945$. Cette calibration ne doit pas être extrapolée aveuglément comme une loi universelle jusqu'à $\sigma=2500$ ou $10000$ ; elle établit néanmoins que le paramètre \texttt{surfaceTensionSigma} est déjà une très bonne approximation de la tension mécanique dans le régime statique résolu. La reconstruction directe du rim décrite ci-dessous contrôle cette hypothèse dans les runs Taylor--Culick eux-mêmes.

Le benchmark x13n utilise
\begin{equation}
  \rho=524288,\qquad H/h=64,\qquad H=0.25.
\end{equation}
En prenant $\sigma_{\rm mech}\simeq\sigma_{\rm declared}$, les références sont donc
\begin{equation}
  U_{\rm TC}(\sigma=10000)=0.390625,
  \qquad
  U_{\rm TC}(\sigma=2500)=0.1953125.
\end{equation}
La correction mécanique x12yl ne modifierait ces vitesses que du facteur $\sqrt{0.99573}\simeq0.99786$, soit environ $0.21\,\%$, très loin des écarts observés.

\subsubsection{Résultats de rétraction : vitesse stationnaire trop faible}

La géométrie de référence longue utilise $L/H=12$ sur une grille $896\times256$, avec $R_{\min}=4h$. Le front est mesuré par le crossing q50 de masse linéique intégrée. Les résultats consolidés sont :
\begin{center}
\begin{tabular}{lcccc}
\toprule
Cas & $L/H$ & $U_{\rm TC}$ & $G_{\rm TC}$ principal & $G_{\rm TC}$ tardif \\
\midrule
$\sigma=10000$, court & 8  & $0.390625$  & $0.8012$ & $0.8152$ \\
$\sigma=10000$, long  & 12 & $0.390625$  & $0.8079$ & $0.8386$ \\
$\sigma=2500$, long   & 12 & $0.1953125$ & $0.7004$ & $0.7334$ \\
\bottomrule
\end{tabular}
\end{center}
Le cas court $\sigma=10000$ a en outre été répété sur trois graines : $G_{\rm TC}=0.8012$, $0.8006$ et $0.7917$, soit une moyenne $\simeq0.7979$ avec une faible dispersion inter-graine. Le défaut n'est donc pas un accident statistique.

L'allongement de la nappe de $L/H=8$ à $12$ n'améliore le gain que d'environ $0.8\,\%$ sur la fenêtre principale et $2.9\,\%$ sur la fenêtre tardive. L'interaction prématurée des deux rims ou le manque de longueur disponible n'est donc pas la cause principale. Le cas $\sigma=2500$ va dans le même sens : diminuer de deux le $U_{\rm TC}$ de référence n'améliore pas le gain, il le dégrade. Une limitation simple par vitesse absolue ou par une échelle acoustique ne suffit donc pas à expliquer le résultat.

Le film prolongé du cas long $\sigma=2500$ confirme qualitativement cette lecture. Une mesure indépendante sur l'enveloppe du film, distincte du q50 et donc non substitutive à la métrologie principale, donne une vitesse presque constante d'environ
\begin{equation}
  U_{\rm film}\simeq0.147,
  \qquad
  U_{\rm film}/U_{\rm TC}\simeq0.754,
\end{equation}
sur plusieurs temps capillaires, avec une très bonne linéarité des deux fronts. Les rims croissent, accumulent la masse de la nappe et finissent par interagir lorsque la portion plane devient courte ; aucune récupération tardive vers $G_{\rm TC}=1$ n'est visible avant cette interaction. La morphologie est donc qualitativement celle d'une rétraction de Taylor--Culick, mais la vitesse sélectionnée reste quantitativement trop faible.

\subsubsection{Contrôles sur x9r, x12a et la longueur de nappe}

Le cutoff de courbure x9r contribue à la dégradation des très faibles $\sigma$, mais il n'explique pas le déficit de base à $\sigma=10000$. Sur le cas court $\sigma=625$, passer de $R_{\min}=4h$ à $3h$ fait passer le gain principal de $0.6118$ à $0.6857$ et le gain tardif de $0.6087$ à $0.7289$ ; la fraction de faces clippées reste alors très élevée. Cette ablation établit la causalité du cutoff dans ce régime sous-résolu, sans autoriser pour autant $R_{\min}<3h$ comme nouveau réglage de production.

À $\sigma=10000$, l'ablation x12a montre au contraire un effet faible : le gain principal passe d'environ $0.8012$ avec x12a à $0.7858$ sans x12a, et le gain tardif de $0.8152$ à $0.8128$. Le refroidissement thermique local x12a n'est donc pas la source du ralentissement Taylor--Culick. De même, l'essai de longueur $L/H=12$ montre que les interactions de fin de nappe n'expliquent qu'une petite partie du déficit.

\subsubsection{Reconstruction locale de la traction du rim}

Pour déterminer si le défaut provenait encore d'une amplitude capillaire insuffisante dans le cas fortement advecté lui-même, un analyseur offline a été construit à partir des dumps \texttt{.smpcd}. Il redépose la phase liquide, reconstruit exactement la géométrie x6c bornée, les crossings x6e $\alpha=0.5$, le masque porteur x5a, les trois passes binomiales et la courbure Scharr p3, puis applique l'interpolation de face et le cutoff x9r. Le contrat temporel est explicite : un \texttt{state\_step\_N} est post-pas, alors que Q6-g-f est pré-transport ; la référence exacte est donc la ligne Q6 du pas $N+1$.

L'ancre canonique \texttt{state\_step\_0}$\leftrightarrow$Q6 step 1 passe tous les contrôles : 1664 faces de frontière porteur, 1664 crossings $\alpha=0.5$, 1664 faces représentées et capillaires, aucun clipping initial et accord de la courbure maximale à l'arrondi près. Les états dynamiques sont analysés même lorsque la cadence des CSV ne contient pas exactement le pas $N+1$ ; ils sont alors étiquetés \texttt{NO\_EXACT\_REFERENCE} et ne sont pas utilisés pour prétendre à une validation solver--solver inexistante.

La résultante géométrique locale est intégrée dans une fenêtre de largeur $W$ autour de chaque rim :
\begin{equation}
  R_{x}^{(j)}(W)
  =\int_{\mathrm{rim}\ j,W}
   \kappa_{\rm eff}\,n_x\,\mathrm ds,
\end{equation}
et la grandeur normalisée
\begin{equation}
  \mathcal T(W)
  =\frac{|R_x^{(L)}|+|R_x^{(R)}|}{4}
  \label{eq:x13n_rim_traction_norm}
\end{equation}
vaut théoriquement 1 lorsque chaque rim fournit la résultante $2\sigma$. À $W=H$, les trois états dynamiques analysés donnent
\begin{equation}
  \mathcal T_{10000}
  =1.0522,\ 1.0205,\ 0.9633,
  \qquad
  \langle\mathcal T_{10000}\rangle=1.0120,
\end{equation}
et
\begin{equation}
  \mathcal T_{2500}
  =1.1663,\ 0.8681,\ 0.9721,
  \qquad
  \langle\mathcal T_{2500}\rangle=1.0022.
\end{equation}
La fenêtre $W=0.75H$ conduit à la même conclusion moyenne. Le clipping x9r agit sur des contributions signées et peut donc augmenter ou diminuer ponctuellement la résultante nette ; il ne se comporte pas comme un simple facteur multiplicatif de perte de traction.

Ces résultats excluent une lecture du type $\sigma_{\rm dynamique}\simeq G_{\rm TC}^2\sigma$. Pour le cas long $\sigma=10000$, $G_{\rm TC}^2\simeq0.653$ alors que la résultante capillaire locale reconstruite vaut environ $1.01$ de la valeur théorique. Pour $\sigma=2500$, $G_{\rm TC}^2\simeq0.491$ alors que la résultante vaut environ $1.00$. Le défaut de vitesse intervient donc après la construction géométrique de $\sigma\kappa$ et de sa résultante locale, dans le transfert dynamique complet vers le rim et la nappe.

\subsubsection{Statut physique après Taylor--Culick}

La conclusion doit rester plus prudente que « la tension superficielle est correcte ». Les validations permettent désormais de distinguer trois niveaux :
\begin{enumerate}[label=(\roman*),leftmargin=2.5em]
  \item la loi mécanique de Laplace et son amplitude statique sont qualifiées par x12yl ;
  \item la dynamique oscillatoire résolue est qualifiée sur les gouttes x13h $n=2,3,4$ ;
  \item la rétraction soutenue avec croissance d'un rim n'est pas quantitativement qualifiée, car $U_{\rm SRC}<U_{\rm TC}$ de manière persistante malgré une résultante capillaire locale correcte.
\end{enumerate}
Ainsi, Taylor--Culick n'invalide pas la loi $\sigma\kappa$ elle-même, mais il interdit de déclarer l'implémentation interfaciale complète universellement qualifiée. Le défaut restant est spécifique à la conversion d'une traction capillaire correcte en flux et stockage de quantité de mouvement dans un rim mobile qui collecte continuellement de la masse. À ce stade, aucune modification de $\sigma$, de x9r ou de la fermeture cinétique ne doit être engagée sans localisation supplémentaire. Le prochain diagnostic prioritaire est un bilan local de quantité de mouvement séparant au mieux la contribution Q6/capillaire, le flux entrant depuis la nappe, le stockage dans le rim, les échanges de la fermeture cinétique x10 et l'action du thermostat.


\subsection{Campagne expérimentale post-x13h : corrections de Taylor--Culick, invalidations et retour au point qualifié}
\label{sec:x13-post-tc-experimental}

Le déficit Taylor--Culick à $\sigma=10000$ a déclenché une campagne volontairement exploratoire. Son objectif n'était pas de renormaliser $\sigma$, déjà qualifiée mécaniquement, mais d'identifier une fermeture cinétique et thermique capable de supprimer la perte de vitesse du rim sans dégrader les autres observables capillaires. Cette campagne est importante à conserver dans le rapport car plusieurs variantes ont donné, prises isolément, des résultats Taylor--Culick très séduisants avant d'être invalidées par une validation croisée. Le résultat final est donc autant méthodologique que numérique.

\subsubsection{x13o : swap normal-only comme ablation, non comme chaîne de production}

Le commit final contient une ablation x13o pilotée par \nolinkurl{MPCD_X10_KINETIC_INTERFACE_ONE_FOR_ONE_NORMAL_ONLY}. Elle exige x10u+x10v et remplace le swap du vecteur vitesse complet par l'échange de la seule composante normale locale, en laissant la composante tangentielle inchangée. Ce flag vaut zéro par défaut. Le message du commit \texttt{7655b81} mentionne ``swap normal only'', mais les runners qualifiés x13k/x13n n'activent pas cette ablation : le log de référence doit afficher \texttt{swapNormalOnly=0} et la sémantique \texttt{full-vector-swap}. Cette distinction doit être conservée afin d'éviter de reconstruire ultérieurement un faux historique à partir du seul titre du commit.

\subsubsection{x13r et x13s : refroidissements locaux trop directs}

Deux premières variantes ont isolé l'hypothèse thermique en supprimant toute correction cinétique post-interface. x13r imposait un refroidissement de type température nulle sur les cellules immédiatement identifiées comme interfaciales. x13s rendait cette fermeture anisotrope : suppression de la composante thermique normale, conservation réglable de la composante tangentielle et possibilité d'étendre la couche de refroidissement de quelques cellules vers l'intérieur du liquide. Ces branches ont été utiles pour montrer que l'agitation thermique près de l'interface participe au défaut Taylor--Culick, mais leur caractère discontinu et fortement dépendant du masque de cellules ne fournissait pas une fermeture générale acceptable. Elles ont donc été supersédées par x13t plutôt que promues en production.

\subsubsection{x13t : refroidissement progressif unifié}

La fermeture x13t a remplacé les seuils abrupts par un facteur thermique progressif. Pour une distance signée $d$ à l'interface et une demi-épaisseur locale $L_{\rm loc}$,
\begin{equation}
 f_\Gamma(d)=\operatorname{smoothstep}\!\left(\frac{d+0.5h}{R_\Gamma}\right),
 \qquad
 f_{\rm thin}=\min\!\left[1,\left(\frac{L_{\rm loc}}{R_c}\right)^2\right],
 \qquad
 f_T=\min(f_\Gamma,f_{\rm thin}),
\end{equation}
avec $R_c/h=25.298221281347036$. Le thermostat \texttt{cell\_relative\_rescale} rescale alors les vitesses relatives autour du centre de masse de cellule ; il ne corrige donc pas directement la vitesse collective locale.

Le sweep Taylor--Culick à $\sigma=10000$ a montré un gain monotone lorsque la couche froide est élargie :
\begin{center}
\begin{tabular}{cc}
\toprule
$R_\Gamma/h$ & $G_{\rm TC}$ \\
\midrule
2.0 & 0.83855 \\
3.0 & 0.86377 \\
4.0 & 0.91597 \\
4.5 & 0.93053 \\
5.0 & 0.95109 \\
\bottomrule
\end{tabular}
\end{center}
Cependant, la corrugation de l'interface augmente avec $R_\Gamma$. Le compromis visuel $R_\Gamma/h=4$ a donc été retenu pour les ablations suivantes. Utilisé seul, x13t donne $G_{\rm TC}\simeq0.916$, mais laisse environ $9.2\,\%$ de masse particulaire dans la population diffuse éloignée de la nappe. Ce résultat montre que le refroidissement améliore la vitesse sans résoudre le confinement cinétique.

\subsubsection{x13u/x13v : rétention one-for-one sous x13t et séparation position/swap}

x13u a constitué le gate de combinaison entre le refroidissement x13t et la relocalisation one-for-one historique x10u sous $r=1$. x13v a ensuite séparé explicitement la relocalisation de position et le swap de vitesse : 
\texttt{MPCD\_X13T\_ONE\_FOR\_ONE\_SIMPLE=1} imposait x10u, tandis que le swap full-vector x10v restait activable ou non par son propre flag. Le fix3 de cette branche a en outre corrigé un défaut d'orchestration : le champ thermique déjà construit par x13t devait survivre au passage de la fermeture cinétique, au lieu d'être invalidé comme si x12a allait le reconstruire. Cette branche n'introduisait donc pas une nouvelle loi de réflexion ; elle servait à dissocier proprement refroidissement, repositionnement et redistribution de vitesse.

La réintroduction des mécanismes historiques de support a séparé clairement vitesse et confinement. À $R_\Gamma/h=4$ :
\begin{center}
\begin{tabular}{lcc}
\toprule
Fermeture & $G_{\rm TC}$ & population lointaine finale \\
\midrule
x13t seul & 0.91597 & $\simeq9.20\,\%$ \\
x13t + x10u position-only & 0.88442 & $\simeq0.025\,\%$ \\
x13t + x10u + x10v full-vector swap & $\simeq0.86284$ & $\simeq0.01\,\%$ \\
\bottomrule
\end{tabular}
\end{center}
Ainsi x10u suffit pratiquement à fermer la fuite, mais coûte environ trois points de gain Taylor--Culick ; x10v apporte peu au confinement supplémentaire et ralentit encore le rim. Cette ablation a motivé la recherche d'une rétention qui ne réinjecte pas la fluctuation balistique sortante.

\subsubsection{x13w : escape $\rightarrow$ inactive $\rightarrow$ reseed local}

x13w a testé une autre sémantique. Lorsqu'un franchissement sortant est confirmé par x10o + CIC + Q2 + x10p, la particule est rendue inactive ; sa quantité de mouvement et son énergie particulières sont abandonnées ; une particule de même masse et de même type est recréée localement côté liquide ; x13t restaure ensuite le niveau thermique relatif. La première version réensemençait avec la moyenne des survivants de cellule. Cette moyenne est biaisée parce que le franchissement retire préférentiellement les vitesses sortantes.

Le fix2, qui corrigeait surtout le choix géométrique de la cellule de réensemencement, donne
\begin{equation}
 G_{\rm TC}=0.86416,
 \qquad
 f_{\rm far2}=0.0537\,\%,
\end{equation}
avec un confinement pratiquement complet. Le fix3 remplace la vitesse des survivants par la moyenne \emph{pré-évasion} $P_0/M_0$ de la cellule. Il remonte à
\begin{equation}
 G_{\rm TC}=0.88248,
 \qquad
 f_{\rm far2}=0.0810\,\%.
\end{equation}
Un rerun ultérieur du même endpoint a donné $G_{\rm TC}\simeq0.894$, ce qui a rappelé que le chemin GPU n'est pas bit à bit reproductible malgré un seed nominal identique. x13w est donc efficace comme mécanisme de confinement, mais son principe même consiste à convertir un franchissement en transport de population vers l'intérieur du support liquide ; cette propriété deviendra le point d'invalidation décisif sur les gouttes fermées.

\subsubsection{x13x : rétention probabiliste et absence de compromis robuste}

Le paramètre de réflexion/rétention $r$ a ensuite été reconnecté à x13w afin d'interpoler entre hard retention ($r=1$) et crossing naturel ($r=0$). Une première série était contaminée par la recapture x10p/x10q ; une ablation corrigée a désactivé cette recapture pour $r<1$. Les résultats principaux sont :
\begin{center}
\begin{tabular}{ccc}
\toprule
$r$ & $G_{\rm TC}$ & far2 final \\
\midrule
1.00 & 0.89438 & $0.0626\,\%$ \\
0.75 & 0.88067 & $5.25\,\%$ \\
0.50 & 0.90689 & $6.57\,\%$ \\
0.25 & 0.92087 & $7.88\,\%$ \\
0.00 & 0.92038 & $9.26\,\%$ \\
\bottomrule
\end{tabular}
\end{center}
La vitesse n'est pas monotone en $r$ et l'ouverture du crossing réintroduit immédiatement plusieurs pourcents de population diffuse. Le paramètre $r$ reste donc un outil diagnostic ; il ne fournit pas une loi d'évaporation ni un compromis de production.

\subsubsection{Le faux test de convergence en $\Delta t$}

Un run à $\Delta t/2$ a donné $G_{\rm TC}=0.80113$, donc un résultat apparemment beaucoup plus mauvais. Ce run ne constitue pas un test de convergence temporelle : dans SRD/MPCD, $\Delta t$ est aussi l'intervalle de collision et entre directement dans la viscosité cinétique et collisionnelle. Le diviser par deux change le fluide. Une compensation par le seul angle de rotation ne permettait pas de retrouver la viscosité de référence dans ce régime. La règle méthodologique est désormais explicite : un raffinement temporel n'est interprétable que si les propriétés constitutives $\nu$, $c_s$, $D$ et les rapports cinétiques pertinents sont rematchés ou recalibrés.

\subsubsection{x13z : changements de grille, gains Taylor--Culick séduisants mais non suffisants}

Un premier maillage alternatif a conservé la densité et plusieurs rapports de similitude en passant approximativement à $1267\times362$ avec $\gamma=4$. Il a donné $G_{\rm TC}=0.95823$, mais avec une asymétrie d'environ $-21\,\%$ et une morphologie fortement corruguée. Ce cas a été rejeté : changer simultanément $\gamma$ et la largeur en cellules de la fermeture thermique modifie à la fois le bruit statistique et la physique subgrid.

Un second design plus conservateur a utilisé $1120\times320$, $h=1/320$, $\gamma=8$, $m=0.64$, $k_BT=0.08$, $\Delta t=0.8\Delta t_0$ et $\alpha_{\rm rot}=1.731256178986674$ rad, de façon à conserver exactement la densité, $k_BT/m$, $\lambda/h$ et une viscosité SRD analytique proche de la représentation historique. Les largeurs numériques ont été gardées en unités de cellules : $R_\Gamma/h=4$, $R_c/h=25.2982$ et $R_{\min}/h=4$. Deux exécutions nominalement identiques ont donné
\begin{equation}
 G_{\rm TC}=0.93936
 \qquad\hbox{et}\qquad
 G_{\rm TC}=0.97498.
\end{equation}
Le gain semblait donc spectaculaire, mais la dispersion de plusieurs points pour le même seed interdisait déjà de parler de convergence. Surtout, la température dynamique finale rapportée à la masse n'était pas strictement identique entre représentations. L'audit de traction locale du rim à largeur $W/H$ appariée n'a pas montré l'augmentation d'environ dix pourcents qui aurait été nécessaire pour expliquer ce gain par une simple hausse de $\sigma_{
m eff}$ ; la cause devait donc être cherchée dans la représentation et la fermeture dynamiques.

\subsubsection{x13za--x13zc : dépendance de grille de la géométrie interfaciale}

Avant d'attribuer le gain Taylor--Culick à la nouvelle fermeture, les gouttes oscillantes historiques ont été comparées sur les deux grilles en conservant la chaîne qualifiée x10u+x10v+x12a. Les gains de fréquence passent de
\begin{equation}
 (1.0031,\ 1.0047,\ 0.9805)
 \quad\hbox{à}\quad
 (1.0363,\ 1.0562,\ 1.0401)
\end{equation}
pour $n=2,3,4$. La grille fine accélère donc génériquement la dynamique capillaire de $3.3$ à $6.1\,\%$, avec un effet croissant avec le mode.

L'analyse statique à $\sigma=10000$ a montré un décalage très régulier du rayon effectif obtenu par aire : $R_{\rm eff}$ est en moyenne $2.73\,\%$ plus petit sur la grille 320. La seule loi $\omega\propto R^{-3/2}$ prédit alors environ $+4.2\,\%$ de fréquence, ce qui explique l'essentiel du décalage commun des gouttes. Le clipping x9r est au contraire plus faible sur la grille fine dans ces gouttes statiques ; il ne peut donc pas être invoqué comme cause générique de l'accélération.

La régression Young--Laplace à $\sigma=10000$ s'est révélée inutilisable : les fluctuations instantanées de pression et de courbure dominent la pente globale, malgré des gouttes visuellement stables. Une seconde tentative à $\sigma=945$ a produit des gains apparents incohérents ($G_\sigma\simeq0.229$ sur 256 et $0.673$ sur 320, avec notamment $R^2=-41.65$ sur la grille fine). Ces nombres ne sont pas des mesures physiques de la tension superficielle : ils ont révélé un défaut plus fondamental du protocole de baseline. Sans tension superficielle, la ``goutte'' $\sigma=0$ se disperse et perd rapidement sa géométrie circulaire. Une longue baseline réduit éventuellement la variance statistique mais introduit un biais physique beaucoup plus grand, car elle ne représente plus le même domaine deep-liquid/deep-gas. Le contrôle $\sigma=0$ libre ne doit donc plus être utilisé comme référence longue de Young--Laplace.

\subsubsection{x13zd : validation croisée décisive de x13t+x13w}

La fermeture qui avait produit les meilleurs Taylor--Culick a finalement été soumise au test qui manquait : gouttes oscillantes $n=2,3,4$ avec exactement x13t + x13w-fix3, x10u/x10v désactivés, sur les deux grilles. Les résultats sont :
\begin{center}
\begin{tabular}{ccccc}
\toprule
Grille & mode & $G_\omega$ historique & $G_\omega$ x13t+x13w & $R^2$ x13t+x13w \\
\midrule
256 & 2 & 1.0031 & 1.0189 & 0.901 \\
256 & 3 & 1.0047 & 1.1020 & 0.994 \\
256 & 4 & 0.9805 & 1.1073 & 0.980 \\
320 & 2 & 1.0363 & 1.2669 & 0.491 \\
320 & 3 & 1.0562 & 1.2111 & 0.827 \\
320 & 4 & 1.0401 & 0.6339 & 0.404 \\
\bottomrule
\end{tabular}
\end{center}
Le rayon effectif moyen pendant le fit chute parallèlement : sur 256 il passe d'environ $0.145$ à $0.133$--$0.136$, et sur 320 d'environ $0.142$ à $0.123$--$0.128$. La masse globale peut rester conservée, mais le support de phase se contracte : le reseed convertit systématiquement des tentatives de sortie en population replacée côté liquide. Sur une nappe ouverte, ce mécanisme peut artificiellement favoriser la collecte dans le rim ; sur une interface fermée, il produit un flux net vers l'intérieur et détruit la dynamique modale. Les $G_{\rm TC}$ proches de l'unité obtenus avec cette fermeture ne peuvent donc pas être retenus comme validation physique.

\subsubsection{Retour au dépôt qualifié et point de référence figé}

À l'issue de cette invalidation, le dépôt a été remis exactement sur \texttt{origin/surf} au commit
\begin{center}
\texttt{7655b81b1b2fd16eecefa8d8b3bebac4cd9f87f1},
\end{center}
puis marqué par le tag
\begin{center}
\texttt{surf-tension-qualified-x13h-20260831}.
\end{center}
La chaîne de référence est donc de nouveau, sans ambiguïté,
\begin{equation}
 \boxed{\text{x10o + CIC + Q2 + x10p/q + x10u + x10v full-vector swap + x12a}.}
\end{equation}
Les développements x13r--x13x ne sont pas présents dans le code commité de cette référence. x13o reste présent comme ablation mais OFF. Après reconstruction du binaire à partir de ce commit, la goutte $n=2$ donne
\begin{equation}
 G_{\omega}=1.00150857,
 \qquad
 R^2=0.961425,
\end{equation}
et Taylor--Culick à $\sigma=10000$ donne
\begin{equation}
 G_{\rm TC}=0.79457172,
 \qquad
 R^2_{\rm half}=0.999551,
 \qquad
 \varepsilon_{\rm sym}=-0.609\,\%.
\end{equation}
Ces nombres constituent le smoke de rollback. Ils ne remplacent pas les campagnes multi-modes/multi-graines antérieures, mais certifient que le tag reproduit bien le régime qualifié historique.

\subsubsection{Règles à conserver pour ne pas répéter cette campagne}

Les résultats négatifs conduisent aux règles suivantes :
\begin{enumerate}[leftmargin=2em]
  \item un gain Taylor--Culick isolé ne valide jamais une fermeture de surface libre ; toute modification doit être recroisée au minimum sur une goutte statique/oscillante fermée ;
  \item une fermeture de confinement doit contrôler explicitement le volume ou l'aire de phase, pas seulement la masse particulaire globale ;
  \item un reseed qui déplace systématiquement une particule vers l'intérieur peut conserver masse et néanmoins créer un transport artificiel de phase ;
  \item $\Delta t$ est un paramètre constitutif du fluide MPCD, pas seulement un pas d'intégration ; un test de convergence exige un rematching du fluide ;
  \item lors d'un changement de grille, il faut séparer les effets de $h$, de $\gamma$, de $m$, de $k_BT$, de $\lambda/h$ et des largeurs subgrid exprimées en cellules ;
  \item une non-reproductibilité bit à bit GPU impose des ensembles de graines ou des répétitions avant d'interpréter un gain de quelques pourcents ;
  \item \texttt{interfaceSpeedRms} mesure une vitesse de cellules interfaciales, pas une amplitude de déformation de la goutte ; il ne doit pas être utilisé seul pour conclure à une instabilité morphologique ;
  \item une intégrale de $|\kappa|$ peut exploser sous bruit local tout en laissant la résultante vectorielle correcte ; elle ne doit pas être assimilée directement à une $\sigma_{\rm eff}$ ;
  \item une baseline $\sigma=0$ n'est exploitable que tant que sa géométrie reste comparable au cas actif ; une goutte libre sans tension superficielle ne fournit pas une référence longue stationnaire ;
  \item le tag \texttt{surf-tension-qualified-x13h-20260831} est la préimage obligatoire de tout futur développement de surface libre, sauf décision explicite de rebasage.
\end{enumerate}

\subsection{Qualification splash et choix du cutoff de petite courbure}

La campagne x12a \texttt{minradius\_splash\_overnight\_fix1} a comparé quatre runs appariés en conservant exactement la même chaîne x10v+x12a :
\begin{equation}
  R/h=40,\qquad
  v_y=-0.35,\qquad
  g_y=-0.005,\qquad
  \sigma=945,\qquad
  k_BT=0.125,
\end{equation}
avec cibles \texttt{wall} et \texttt{puddle}, chacune exécutée pour
\texttt{surfaceTensionMinRadiusCells}=3 puis 4 pendant 20000 pas. Les films et champs enregistrés montrent que le passage à 4 n'annule pas les mécanismes d'impact : l'étalement initial, la lamelle, la couronne/splash et les changements de topologie restent présents. Les différences observées sont quantitatives plutôt que qualitatives.

Ce résultat complète les tests de petites gouttes, où $R_{\min}=4h$ améliore nettement la stabilité des structures sous-résolues. La valeur 4 a donc été retenue dans les runners JFM et x12yl actuels. Elle ne doit toutefois pas devenir une constante physique cachée : x12cal conserve 3, et le défaut de \texttt{SimulationParams} reste 0. Le paramètre représente explicitement une régularisation de résolution de la courbure utilisée dans $\sigma\kappa$.

\subsection{Benchmark JFM 524 : impact de goutte et obstacle déclencheur de splash}

La campagne x12b--x12d prépare une comparaison externe avec l'expérience de Josserand, Lemoyne, Troeger et Zaleski \cite{JosserandLemoyneTroegerZaleski2005}, dans laquelle un petit obstacle sur une surface sèche déclenche localement l'éjection d'une nappe liquide. Le runner x12d est la version de mesure actuellement retenue ; x12b et x12c sont des étapes de construction et de réduction du domaine, sans nouvelle physique C++.

La transposition numérique fixe
\[
  D/h=320,\qquad
  R/h=160,\qquad
  h=\frac{1}{256},
\]
sur une grille par défaut $3200\times1024$ de taille $L_x=12.5$, $L_y=4$. Le cas de référence de l'article est approché par
\[
  \frac{H_{\rm obs}}{D}
  =\frac{6}{320}
  =0.01875,
  \qquad
  \frac{d_{\rm obs}}{D}
  =\frac{173}{320}
  =0.540625,
\]
à comparer aux rapports $0.07/3.7\simeq0.018919$ et $2.0/3.7\simeq0.540541$ utilisés comme cibles dans le runner. Le nombre de Weber est réglé à
\[
  We\simeq514
\]
avec
\[
  \sigma=392.149185,
\]
et la gravité est choisie pour reproduire le rapport de Froude de l'impact aussi près que le permet l'échelle numérique.

L'obstacle est représenté par le champ Darcy/$\chi$ déjà qualifié : $\chi=0$ dans le step solide et $\chi=1$ dans le fluide. La pénalisation utilise $\alpha_{\rm Darcy}=8\times10^5$, donc un régime $\lambda\rightarrow1$ sur un pas. Une limitation importante est explicitement conservée : \texttt{darcyChiCollisionVpEnable=false}, car la fermeture de contact x9m rejette encore la coexistence directe de la géométrie de paroi du domaine et de la géométrie $\chi$ pour l'angle de contact. Le contact à $90^\circ$ reste donc imposé sur la paroi physique inférieure, mais le benchmark ne revendique pas encore une équivalence complète du mouillage sur l'obstacle $\chi$.

Le second écart majeur est hydrodynamique. L'expérience JFM 524 correspond à une goutte millimétrique avec $We\simeq514$ et $Re\simeq12200$, alors que le fluide SRC actuellement qualifié à $D/h=320$ et $U_{\rm impact}\simeq0.35$ donne seulement
\[
  Re\simeq649.
\]
Le runner documente volontairement cette différence plutôt que de la masquer. À ce stade, x12d est donc un \emph{benchmark de mesure géométrique et inertio-capillaire à Weber imposé}, destiné à quantifier angle, vitesse et trajectoire de la nappe déclenchée par l'obstacle. Il ne constitue pas encore une reproduction quantitative complète de l'expérience à Reynolds identique.

\subsection{Statut actuel, coût et chantiers encore ouverts}

La chaîne de travail actuelle pour la surface libre liquide--vide est donc
\begin{equation}
  \boxed{
  \text{Q6-g-f + x9 Laplace}
  \;+\;
  \text{x10o + CIC + Q2 + x10p/q + x10u + x10v}
  \;+\;
  \text{x12a}.
  }
\end{equation}
Elle est utilisée avec resampling, refill vide et kick viriel désactivés dans les calibrateurs capillaires et le benchmark JFM, afin de ne pas confondre les mécanismes.

Le statut de qualification doit maintenant être énoncé par régime plutôt que résumé par un seul facteur « tension superficielle effective » :
\begin{itemize}[leftmargin=2em]
  \item \textbf{statique mécanique :} x12yl valide la loi de Laplace à $G_\sigma=0.99573$ pour $\sigma_{\rm declared}=945$, avec $R^2=0.999978$ ;
  \item \textbf{dynamique oscillatoire résolue :} la requalification x13h valide les gouttes 2-D $n=2,3,4$ avec $G_\omega\simeq0.992$, $0.999$ et $0.987$ ;
  \item \textbf{rim en rétraction soutenue :} Taylor--Culick reste non qualifié quantitativement, avec $G_{\rm TC}\simeq0.81$ à $\sigma=10000$ et $\simeq0.70$ à $\sigma=2500$ sur les fenêtres principales longues ;
  \item \textbf{traction capillaire locale du rim :} la reconstruction offline validée donne pourtant $\langle\mathcal T(H)\rangle\simeq1.012$ à $\sigma=10000$ et $1.002$ à $\sigma=2500$, ce qui écarte une simple perte multiplicative d'amplitude de $\sigma\kappa$ ;
  \item \textbf{changements de topologie :} dripping, pincement, impact wall/puddle, lamelle et splash sont obtenus dans les cas de qualification/démonstration ;
  \item \textbf{benchmark JFM :} la géométrie et le Weber sont transposés, mais la similitude de Reynolds reste insuffisante pour une validation quantitative complète.
\end{itemize}

Les avantages architecturaux restent inchangés : la capillarité est imposée comme contrainte normale Q6 plutôt que comme kick particulaire ad hoc ; la géométrie cinétique CIC/Q2 est distincte du champ Q6 ; x10u ne crée ni split ni merge ; x10v est localement conservatif en quantité de mouvement et énergie pour les paires de masses égales ; x12a agit sur la fermeture thermique locale sans modifier directement $\sigma\kappa$.

Les limitations à ne pas masquer sont désormais plus précisément localisées :
\begin{itemize}[leftmargin=2em]
  \item les anciennes ablations $n=4$ montrent une sensibilité réelle à x10u/x10v, mais la requalification x13h interdit d'en déduire un ralentissement universel de x10v ;
  \item Taylor--Culick met en évidence un défaut spécifique au régime d'interface translatée qui collecte continuellement de la masse ; son origine exacte dans le couplage Q6--streaming--fermeture cinétique--thermostat n'est pas encore démontrée ;
  \item x9r devient causalement important lorsque la courbure locale est sous-résolue, comme à $\sigma=625$, mais il n'explique pas le déficit de base à $\sigma=10000$ ;
  \item le CIC cinétique est séparé en données mais pas encore en orchestration : il dépend toujours du chemin Q6 dans le snapshot courant ;
  \item x12a est une loi subgrid numérique calibrée sur l'épaisseur locale, non une thermodynamique de changement de phase ;
  \item $r<1$ et \texttt{phaseInterfaceEvaporationTargetType} ne constituent toujours pas une loi d'évaporation/condensation qualifiée ;
  \item le mouillage dynamique et l'angle de contact sur obstacles $\chi$ restent incomplets ;
  \item le benchmark JFM actuel reproduit $We$ et la géométrie mais pas encore $Re$.
\end{itemize}

La campagne expérimentale postérieure, documentée à la section~\ref{sec:x13-post-tc-experimental}, a effectivement exploré plusieurs modifications thermiques, cinétiques et de grille. Certaines ont porté $G_{\rm TC}$ près de l'unité, mais la validation croisée a montré qu'elles modifiaient artificiellement le support de phase ou changeaient le fluide constitutif. La décision de production est donc de conserver x10u+x10v full-vector+x12a au tag \texttt{surf-tension-qualified-x13h-20260831}. Taylor--Culick reste un défaut quantitatif connu de cette chaîne, et non une justification pour réintroduire les branches x13t/x13w invalidées.



\section{Cycle 0493x13 : calibrateurs de transport, qualification haut Reynolds et domaine de validité des chemins SRC/Q6-g-f}
\label{sec:x13-transport}

\subsection{Motivation : transformer les paramètres numériques en propriétés de fluide utilisables}

Le cycle 0493x13 a été lancé à partir d'un constat devenu critique dans les validations de von Kármán et de splash : une simulation ne peut être comparée proprement à une référence externe que si les propriétés de transport du fluide SRC sous-jacent sont connues avec une incertitude compatible avec l'écart que l'on cherche à interpréter. Les campagnes antérieures disposaient déjà du calibrateur 0493w1, fondé sur Taylor--Green pour la viscosité, une perturbation longitudinale pour la vitesse du son et l'amortissement acoustique, et un MSD pour l'autodiffusion. Ce socle a permis d'identifier des régimes très visqueux ou très compressibles, mais les campagnes x13 ont montré plusieurs limites méthodologiques : dépendance de la viscosité apparente à la longueur d'onde, bruit important des fits individuels, sensibilité du fit longitudinal à l'amplitude, et non-reproductibilité trajectorielle GPU même lorsque l'état initial et le seed nominal sont identiques.

L'objectif du cycle x13 n'est donc pas seulement d'obtenir une nouvelle valeur de viscosité. Il est de construire un \emph{banc de caractérisation constitutive} qui permette, pour un jeu de paramètres SRC donné, de déterminer
\begin{equation}
  \nu_T,\qquad
  \nu_L,\qquad
  c_s,\qquad
  D_{\rm self},
\end{equation}
de vérifier leur caractère local aux échelles hydrodynamiques, de quantifier leur variabilité statistique et de définir le domaine de densité et de Mach dans lequel ces coefficients peuvent être utilisés pour dimensionner une application.

Pour une grille isotrope,
\begin{equation}
  h=\frac{L_x}{N_x}=\frac{L_y}{N_y},
\end{equation}
la signature locale du fluide est plus utilement décrite par
\begin{equation}
  \boxed{
  \Theta_{\rm SRC,local}
  =
  \left(
  \gamma,\alpha_{\rm rot},\frac{\lambda}{h},
  \mathcal T,\mathcal G,\mathcal S_\pm
  \right),
  }
  \label{eq:x13-local-signature}
\end{equation}
où $\gamma$ est l'occupation moyenne, $\alpha_{\rm rot}$ l'angle de rotation SRC, $\mathcal T$ la politique de thermostat, $\mathcal G$ le décalage aléatoire de grille et $\mathcal S_\pm$ la règle de signe de rotation. Dans la convention de travail du cycle x13,
\begin{equation}
  \frac{\lambda}{h}
  =
  \sqrt{\frac{\pi k_BT}{2m}}\frac{\Delta t}{h}.
  \label{eq:x13-lambda}
\end{equation}
Ainsi, $N_x,N_y,L_x,L_y$ n'agissent pas comme quatre paramètres indépendants sur le fluide local : ils fixent d'abord $h$ et les longueurs d'onde hydrodynamiques disponibles. À $h$ et microphysique fixés, augmenter simultanément la taille macroscopique et le nombre de cellules augmente $L/h$ sans changer le fluide local ; à domaine fixé, augmenter $N$ change $h$ et donc le fluide si $\Delta t$, $k_BT$ et les autres paramètres ne sont pas rééchelonnés.

Cette distinction conduit à une lecture simple du problème inverse. Pour une longueur caractéristique $L$ et une vitesse $U$,
\begin{equation}
  Re=\frac{UL}{\nu_T},
  \qquad
  Ma=\frac{U}{c_s}.
\end{equation}
En définissant la portée intrinsèque par cellule
\begin{equation}
  H_h=\frac{c_s h}{\nu_T},
  \label{eq:x13-Hh}
\end{equation}
on obtient
\begin{equation}
  \boxed{
  Re=Ma\,\frac{L}{h}\,H_h.
  }
  \label{eq:x13-Re-Ma-Hh}
\end{equation}
Le gain à grand Reynolds peut donc provenir soit d'une augmentation de $L/h$, soit d'une augmentation du Mach admissible, soit d'un fluide dont $H_h$ est plus élevé. Le cycle x13 a précisément cherché à augmenter $H_h$ sans quitter le domaine hydrodynamique local.

\subsection{Calibrateur standalone général et guide utilisateur sommaire}
\label{sec:x13-standalone-calibrator}

La campagne x13 a conduit à regrouper la caractérisation du fluide dans un calibrateur standalone unique, \texttt{calibrate\_fluid\_0493w1\_standalone.sh}. Son objectif est de supprimer l'ambiguïté entre ``fluide SRC sous-jacent'' et ``chemin numérique effectivement utilisé''. Le calibrateur n'appelle ni \texttt{run\_common}, ni runner x7n, ni helper \texttt{suite\_*} : il écrit directement les paramètres périodiques, génère les états de calibration, sélectionne les flags CUDA et analyse les résultats. Le même fichier peut donc caractériser
\begin{equation}
  \texttt{CALIBRATION\_PATH=src}
  \qquad\text{ou}\qquad
  \texttt{CALIBRATION\_PATH=src-q6-g-f}.
\end{equation}

Pour \texttt{src}, les trois observables standards sont qualifiables : Taylor--Green pour la viscosité transverse, onde longitudinale pour $c_s$ et l'atténuation, et MSD pour $D_{\rm self}$. Pour \texttt{src-q6-g-f}, Taylor--Green et MSD restent des mesures physiques du chemin complet, mais une vitesse du son ordinaire n'est \emph{pas} qualifiée : la projection modifie précisément les modes compressifs longitudinaux. Une expérience \texttt{sound} peut être conservée comme diagnostic de réponse longitudinale résiduelle, mais sa sortie doit être marquée \texttt{NOT\_APPLICABLE}/non qualifiée pour $c_s$.

Le profil Q6-g-f par défaut du calibrateur reprend le chemin de production VK : \texttt{prestream\_single\_fused}, projection CUDA \texttt{auto\_fv\_cg}, tolérance $10^{-5}$, $\tau_\rho=0.25$, seuils de compression/traction $3/6$ particules, gain unité, \texttt{minFill=0.10}, CG x7j résident, single-block 0407 désactivé, resampling et empty-refill désactivés. Ces paramètres restent surchargeables par variables d'environnement lorsque l'application de production utilise explicitement un autre profil.

Le mode d'emploi minimal est :
\begin{verbatim}
# Vérification interne, sans simulation de production
bash scripts/calibrate_fluid_0493w1_standalone.sh --check

# SRC complet : nu, cs, Dself, Sc
CALIBRATION_PATH=src \
CALIBRATION_EXPERIMENTS="tg sound msd" \
RUN_ROOT=runs/calibrations/fluid_src \
bash scripts/calibrate_fluid_0493w1_standalone.sh

# Q6-g-f : nu, Dself, Sc ; pas de qualification cs
CALIBRATION_PATH=src-q6-g-f \
CALIBRATION_EXPERIMENTS="tg msd" \
RUN_ROOT=runs/calibrations/fluid_q6gf \
bash scripts/calibrate_fluid_0493w1_standalone.sh
\end{verbatim}

Une caractérisation multi-graines se demande directement par
\begin{verbatim}
SEEDS=493201,593201,693201,...
\end{verbatim}
sans boucle shell externe. Pour une campagne de viscosité seule, \texttt{CALIBRATION\_EXPERIMENTS=tg} évite les coûts acoustiques et MSD. Les paramètres locaux essentiels à archiver avec le résultat sont $h$, $\gamma$, $\Delta t$ ou $\lambda/h$, $k_BT$, $m$, $\alpha_{\rm rot}$, thermostat, signe de rotation, grid shift, chemin numérique et profil Q6-g-f éventuel.

Les sorties consolidées sont placées sous \texttt{RUN\_ROOT/analysis/}. Le contrat courant comprend au minimum \texttt{fluid\_characterization.csv}, \texttt{fluid\_characterization.json}, \texttt{README\_RESULTS.md} et, selon les expériences demandées, les tables détaillées Taylor--Green, acoustiques et MSD. Le résumé rapporte la moyenne, l'écart-type et le coefficient de variation de $\nu$, la statistique de $D_{\rm self}$ et $Sc$, ainsi que $Re$ lorsque \texttt{CHARACTERISTIC\_U} et \texttt{CHARACTERISTIC\_L} sont fournis. Cette organisation permet d'utiliser exactement le même outil pour le présweep haut Reynolds, la qualification d'un fluide de production et le contrôle périodique de non-régression d'un chemin Q6-g-f.

\subsection{Lien avec le calibrateur de tension superficielle : propriété constitutive versus dynamique dispersive}

Les développements x12 sur la tension superficielle et x13 sur la viscosité convergent vers la même règle méthodologique. Pour $\sigma$, le calibrateur x12yl définit la propriété mécanique/statique par la loi de Laplace,
\begin{equation}
  \Delta p_{\rm cap}
  =
  G_\sigma\,\sigma_{\rm declared}
  \langle\kappa\rangle_{\rm active},
\end{equation}
et la campagne finale a donné
\begin{equation}
  \sigma_{\rm eff}=940.965,\qquad
  G_\sigma=0.99573
\end{equation}
pour $\sigma_{\rm declared}=945$. En parallèle, x12cal mesure une dispersion dynamique d'ondes capillaires dépendante de $kh$. Le ralentissement des modes courts n'est donc pas interprété comme une tension superficielle mécanique plus faible : il constitue une erreur dispersive de la chaîne cinétique et doit être caractérisé séparément.

Pour $\nu$, la même séparation est désormais retenue. La viscosité utilisée dans $Re$ doit être une viscosité transverse de longue longueur d'onde, statistiquement convergée. Un amortissement apparent mesuré à $k\lambda$ fini, sur un mode couplé ou sur une seule trajectoire bruyante ne constitue pas à lui seul la valeur constitutive du fluide. Le calibrateur x13 mesure donc le coefficient de transport et, séparément, la dépendance en longueur d'onde, en densité et en amplitude. Cette distinction est importante pour les futures simulations multiphysiques : $\sigma_{\rm eff}$ et $\nu_T$ doivent être calibrés indépendamment puis utilisés pour construire $We$, $Re$, $Oh$ et les autres groupes sans dimension ; il ne faut pas corriger artificiellement l'un pour compenser une erreur dynamique de l'autre.

\subsection{x13a : pré-balayage analytique et définition de la portée intrinsèque}

Le pré-balayage x13a a exploré des valeurs de $\lambda/h$ nettement plus grandes que dans les premières configurations historiques. Pour orienter la recherche, l'estimation analytique SRD
\begin{equation}
  f_\gamma=\frac{\gamma-1+e^{-\gamma}}{\gamma},
  \qquad
  q=f_\gamma\left(1-\cos\alpha_{\rm rot}\right),
\end{equation}
\begin{equation}
  \nu_{\rm SRD}
  =
  k_BT\,\Delta t\left(\frac{1}{q}-\frac12\right)
  +
  \frac{h^2q}{12\Delta t}
  \label{eq:x13-nu-srd}
\end{equation}
a été utilisée comme guide. Les points A0--A6 couvraient approximativement $\lambda/h=0.23$ à 3, avec des angles allant de $90^\circ$ à $175^\circ$.

La première leçon est que la formule analytique reste utile pour repérer des zones prometteuses mais qu'elle ne doit pas être traitée comme une calibration. À partir de certains angles et libres parcours, la viscosité mesurée pouvait différer fortement de~\eqref{eq:x13-nu-srd}, alors que la dynamique particulaire restait parfaitement exécutable. Inversement, un statut \texttt{INVALID} de l'ancien fit hydrodynamique signifiait parfois seulement que le protocole de mesure n'était plus adapté à la longueur d'onde ou au rapport signal/bruit. Il ne faut donc pas confondre \emph{fit invalide} et \emph{fluide physiquement invalide}.

La seconde leçon est que le problème de grand Reynolds se formule plus proprement par $H_h$ que par une viscosité seule. Un fluide peut avoir une viscosité faible mais une vitesse du son faible dans les mêmes proportions, ce qui ne donne aucun gain à Mach fixé. Le bon objectif de recherche devient la maximisation de $H_h=c_s h/\nu_T$ sous contrainte de localité hydrodynamique.

\subsection{x13b : cisaillement transverse pur, apport initial et limites métrologiques}

Le calibrateur historique Taylor--Green utilise un champ bidimensionnel divergence-free mais combine deux composantes spatiales et peut devenir moins lisible lorsque le bruit stochastique ou des corrections additionnelles sont importants. x13b a introduit un mode de cisaillement transverse pur,
\begin{equation}
  u_x(y,0)=U_0\sin(ky),
  \qquad
  u_y(y,0)=0.
  \label{eq:x13-pure-shear}
\end{equation}
Ce choix possède deux avantages. Le champ est exactement solénoïdal et, surtout, le terme advectif non linéaire s'annule. Pour un fluide newtonien local,
\begin{equation}
  A(t)=A_0\exp(-\nu_T k^2t),
\end{equation}
d'où
\begin{equation}
  \boxed{
  \nu_T
  =
  -\frac{1}{k^2}
  \frac{\mathrm d}{\mathrm dt}\ln|A(t)|.
  }
  \label{eq:x13-nuT}
\end{equation}
Le fit porte donc directement sur une relaxation de quantité de mouvement transverse.

Les amplitudes $U_0=0.05$, $0.025$ et $0.0125$ ont été comparées. Les amplitudes les plus faibles n'ont pas révélé une meilleure physique linéaire ; elles ont principalement réduit le rapport signal/bruit et augmenté la sensibilité à la fenêtre de fit. Le choix $U_0=0.05$ a donc été conservé pour les campagnes constitutives. Comme le mode~\eqref{eq:x13-pure-shear} ne possède pas de non-linéarité advective, ce niveau d'amplitude est beaucoup moins problématique qu'il ne le serait dans un vortex ou un écoulement avec convection propre.

Le calibrateur publie la viscosité, la qualité exponentielle, le nombre de points utilisés, l'étendue logarithmique du signal et le temps mur. Les campagnes x13 ont ensuite ajouté une comparaison systématique entre longueurs d'onde pour vérifier la limite $k\lambda\rightarrow0$.

Un réaudit tardif a toutefois montré que cet estimateur ne peut plus être traité comme l'autorité constitutive absolue de $\nu_T$. Les variantes V1/V2/V2.1 ont produit une dépendance d'échelle anormale, avec des viscosités de l'ordre de $0.6$--$0.7\times10^{-3}$ aux longueurs proches de $100h$--$128h$, mais pouvant monter vers $1.1$--$1.3\times10^{-3}$ autour de $200h$--$256h$ et jusqu'à environ $2.0\times10^{-3}$ à $400h$. Le contrôle V2.1 au plus grand mode était en outre seulement \texttt{REVIEW}. Cette dérive est incompatible avec l'interprétation d'un plateau constitutif simple et a motivé le retour au Taylor--Green historique comme métrologie primaire. Le cisaillement transverse pur reste utile comme diagnostic de relaxation et de localité, mais ses valeurs absolues ne doivent plus, seules, dimensionner un Reynolds.

\subsection{x13b--x13d : correction du calibrateur acoustique et fit longitudinal amorti}

La mesure acoustique a également été profondément révisée. Le premier problème concernait la génération d'une faible perturbation de densité à faible $\gamma$. Arrondir indépendamment l'occupation de chaque cellule quantifie artificiellement les amplitudes et peut supprimer une perturbation faible. Le générateur x13b-C a donc utilisé une répartition fractionnaire avec résidu conservatif de colonne, préservant exactement la population totale.

Le second problème était plus fondamental : l'estimateur cumulatif utilisé historiquement pour $\nu_L$ restait instable même lorsque $c_s$ était bien déterminé. Les séries longitudinales ont donc été ajustées directement par
\begin{equation}
  \rho_k(t)
  =
  C
  +
  e^{-\beta t}
  \left[
  A\cos(\omega t)+B\sin(\omega t)
  \right].
  \label{eq:x13-damped-mode}
\end{equation}
On en déduit
\begin{equation}
  \boxed{
  \nu_L=\frac{2\beta}{k^2},
  \qquad
  c_s=\frac{\sqrt{\omega^2+\beta^2}}{k}.
  }
  \label{eq:x13-cs-nuL}
\end{equation}
Cette formulation a stabilisé l'amortissement longitudinal sur les points de travail G08 et G10.

Le coût de l'analyse bootstrap est d'abord devenu excessif, car le fit global testait des dizaines de milliers de couples $(\beta,\omega)$ pour chaque rééchantillonnage. x13d-C-fastfit a remplacé cette recherche par une stratégie en deux niveaux : une recherche globale poolée par groupe, puis une recherche locale en trois niveaux autour de l'optimum pour chaque bootstrap, avec fallback global uniquement si l'optimum touche le bord du domaine. Le nombre de candidats par fit a été réduit d'environ un facteur 25 sans modifier l'estimateur physique. Cette optimisation a rendu possible l'usage systématique de 500 bootstraps sur plusieurs groupes.

\subsection{x13c : qualification de l'occupation moyenne et choix de \texorpdfstring{$\gamma=8$}{gamma=8}}

Une campagne multi-graines a ensuite comparé $\gamma=6,8,10,14,20$ avec
\begin{equation}
  \alpha_{\rm rot}=120^\circ,
  \qquad
  \lambda/h=0.48.
\end{equation}
La comparaison a été menée aux longueurs d'onde 64 et 128, puis complétée à 256 pour les points les plus importants. Le résultat a corrigé une intuition issue de runs isolés : abaisser $\gamma$ ne produit pas une diminution spectaculaire et monotone de $\nu_T$. G08, G10, G14 et G20 donnent des viscosités du même ordre, avec une dispersion inter-graines non négligeable. G06 devient sensiblement plus bruité.

Le principal bénéfice de la réduction de $\gamma$ est donc le coût particulaire. À nombre de cellules fixé, passer de $\gamma=20$ à 8 réduit la population d'un facteur 2.5. Le gain mur mesuré sur les petites boîtes est inférieur à 2.5, en partie parce que les petits domaines n'occupent pas totalement le GPU, mais la réduction mémoire et le potentiel de throughput sur de grandes grilles restent importants. G08 a été retenu comme meilleur compromis entre coût, stabilité statistique et transport.

\subsection{x13d : convergence de longue longueur d'onde et premier fluide G08 qualifié}

La campagne x13d a ajouté $N_\lambda=256$ pour les points G08, G14 et G20. Pour G08 à $\lambda/h=0.48$, les valeurs obtenues étaient
\begin{equation}
  \nu_T(N_\lambda=128)\simeq5.37\times10^{-4},
  \qquad
  \nu_T(N_\lambda=256)\simeq5.33\times10^{-4},
\end{equation}
soit un écart relatif d'environ $0.69\,\%$, avec quatre runs valides sur quatre et un statut \texttt{LOCAL\_PASS}. Le calibrateur longitudinal donnait
\begin{equation}
  c_s\simeq0.35545,
  \qquad
  \nu_L\simeq1.028\times10^{-3}.
\end{equation}
Le point G08 0.48 définissait donc
\begin{equation}
  H_h\simeq2.61.
\end{equation}

Cette étape a établi deux règles de qualification. Premièrement, une viscosité destinée au Reynolds doit être vérifiée à au moins deux longueurs d'onde lorsque $\lambda/h$ est augmenté. Deuxièmement, l'accord avec~\eqref{eq:x13-nu-srd} est rassurant lorsqu'il existe, mais la qualification repose sur la convergence mesurée, non sur la formule analytique.

\subsection{x13e : sweep Mach et distinction entre régime non linéaire résolu et domaine constitutif}

Le sweep compressible x13e a imposé une onde longitudinale de vitesse à densité initiale uniforme,
\begin{equation}
  u_x(x,0)=U_0\sin(kx),
  \qquad
  U_0=Ma\,c_s.
  \label{eq:x13-mach-init}
\end{equation}
Les Mach $0.05$, $0.1$, $0.2$, $0.3$, $0.4$, $0.5$, $0.7$, $0.9$ et 1.0 ont été comparés à $N_\lambda=128$ et 256 avec plusieurs seeds. L'analyse suit la célérité effective, la fraction d'énergie transférée vers les harmoniques $m=2,3,4$ et les extrema de densité.

Aucune rupture numérique grossière n'a été observée jusqu'à $Ma=1$. Jusqu'à $Ma\simeq0.2$, les coefficients linéaires restent directement applicables. À partir de $Ma\simeq0.3$, la génération d'harmoniques devient significative ; elle est interprétée comme une non-linéarité compressible physique tant que les résultats convergent avec la longueur d'onde. À $Ma=0.5$, la fraction harmonique au pic du premier cycle est déjà d'environ $26\,\%$, mais la célérité à $N_\lambda=256$ diffère seulement d'environ $1.4\,\%$ du $c_s$ linéaire et les longueurs d'onde 128/256 restent cohérentes.

Le sweep a donc révélé une distinction essentielle : une onde peut être \emph{numériquement résolue} à grand Mach tout en traversant des densités où le coefficient de transport local n'est plus qualifié. Les extrema de densité, et non une limite arbitraire de Mach seule, doivent alors être confrontés à une carte $\nu_T(\gamma)$.

\subsection{x13f : optimisation locale de \texorpdfstring{$(\alpha_{\rm rot},\lambda/h)$}{(alpha,lambda/h)} à \texorpdfstring{$\gamma=8$}{gamma=8}}

Le sweep local x13f a exploré
\begin{equation}
  \alpha_{\rm rot}
  \in
  \{120^\circ,130^\circ,140^\circ,150^\circ\},
\end{equation}
\begin{equation}
  \lambda/h
  \in
  \{0.48,0.56,0.64,0.72,0.80\}.
\end{equation}
L'écran à $N_\lambda=128$ a été suivi d'une qualification à 256 sur une shortlist.

Le paysage obtenu est nettement différent de la simple prédiction SRD. À grand angle, augmenter $\lambda/h$ devient rapidement contre-productif. Par exemple, à $150^\circ$ et $\lambda/h=0.80$, la formule analytique prédit une viscosité de l'ordre de $3.94\times10^{-4}$, alors que la mesure est supérieure à $1.8\times10^{-3}$. Le bon corridor reste proche de $\alpha_{\rm rot}=120^\circ$.

Sur cette ligne, $\lambda/h=0.56$ et 0.72 ont passé la qualification de localité, tandis que 0.64 montrait une variation plus importante entre 128 et 256 et 0.80 devenait trop dispersé. Le point 0.72 était le meilleur finaliste, avec une viscosité à $N_\lambda=256$ de l'ordre de $5.21\times10^{-4}$ dans cette campagne.

Un avantage supplémentaire est indépendant de la petite variation de viscosité :
\begin{equation}
  \frac{\Delta t_{0.72}}{\Delta t_{0.48}}=1.5.
  \label{eq:x13-dt-gain}
\end{equation}
À durée physique donnée, le nouveau fluide nécessite donc idéalement un tiers de pas de temps en moins. Sur les benchmarks de reproductibilité x13g, le gain de throughput physique mesuré était de l'ordre de 1.3--1.4, inférieur au facteur 1.5 théorique mais déjà significatif.

\subsection{x13g : découverte de la non-reproductibilité bit à bit et fiabilité statistique de \texorpdfstring{$\nu$}{nu}}

L'écart entre les valeurs de viscosité du même point G08 0.48 dans x13d et x13f a motivé une campagne spécifique de reproductibilité. Pour 0.48 et 0.72, x13g a exécuté d'une part six répétitions strictes du même état initial avec le même seed nominal, d'autre part huit réalisations indépendantes appariées.

La conclusion est importante pour toute interprétation future :
\begin{quote}
un état initial identique et un seed nominal identique ne garantissent pas une trajectoire GPU bit à bit identique sur le chemin courant.
\end{quote}
Les six répétitions strictes produisent six hashes finaux différents pour chacun des deux fluides. Le coefficient de variation des viscosités individuelles atteint environ $16\,\%$ à 0.48 et $10\,\%$ à 0.72. Une partie de la dispersion observée entre anciennes calibrations ne provenait donc pas nécessairement d'un changement de physique ou de code : elle peut être due à la stochasticité effective et à l'ordre d'exécution non déterministe du GPU.

Sur les huit réalisations indépendantes, les moyennes étaient
\begin{equation}
  \overline\nu_T(0.48)=5.83\times10^{-4},
  \qquad
  \overline\nu_T(0.72)=5.07\times10^{-4}.
\end{equation}
La réduction moyenne est de l'ordre de $13\,\%$. Un bootstrap apparié donnait un intervalle 95\,\% approximatif
\begin{equation}
  [4.5\,\%,19.7\,\%]
\end{equation}
pour la réduction relative de 0.72 par rapport à 0.48. Même une comparaison plus conservatrice des deux groupes comme ensembles indépendants restait compatible avec un gain statistiquement significatif.

La deuxième conclusion concerne l'estimateur lui-même. Le fit de la trajectoire d'amplitude moyenne d'ensemble est beaucoup plus régulier que les fits individuels et produit des $R^2$ très proches de l'unité. Dans x13g, un fit d'ensemble particulièrement favorable donnait environ $4.7\times10^{-4}$ pour 0.72, mais les campagnes suivantes ont montré qu'il serait trop optimiste de l'utiliser seul pour le dimensionnement. La politique retenue est donc :
\begin{enumerate}[leftmargin=2em]
  \item utiliser plusieurs réalisations pour toute propriété constitutive ;
  \item publier la moyenne et la dispersion des fits individuels ;
  \item utiliser, lorsque l'observable le permet, un fit de la réponse moyenne d'ensemble ;
  \item pour le dimensionnement, retenir une valeur conservatrice compatible avec plusieurs campagnes, pas l'estimation la plus favorable.
\end{enumerate}

Cette découverte modifie aussi la lecture des validations de von Kármán. Une fréquence de shedding peut être déterminée avec quelques pourcents de précision alors que le Reynolds auquel on la rattache dépend linéairement de $1/\nu$. Si $\nu$ provient d'un seul run, l'incertitude sur $Re$ peut être du même ordre ou supérieure à l'écart de fréquence comparé à la littérature. Pour les futures validations VK, la viscosité devra donc être calibrée par plusieurs runs sur le même chemin physique que la simulation cible ; une valeur SRC ne doit pas être utilisée pour attribuer précisément un Reynolds à Q6-g-f, et réciproquement.

\subsection{Clôture de la campagne haut Reynolds : réaudit Taylor--Green historique et choix de conception}
\label{sec:x13-highre-tg-reaudit}

La découverte des anomalies d'échelle du cisaillement transverse a conduit à réauditer la campagne haut Reynolds avec le Taylor--Green historique 0493w1, sans modifier le solveur. Ce choix est méthodologiquement important : l'objectif n'était plus de trouver un nouvel estimateur donnant la viscosité la plus faible, mais de revenir à une métrologie déjà utilisée historiquement et de vérifier qu'elle reproduit exactement ses propres références avant de comparer les nouveaux fluides.

Le Taylor--Green canonique utilise
\begin{equation}
 u_x=A\sin(k_xx)\cos(k_yy),\qquad
 u_y=-A\cos(k_xx)\sin(k_yy),
\end{equation}
avec
\begin{equation}
 A(t)=A_0\exp[-\nu(k_x^2+k_y^2)t].
\end{equation}
Le banc de réaudit emploie une boîte $64\times64$, une durée physique $t=4$, trois graines par point et une amplitude $A_0=0.05$. Les graines historiques A0 et A1 reproduisent les valeurs de référence 0493w1 à environ $3\times10^{-4}\,\%$ relatif, ce qui valide la chaîne générateur--solver--analyseur et exclut une dérive silencieuse du calibrateur.

Les résultats consolidés sont :
\begin{table}[htbp]
\centering
\small
\caption{Réaduit Taylor--Green de la campagne haut Reynolds. $CV$ est la dispersion inter-graines des viscosités individuelles ; $H_h=c_sh/\nu$ utilise la célérité acoustique qualifiée ou son équivalent du même fluide.}
\label{tab:x13-highre-tg}
\begin{tabular}{lccccc}
\toprule
Point & $\gamma$ & $(\alpha_{\rm rot},\lambda/h)$ & $\overline\nu_T$ & $CV$ & $H_h$ \\
\midrule
A0 & 20 & $(90^\circ,0.2269)$ & $6.6810\times10^{-4}$ & $9.25\,\%$ & $2.08$ \\
A1 & 20 & $(120^\circ,0.48)$ & $5.2882\times10^{-4}$ & $7.83\,\%$ & $2.61$ \\
G08-0.48 & 8 & $(120^\circ,0.48)$ & $5.8111\times10^{-4}$ & $17.42\,\%$ & $2.39$ \\
G08-0.72 & 8 & $(120^\circ,0.72)$ & $5.6029\times10^{-4}$ & $17.59\,\%$ & $2.48$ \\
\bottomrule
\end{tabular}
\end{table}

A0 et A1 sont classés \texttt{PASS}. Les deux points G08 restent \texttt{REVIEW} à cause de leur dispersion inter-graines plus forte, et non parce que leurs fits centraux seraient incohérents. Le passage G08 0.48 $\rightarrow$ 0.72 diminue la moyenne de seulement $3.58\,\%$, donc reste dans une bande de $5\,\%$ très inférieure au $CV$ individuel. Il n'est plus justifié d'affirmer que 0.72 possède une viscosité constitutivement plus faible que 0.48 ; les deux points doivent être considérés comme approximativement équivalents en $\nu_T$ à la précision actuelle.

Le choix de conception se sépare alors en deux objectifs. A1 fournit la meilleure portée hydrodynamique brute parmi les points \texttt{PASS}, avec $H_h\simeq2.61$. G08-0.72 reste le meilleur compromis coût/portée pour les très grandes grilles : $\gamma=8$ réduit la population d'un facteur 2.5 par rapport à A1 et $\lambda/h=0.72$ augmente $\Delta t$ d'un facteur 1.5 par rapport à 0.48. Sur le proxy coût particule--pas, G08-0.72 représente environ $27\,\%$ du coût d'A1, soit un facteur proche de 3.7, alors que sa viscosité moyenne n'est qu'environ $6\,\%$ plus élevée. Le gain de 0.72 est donc avant tout un gain de coût temporel, pas une réduction démontrée de viscosité.

Cette campagne clôt le présweep haut Reynolds sur une base plus robuste. Les valeurs absolues de viscosité issues du cisaillement transverse x13b--x13h doivent être lues comme diagnostics historiques de localité et de tendance ; pour le dimensionnement d'une application, la valeur constitutive autoritaire est désormais le Taylor--Green multi-graines exécuté sur le même chemin numérique que la production. L'acoustique x13h reste indépendante de cette révision et conserve $c_s\simeq0.35512$ et $\nu_L\simeq1.007\times10^{-3}$ pour G08-0.72.

La validation croisée sur le fluide du benchmark VK confirme ce principe. Le même point A0 donne $\overline\nu_{\rm SRC}=6.68103\times10^{-4}$ sous SRC et, sur huit graines sous le chemin Q6-g-f de production, $\overline\nu_{\rm Q6GF}=6.65127\times10^{-4}$ avec $CV=4.90\,\%$. L'écart central de $-0.45\,\%$ est négligeable face à la variabilité. Le chemin projeté ne doit donc ni être supposé iso-visqueux par principe, ni être recalibré avec un outil SRC-only : il doit être mesuré directement, ce que permet désormais le calibrateur standalone général.

\subsection{x13h-A : qualification acoustique du point final \texorpdfstring{$\lambda/h=0.72$}{lambda/h=0.72}}

Le point final
\begin{equation}
  \gamma=8,
  \qquad
  \alpha_{\rm rot}=120^\circ,
  \qquad
  \lambda/h=0.72
  \label{eq:x13-final-micro}
\end{equation}
a ensuite été qualifié indépendamment.

Le bloc acoustique x13h-A a utilisé deux amplitudes de perturbation, 0.04 et 0.08, avec 30 réalisations par amplitude et 500 bootstraps sur le fit poolé. Les valeurs obtenues sont
\begin{equation}
  c_s(0.04)=0.354109,
  \qquad
  c_s(0.08)=0.356133,
\end{equation}
soit seulement $0.57\,\%$ d'écart, et
\begin{equation}
  \nu_L(0.04)=9.87\times10^{-4},
  \qquad
  \nu_L(0.08)=1.027\times10^{-3},
\end{equation}
soit environ $4\,\%$ d'écart. Le résultat consolidé est
\begin{equation}
  \boxed{
  c_s=0.35512,
  \qquad
  \nu_L=1.007\times10^{-3}.
  }
  \label{eq:x13-final-acoustic}
\end{equation}
Les fits poolés ont $R^2>0.999$.

À faible amplitude, plusieurs fits individuels sont \texttt{REVIEW} ou \texttt{INVALID}, alors que le fit d'ensemble reste excellent. À amplitude 0.08, presque tous les fits individuels passent. Ce résultat confirme que l'échec d'un fit individuel sur une perturbation faible ne doit pas être interprété comme une instabilité du fluide ; il traduit souvent simplement un signal trop proche du bruit stochastique.

Il est également notable que la réponse acoustique de 0.72 est presque identique à celle de 0.48 : $c_s$ varie de moins de $0.1\,\%$ et $\nu_L$ de quelques pourcents. Le gain de coût apporté par~\eqref{eq:x13-dt-gain} n'a donc pas été payé par une dégradation visible de la propagation acoustique.

\subsection{x13h-B : viscosité en fonction de la densité et frontière de localité}

Le bloc densité a mesuré $\nu_T(\gamma)$ pour
\begin{equation}
  \gamma\in\{3,4,6,8,12,16\},
\end{equation}
avec $N_\lambda=128$ pour tous les points et $N_\lambda=256$ pour les occupations basses et le point de référence. Six réalisations ont été utilisées par configuration. Les valeurs de longue longueur d'onde retenues sont résumées dans le tableau~\ref{tab:x13-density-map}.

\begin{table}[htbp]
\centering
\small
\caption{Carte diagnostique issue du cisaillement transverse x13h à $\alpha_{\rm rot}=120^\circ$ et $\lambda/h=0.72$. Les tendances de densité/localité restent informatives, mais l'échelle absolue de $\nu_T$ est supersédée pour le dimensionnement par le réaudit Taylor--Green de la section~\ref{sec:x13-highre-tg-reaudit}.}
\label{tab:x13-density-map}
\begin{tabular}{cccc}
\toprule
$\gamma$ & $\nu_T$ retenu & ratio à $\gamma=8$ & statut \\
\midrule
3  & $7.44\times10^{-4}$ & 1.44 & \texttt{LOCAL\_REVIEW} \\
4  & $6.43\times10^{-4}$ & 1.25 & \texttt{LOCAL\_PASS} \\
6  & $4.85\times10^{-4}$ & 0.94 & \texttt{LOCAL\_REVIEW} \\
8  & $5.15\times10^{-4}$ & 1.00 & \texttt{LOCAL\_PASS} \\
12 & $4.58\times10^{-4}$ & 0.89 & \texttt{SCREEN\_PASS} \\
16 & $4.79\times10^{-4}$ & 0.93 & \texttt{SCREEN\_PASS} \\
\bottomrule
\end{tabular}
\end{table}

Le point $\gamma=3$ présente un écart d'environ $24\,\%$ entre $N_\lambda=128$ et 256 et constitue la frontière basse densité actuelle. À $\gamma=4$, l'écart est de l'ordre de $6\,\%$ et la localité reste acceptable. Le cas $\gamma=6$ est plus bruité et classé \texttt{REVIEW}, mais il est encadré par 4 et 8 qui passent ; il n'est pas interprété comme une rupture monotone du transport.

À $\gamma=8$, les six fits individuels de cette campagne de cisaillement à $N_\lambda=256$ donnaient une moyenne voisine de $5.03\times10^{-4}$ et le fit d'ensemble environ $5.15\times10^{-4}$. Après le réaudit métrologique de la section~\ref{sec:x13-highre-tg-reaudit}, ces nombres ne sont plus utilisés comme échelle absolue de conception. La valeur Taylor--Green multi-graines recommandée pour G08-0.72 devient
\begin{equation}
  \boxed{
  \nu_T^{\rm design}
  \simeq
  5.60\times10^{-4},
  }
  \label{eq:x13-final-nuT}
\end{equation}
avec un statut statistique \texttt{REVIEW} ($CV\simeq17.6\,\%$). La carte ci-dessus reste utile pour localiser la frontière basse densité, mais pas pour attribuer seule un Reynolds précis.

En combinant~\eqref{eq:x13-final-acoustic} et~\eqref{eq:x13-final-nuT},
\begin{equation}
  \boxed{
  H_h
  =
  \frac{c_s h}{\nu_T}
  \simeq
  2.48.
  }
  \label{eq:x13-final-Hh}
\end{equation}

\subsection{x13h-C : domaine Mach final du fluide 0.72}

Le sweep Mach final a utilisé
\begin{equation}
  Ma\in\{0.2,0.5,0.7,0.9\}
\end{equation}
à $N_\lambda=128$ et 256. Les ondes restent convergées en longueur d'onde jusque 0.9 ; la limite d'usage vient principalement des densités visitées.

En convertissant les extrema de densité en occupation effective par
\begin{equation}
  \gamma_{\rm eff}
  \simeq
  8\frac{\rho}{\rho_0},
\end{equation}
on obtient le tableau~\ref{tab:x13-mach-envelope}.

\begin{table}[htbp]
\centering
\small
\caption{Enveloppe de densité visitée dans le sweep Mach x13h.}
\label{tab:x13-mach-envelope}
\begin{tabular}{cccc}
\toprule
$Ma$ & $\rho_{\min}/\rho_0$ & $\rho_{\max}/\rho_0$ & enveloppe $\gamma_{\rm eff}$ \\
\midrule
0.2 & 0.596 & 1.457 & 4.77--11.66 \\
0.5 & 0.479 & 1.830 & 3.83--14.64 \\
0.7 & 0.387 & 2.162 & 3.09--17.30 \\
0.9 & 0.297 & 2.623 & 2.38--20.98 \\
\bottomrule
\end{tabular}
\end{table}

La lecture recommandée est
\begin{equation}
  \boxed{
  \begin{array}{ll}
  Ma\lesssim0.2
  &:\ \text{régime linéaire pleinement qualifié},\\[0.2em]
  0.2<Ma\lesssim0.5
  &:\ \text{compressible non linéaire qualifié},\\[0.2em]
  Ma\simeq0.7
  &:\ \text{frontière constitutive / review},\\[0.2em]
  Ma\simeq0.9
  &:\ \text{stress test résolu mais hors enveloppe qualifiée}.
  \end{array}
  }
  \label{eq:x13-mach-status}
\end{equation}
À $Ma=0.5$, la raréfaction minimale correspond à $\gamma_{\rm eff}\simeq3.8$, très proche du point $\gamma=4$ qui est encore \texttt{LOCAL\_PASS}. Le régime est donc considéré comme un vrai régime compressible de production crédible. À $Ma=0.7$, les raréfactions atteignent $\gamma\simeq3$, précisément la frontière \texttt{LOCAL\_REVIEW}. À 0.9, la simulation reste numériquement résolue mais quitte nettement la carte de transport qualifiée.

\subsection{Autodiffusion et nombre de Schmidt : outil disponible mais mesure à refaire sur le point final si nécessaire}

Le calibrateur historique 0493w1 conserve un module d'autodiffusion par déplacement quadratique moyen,
\begin{equation}
  \left\langle
  \left|
  \bm x(t)-\bm x(0)
  \right|^2
  \right\rangle
  \simeq
  4D_{\rm self}t
\end{equation}
en deux dimensions aux temps suffisamment longs. Le nombre de Schmidt est ensuite
\begin{equation}
  Sc=\frac{\nu_T}{D_{\rm self}}.
\end{equation}

Le cycle x13 a concentré le budget de calcul sur $\nu_T$, $c_s$, $\nu_L$, la localité et le domaine Mach. $D_{\rm self}$ n'a donc pas été remesuré systématiquement pour chaque candidat du sweep $(\alpha,\lambda/h)$, notamment pour le point final 0.72. Il ne faut pas supposer que le passage 0.48 $\rightarrow$ 0.72 conserve exactement $Sc$. Pour toute future simulation où le mélange, la diffusion d'espèces, $Pe$ ou $Sc$ jouent un rôle, un run MSD sur le jeu final doit être ajouté à la calibration standard.

\subsection{Rôle de \texorpdfstring{$k_BT$}{kBT} : pas de levier Reynolds indépendant démontré}

Le cycle x13 n'a pas balayé $k_BT$ indépendamment. Cette absence est volontairement documentée car il serait tentant de réduire $k_BT$ pour diminuer la viscosité. À $(\gamma,\alpha_{\rm rot},\lambda/h)$ fixé,~\eqref{eq:x13-lambda} impose cependant
\begin{equation}
  \Delta t
  \propto
  \frac{1}{\sqrt{k_BT/m}}.
\end{equation}
Dans la loi SRD, les deux contributions à la viscosité évoluent alors, au premier ordre, comme
\begin{equation}
  \nu_T\propto h\sqrt{\frac{k_BT}{m}},
\end{equation}
tandis que
\begin{equation}
  c_s\propto\sqrt{\frac{k_BT}{m}}.
\end{equation}
Par conséquent,
\begin{equation}
  H_h=\frac{c_sh}{\nu_T}
\end{equation}
reste approximativement invariant. Baisser $k_BT$ à vitesse $U$ fixée baisse la viscosité, mais baisse aussi $c_s$ et augmente $Ma$ : le gain apparent de Reynolds est payé par une compressibilité plus forte.

Cette invariance d'échelle n'a pas encore été testée explicitement sur le binaire courant. Un test futur peu coûteux peut utiliser le point 0.72 avec $k_BT$ divisé ou multiplié par quatre et $\Delta t$ multiplié ou divisé par deux, de sorte que $\lambda/h$ reste constant. On devrait alors vérifier
\begin{equation}
  c_s\propto\sqrt{k_BT},
  \qquad
  \nu_T\propto\sqrt{k_BT},
  \qquad
  D_{\rm self}\propto\sqrt{k_BT},
\end{equation}
et l'invariance, à l'incertitude statistique près, de $H_h$ et $Sc$. Tant que ce contrôle n'est pas fait, $k_BT$ doit être considéré comme un paramètre d'échelle, non comme un bouton indépendant d'optimisation de Reynolds.

\subsection{Recommandations pratiques pour les futures simulations}

Le point de référence économique actuel est
\begin{equation}
  \boxed{
  \gamma=8,
  \qquad
  \alpha_{\rm rot}=120^\circ,
  \qquad
  \lambda/h=0.72.
  }
  \label{eq:x13-reference}
\end{equation}
Dans les unités de calibration,
\begin{equation}
  \boxed{
  c_s\simeq0.3551,
  \qquad
  \nu_T^{\rm design}\simeq5.60\times10^{-4},
  \qquad
  \nu_L\simeq1.01\times10^{-3},
  \qquad
  H_h\simeq2.48.
  }
  \label{eq:x13-reference-values}
\end{equation}
La relation~\eqref{eq:x13-Re-Ma-Hh} devient approximativement
\begin{equation}
  Re\simeq2.48\,Ma\,\frac{L}{h}.
\end{equation}
Pour $Re=10^4$, il faut donc environ
\begin{equation}
  \frac{L}{h}\simeq8080
  \quad\text{à}\quad Ma=0.5,
\end{equation}
\begin{equation}
  \frac{L}{h}\simeq5770
  \quad\text{à}\quad Ma=0.7,
\end{equation}
et
\begin{equation}
  \frac{L}{h}\simeq4490
  \quad\text{à}\quad Ma=0.9.
\end{equation}
La première valeur se situe dans le domaine compressible actuellement qualifié. La seconde touche la frontière basse densité, et la troisième n'est qu'un stress test. Ces nombres représentent la résolution de la \emph{longueur caractéristique} $L$, non la taille totale de la grille. Un cylindre ou un obstacle de hauteur $H$ résolu par plusieurs milliers de cellules dans un domaine de plusieurs $H$ de long et de haut conduit à une grille globale beaucoup plus grande.

Pour une nouvelle application, le protocole recommandé est désormais :
\begin{enumerate}[leftmargin=2em]
  \item choisir un point local $(\gamma,\alpha_{\rm rot},\lambda/h)$ déjà qualifié ;
  \item fixer $h$ et $L/h$ à partir de $Re$, $Ma$ et du coût de grille acceptable ;
  \item relancer un calibrateur transverse multi-runs sur le binaire et le chemin physique effectivement utilisés, avec au moins une vérification de longueur d'onde si le point n'est pas exactement celui déjà qualifié ;
  \item utiliser un fit d'ensemble pour la valeur constitutive et conserver la dispersion des runs individuels comme incertitude ;
  \item mesurer $c_s$ et $\nu_L$ par mode amorti si la compressibilité ou les ondes sont importantes ;
  \item mesurer $D_{\rm self}$ par MSD si $Sc$ ou $Pe$ intervient ;
  \item pour une simulation capillaire, combiner $\nu_T$ avec la $\sigma_{\rm eff}$ mécanique de x12yl et ne pas absorber la dispersion x12cal dans une tension superficielle artificielle ;
  \item vérifier que l'enveloppe de densité reste dans une plage où $\nu_T(\gamma)$ est locale ; pour le point~\eqref{eq:x13-reference}, $\gamma\gtrsim4$ est actuellement la zone sûre et $\gamma\simeq3$ la frontière de revue ;
  \item enregistrer le hash du binaire, le HEAD, les paramètres locaux et la statistique de calibration. Un seed nominal identique n'est pas une garantie de reproductibilité bit à bit.
\end{enumerate}

Pour les validations de von Kármán, la règle est maintenant vérifiée sur le cas Zovatto lui-même. La fréquence de shedding est déterminée à mieux de $1\,\%$ entre POD et sondes, tandis que la viscosité Q6-g-f mesurée sur huit graines possède un $CV$ de $4.90\,\%$. Cette dispersion se transmet directement à $Re$, mais la moyenne $\nu=6.65127\times10^{-4}$ place la fenêtre pleinement développée du run cluster à $Re_H=282.125$, donc pratiquement au point bibliographique $280$. La variabilité de $\nu$ doit être conservée comme barre d'abscisse lorsqu'une dépendance fine $T^*(Re)$ est discutée ; elle ne justifie plus de déplacer la valeur centrale suffisamment pour expliquer le résidu de période désormais réduit à environ $4$--$5\,\%$.

\subsection{Statut du chantier de caractérisation}

Le cycle x13 atteint l'objectif du chantier de propriétés de transport. Le projet dispose maintenant d'une vision opérationnelle de l'application
\begin{equation}
  \left(
  N_x,N_y,L_x,L_y,\Delta t,k_BT,m,\gamma,\alpha_{\rm rot},
  \mathcal T,\mathcal G,\mathcal S_\pm
  \right)
  \longrightarrow
  \left(
  \nu_T,\nu_L,c_s,D_{\rm self}
  \right),
\end{equation}
à condition de comprendre que $N_x,N_y,L_x,L_y$ agissent d'abord via $h$ et les longueurs d'onde disponibles, et que $(\gamma,\alpha_{\rm rot},\lambda/h)$ définissent la famille locale de fluide. Cette application n'est pas une formule analytique universelle : les lois SRD donnent des estimations de positionnement, puis les calibrateurs numériques certifient les coefficients, leur incertitude et leur domaine de validité.

Le banc de caractérisation de référence comprend désormais :
\begin{itemize}[leftmargin=2em]
  \item le Taylor--Green historique multi-graines comme métrologie primaire de $\nu_T$ destinée au Reynolds ;
  \item le cisaillement transverse pur comme diagnostic secondaire de relaxation/localité, non comme autorité absolue de viscosité ;
  \item le mode longitudinal amorti pour $c_s$ et $\nu_L$ sous SRC non projeté ;
  \item le MSD périodique pour $D_{\rm self}$ et $Sc$ ;
  \item le calibrateur standalone général, path-aware, pour exécuter exactement ces mesures sous \texttt{src} ou \texttt{src-q6-g-f} sans dépendance à un runner de campagne ;
  \item les comparaisons de longueur d'onde pour la localité en $k\lambda$ et les ensembles multi-runs pour la fiabilité statistique ;
  \item la carte de densité et le sweep Mach comme diagnostics de domaine de validité, avec prudence sur l'échelle absolue des viscosités issue du cisaillement historique ;
  \item les calibrateurs x12yl/x12cal pour séparer tension superficielle mécanique et dispersion capillaire.
\end{itemize}

Pour les futurs développements, ces calibrateurs doivent constituer le point de départ avant toute conclusion sur un Reynolds effectif, une fréquence de shedding, une diffusion, une compressibilité ou une dynamique capillaire. La priorité n'est plus de chercher à l'aveugle un jeu de paramètres ``plus rapide'' ou ``moins visqueux'', mais de choisir un fluide déjà cartographié, le recalibrer brièvement dans l'environnement de production, puis dimensionner l'application à partir de ses coefficients effectivement mesurés.




\section{Cycle 0493x14 : couplage liquide--gaz explicite, pression interfaciale, capillarité et transfert de quantité de mouvement}
\label{sec:x14-liquid-gas}

\subsection{Objectif et principe de séparation des mécanismes}

Le cycle 0493x14 étend la chaîne de surface libre qualifiée à une véritable phase gazeuse particulaire explicite, sans revenir sur la fermeture liquide validée. Le point de départ reste donc
\begin{equation}
  \text{Q6-g-f}
  + \text{x9 Laplace}
  + \text{x10o}
  + \text{CIC}
  + \text{Q2}
  + \text{x10p/q}
  + \text{x10u}
  + \text{x10v full-vector}
  + \text{x12a},
\end{equation}
avec resampling, refill vide et kick viriel désactivés dans les calibrateurs interfaciaux. Le contrat de développement x14 est de conserver cette chaîne liquide inchangée et d'ajouter uniquement les briques nécessaires au côté gaz.

La difficulté est de ne pas confondre quatre échanges physiques différents :
\begin{enumerate}[label=(\roman*)]
  \item l'imperméabilité cinétique normale de l'interface ;
  \item la pression thermodynamique du gaz qui entre dans la condition de contrainte normale ;
  \item la partie cinétique hors équilibre de l'impulsion normale associée aux impacts/réflexions gazeux ;
  \item le transfert tangentiel visqueux entre les deux fluides.
\end{enumerate}
La tension superficielle n'est pas réimplémentée dans x14. Elle reste portée par x9 comme saut de Laplace dans le potentiel de Dirichlet de Q6-g-f. Avec une phase gaz explicite, la condition complète reste
\begin{equation}
  p_L\big|_\Gamma = p_G + \sigma\kappa,
\end{equation}
soit, avec la convention du solveur,
\begin{equation}
  \phi_\Gamma =
  \frac{\Delta t}{\rho_{L,\mathrm{ref}}}
  \left(p_G + \sigma\kappa - p_{\mathrm{ref}}\right).
  \label{eq:x14-dirichlet}
\end{equation}
La contribution capillaire et la pression gazeuse se composent donc dans le même terme de frontière ; aucune force CSF et aucun kick capillaire particulaire supplémentaire ne sont introduits.

Les cas de travail x14 utilisent deux espèces enregistrées. La phase liquide est de type~1, de masse particulaire $m_L=1$, avec $q6Strength=1$ et une température propre. La phase gaz est de type~2, de masse $m_G=0.1$, avec $q6Strength=0$ et sa propre température. Le gaz participe aux collisions SRC communes mais ne reçoit aucune correction Q6 directe. Cette construction conserve ainsi une collision mésoscopique commune dans les cellules mixtes, qui fournit naturellement le principal mécanisme d'échange tangentiel.

\subsection{Thermostat séparé par espèce : x14a--x14j}

Une phase gazeuse explicite de masse et de température différentes ne peut pas être qualifiée avec un thermostat unique appliqué indistinctement à tout le mélange. Le premier ajout structurel de x14 est donc un thermostat par espèce, opt-in et désactivé par défaut. Le paramètre
\begin{center}
\small
\nolinkurl{speciesThermostatEnable = true}
\end{center}
active cette séparation. Chaque espèce enregistrée reçoit alors une cible dynamique
\begin{center}
\small
\nolinkurl{speciesKThermostatTargetKBT = value},
\end{center}
où $K$ est l'indice de déclaration ; une valeur négative signifie héritage de la cible globale. Lorsque le mode est actif, le parseur exige un thermostat global activé, un registre d'espèces non vide, \nolinkurl{speciesRequireRegisteredTypes=true} et une cible positive finie résolue pour chaque espèce.

L'implémentation x14d conserve volontairement la \emph{collision SRC commune} au mélange. Elle ne remplace donc pas la collision par deux fluides qui ne se verraient plus. Seule l'étape de remise à température est séparée par type après collision. Dans le chemin CUDA résident, les espèces réutilisent le même workspace cellulaire successivement : la mémoire ne croît pas avec le nombre d'espèces. Le pont x14g utilise les identifiants de cellule exacts post-streaming/post-grid-shift du chemin SRC persistant, au lieu de réutiliser un état de cellule préstream devenu obsolète. Cette distinction est importante pour ne pas introduire artificiellement de diffusion ou de défaut de température au voisinage d'une interface mobile.

Enfin, x14m rend ce thermostat compatible avec x12a : le refroidissement thermique local de petites structures reste appliqué uniquement au côté A liquide, tandis que le gaz conserve sa propre cible $k_BT_G$. La qualification x14j sur goutte à deux températures a servi de cas d'intégration pour cette architecture avant d'activer les fermetures cinétiques bilatérales.

\subsection{Géométrie cinétique bilatérale et réflexion gazeuse : x14k--x14m}

La chaîne x10 historique était volontairement qualifiée pour une interface liquide--vide. x14k introduit le paramètre
\begin{center}
\small
\nolinkurl{phaseInterfaceKineticBilateralRelocation = true},
\end{center}
qui autorise une seule espèce gazeuse explicite du côté B à consommer la même paroi cinétique mobile x10o+CIC+Q2, avec sens de phase inversé. Ce gate ne change à lui seul ni vitesse, ni masse, ni type : il étend uniquement la géométrie de relocalisation positionnelle au second côté de l'interface. Pour éviter les configurations ambiguës, le mode bilatéral exige dans le snapshot courant des sélecteurs explicites \nolinkurl{type:<id>} pour A et B, une phase A liquide unique projetée, une phase B gazeuse unique non projetée, et
\begin{equation}
  \texttt{phaseInterfaceKineticReflectionFraction}=1.
\end{equation}
Il est incompatible avec une cible d'évaporation active. Le chemin historique $0<r<1$ reste donc une ablation de transmission/évaporation non qualifiée et n'est pas utilisé pour le couplage liquide--gaz courant.

x14l ajoute ensuite la fermeture cinétique du gaz proprement dite, commandée par
\begin{center}
\small
\nolinkurl{MPCD_X14L_GAS_SPECULAR_REFLECTION=1}.
\end{center}
Lorsqu'une particule gazeuse tente de franchir l'interface matérielle, sa vitesse est réfléchie spéculairement dans le repère local de l'interface : seule la composante normale relative change de signe. La composante tangentielle reste inchangée. Le côté liquide conserve exactement x10u et la fermeture cinétique déjà qualifiée ; x14l évite que la correction positionnelle x10u n'écrase la vitesse spéculaire calculée pour le gaz.

x14m constitue le premier assemblage complet : liquide x10o+CIC+Q2+x10p/q+x10u+x10v+x12a inchangé, gaz x14l spéculaire, thermostat séparé par espèce. Le mécanisme x10w reste désactivé. Cette étape fixe une propriété importante pour la suite : la réflexion gaz assure l'imperméabilité normale, mais elle ne crée volontairement aucun frottement tangentiel ad hoc. Le transfert tangent doit donc émerger essentiellement des collisions SRC communes dans les cellules interfaciales.

\subsection{Pression gazeuse et correction de volume accessible : x14n--x14s}

Le couplage normal macroscopique du gaz est porté par x6g. Les ablations x14n et x14o ont d'abord séparé la fermeture cinétique et la pression : x14n coupe la fermeture gazeuse, tandis que x14o utilise une pression gazeuse constante choisie de façon à annuler la contribution de jauge. La tension superficielle x9 reste active dans ces ablations, ce qui permet de ne pas confondre $p_G$ et $\sigma\kappa$.

Avec une phase gazeuse particulaire, la formule brute
\begin{equation}
  p_G^{\mathrm{raw}} = \frac{N_G k_B T_G}{A_c}
\end{equation}
sous-estime la pression locale lorsque le nombre de particules est compté dans une cellule dont seule une fraction est géométriquement accessible au gaz. L'analyse offline x14r a relié ce défaut à l'alpha Q6 du côté gaz et conduit à la variante x14s
\begin{center}
\small
\nolinkurl{MPCD_Q6_PHASE_GAS_PRESSURE_MODE_0493X6G=eos_accessible_volume}.
\end{center}
En posant
\begin{equation}
  g = 1-\alpha_{\mathrm{Q6}},
\end{equation}
la fraction accessible utilisée par le code est
\begin{equation}
  f_{\mathrm{Q6}} =
  \min\!\left[
    1,\,
    \max\!\left(
      0.05,\,
      \frac{g-g_0}{g_1-g_0}
    \right)
  \right],
  \qquad
  g_0=0.37396789550781245,\quad
  g_1=0.9153533203125,
  \label{eq:x14-accessible-volume}
\end{equation}
et la pression devient
\begin{equation}
  p_G^{\mathrm{acc}} =
  \frac{p_G^{\mathrm{raw}}}{f_{\mathrm{Q6}}}.
  \label{eq:x14-gas-pressure}
\end{equation}
Le plancher $0.05$ est un garde défensif ; sur une trace gaz x6f valide, la fraction issue du fit est déjà sensiblement plus grande. Dans les runners x14, la valeur globale de $k_BT$ consommée par l'EOS x6g est explicitement alignée sur la cible $k_BT_G$ de l'espèce gazeuse, tandis que le thermostat séparé impose indépendamment $T_L$ et $T_G$ aux particules. La correction n'est évaluée que sur les faces d'interface $\alpha=0.5$ déjà préparées par x6f/x6g. Elle n'ajoute donc ni passe $O(N_{\mathrm{cell}})$ globale, ni passe particulaire, ni buffer résident supplémentaire. Les modes historiques \nolinkurl{eos} et \nolinkurl{constant} restent inchangés.

L'intérêt de cette formulation est également architectural : x14s corrige uniquement l'estimation thermodynamique $p_G$ ; la capillarité continue d'ajouter $\sigma\kappa$ dans la même condition \eqref{eq:x14-dirichlet}. La modification de l'EOS gazeuse ne change donc ni le calcul de courbure x9, ni le CG, ni B1.

\subsection{Séparation pression thermodynamique / impulsion cinétique normale : x14t--x14v}

La qualification a ensuite été conduite en deux étapes pour éviter tout double comptage. Le piston plan x14t impose un chargement normal gazeux sans ajouter de mécanisme impulsionnel direct. Il montre que x6g, notamment avec la correction de volume accessible x14s, transmet correctement la pression thermodynamique moyenne au liquide. Ce test qualifie le terme d'équilibre qui figure déjà dans le Dirichlet Q6-g-f.

Le benchmark x14u impose au contraire un gaz incident avec quantité de mouvement normale dirigée. Avec x6g réglé en mode constant pour supprimer la variation thermodynamique, la seule réflexion x14l rend le gaz imperméable mais ne restitue pratiquement pas au liquide le moment perdu : le gain de transfert direct reste de l'ordre de $G_{\mathrm{mom}}\simeq0.07$ vers le pas~20. Cette expérience localise donc un terme physique manquant, distinct de $p_G$ : la partie cinétique hors équilibre du flux d'impulsion normal.

x14v ajoute exactement cette composante, avec le gate
\begin{center}
\small
\nolinkurl{MPCD_X14V_GAS_KINETIC_EXCESS_KICK=1}.
\end{center}
Le principe est
\begin{equation}
  \mathbf J_{\mathrm{excess}}
  =
  \mathbf J_{\mathrm{refl,\,gas}}
  -
  \mathbf J_{\mathrm{thermo,\,x6g}},
  \label{eq:x14v-excess}
\end{equation}
où $\mathbf J_{\mathrm{refl,\,gas}}$ est l'impulsion réelle agrégée des réflexions spéculaires et $\mathbf J_{\mathrm{thermo,\,x6g}}$ la traction d'équilibre déjà représentée par x6g. Seul le résidu est transféré collectivement au liquide. Cette soustraction est essentielle : ajouter directement toute l'impulsion de réflexion par-dessus x6g doublerait la partie thermodynamique de la contrainte normale.

L'implémentation respecte les contraintes de coût du chemin de production. L'impulsion gazeuse est accumulée par propriétaire d'interface dans le stockage déjà disponible ; la masse liquide CIC post-x10u est maintenue par mise à jour signée, sans seconde passe particulaire ; un kernel cellule remplace des remises à zéro déjà nécessaires à x10v pour former l'excès et le distribuer sur le support CIC liquide ; enfin le kick est fusionné dans le redépôt de moments post-cinétique existant. Il n'y a donc ni nouveau buffer résident permanent, ni transfert hôte, ni nouvelle traversée $O(N_p)$. Les lois x10u, x10v et x12a ne sont pas modifiées.

Le benchmark x14u miroir \emph{bottom-impact/top-impact} ferme quantitativement ce bilan. Avec x14v et x6g constant, le gain de moment vaut environ
\begin{equation}
  G_{\mathrm{mom}}
  =
  1.000,\ 1.001,\ 0.988,\ 0.978
\end{equation}
aux pas $1,5,10,20$, soit une moyenne des pas $1\ldots20$ proche de $0.988$. Aux pas plus tardifs il reste voisin de l'unité. Avec x14v et x6g \nolinkurl{eos_accessible_volume}, le même gain reste voisin de un, sans signature significative de double comptage. À ce stade, l'imperméabilité normale, la pression thermodynamique $p_G$ et la composante cinétique normale excédentaire sont donc qualifiées séparément puis en composition.

Le coût mesuré est inférieur à la résolution du chronométrage utilisé. Sur une goutte de 500 pas sans LiveVis ni recording, les temps pairés sont $27.58$~s x14v OFF et $27.15$~s x14v ON pour le temps écoulé, avec des temps utilisateur pratiquement identiques. Ce résultat est cohérent avec l'absence de nouvelle passe particulaire globale.

\subsection{Non-régression de la goutte statique avec gaz explicite}

Une goutte statique $400\times400$, $\gamma=20$, $R/h=40$, $\sigma=2560$, $k_BT_L=0.02$, $k_BT_G=0.08$ et $\Delta t=0.002$ a été exécutée avec x6g accessible-volume, x14l et x14v, en gardant toute la chaîne liquide. Au pas~500,
\begin{equation}
  \langle\kappa\rangle = 6.4136905,
  \qquad
  \kappa_{\mathrm{exact}}=6.4,
\end{equation}
soit une erreur de Laplace de $+0.214\,\%$. Le rayon $R_{50}/h=40.7039$ et l'épaisseur $10$--$90\,\%$ vaut $5.0478h$. Les contrôles de contamination donnent zéro liquide à $R_{50}+4h$ et $R_{50}+8h$, et zéro gaz à $R_{50}-4h$. Cette expérience ne remplace pas une campagne statistique longue mais exclut une non-régression grossière de la géométrie, de la loi de Laplace ou du confinement des phases après ajout de x14v.

\subsection{Transfert tangentiel : benchmark Couette biphasique x14w}

La dernière interaction élémentaire est le transfert tangent. Le benchmark x14w est volontairement construit sans nouvelle physique de paroi ni nouveau diagnostic runtime. La géométrie est
\begin{equation}
  G_{\mathrm{bottom}}\;|\;L\;|\;G_{\mathrm{top}},
\end{equation}
avec $x$ périodique et deux parois solides en $y$. Les deux parois physiques touchent le gaz et non le liquide. Sur la grille de référence $128\times128$, $L_x=L_y=0.5$ et $h=1/256$, les couches ont respectivement $32h$, $64h$ et $32h$. La largeur liquide dépasse $2R_c$ avec $R_c/h\simeq25.3$, de sorte que x12a n'est pas artificiellement déclenché dans son coeur.

Les parois utilisent le mécanisme thermique/collisionnel déjà employé dans les validations de canal :
\begin{equation}
  U_{w,\mathrm{bottom}}=-U_w,
  \qquad
  U_{w,\mathrm{top}}=+U_w,
\end{equation}
avec masse virtuelle $m_G$, température $T_G$ et occupation virtuelle $\gamma=20$. La réflexion géométrique solide ne force pas directement la composante tangentielle ; la contrainte tangentielle pariétale est donc transmise par les collisions avec particules virtuelles. À l'interface liquide--gaz, x14l laisse également la composante tangentielle inchangée : le transfert recherché doit principalement provenir du SRC commun dans les cellules mixtes.

L'analyse est entièrement offline à partir des dumps \texttt{.smpcd}. Pour éliminer les dérives communes, elle construit le profil antisymétrique
\begin{equation}
  u_a(z)=\frac{u_{\mathrm{top}}(z)-u_{\mathrm{bottom}}(z)}{2},
\end{equation}
puis ajuste loin des murs et de l'interface
\begin{equation}
  u_L(z)=a_Lz+b_L,
  \qquad
  u_G(z)=a_Gz+b_G.
\end{equation}
Les deux conditions physiques visées sont
\begin{equation}
  u_L^\Gamma=u_G^\Gamma,
  \qquad
  \mu_L a_L=\mu_G a_G.
  \label{eq:x14-couette-targets}
\end{equation}
Le slip de paroi gaz est suivi séparément afin de ne pas attribuer à l'interface un défaut de couplage aux murs.

Le premier run à rapport signal/bruit suffisant a utilisé $U_w=0.075$ pendant 18000 pas. Le moyennage temporel des profils tardifs $t=8\ldots36$ donne
\begin{center}
\begin{tabular}{lcc}
\toprule
Grandeur & Liquide & Gaz\\
\midrule
pente $a$ & $0.11756$ & $0.49910$\\
$R^2$ linéaire & $0.97694$ & $0.90314$\\
\bottomrule
\end{tabular}
\end{center}
et
\begin{equation}
  \frac{a_L}{a_G}=0.23555.
\end{equation}
Pour les paramètres SRC des deux phases, l'estimation analytique SRD indépendante donne
\begin{equation}
  \frac{\mu_G}{\mu_L}\simeq0.23769,
\end{equation}
soit un écart de l'ordre de $0.9\,\%$ sur la continuité de contrainte tangentielle relativement à cette référence analytique. Cette comparaison est un contrôle indépendant utile, mais elle ne remplace pas une calibration Taylor--Green path-aware des deux viscosités si une métrologie sub-pourcent devient nécessaire. Le slip interfacial extrapolé vaut environ
\begin{equation}
  \frac{\Delta u_\Gamma}{U_w}\simeq-5.7\,\%,
\end{equation}
et le slip gaz aux parois environ $+3.6\,\%$ de $U_w$. Le rapport des pentes et la linéarité des profils sont donc déjà \emph{PASS-like}; en revanche la continuité de vitesse à l'interface reste à confirmer quantitativement sur une seconde graine, car le gaz conserve un bruit thermique important. Le run de confirmation inter-graine est maintenu à $U_w=0.075$ sans aucune modification C++/CUDA.

Cette lecture est importante pour le statut du développement. Le Couette ne fournit à ce stade aucun motif pour ajouter un frottement tangentiel artificiel à x14l. Le mécanisme SRC commun transmet déjà une contrainte tangentielle compatible avec le rapport de viscosités attendu, tandis que le petit slip résiduel doit d'abord être séparé de la variance statistique et du slip pariétal.


\subsection{Interface courbe : séparation de la loi locale et de la résultante globale}

La goutte oscillante diphasique $n=2$ a servi de banc intégré pour vérifier simultanément la pression gazeuse, la capillarité, la réflexion x14l, le transfert cinétique x14v et la projection liquide. La référence inertielle utilisée est celle d'un cylindre bidimensionnel à deux fluides,
\begin{equation}
  \omega_{n}^{2}=\frac{n(n^2-1)\sigma}{(\rho_L+\rho_G)R^3},
\end{equation}
avec $R=0.15625$, $\rho_L=1310720$, $\rho_G=131072$, $\rho_G/\rho_L=0.1$ et $\sigma=2560$. Pour $n=2$, elle donne $\omega_{2,\mathrm{th}}=1.6711455$ et $T_{2,\mathrm{th}}=3.7598075$. Cette expérience a montré qu'une fermeture pouvait conserver correctement le mode de forme $n=2$ tout en laissant apparaître une translation parasite $n=1$ : les deux propriétés doivent donc être qualifiées séparément.

Le point de départ x14v écrit l'impulsion transmise au liquide sous la forme
\begin{equation}
  \bm J_{\mathrm{excess}}=\bm J_{\mathrm{raw}}-\bm J_{\mathrm{thermo}},
\end{equation}
où $\bm J_{\mathrm{raw}}$ est l'impulsion réellement perdue par le gaz lors des réflexions x14l et $\bm J_{\mathrm{thermo}}$ la traction d'équilibre déjà portée par x6g. La difficulté mise en évidence par les interfaces courbes ne porte pas sur cette séparation physique, mais sur la représentation discrète de $\bm J_{\mathrm{thermo}}$. Q6 agit sur des faces de cellules x6f/x6g, alors que la géométrie cinétique et le scatter liquide utilisent les segments x10n et les poids CIC. Plusieurs formulations ont donc été comparées en conservant la chaîne liquide inchangée.

La formulation locale finalement retenue est x14ad. Les normales et longueurs de segments restent celles de x10n, qui avaient préservé la dynamique de forme. La pression de jauge locale n'est plus prise dans une cellule gazeuse voisine choisie géométriquement : elle est reconstruite comme la moyenne arithmétique des deux faces x6g représentées qui terminent le même segment de marching-squares ; si une seule face est représentée elle est utilisée seule, et l'ancien échantillonnage n'intervient qu'en fallback lorsqu'aucune des deux ne l'est. Les identifiants des arêtes sont empaquetés dans le scratch d'un octet par cellule déjà alloué par x10m. Aucune nouvelle passe globale ni aucun nouveau buffer $O(N)$ n'est ajouté.

La résultante globale de cette traction locale est corrigée par la projection conservative introduite en x14ac et conservée dans x14ad. Si $\Delta p_s$ désigne la pression de jauge nominale sur un segment $s$, $L_s$ sa longueur et $\bm m_s$ sa normale de force, on cherche la correction de norme quadratique minimale
\begin{equation}
  \Delta p_s^*=\Delta p_s+\bm m_s\cdot\bm\lambda,
\end{equation}
avec
\begin{equation}
  \bm M\bm\lambda=\bm F_{\mathrm{target}}-\bm F_{\mathrm{local}},
  \qquad
  \bm M=\sum_s L_s\,\bm m_s\bm m_s^{\mathsf T}.
  \label{eq:x14-resultant-projection}
\end{equation}
Le scratch nécessaire contient sept doubles. Il est rempli dans le kernel x10n déjà exécuté et résolu dans le kernel x14v existant ; les géométries de rang insuffisant retombent volontairement sur l'absence de projection plutôt que d'amplifier une composante non résolue. x14ad a ainsi fourni la meilleure combinaison obtenue entre traction locale cohérente avec x6g et préservation du mode de forme, mais une dérive barycentrique résiduelle $|dU/dt|\simeq4.74\times10^{-4}$, équivalente à une erreur globale d'environ $0.39$ unité d'impulsion par pas sur le benchmark, subsistait.

\subsection{Diagnostics causaux x14ae--x14af}

La dérive résiduelle a été traitée par diagnostics discriminants plutôt que par ajout immédiat d'une nouvelle loi. x14ae teste l'hypothèse d'une perte terminale dans le scatter CIC : une impulsion valide pourrait ne trouver aucune masse liquide après les quatre noeuds CIC puis le fallback jusqu'au rayon 2. Le diagnostic cumule uniquement le nombre de telles pertes et leurs deux composantes d'impulsion, sans modifier le scatter. Sur le run x14ad de référence, \texttt{terminalNoSupportCount=0} et $J_{x,\mathrm{lost}}=J_{y,\mathrm{lost}}=0$ ; l'hypothèse d'une quantité de mouvement perdue faute de support liquide est donc rejetée.

x14af mesure ensuite, à chaque pas contrôlé, les résultantes nécessaires au bilan global sans modifier la physique. Sur 200 pas avec \texttt{SUMMARY\_EVERY=1}, le bilan se ferme au niveau de l'arrondi numérique selon
\begin{equation}
  \Delta\bm P_{\mathrm{tot}}
  =\bm J_{Q6}^{\mathrm{applied}}-\bm J_{\mathrm{thermo}}^{\mathrm{x14ad}},
  \label{eq:x14af-balance}
\end{equation}
avec un résidu RMS de l'ordre de $8\times10^{-12}$. Ce résultat localise précisément le défaut : la résultante thermodynamique reconstruite sur l'interface x14ad diffère légèrement de l'impulsion que le chemin discret CG $\rightarrow$ corrections faces/cellules $\rightarrow$ B1/RT0 applique réellement aux particules liquides. Il n'est donc ni nécessaire ni justifié de modifier la réflexion x14l, le scatter CIC ou la loi locale de pression.

\subsection{Fermeture conservative x14ai et correctif x14ai-fix1}

x14ai remplace seulement la cible globale de l'équation~\eqref{eq:x14-resultant-projection} par la résultante Q6 réellement appliquée :
\begin{equation}
  \bm F_{\mathrm{target}}
  =\frac{\bm J_{Q6}^{\mathrm{applied}}}{\Delta t},
  \qquad
  \bm M\bm\lambda
  =\bm F_{\mathrm{target}}-\bm F_{\mathrm{local}}.
  \label{eq:x14ai-target}
\end{equation}
La loi locale x14ad, x6g, x14l et l'ensemble x10o+CIC+Q2+x10p/q+x10u+x10v+x12a restent inchangés. Pour une composante liquide fermée et isolée des frontières Q6 externes, l'identité recherchée devient alors
\begin{equation}
  \bm J_{L,\mathrm{interface}}
  =\bm J_{Q6}+\bm J_{\mathrm{raw}}-\bm J_{\mathrm{thermo}}
  =\bm J_{\mathrm{raw}},
\end{equation}
c'est-à-dire que le mode de translation interfacial est porté uniquement par l'échange cinétique réel avec le gaz, tandis que la distribution locale de pression reste celle de x14ad.

L'implémentation respecte le contrat de coût du chantier. Deux doubles persistants, soit 16 octets, accumulent dans B1 la quantité $m\,\Delta\bm v$ ; la réduction se fait par warp/bloc et la cible est consommée dans le kernel x14v déjà présent. Il n'y a ni nouveau kernel, ni nouvelle passe particulaire ou cellule, ni transfert device--host par pas, ni second CG, ni buffer $O(N_p)$ ou $O(N_{\mathrm{cell}})$. Le timing apparié sur trois paires de runs de 400 pas donne $0.06855\,\mathrm{s/pas}$ avec x14ai OFF et $0.067925\,\mathrm{s/pas}$ avec x14ai ON ; le surcoût médian apparié vaut $-0.695\,\%$, inférieur à la dispersion de mesure, et le verdict est \texttt{COST\_NEGLIGIBLE\_LE\_2PCT}.

Le premier x14ai révélait toutefois deux sauts discrets du moment total, vers les pas 710 et 1260. L'audit de B1 a montré que les particules reçoivent en domaine périodique
\begin{equation}
  \Delta\bm v_{\mathrm{applied}}
  =\Delta\bm v_{\mathrm{RT0}}-\bm C_{\mathrm{periodic}},
\end{equation}
alors que la première réduction x14ai cumulait seulement $m\Delta\bm v_{\mathrm{RT0}}$. x14ai-fix1 corrige la cible une fois globalement,
\begin{equation}
  \bm J_{Q6}^{\mathrm{actual}}
  =\sum_i m_i\Delta\bm v_{\mathrm{RT0},i}
   -M_{\mathrm{active}}\bm C_{\mathrm{periodic}},
  \label{eq:x14ai-fix1}
\end{equation}
sans ajouter d'opération par particule. Le marqueur runtime devient\begin{center}\small\nolinkurl{deviceAppliedQ6ResultantClosure=B1-exact-post-periodic-device-target}.\end{center}

La portée de cette fermeture est volontairement limitée. Elle est candidate de production pour une composante liquide fermée qui ne touche ni paroi, ni inlet/outlet, ni autre frontière externe Q6. Dans ces autres topologies, $\bm J_{Q6}^{\mathrm{applied}}$ peut contenir une réaction extérieure physique ; la compenser comme une erreur d'interface serait incorrect. Cette restriction fait partie de la définition de x14ai-fix1 et non d'une simple réserve de validation.

\subsection{Qualification intégrée x14ai-fix1 : oscillation $n=2$ et traînée gazeuse}

Le run de goutte oscillante x14ai-fix1 reprend exactement le cas $400\times400$, $R/h=40$, $\epsilon=0.04$, $\sigma=2560$, $m_L=1$, $m_G=0.1$, $k_BT_L=0.02$, $k_BT_G=0.08$ et $\Delta t=0.002$. Sur $0<t\le4$, le fit amorti donne
\begin{equation}
  \omega=1.6375276,
  \qquad G_\omega=0.9798833,
  \qquad \beta=0.1694319,
  \qquad R^2=0.997111.
\end{equation}
L'écart de fréquence est donc $-2.01\,\%$ relativement à la théorie à deux fluides. Il n'y a encore que deux passages par zéro, de sorte que la fréquence par zero-crossing n'est pas utilisée comme métrique de qualification. La population de chaque phase reste constante, les températures moyennes restent proches des cibles ($\langle k_BT_L\rangle=0.01929$, $\langle k_BT_G\rangle=0.07961$), aucune troncature de carrier ni face d'interface non couverte n'est observée, et le clipping de courbure reste rare : fraction maximale $0.815\,\%$ des faces capillaires sur un échantillon.

Le résultat principal est la fermeture du mode global. En sommant les moments liquide et gaz à partir de \texttt{species\_runtime}, on obtient
\begin{equation}
  \bm P_{\mathrm{tot}}(0)=\bm 0,
  \qquad
  \bm P_{\mathrm{tot}}(t=4)
  =(-3.94\times10^{-11},-1.00\times10^{-11}),
\end{equation}
avec une dérive vectorielle maximale $5.08\times10^{-11}$. Les sauts du premier x14ai ont disparu au niveau de l'arrondi numérique. Le centre de masse liquide se déplace néanmoins de $8.87\times10^{-4}$, soit $0.227h=0.00568R$, entre le début et la fin du run. Cette translation est plus grande que dans le premier x14ai avant fix1 ; elle ne doit pas être confondue avec une violation de conservation globale, car les vitesses moyennes finales liquide et gaz se compensent dans le moment total. Une prolongation de ce même cas jusqu'au pas global 6000 est prévue pour déterminer si ce déplacement relatif est borné/stochastique ou séculaire et pour obtenir plusieurs périodes utiles à l'estimation robuste de $\omega$ et de $\beta$.

Le second test intégré est x14ah, un canal périodique en $x$ avec parois solides en $y$, sans inlet/outlet ni force volumique. Un profil de Poiseuille gazeux initial entraîne une goutte initialement au repos. Ce benchmark ne possède volontairement pas de coefficient de traînée analytique imposé : il vérifie que la fermeture conservative ne supprime pas la vraie interaction gaz--liquide. Avec x14ai-fix1, le gaz reste dirigé vers $+x$ pendant tout le run, avec $U_{G,x}=0.03815$ à $t=8$, et la goutte atteint
\begin{equation}
  \Delta x=0.12480=0.7987R,
  \qquad
  U_{L,x}^{\mathrm{late}}=0.01814.
\end{equation}
Le déplacement transverse final vaut seulement $-6.52\times10^{-4}= -0.00417R$, et le rapport d'axes maximal est $1.08065$. Par rapport au benchmark x14ah antérieur sans fermeture device-side, qui donnait environ $\Delta x/R=0.51$, $U_{L,x}\simeq0.0107$ et $\Delta y/R\simeq-0.07$, la réponse aval augmente tandis que la dérive transverse absolue est réduite d'environ $94\,\%$, soit un facteur proche de 17. x14ai-fix1 ne projette donc pas artificiellement la traînée physique.

Le bilan longitudinal du canal donne $\Delta P_{G,x}\simeq-2753.25$, $\Delta P_{L,x}\simeq+1991.63$ et $\Delta P_{\mathrm{tot},x}\simeq-761.61$ sur $0\le t\le8$. Environ $72\,\%$ de la perte de moment gazeux apparaît donc dans le liquide ; le reliquat est compatible avec la réaction visqueuse des parois fixes. Le compact ne contient toutefois pas de diagnostic direct de l'impulsion pariétale permettant de fermer ce dernier terme pas à pas, et il ne faut donc pas présenter ce benchmark comme une mesure quantitative d'un coefficient de traînée. Son statut est celui d'une validation intégrée du signe, de la persistance et de l'ordre de grandeur du transfert de quantité de mouvement.

\subsection{Méthodes écartées, ablations et diagnostics non retenus}

La séquence x14y--x14ai a été construite de façon à rejeter un mécanisme à la fois. Les variantes suivantes restent utiles pour reproductibilité ou diagnostic, mais ne font pas partie de la fermeture candidate :
\begin{center}
\small
\begin{tabular}{p{0.10\textwidth}p{0.34\textwidth}p{0.46\textwidth}}
\toprule
Jalon & Principe & Statut et motif\\
\midrule
x14y & Supprimer la soustraction de $p_G$ dans x14v et transférer $\bm J_{\rm raw}$ directement. & Ablation uniquement. Elle réintroduit le double comptage de la pression d'équilibre déjà portée par x6g ; non retenue comme loi physique.\\
x14z & Fermer géométriquement la résultante du seul $p_{\rm ref}$ uniforme sur la polyligne x10n. & Rejeté comme explication du défaut $n=1$ : la fermeture géométrique de $p_{\rm ref}$ est déjà quasi exacte et ne traite pas le décalage Q6/B1 observé.\\
x14aa & Reporter toute la traction thermodynamique absolue sur les faces x6g représentées. & Non retenu. La résultante est alignée sur la géométrie Q6, mais le grand terme uniforme $p_{\rm ref}$ quitte la géométrie locale x10n/CIC qui préserve la dynamique de forme.\\
x14ab & Garder $p_{\rm ref}$ sur x10n et appliquer seulement la pression de jauge sur les faces x6g. & Non retenu. La séparation référence/jauge est plus propre, mais la traction variable reste distribuée sur les normales axiales des faces Q6 plutôt que sur les normales locales x10n.\\
x14ac & Conserver l'échantillonnage local historique et corriger seulement sa résultante globale par projection minimum-$L^2$. & Principe de projection conservé, mais méthode locale supplantée par x14ad : l'association locale pression/segment restait moins cohérente avec les faces x6g réellement représentées.\\
x14ae & Compter les impulsions x14v perdues après échec du support CIC et du fallback rayon 2. & Diagnostic seulement. Résultat nul sur le cas discriminant ; l'hypothèse de perte terminale est rejetée.\\
x14af & Mesurer les résultantes Q6 appliquée, thermodynamique et x14v afin de fermer le bilan global. & Diagnostic seulement. Il identifie au roundoff le mismatch $\bm J_{Q6}^{\rm applied}-\bm J_{\rm thermo}^{\rm x14ad}$ comme source du défaut.\\
x14ai initial & Utiliser la somme B1/RT0 $\sum m\Delta\bm v_{\rm RT0}$ comme cible globale. & Concept conservé mais implémentation supplantée par fix1 : elle oubliait la correction périodique uniforme appliquée ensuite par B1 et produisait des sauts de $\bm P_{\rm tot}$.\\
x14ag & Benchmark de traînée avec inlet/outlet. & Benchmark abandonné : les conditions ouvertes entraînaient une évolution du gaz jusqu'à inversion du courant, ce qui empêchait d'isoler proprement le transfert sur la goutte. Remplacé par x14ah périodique-$x$.\\
\bottomrule
\end{tabular}
\end{center}

\subsection{Paramètres et gates x14 retenus}

Les entrées physiques et les gates runtime qui définissent la chaîne candidate sont résumées ci-dessous. Les gates de prototypes x14z/x14aa/x14ab/x14ac et les diagnostics x14ae/x14af restent OFF en production.
\begin{center}
\footnotesize
\begin{tabular}{p{0.43\textwidth}p{0.47\textwidth}}
\toprule
Entrée & Rôle\\
\midrule
\nolinkurl{speciesThermostatEnable} &
active la remise à température séparée par espèce après collision SRC commune ; défaut \texttt{false}.\\
\nolinkurl{speciesKThermostatTargetKBT} &
cible $k_BT$ de l'espèce $K$ ; valeur négative = héritage global.\\
\nolinkurl{phaseInterfaceKineticBilateralRelocation} &
autorise l'extension positionnelle x10u à une paire explicite liquide/gaz ; défaut \texttt{false}.\\
\nolinkurl{phaseInterfaceKineticReflectionFraction} &
reste fixé à $1$ dans la chaîne x14 bilatérale courante.\\
\nolinkurl{phaseInterfaceASelector}, \nolinkurl{phaseInterfaceBSelector} &
sélectionnent explicitement le type liquide A projeté et le type gaz B non projeté.\\
\nolinkurl{MPCD_Q6_PHASE_GAS_PRESSURE_MODE_0493X6G=eos_accessible_volume} &
corrige l'EOS gazeuse par la fraction de volume accessible dérivée de l'alpha Q6.\\
\nolinkurl{MPCD_X14L_GAS_SPECULAR_REFLECTION} &
active la réflexion spéculaire normale du gaz, tangentielle inchangée.\\
\nolinkurl{MPCD_X14V_GAS_KINETIC_EXCESS_KICK} &
active le transfert au liquide de l'impulsion gazeuse excédentaire.\\
\nolinkurl{MPCD_X14V_SUBTRACT_X6G_THERMODYNAMIC_TRACTION} &
doit rester à son défaut $1$ : soustrait la traction thermodynamique déjà portée par x6g.\\
\nolinkurl{MPCD_X14V_X6G_LOCAL_FACE_GAUGE_PROJECTION} &
active la loi locale x14ad : pression de jauge des faces x6g reportée sur les normales x10n avec projection de résultante.\\
\nolinkurl{MPCD_X14V_DEVICE_APPLIED_Q6_RESULTANT_CLOSURE} &
active x14ai-fix1 : remplace la cible globale par la résultante B1 réellement appliquée après correction périodique ; défaut \texttt{0}, candidate seulement pour composante liquide fermée/isolée.\\
\bottomrule
\end{tabular}
\end{center}

Les contraintes de validation restent strictes. x14v requiert x6g actif, la géométrie x10o + CIC + Q2 + x10u + x10v et le côté gaz x14l. x14ai-fix1 requiert en plus x14ad local-face et interdit simultanément x14z. Les quatre géométries thermodynamiques x14aa/x14ab/x14ac/x14ad sont mutuellement exclusives. Les défauts des nouvelles gates préservent le comportement historique ; les runners de qualification activent explicitement les seules branches nécessaires.


\subsection{Validation externe x14at : profondeur de cavité sous jet gazeux incident, comparaison à Sato et al.}
\label{sec:x14at-sato}

La campagne 0493x14at constitue la première comparaison externe directe de la chaîne liquide--gaz x14 avec une mesure de déformation de surface libre sous jet gazeux incident. La référence retenue est l'expérience eau--air de Sato et al.~\cite{Sato2020}, dans laquelle une lance circulaire de diamètre $d=10\,\mathrm{mm}$ souffle verticalement sur un bain d'eau de profondeur $10d$, dans une cuve de diamètre $20d$, pour des hauteurs de lance $H/d=0.8$ et $1.7$. La profondeur $h_j$ y est mesurée entre le niveau libre initial avant soufflage et le point le plus bas de la cavité. Les auteurs indiquent en outre une dispersion expérimentale typique de l'ordre de $\pm20\,\%$ sur cette mesure et proposent, dans le régime de faible débit (stage~A), la corrélation
\begin{equation}
  \frac{h_j}{d}=1.30\,Fr'_m,
  \qquad Fr'_m<1.17,
  \label{eq:x14at-sato-stageA}
\end{equation}
avec $Fr'_m$ construit à partir de la vitesse du jet au voisinage du bain. Dans les configurations proches de la surface considérées par Sato, cette vitesse reste pratiquement égale à la vitesse de sortie de lance.

Le runner x14at a été construit pour reproduire les grandeurs de similitude qui organisent directement ce benchmark : $L_x/D=20$, profondeur liquide $=10D$, $H/D=0.8$ ou $1.7$, diamètre de jet $D=20h$, et nombre de Bond liquide au diamètre de buse $Bo_D\simeq13.584$. Le gaz explicite utilise la chaîne x6g+x9+x10o/CIC/Q2+x10u/x10v+x12a+x14l+x14v+x14ad ; x14ai est forcé OFF parce que le bain touche des frontières externes, hors du domaine de validité de cette fermeture globale. La comparaison quantitative est effectuée avec le \emph{Froude modifié mesuré} dans le jet incident, et non avec sa cible nominale, afin de comparer le chargement effectivement reçu par le liquide.

Cette similitude reste volontairement partielle. La simulation est bidimensionnelle et utilise $\rho_G/\rho_L=0.1$, alors que le rapport eau--air expérimental est de l'ordre de $1.2\times10^{-3}$. Les viscosités calibrées du modèle donnent, sur les quatre forcings de la campagne, $Re_L\simeq22$--$46$, $Re_G\simeq20$--$43$ et $Oh_L\simeq0.305$ ; ces nombres ne sont pas ceux de l'expérience eau--air. Il faut donc parler ici de \emph{similitude de benchmark} sur la géométrie, $Bo_D$ et $Fr'_m$, et non de similitude dynamique complète en Reynolds, Ohnesorge, Mach et rapport de densités. Cette réserve est particulièrement importante hors du régime stage~A, mais elle n'empêche pas de tester si la fermeture x14 produit le bon ordre de grandeur de cavité lorsque le chargement inertio-gravitaire imposé est équivalent.

\paragraph{Cas quantitatifs exploitables à $H/D=0.8$.}
Pour cette hauteur, le protocole long a montré que la fenêtre $3000$--$4000$ fournit une moyenne représentative : les sentinelles prolongées à 6000 pas donnent un biais de moyenne inférieur à environ $2.2\,\%$. Les quatre résultats sont comparés ci-dessous à la valeur $1.30\,Fr'_m$ évaluée au $Fr'_m$ \emph{mesuré} :
\begin{center}
\small
\begin{tabular}{cccc}
\toprule
$Fr'_{m,\mathrm{mes}}$ & $(h/D)_{\mathrm{SRC}}$ & $(h/D)_{\mathrm{Sato}}=1.30Fr'_{m,\mathrm{mes}}$ & écart relatif \\
\midrule
0.3285 & 0.3180 & 0.4271 & $-25.5\,\%$ \\
0.4491 & 0.4605 & 0.5838 & $-21.1\,\%$ \\
\textbf{0.4880} & \textbf{0.6039} & \textbf{0.6344} & \textbf{$-4.8\,\%$} \\
\textbf{0.5859} & \textbf{0.6861} & \textbf{0.7616} & \textbf{$-9.9\,\%$} \\
\bottomrule
\end{tabular}
\end{center}
Le résultat principal est donc concentré dans la plage $Fr'_m\simeq0.49$--$0.59$ : à chargement modifié équivalent, la cavité prédite par SRC diffère de la corrélation expérimentale de seulement $4.8\,\%$ et $9.9\,\%$. Ces écarts sont inférieurs à la dispersion expérimentale d'environ $\pm20\,\%$ rapportée par Sato. Avec $D=20h$, les deux cas correspondent respectivement à des profondeurs SRC d'environ $12.1h$ et $13.7h$, contre $12.7h$ et $15.2h$ pour la corrélation expérimentale. On obtient donc non seulement la bonne tendance, mais aussi le bon ordre de grandeur géométrique de la cavité relativement au diamètre de buse.

Sur l'ensemble des quatre points $H/D=0.8$, le fit contraint par l'origine donne
\begin{equation}
  \frac{h}{D}=1.1313\,Fr'_{m,\mathrm{mes}},
  \qquad R^2=0.894,
\end{equation}
contre $1.30\,Fr'_m$ pour la référence tridimensionnelle. Les deux plus faibles forcings sont moins concluants : les écarts atteignent $21$--$26\,\%$, et le cas le plus faible ne forme qu'une cavité de quelques cellules de profondeur. Ils sont conservés comme points de tendance mais ne constituent pas les ancres de la validation quantitative.

\paragraph{Pourquoi les cas $H/D=1.7$ ne sont pas utilisés comme validation quantitative.}
Les runs $H/D=1.7$ ont été prolongés jusqu'au pas global 8000. La branche tardive reste nettement plus basse que la corrélation de Sato ; par exemple, sur $7000$--$8000$, les profondeurs sont de l'ordre de $h/D=0.116$, $0.242$, $0.345$ et $0.459$ pour des $Fr'_m$ mesurés d'environ $0.407$, $0.517$, $0.709$ et $0.795$. Cette différence ne doit toutefois pas être interprétée comme un échec physique du couplage x14, car le runner courant n'a pas conservé une géométrie de lance strictement identique entre les deux hauteurs : la longueur de buse est dérivée de $H/D$ pour maintenir la sortie à la bonne distance du bain. Avec la grille de référence, cela revient à utiliser environ $40h$ de buse à $H/D=0.8$ mais seulement $22h$ à $H/D=1.7$. L'expérience de Sato déplace au contraire une même lance. Le développement du jet dans la buse et à sa sortie n'est donc pas contrôlé de façon strictement homologue entre ces deux séries.

À cette réserve géométrique s'ajoute l'absence déjà signalée de similitude complète en $Re$, $Oh$ et rapport de densités. La série $H/D=1.7$ est par conséquent classée \texttt{REVIEW} pour la validation externe : elle démontre la robustesse temporelle du calcul et la monotonie de la réponse au forcing, mais elle ne fournit pas une comparaison expérimentale propre de la profondeur. Une campagne future destinée à fermer ce point devra conserver la longueur et la géométrie de la buse fixes et modifier la hauteur en déplaçant la buse ou en agrandissant le domaine gazeux, puis comparer à nouveau les profondeurs au $Fr'_m$ incident réellement mesuré.

\paragraph{Statut de validation x14at.}
La conclusion exploitable est donc volontairement ciblée :
\begin{quote}
\textbf{pour $H/D=0.8$ et $Fr'_m\simeq0.49$--$0.59$, dans le régime stage~A et avec la similitude géométrique/$Bo_D$/$Fr'_m$ imposée, la profondeur de cavité SRC est quantitativement comparable à celle de Sato, avec un écart de $5$--$10\,\%$, inférieur à la dispersion expérimentale rapportée.}
\end{quote}
Ce résultat constitue une validation externe de l'ordre de grandeur de la réponse normale intégrée de la chaîne liquide--gaz x14. Il ne qualifie pas encore la dépendance à la hauteur de lance ni une similitude eau--air complète ; ces deux questions restent séparées du résultat positif obtenu sur les cas $H/D=0.8$.

\subsection{Statut physique du couplage liquide--gaz au 7 septembre 2026}

Le statut peut être formulé sans requalifier les briques déjà fermées :
\begin{itemize}
  \item la tension superficielle reste la fermeture x9 qualifiée, composée avec la pression gazeuse par $p_L=p_G+\sigma\kappa$ ;
  \item le thermostat séparé permet des masses et températures liquide/gaz différentes sans supprimer la collision SRC commune ;
  \item l'imperméabilité normale du gaz est qualifiée par x14l ;
  \item la pression thermodynamique normale $p_G\rightarrow L$ est qualifiée par x6g+x14s ;
  \item le transfert cinétique normal hors équilibre est porté par x14v ; x14ad est retenu pour sa distribution locale de traction sur les normales x10n ;
  \item x14ai-fix1 ferme le mode de résultante global au roundoff sur la goutte périodique tout en laissant subsister la traînée réelle dans x14ah ;
  \item le domaine de qualification de x14ai-fix1 est limité aux composantes liquides fermées et isolées des frontières Q6 externes ;
  \item la goutte statique ne montre pas de régression mécanique évidente ;
  \item la campagne externe x14at reproduit quantitativement la profondeur de cavité de Sato pour $H/D=0.8$ et $Fr'_m\simeq0.49$--$0.59$, avec des écarts de $4.8\,\%$ et $9.9\,\%$ ; la série $H/D=1.7$ reste \texttt{REVIEW} car la similitude géométrique de la lance n'est pas stricte ;
  \item le transfert de contrainte tangentielle est \emph{PASS-like} sur le Couette à fort signal, sans motif pour ajouter un frottement tangent ad hoc ; la continuité de vitesse interfaciale sub-pourcent n'est pas encore déclarée fermée.
\end{itemize}

La prochaine opération de qualification déjà définie n'est pas un nouveau développement : il s'agit de prolonger exactement la goutte oscillante $n=2$ x14ai-fix1 du pas 2000 au pas global 6000. Ce run doit vérifier simultanément la conservation du moment sur une durée triple, fournir assez de cycles pour une estimation robuste de $G_\omega$ et $\beta$, et déterminer si le déplacement relatif du centre de masse liquide est borné ou séculaire. Aucune modification de la chaîne liquide, aucun nouveau diagnostic et aucun terme tangent artificiel ne sont justifiés avant ce résultat.


\section{Audit documentaire du point de référence surface libre au 31 août 2026}
\label{sec:inventory-audit-7655}

Le point de référence documentaire est le commit \nolinkurl{7655b81b1b2fd16eecefa8d8b3bebac4cd9f87f1}, tag \nolinkurl{surf-tension-qualified-x13h-20260831}. L'audit a comparé les deux inventaires consolidés du 26 août avec le parseur de paramètres et les lectures directes de variables d'environnement du code source correspondant. Le snapshot de travail ultérieur a été utilisé uniquement comme sur-ensemble pour détecter les noms apparus pendant les expériences ; chaque nom absent de l'inventaire a ensuite été classé selon sa présence ou son absence dans le code commité.

Pour l'inventaire des paramètres, \nolinkurl{src/params_io_base.cpp} contient 442 clés de parseur distinctes. 430 possèdent une ligne de nom canonique directe dans l'inventaire et les 12 restantes sont explicitement couvertes comme aliases de champs canoniques. Il n'existe donc aucune clé parseur non documentée au commit de référence. Deux exemples importants sont les couples d'alias suivants :
\begin{center}
\small
\nolinkurl{kineticReflectionFraction} $\rightarrow$ \nolinkurl{phaseInterfaceKineticReflectionFraction}\\
\nolinkurl{evaporationTargetType} $\rightarrow$ \nolinkurl{phaseInterfaceEvaporationTargetType}.
\end{center}
Ils ne nécessitent pas des entrées canoniques supplémentaires.

Pour les variables d'environnement lues par \texttt{getenv} ou par les helpers typés des sources C++/CUDA, le snapshot expérimental contient 224 noms littéraux. Douze n'étaient pas présents dans l'inventaire du 26 août. Onze appartiennent aux branches expérimentales postérieures et sont absents du commit qualifié :
\begin{center}
\begin{minipage}{0.92\textwidth}
\small\raggedright
\nolinkurl{MPCD_X10_KINETIC_INTERFACE_FREE_CROSSING_CELLS},\linebreak
\nolinkurl{MPCD_X10_KINETIC_INTERFACE_ORPHAN_TURNOVER},\linebreak
\nolinkurl{MPCD_X13R_INTERFACE_CELL_COOLING},\linebreak
\nolinkurl{MPCD_X13S_INTERFACE_ANISOTROPIC_COOLING},\linebreak
\nolinkurl{MPCD_X13S_INTERFACE_DEPTH_CELLS},\linebreak
\nolinkurl{MPCD_X13S_TANGENTIAL_VELOCITY_FRACTION},\linebreak
\nolinkurl{MPCD_X13T_INTERFACE_RECOVERY_CELLS},\linebreak
\nolinkurl{MPCD_X13T_ONE_FOR_ONE_SIMPLE},\linebreak
\nolinkurl{MPCD_X13T_THIN_RADIUS_CELLS},\linebreak
\nolinkurl{MPCD_X13T_UNIFIED_PROGRESSIVE_INTERFACE_COOLING} et\linebreak
\nolinkurl{MPCD_X13W_ESCAPE_RESEED}.
\end{minipage}
\end{center}
Le douzième nom, \nolinkurl{MPCD_X10_KINETIC_INTERFACE_ONE_FOR_ONE_NORMAL_ONLY}, est au contraire bien présent dans le commit \texttt{7655b81}. Il a été ajouté à l'inventaire révisé avec le statut ``ablation x13o, OFF dans la chaîne qualifiée''. Après cette correction, les 213 noms de cet ensemble appartenant au code commité sont couverts par l'inventaire.

L'audit vérifie aussi le contrat de production : les runners x13k et x13n du commit activent \texttt{ONE\_FOR\_ONE=1}, \texttt{ONE\_FOR\_ONE\_SWAP=1} et x12a, mais ne demandent pas le normal-only. Le code sélectionne alors explicitement la sémantique \nolinkurl{position-mirror+local-equal-mass-full-vector-swap}. Le titre du commit ne doit donc jamais être utilisé comme source unique pour décrire la fermeture active.

Les fichiers d'inventaire joints à cette version du rapport sont :
\begin{center}
\small
\nolinkurl{src_mpcd_env_flags_inventory_commit_7655b81_310826.csv}\\
\nolinkurl{src_mpcd_params_inventory_commit_7655b81_310826.csv}.
\end{center}
Le premier contient 507 entrées après ajout de x13o ; le second conserve 848 entrées et passe l'audit 442/442 des clés parseur.


\subsection{Extension de l'audit aux ajouts x14 du snapshot du 4 septembre 2026}

Le snapshot \nolinkurl{snap_040926.zip} contient toujours 444 clés littérales directes reconnues par le parseur, contre 442 au point de référence du 31 août. Les deux clés canoniques ajoutées par x14 restent
\begin{center}
\small
\nolinkurl{speciesThermostatEnable}
\qquad\text{et}\qquad
\nolinkurl{phaseInterfaceKineticBilateralRelocation}.
\end{center}
La famille dynamique \nolinkurl{speciesKThermostatTargetKBT} reste indexée comme les déclarations \nolinkurl{speciesK}. Aucun jalon x14y--x14ai-fix1 n'ajoute de nouvelle clé \texttt{params.kv} : ces expériences changent uniquement des gates runtime du backend et réutilisent la géométrie, les sélecteurs de phases et les paramètres thermiques déjà documentés. L'inventaire paramètres conserve donc 854 lignes et couvre les 444/444 clés littérales directes du parseur, dont 12 par leurs aliases canoniques ; les lignes x14 concernées sont seulement réannotées pour tracer l'audit du 4 septembre.

L'audit des lectures runtime a été refait sur les appels \texttt{getenv} et helpers \texttt{env\_*} littéraux réellement présents dans \nolinkurl{src/} et \nolinkurl{include/}. Il trouve 354 noms runtime préfixés \nolinkurl{MPCD_} ou \nolinkurl{SRC_MPCD_} distincts, tous couverts après ajout de neuf gates x14 absentes de l'inventaire du 3 septembre :
\begin{center}
\begin{minipage}{0.94\textwidth}
\small\raggedright
\nolinkurl{MPCD_X14V_SUBTRACT_X6G_THERMODYNAMIC_TRACTION},\linebreak
\nolinkurl{MPCD_X14V_REFERENCE_PRESSURE_GEOMETRIC_CLOSURE},\linebreak
\nolinkurl{MPCD_X14V_X6G_FACE_THERMO_TRACTION},\linebreak
\nolinkurl{MPCD_X14V_X6G_GAUGE_FACE_THERMO_TRACTION},\linebreak
\nolinkurl{MPCD_X14V_X6G_GAUGE_RESULTANT_PROJECTION},\linebreak
\nolinkurl{MPCD_X14V_X6G_LOCAL_FACE_GAUGE_PROJECTION},\linebreak
\nolinkurl{MPCD_X14V_SCATTER_LOSS_DIAGNOSTIC},\linebreak
\nolinkurl{MPCD_X14V_GLOBAL_BALANCE_DIAGNOSTIC},\linebreak
\nolinkurl{MPCD_X14V_DEVICE_APPLIED_Q6_RESULTANT_CLOSURE}.
\end{minipage}
\end{center}
Le premier de ces flags a un défaut $1$ et doit rester actif dans la chaîne physique x14v ; les variantes x14z/x14aa/x14ab/x14ac et les diagnostics x14ae/x14af ont un défaut $0$. Le gate x14ad et le gate x14ai ont également un défaut $0$ afin de préserver la non-régression ; ils sont activés explicitement ensemble dans le chemin candidat x14ai-fix1. Les modes x14aa/x14ab/x14ac/x14ad sont mutuellement exclusifs dans le code.

Les inventaires associés à cette extension sont
\begin{center}
\small
\nolinkurl{src_mpcd_params_inventory_snapshot_040926_x14ai.csv}\\
\nolinkurl{src_mpcd_env_flags_inventory_snapshot_040926_x14ai.csv}.
\end{center}
L'inventaire environnement contient désormais 524 lignes au total. Ce nombre inclut les aliases et contrôles de runners déjà présents dans l'inventaire documentaire ; le contrôle source proprement dit porte sur les 354 lectures runtime préfixées littérales et passe 354/354. L'inventaire paramètres reste à 854 lignes et passe 444/444 sur les clés littérales directes du parseur, en plus des familles dynamiques explicitement documentées.


\section{Conclusion}

Le statut 0493x7q modifie substantiellement la conclusion du cycle 0493x7e. Q6-g-f n'est plus seulement une fermeture de dam-break : le chemin signé est maintenant qualifié sur plusieurs familles de conditions limites et son défaut de moment uniforme a été identifié puis corrigé au niveau exact où il naissait, c'est-à-dire dans l'échantillonnage particulaire de la reconstruction B1. Le probe périodique démontre la neutralité du gradient de pression sur le mode $k=0$ au roundoff ; Poiseuille retrouve un profil et une viscosité cohérents ; Taylor--Green, bend-pipe et same-face IO ne montrent pas de régression structurelle ; le dam-break conserve son chemin partiel historique et sa dynamique qualitative.


Le cycle 0493x8 ajoute une validation externe quantitative à cette qualification interne, désormais reprise sur le domaine bibliographique complet. Les inlets segmentés possèdent un profil de Poiseuille réellement local au segment, tandis que la sortie \texttt{neumann} combine une continuation cinétique locale et une sortie de pression passive conditionnée pour le CG. Le run cluster CCUB $4400\times400$, $15D$ en amont et $40D$ en aval, a été prolongé jusqu'à $75000$ pas en trois segments restartés. L'inspection des champs montre que l'allée n'est pleinement développée sur la longue zone aval qu'au troisième segment ; la comparaison finale utilise donc les pas $59001\ldots74901$. Sur cette fenêtre, $f_{\rm POD}=0.0969335$ et $f_{\rm probe}=0.0974010$ diffèrent de seulement $0.481\,\%$, avec $R^2_{\rm phase}=0.999923$, $R^2_{\rm probe}=0.9416$ et une capture POD de $73.1\,\%$. Le débit bulk à $-3D$ vaut $0.120095201$, soit $-0.6158\,\%$ de la cible, et les dispersions amont de $Q_v$, $\bar\rho$ et $J_\rho$ restent sous $0.224\,\%$. La calibration constitutive Q6-g-f à huit graines donne $\nu=6.65127\times10^{-4}$, $CV=4.90\,\%$ et $Re_H=282.125$ avec la vitesse mesurée. Le run est donc effectivement centré sur le point $Re_H\simeq280$ de Zovatto--Pedrizzetti. La période $T^*=0.791017$ reste $4.70\,\%$ sous la valeur $0.83$ : la réserve de viscosité, de débit et d'établissement spatial est levée, et le résidu devient un problème de précision numérique/fermeture et de durée statistique. La fenêtre stricte ne contenant que $3.085$ cycles, le résultat doit être lu comme $T^*\simeq0.79$ avec un écart de l'ordre de $4$--$5\,\%$, et non comme une métrologie sub-pourcent. L'ancien accord à environ $2\,\%$ obtenu en domaine réduit ne doit plus être présenté comme le résultat quantitatif final.

La chaîne de production retenue est donc \texttt{prestream\_single\_fused} + \texttt{free\_surface\_masked} + stencil x6f/x6g si gaz présent + B1 + restauration de densité signée cohérente + CG x7j résident + fermeture périodique exacte x7q lorsque son domaine d'application est satisfait. Le resampling et le viriel restent des modules séparés et opt-in : ils ne sont pas requis pour reproduire les validations finales de cette chaîne.


Le cycle x14 complète maintenant cette hiérarchie en remplaçant le côté ``gaz/vide'' abstrait par une phase gazeuse particulaire explicite lorsque l'application le demande. La capillarité reste portée par x9 et la pression gazeuse par x6g ; x14 ne superpose pas une seconde loi de tension superficielle. Le thermostat par espèce conserve une collision SRC commune au mélange, x14l ferme l'imperméabilité normale du gaz sans modifier sa vitesse tangentielle, x14s corrige l'EOS interfaciale pour le volume réellement accessible, et x14v restitue au liquide la seule impulsion cinétique normale qui n'était pas déjà comptée par x6g. La campagne x14y--x14af a ensuite séparé le problème local du problème global : x14ad est retenu pour échantillonner la pression de jauge x6g sur les faces représentées tout en appliquant la traction sur les normales x10n, et x14af démontre que le défaut de translation résiduel est exactement le mismatch entre cette résultante thermodynamique et la résultante Q6 effectivement transmise par B1/RT0. x14ai-fix1 ferme ce seul mode de résultante avec la quantité de mouvement B1 post-périodique réellement appliquée. Sur la goutte $n=2$, $G_\omega=0.97988$, $R^2=0.9971$ et la dérive du moment total reste sous $5.1\times10^{-11}$ ; sur x14ah, la goutte parcourt $0.799R$ vers l'aval et la dérive transverse est réduite d'environ un facteur 17 sans suppression de la traînée. Le coût de x14ai est non résolu sous le bruit de chronométrage et classé inférieur à $2\,\%$. Cette fermeture devient donc candidate pour les composantes liquides fermées et isolées des frontières Q6 externes, mais elle ne doit pas être généralisée aux ménisques, nappes attachées ou liquides touchant une ouverture sans qualification spécifique. Le Couette x14w indique en parallèle que la contrainte tangentielle peut être transmise par le SRC commun sans terme de frottement interfacial ajouté : $a_L/a_G=0.23555$ pour une référence $\mu_G/\mu_L\simeq0.23769$ ; la fermeture sub-pourcent de la continuité de vitesse interfaciale reste une question distincte.


La campagne x14at ajoute à cette qualification interne une validation externe ciblée sur la profondeur de cavité créée par un jet gazeux. Pour $H/D=0.8$, $D=20h$, une profondeur de bain $10D$ et un nombre de Bond de buse calé sur l'expérience de Sato, les deux cas dont le Froude modifié incident mesuré vaut $Fr'_m=0.488$ et $0.586$ donnent respectivement $h/D=0.604$ et $0.686$, contre $0.634$ et $0.762$ pour la corrélation expérimentale $h/D=1.30Fr'_m$. Les écarts, $-4.8\,\%$ et $-9.9\,\%$, sont inférieurs à la dispersion expérimentale rapportée. Cette validation porte sur une similitude partielle géométrie/$Bo_D$/$Fr'_m$ en stage~A, et non sur une similitude complète en Reynolds, Ohnesorge ou rapport de densités. La série $H/D=1.7$ n'est pas utilisée comme validation quantitative, car le runner courant modifie la longueur de buse lorsqu'il change la hauteur ; une reprise à géométrie de lance fixe est nécessaire pour qualifier l'effet de $H/D$ proprement dit.

Le jalon 0493w8 de la branche \texttt{surf}, au commit \texttt{caa5b6b}, a ajouté un jalon majeur à cette hiérarchie : Q6 n'est plus limité à une correction barycentrique commune dans les mélanges. Le mode \texttt{independent\_masked} fournit une projection distincte par espèce, strictement sélective, CUDA résidente et compatible avec les familles de conditions limites qualifiées du backend Q6 résident. Les smokes 0493w5--0493w7 et la qualification Taylor--Green 0493w8 établissent la non-régression monoespèce, la neutralité des étiquettes de type sous SRC, l'accord de transport d'un fluide dual identique et l'absence de correction directe du gaz lorsque sa force Q6 vaut zéro.

Le cycle 0493x a depuis levé une part importante de cette limite sur le cas de surface libre fermé. La difficulté temporelle a été séparée de la difficulté de support : Q6-g projette la vitesse tentative incluant la force avant le streaming, tandis que \texttt{free\_surface\_masked} et le stencil x6f distinguent le carrier numérique de l'interface physique $\alpha=0.5$. Le couplage x6g impose $p_l|_\Gamma=p_g$ sans modifier la matrice CG. La projection Q6-g-f ajoute la reconstruction face--particule B1 et une restauration de densité directement dans le RHS, paramétrée par la constante de temps $\tau_\rho$. Les runs longs appariés $300\times150/5000$ et $600\times300/10000$ atteignent $t=25$ avec une erreur volumique finale d'environ $5$--$6\,\%$, sans dérive séculaire, avec ou sans pression gazeuse. La calibration montre que ces deux mailles ne conservent pas le même Reynolds ; le résultat qualifie donc la robustesse de la fermeture volumique, non la convergence hydrodynamique complète du dam-break. La limite immédiate n'est plus l'incompressibilité du liquide ni la définition d'une pression gazeuse interfaciale, mais la reconstruction suffisamment régulière de la courbure et la qualification du saut capillaire.



Le dernier développement GPU ajoute un jalon distinct. La branche \texttt{SRC\_GPU} dispose maintenant d'un chemin \textbf{SRC/MPCD classic full CUDA résident}. Ici, \emph{SRC classic} signifie explicitement advection/streaming, décalage aléatoire de grille, collision/rotation SRC et thermostat. Les validations consolidées couvrent les familles périodique, wall-simple, solide rectangle, solide circulaire, piston/mobile wall, inlet/outlet full-face, inlet/outlet segmenté, et cercle avec inlet/outlet pour le cas Von Karman. Les ouvertures CUDA disposent en outre de régimes explicites \texttt{neumann}, \texttt{equilibrium\_flux} et \texttt{forced\_flux}, utilisables en full-face ou segmenté dans le périmètre SRC classic-only. Les modules de fermeture restent activables séparément : Q6 CUDA résident est désormais qualifié en mono-espèce et, via \texttt{independent\_masked}, en multi-espèces sur les familles périodique, canal à parois, entrée/sortie pleine face, entrée/sortie segmentée et Darcy--Brinkman ; le resampling et le viriel conservent leurs propres domaines de qualification et leurs garde-fous.

Les développements post-0143 ajoutent à cette hiérarchie un autre résultat important. Les conditions limites segmentées permettent désormais plusieurs inlets/outlets sur une même face solide, ce qui généralise fortement les cas de cuve, injection localisée ou injection multi-type marquée. La suppression de \texttt{method} rend les ablations propres : Q6 et resampling sont des briques indépendantes, activées par \texttt{projectionEnable} et \texttt{resamplingEnable}. Enfin, la réponse de capacité fermée fournit une voie continue entre régime quasi-incompressible et régime faiblement compressible : une surmasse globale affaiblit Q6, augmente le viriel effectif, modifie la cible de remap pour conserver la compression, puis projette \(P_{\rm kin}+P_{\rm vir}\) en charge pariétale. La validation statique de boîte uniformément surcompressée montre que la pression pariétale est uniforme par face et que la force nette est nulle, ce qui constitue un jalon majeur pour un futur couplage à des parois déformables ou rompables.



Le dernier cycle CUDA ajoute maintenant une réponse partielle mais validée à cette difficulté. Le resampling post-SRC CUDA ne modifie pas la collision SRC ; il reconditionne la représentation particulaire après que le pas physique a produit l'état fluide. Les modules 0295--0299 établissent une chaîne survey, reconditionnement de masse, split/merge local, restauration $M/P/K_{\mathrm{rel}}$ et filtrage boundary-aware. Les campagnes 0300--0302 montrent que le backward step, qui perd jusqu'à $191$ cellules humides vides à $U_{\mathrm{in}}=0.60$ en SRC classic, n'en conserve plus aucune avec le réglage nominal \texttt{12:20:32}, alors que les erreurs locales de conservation restent au niveau du roundoff. Ce résultat donne une première validation physique du resampling CUDA comme contrôle de support particulaire, distinct de Q6 et complémentaire du chemin SRC classic résident.


Le cycle 0490a--0490p généralise maintenant cette chaîne aux mélanges. La conservation est imposée par \texttt{type}, depuis le split/merge local jusqu'au refill, au plan donneur--receveur et à la fermeture de masse. Le chemin CUDA résident élimine le plan CPU, les vecteurs d'opérations hôte, les copies complètes de rollback, les scans de particules pour les dépôts et le pool, puis le dernier miroir CPU wet/poor/rich. La validation à 10000 pas et l'audit \texttt{remaining\_cpu\_scope=none} constituent un jalon d'architecture : le resampling multi-espèces peut désormais rester sur l'état CUDA partagé dans le domaine périodique qualifié. La projection Q6 dépendante de la phase est disponible via \texttt{independent\_masked}. Les difficultés alors identifiées sur la condition physique d'interface et les supports partiels ont depuis été traitées, pour le cas fermé de surface libre, par \texttt{free\_surface\_masked}, la géométrie x6c, le stencil x6f et la condition $p_l=p_g$ x6g. L'élargissement du resampling multi-espèces aux frontières et aux solides reste en revanche un chantier distinct, à conduire sans rompre la résidence CUDA ni la conservation par espèce.

Le projet SRC/MPCD quasi-incompressible aboutit maintenant à une hiérarchie plus claire des méthodes. La redistribution dure et la fermeture virielle historique ont joué un rôle exploratoire, mais ne constituent plus le socle principal de validation. La génération Q6/Q9 a fourni une formulation plus robuste de l'incompressibilité : collision SRC classique, projection de vitesse, correction basse fréquence du flux de masse, thermostat corrigé, diagnostics de conservation et traitement explicite des parois par particules virtuelles.

La difficulté décisive n'était toutefois pas seulement la divergence du champ de vitesse ou du flux de masse. Les cas de poches vides, de remplissage et d'inlet/outlet ont montré que la méthode doit garantir un support particulaire local suffisant. Une cellule pauvre ne doit pas être traitée comme un puits de pression ou une anomalie eulérienne à corriger par force ; elle doit être traitée comme un défaut de représentation particulaire locale.

La formulation SRC/MPCD pondérée avec resampling fournit cette réponse. Dans sa forme stabilisée, elle combine :
\begin{equation}
  \text{rôles particulaires}
  + \text{pool préalloué}
  + \text{garde }N_{\min}/N_{\mathrm{target}}/N_{\max}
  + \text{masses pondérées}
  + \text{remap }M,P,E
  + \text{Q6}.
\end{equation}
Le point essentiel est l'ordre logique : on maintient d'abord la population active des cellules mouillées, puis on impose la masse et le moment par remap pondéré. La correction de masse seule est insuffisante, car elle peut masquer un dépeuplement local par des particules trop lourdes.

Le portage C++/OpenMP constitue une étape majeure : il démontre que la méthode peut être organisée dans une architecture compatible avec une exécution parallèle, des états persistants \texttt{.smpcd}, des scripts MATLAB de préparation et d'analyse, et des scripts \texttt{bash} de lancement reproductibles. Les validations récentes indiquent que Taylor--Green, Poiseuille \texttt{wallVP}, le cylindre périodique, la marche, l'obstacle ouvert et le remplissage d'un domaine initialement inactif se comportent de façon cohérente. Le cas Poiseuille montre que la physique de canal est préservée ; le cas injection/fill montre que la couche wet/dry et le pool actif/inactif sont opérationnels ; le cas obstacle ouvert a joué un rôle utile en révélant la nécessité du garde de population maintenant réintroduit.

Le statut actuel peut donc être résumé ainsi : le socle SRC CUDA résident est validé sur les principales familles de conditions limites ; Q6 monoespèce et Q6 multi-espèces sont préservés ; Q6-g-f emploie le séquencement force-aware \texttt{prestream\_single\_fused}, le masque de support liquide \texttt{free\_surface\_masked}, le stencil x6f, la condition gazeuse x6g lorsqu'une phase gaz est enregistrée, B1 et la restauration de densité signée cohérente. Le CG x7j maintient le solveur elliptique dans le chemin CUDA résident et x7q ferme exactement le moment $k=0$ dans les directions périodiques du domaine monophase complet. Le cycle x8 ajoute un inlet segmenté Poiseuille local et une sortie Neumann cinétique--pression passive, puis valide leur combinaison sur un cylindre confiné. Le cycle x9 ajoute la courbure p3/Scharr, le saut de Laplace $p_A-p_B=\sigma\kappa$, les sélecteurs de phases et le cutoff de courbure x9r.

Les cycles x10--x12 ont depuis révisé la fermeture cinétique de cette surface libre. La chaîne de travail n'est plus la réflexion x10q seule : elle combine x10o, un alpha CIC cinétique distinct de x6c, une reconstruction Q2, x10p/x10q, la relocalisation one-for-one x10u, le swap local conservatif x10v et le refroidissement thermique local x12a. Les variantes x10j/k/m/n/r/s/t et x10w sont conservées comme ablations mais restent désactivées dans les runners x12 courants. La calibration mécanique x12yl donne, pour $\sigma_{\rm declared}=945$, $\sigma_{\rm eff}=940.965$, $G_\sigma=0.99573$ et $R^2=0.999978$. La requalification x13h complète désormais ce résultat : les gouttes oscillantes 2-D $n=2,3,4$ ont des fréquences proches de la théorie, alors que Taylor--Culick reste trop lent ($G_{\rm TC}\simeq0.81$ à $\sigma=10000$ et $\simeq0.70$ à $\sigma=2500$) malgré une résultante capillaire locale du rim proche de $2\sigma$. Le défaut n'est donc plus formulé comme une simple mauvaise amplitude de tension superficielle ni comme un ralentissement universel de x10v : il concerne le couplage dynamique soutenu qui convertit la traction capillaire en flux et stockage de quantité de mouvement dans le rim. Les campagnes splash $R/h=40$ confirment par ailleurs que le choix $R_{\min}=4h$ conserve les mécanismes d'impact, tandis que x12d prépare la comparaison JFM 524 à $We\simeq514$ avec géométrie d'obstacle dimensionnée.

La campagne x13r--x13zd menée ensuite doit être conservée comme historique d'invalidation : refroidissements local/anisotrope/progressif, variantes de rétention, reseed et changements de grille peuvent améliorer Taylor--Culick tout en contractant artificiellement le support d'une goutte fermée. Le test x13zd a invalidé la fermeture x13t+x13w comme chemin général. Le dépôt a donc été ramené au commit \texttt{7655b81}, tag \texttt{surf-tension-qualified-x13h-20260831}; le smoke de rollback retrouve $G_{\omega,n=2}=1.00151$ et $G_{\rm TC}=0.79457$.


Le complément Darcy/chiVP modifie également le statut des géométries complexes. Le champ $\chi$ peut maintenant être traité comme une donnée géométrique externe de même rang pratique que l'état particulaire initial : il décrit la forme, désactive les particules initiales dans le solide, pilote le frein Darcy moyen, oriente le bain de réémission et fournit les moments virtuels de collision d'interface. La correction 0426 des fastflags CUDA résidents montre que cette généralité n'impose pas de pénalité majeure de performance lorsque les scripts sélectionnent le backend résident optimisé. Cette observation est importante pour les futurs chantiers d'optimisation topologique, car elle rend réaliste une boucle externe qui modifie $\chi$ sans recompilation du solveur.

La conclusion méthodologique est maintenant plus précise : rendre SRC/MPCD quasi-incompressible ne consiste ni à projeter un champ de vitesse isolé, ni à corriger toute difficulté par resampling ou par viriel. Quatre niveaux doivent rester distincts : la vitesse de transport et son séquencement force--projection, la qualité du support particulaire, la loi de contrainte mécanique aux interfaces/parois, et l'échange cinétique éventuel de masse au travers d'une interface. Le cycle 0493x a séparé les premiers niveaux sur le dam-break ; x9 a ajouté la contrainte capillaire via Q6-g-f ; x10--x12 montrent qu'une fermeture de support peut conserver masse, impulsion et énergie tout en introduisant encore une dispersion propre qu'il faut mesurer séparément. x12yl établit désormais la valeur mécanique de $\sigma$, tandis que x12cal quantifie la dynamique dispersive de la chaîne complète.

Le cycle 0493x13 modifie en outre le statut du verrou Reynolds. Le fluide n'est plus traité comme une boîte noire dont la viscosité est extraite d'un run unique : le calibrateur standalone path-aware, le Taylor--Green multi-graines, le longitudinal amorti et le MSD permettent de construire $\nu_T$, $\nu_L$, $c_s$ et $D_{\rm self}$ sur le chemin numérique réellement utilisé. Le réaudit final a déclassé le cisaillement transverse pur comme métrologie absolue en raison de sa dépendance d'échelle et rétabli Taylor--Green comme référence constitutive de $\nu_T$. A1 ($\gamma=20$, $120^\circ$, $\lambda/h=0.48$) fournit la meilleure portée brute parmi les points \texttt{PASS}, tandis que G08-0.72 reste la référence économique pour les très grandes grilles : $c_s\simeq0.3551$, $\nu_T^{\rm design}\simeq5.60\times10^{-4}$, $\nu_L\simeq1.01\times10^{-3}$ et $H_h\simeq2.48$, avec une dispersion de viscosité qui impose encore un statut \texttt{REVIEW}. Sa propagation acoustique reste qualifiée pour un usage compressible non linéaire jusqu'à environ $Ma=0.5$ ; les régimes plus rapides atteignent progressivement la frontière basse densité. La calibration Q6-g-f à huit graines du fluide VK démontre en parallèle que SRC et Q6-g-f ont la même viscosité transverse à la précision actuelle ($-0.45\,\%$ d'écart central) et fixe la fenêtre pleinement développée du benchmark cluster à $Re_H=282.125$. Cette caractérisation rend désormais possible un dimensionnement raisonné des cas JFM, VK et des futures applications à grand Reynolds avant d'engager les très grandes grilles.

La priorité de maintenance est désormais la stabilité du point de référence, pas la poursuite d'une optimisation locale de Taylor--Culick. Tout futur chantier de surface libre doit repartir du tag \texttt{surf-tension-qualified-x13h-20260831}, conserver les validations gouttes oscillantes et Taylor--Culick comme tests couplés, et démontrer explicitement la conservation du support de phase avant d'être comparé à la référence. Les études encore ouvertes sont la convergence capillaire à fluide constitutif rematché, la qualification JFM à similitude Reynolds/Weber plus complète, le mouillage sur géométrie $\chi$ et, séparément, une véritable thermodynamique d'évaporation/condensation. Les branches x13t/x13w/x13x ne sont pas des préimages de production.




\begin{thebibliography}{99}

\bibitem{MalevanetsKapral1999}
A. Malevanets and R. Kapral,
``Mesoscopic model for solvent dynamics,''
\emph{The Journal of Chemical Physics}, vol. 110, no. 17, pp. 8605--8613, 1999.
doi: \href{https://doi.org/10.1063/1.478857}{10.1063/1.478857}.

\bibitem{GompperIhleKrollWinkler2009}
G. Gompper, T. Ihle, D. M. Kroll, and R. G. Winkler,
``Multi-Particle Collision Dynamics: A Particle-Based Mesoscale Simulation Approach to the Hydrodynamics of Complex Fluids,''
in \emph{Advanced Computer Simulation Approaches for Soft Matter Sciences III},
Advances in Polymer Science, vol. 221, Springer, Berlin, Heidelberg, pp. 1--87, 2009.
See also arXiv:0808.2157.

\bibitem{IhleTuzelKroll2006}
T. Ihle, E. Tüzel, and D. M. Kroll,
``Consistent particle-based algorithm with a non-ideal equation of state,''
\emph{Europhysics Letters}, vol. 73, no. 5, pp. 664--670, 2006.

\bibitem{ZantopStark2021a}
A. W. Zantop and H. Stark,
``Multi-particle collision dynamics with a non-ideal equation of state. I,''
\emph{The Journal of Chemical Physics}, vol. 154, 024105, 2021.

\bibitem{ZantopStark2021b}
A. W. Zantop and H. Stark,
``Multi-particle collision dynamics with a non-ideal equation of state. II,''
\emph{The Journal of Chemical Physics}, vol. 155, 134904, 2021.

\bibitem{Chorin1968}
A. J. Chorin,
``Numerical solution of the Navier--Stokes equations,''
\emph{Mathematics of Computation}, vol. 22, no. 104, pp. 745--762, 1968.

\bibitem{SolenthalerPajarola2009}
B. Solenthaler and R. Pajarola,
``Predictive-corrective incompressible SPH,''
\emph{ACM Transactions on Graphics}, vol. 28, no. 3, article 40, 2009.

\bibitem{DeCourcy2024}
J. J. De Courcy, T. C. S. Rendall, L. Constantin, B. Titurus, and J. E. Cooper,
``Incompressible $\delta$-SPH via artificial compressibility,''
\emph{Computer Methods in Applied Mechanics and Engineering}, vol. 420, article 116700, 2024.

\bibitem{Yang2021}
K. Yang, T. Aoki, and co-authors,
``Weakly compressible Navier--Stokes solver based on evolving pressure projection method for two-phase flow simulations,''
\emph{Journal of Computational Physics}, vol. 431, article 110113, 2021.

\bibitem{WinklerHuang2009}
R. G. Winkler and C.-C. Huang,
``Stress tensors of multiparticle collision dynamics fluids,''
\emph{The Journal of Chemical Physics}, vol. 130, 074907, 2009.

\bibitem{ImperioPaddingBriels2011}
A. Imperio, J. T. Padding, and W. J. Briels,
``Force calculation on walls and embedded particles in multi-particle collision dynamics simulations,''
\emph{Physical Review E}, vol. 83, 046704, 2011.

\bibitem{ValizadehMonaghan2015}
A. Valizadeh and J. J. Monaghan,
``A study of solid wall models for weakly compressible SPH,''
\emph{Journal of Computational Physics}, vol. 300, pp. 5--19, 2015.

\bibitem{Islam2024}
M. R. I. Islam,
``SPH-based framework for modelling fluid--structure interaction problems with finite deformation and fracturing,''
\emph{Ocean Engineering}, vol. 294, article 116722, 2024.

\bibitem{Singh2025}
V. Singh, C. Peng, and M. R. I. Islam,
``Three-dimensional SPH modeling of brittle fracture under hydrodynamic loading,''
\emph{Computer Methods in Applied Mechanics and Engineering}, 2025.


\bibitem{Brinkman1949}
H. C. Brinkman,
``A calculation of the viscous force exerted by a flowing fluid on a dense swarm of particles,''
\emph{Applied Scientific Research}, vol. 1, pp. 27--34, 1949.
doi: \href{https://doi.org/10.1007/BF02120313}{10.1007/BF02120313}.

\bibitem{BorrvallPetersson2003}
T. Borrvall and J. Petersson,
``Topology optimization of fluids in Stokes flow,''
\emph{International Journal for Numerical Methods in Fluids}, vol. 41, no. 1, pp. 77--107, 2003.
doi: \href{https://doi.org/10.1002/fld.426}{10.1002/fld.426}.

\bibitem{ZovattoPedrizzetti2001}
L. Zovatto and G. Pedrizzetti,
``Flow about a circular cylinder between parallel walls,''
\emph{Journal of Fluid Mechanics}, vol. 440, pp. 1--25, 2001.
doi: \href{https://doi.org/10.1017/S0022112001004608}{10.1017/S0022112001004608}.

\bibitem{SahinOwens2004}
M. Sahin and R. G. Owens,
``A numerical investigation of wall effects up to high blockage ratios on two-dimensional flow past a confined circular cylinder,''
\emph{Physics of Fluids}, vol. 16, no. 5, pp. 1305--1320, 2004.
doi: \href{https://doi.org/10.1063/1.1668285}{10.1063/1.1668285}.


\bibitem{Sato2020}
S. Sato, M. Ando, J. Okada, Y. Ueda, and M. Iguchi,
``Prediction of Plunging Depth Induced by Top Lance Gas Blowing onto a Low-Melting-Point Metal Bath,''
\emph{ISIJ International}, vol. 60, pp. 1675--1683, 2020.
doi: \href{https://doi.org/10.2355/isijinternational.ISIJINT-2019-790}{10.2355/isijinternational.ISIJINT-2019-790}.

\bibitem{JosserandLemoyneTroegerZaleski2005}
C. Josserand, L. Lemoyne, R. Troeger, and S. Zaleski,
``Droplet impact on a dry surface: triggering the splash with a small obstacle,''
\emph{Journal of Fluid Mechanics}, vol. 524, pp. 47--56, 2005.
doi: \href{https://doi.org/10.1017/S0022112004002393}{10.1017/S0022112004002393}.

\end{thebibliography}

\end{document}
